%% =====================================================================
%%  2026-11-05-途中経過：DAI に出す研究（少数のターゲットラベルでのマルチモーダル感情認識の転移）の途中経過
%%  日本語版 main_ja.tex と英語版 main_en.tex が同じこのファイルを読む．文は \T{日本語}{English} で書く
%%  数値は figures/numbers.tex から \N{key} で引く（scripts/export_figdata.py が結果ファイルから書く．手で数を書かない）
%%  図の値は figures/*.csv（同じ台本が書く．各 CSV の先頭に出どころ）
%%  模式図（設定，手法の流れ，位置と向き，重みづけ，LOO）は実験の値ではない（図の下に書く）
%% =====================================================================
\input{preamble}   % common/preamble.tex（日本語・テーマ・パッケージ・マクロ）
\input{merda}      % common/merda.tex（上端の帯，補足のページを指す命令，色，図の部品）
% scripts/export_figdata.py が書いた．手で直さない．
\expandafter\gdef\csname NUM@advAFEWAVTfive\endcsname{+1.3}
\expandafter\gdef\csname NUM@advAFEWAVTtwofive\endcsname{+0.9}
\expandafter\gdef\csname NUM@advAFEWAfive\endcsname{+21.7}
\expandafter\gdef\csname NUM@advAFEWAtwofive\endcsname{+24.2}
\expandafter\gdef\csname NUM@advAFEWTfive\endcsname{+27.3}
\expandafter\gdef\csname NUM@advAFEWTtwofive\endcsname{+31.1}
\expandafter\gdef\csname NUM@advAFEWVfive\endcsname{+1.7}
\expandafter\gdef\csname NUM@advAFEWVtwofive\endcsname{+3.1}
\expandafter\gdef\csname NUM@advIEMOCAPAVTfive\endcsname{+1.2}
\expandafter\gdef\csname NUM@advIEMOCAPAVTtwofive\endcsname{+0.1}
\expandafter\gdef\csname NUM@advIEMOCAPAfive\endcsname{+7.2}
\expandafter\gdef\csname NUM@advIEMOCAPAtwofive\endcsname{+7.0}
\expandafter\gdef\csname NUM@advIEMOCAPTfive\endcsname{+25.0}
\expandafter\gdef\csname NUM@advIEMOCAPTtwofive\endcsname{+23.6}
\expandafter\gdef\csname NUM@advIEMOCAPVfive\endcsname{+17.6}
\expandafter\gdef\csname NUM@advIEMOCAPVtwofive\endcsname{+19.2}
\expandafter\gdef\csname NUM@advMELDAVTfive\endcsname{+2.5}
\expandafter\gdef\csname NUM@advMELDAVTtwofive\endcsname{+1.6}
\expandafter\gdef\csname NUM@advMELDAfive\endcsname{+8.4}
\expandafter\gdef\csname NUM@advMELDAtwofive\endcsname{+9.5}
\expandafter\gdef\csname NUM@advMELDTfive\endcsname{+0.3}
\expandafter\gdef\csname NUM@advMELDTtwofive\endcsname{+0.8}
\expandafter\gdef\csname NUM@advMELDVfive\endcsname{+12.6}
\expandafter\gdef\csname NUM@advMELDVtwofive\endcsname{+14.5}
\expandafter\gdef\csname NUM@advMOSEIAVTfive\endcsname{+0.1}
\expandafter\gdef\csname NUM@advMOSEIAVTtwofive\endcsname{+0.9}
\expandafter\gdef\csname NUM@advMOSEIAfive\endcsname{+6.8}
\expandafter\gdef\csname NUM@advMOSEIAtwofive\endcsname{+7.5}
\expandafter\gdef\csname NUM@advMOSEITfive\endcsname{+8.0}
\expandafter\gdef\csname NUM@advMOSEITtwofive\endcsname{+8.8}
\expandafter\gdef\csname NUM@advMOSEIVfive\endcsname{+5.0}
\expandafter\gdef\csname NUM@advMOSEIVtwofive\endcsname{+7.2}
\expandafter\gdef\csname NUM@agreeAFEWAVTfive\endcsname{89}
\expandafter\gdef\csname NUM@agreeAFEWAVTtwofive\endcsname{92}
\expandafter\gdef\csname NUM@agreeAFEWAfive\endcsname{66}
\expandafter\gdef\csname NUM@agreeAFEWAtwofive\endcsname{63}
\expandafter\gdef\csname NUM@agreeAFEWTfive\endcsname{55}
\expandafter\gdef\csname NUM@agreeAFEWTtwofive\endcsname{54}
\expandafter\gdef\csname NUM@agreeAFEWVfive\endcsname{65}
\expandafter\gdef\csname NUM@agreeAFEWVtwofive\endcsname{73}
\expandafter\gdef\csname NUM@agreeIEMOCAPAVTfive\endcsname{87}
\expandafter\gdef\csname NUM@agreeIEMOCAPAVTtwofive\endcsname{89}
\expandafter\gdef\csname NUM@agreeIEMOCAPAfive\endcsname{71}
\expandafter\gdef\csname NUM@agreeIEMOCAPAtwofive\endcsname{74}
\expandafter\gdef\csname NUM@agreeIEMOCAPTfive\endcsname{55}
\expandafter\gdef\csname NUM@agreeIEMOCAPTtwofive\endcsname{58}
\expandafter\gdef\csname NUM@agreeIEMOCAPVfive\endcsname{63}
\expandafter\gdef\csname NUM@agreeIEMOCAPVtwofive\endcsname{66}
\expandafter\gdef\csname NUM@agreeMELDAVTfive\endcsname{82}
\expandafter\gdef\csname NUM@agreeMELDAVTtwofive\endcsname{83}
\expandafter\gdef\csname NUM@agreeMELDAfive\endcsname{58}
\expandafter\gdef\csname NUM@agreeMELDAtwofive\endcsname{61}
\expandafter\gdef\csname NUM@agreeMELDTfive\endcsname{52}
\expandafter\gdef\csname NUM@agreeMELDTtwofive\endcsname{56}
\expandafter\gdef\csname NUM@agreeMELDVfive\endcsname{53}
\expandafter\gdef\csname NUM@agreeMELDVtwofive\endcsname{51}
\expandafter\gdef\csname NUM@agreeMOSEIAVTfive\endcsname{80}
\expandafter\gdef\csname NUM@agreeMOSEIAVTtwofive\endcsname{82}
\expandafter\gdef\csname NUM@agreeMOSEIAfive\endcsname{58}
\expandafter\gdef\csname NUM@agreeMOSEIAtwofive\endcsname{56}
\expandafter\gdef\csname NUM@agreeMOSEITfive\endcsname{45}
\expandafter\gdef\csname NUM@agreeMOSEITtwofive\endcsname{49}
\expandafter\gdef\csname NUM@agreeMOSEIVfive\endcsname{49}
\expandafter\gdef\csname NUM@agreeMOSEIVtwofive\endcsname{48}
\expandafter\gdef\csname NUM@cagreeAFEWAVTfive\endcsname{91}
\expandafter\gdef\csname NUM@cagreeAFEWAVTtwofive\endcsname{92}
\expandafter\gdef\csname NUM@cagreeAFEWAfive\endcsname{86}
\expandafter\gdef\csname NUM@cagreeAFEWAtwofive\endcsname{88}
\expandafter\gdef\csname NUM@cagreeAFEWTfive\endcsname{86}
\expandafter\gdef\csname NUM@cagreeAFEWTtwofive\endcsname{89}
\expandafter\gdef\csname NUM@cagreeAFEWVfive\endcsname{76}
\expandafter\gdef\csname NUM@cagreeAFEWVtwofive\endcsname{82}
\expandafter\gdef\csname NUM@cagreeIEMOCAPAVTfive\endcsname{86}
\expandafter\gdef\csname NUM@cagreeIEMOCAPAVTtwofive\endcsname{89}
\expandafter\gdef\csname NUM@cagreeIEMOCAPAfive\endcsname{81}
\expandafter\gdef\csname NUM@cagreeIEMOCAPAtwofive\endcsname{84}
\expandafter\gdef\csname NUM@cagreeIEMOCAPTfive\endcsname{80}
\expandafter\gdef\csname NUM@cagreeIEMOCAPTtwofive\endcsname{83}
\expandafter\gdef\csname NUM@cagreeIEMOCAPVfive\endcsname{77}
\expandafter\gdef\csname NUM@cagreeIEMOCAPVtwofive\endcsname{81}
\expandafter\gdef\csname NUM@cagreeMELDAVTfive\endcsname{55}
\expandafter\gdef\csname NUM@cagreeMELDAVTtwofive\endcsname{58}
\expandafter\gdef\csname NUM@cagreeMELDAfive\endcsname{51}
\expandafter\gdef\csname NUM@cagreeMELDAtwofive\endcsname{53}
\expandafter\gdef\csname NUM@cagreeMELDTfive\endcsname{45}
\expandafter\gdef\csname NUM@cagreeMELDTtwofive\endcsname{47}
\expandafter\gdef\csname NUM@cagreeMELDVfive\endcsname{52}
\expandafter\gdef\csname NUM@cagreeMELDVtwofive\endcsname{54}
\expandafter\gdef\csname NUM@cagreeMOSEIAVTfive\endcsname{48}
\expandafter\gdef\csname NUM@cagreeMOSEIAVTtwofive\endcsname{53}
\expandafter\gdef\csname NUM@cagreeMOSEIAfive\endcsname{45}
\expandafter\gdef\csname NUM@cagreeMOSEIAtwofive\endcsname{49}
\expandafter\gdef\csname NUM@cagreeMOSEITfive\endcsname{44}
\expandafter\gdef\csname NUM@cagreeMOSEITtwofive\endcsname{48}
\expandafter\gdef\csname NUM@cagreeMOSEIVfive\endcsname{43}
\expandafter\gdef\csname NUM@cagreeMOSEIVtwofive\endcsname{46}
\expandafter\gdef\csname NUM@cdisAFEWAVTfive\endcsname{62}
\expandafter\gdef\csname NUM@cdisAFEWAVTtwofive\endcsname{66}
\expandafter\gdef\csname NUM@cdisAFEWAfive\endcsname{59}
\expandafter\gdef\csname NUM@cdisAFEWAtwofive\endcsname{65}
\expandafter\gdef\csname NUM@cdisAFEWTfive\endcsname{58}
\expandafter\gdef\csname NUM@cdisAFEWTtwofive\endcsname{61}
\expandafter\gdef\csname NUM@cdisAFEWVfive\endcsname{67}
\expandafter\gdef\csname NUM@cdisAFEWVtwofive\endcsname{71}
\expandafter\gdef\csname NUM@cdisIEMOCAPAVTfive\endcsname{61}
\expandafter\gdef\csname NUM@cdisIEMOCAPAVTtwofive\endcsname{65}
\expandafter\gdef\csname NUM@cdisIEMOCAPAfive\endcsname{62}
\expandafter\gdef\csname NUM@cdisIEMOCAPAtwofive\endcsname{65}
\expandafter\gdef\csname NUM@cdisIEMOCAPTfive\endcsname{59}
\expandafter\gdef\csname NUM@cdisIEMOCAPTtwofive\endcsname{62}
\expandafter\gdef\csname NUM@cdisIEMOCAPVfive\endcsname{61}
\expandafter\gdef\csname NUM@cdisIEMOCAPVtwofive\endcsname{64}
\expandafter\gdef\csname NUM@cdisMELDAVTfive\endcsname{42}
\expandafter\gdef\csname NUM@cdisMELDAVTtwofive\endcsname{43}
\expandafter\gdef\csname NUM@cdisMELDAfive\endcsname{41}
\expandafter\gdef\csname NUM@cdisMELDAtwofive\endcsname{43}
\expandafter\gdef\csname NUM@cdisMELDTfive\endcsname{44}
\expandafter\gdef\csname NUM@cdisMELDTtwofive\endcsname{46}
\expandafter\gdef\csname NUM@cdisMELDVfive\endcsname{38}
\expandafter\gdef\csname NUM@cdisMELDVtwofive\endcsname{38}
\expandafter\gdef\csname NUM@cdisMOSEIAVTfive\endcsname{38}
\expandafter\gdef\csname NUM@cdisMOSEIAVTtwofive\endcsname{39}
\expandafter\gdef\csname NUM@cdisMOSEIAfive\endcsname{38}
\expandafter\gdef\csname NUM@cdisMOSEIAtwofive\endcsname{39}
\expandafter\gdef\csname NUM@cdisMOSEITfive\endcsname{37}
\expandafter\gdef\csname NUM@cdisMOSEITtwofive\endcsname{38}
\expandafter\gdef\csname NUM@cdisMOSEIVfive\endcsname{38}
\expandafter\gdef\csname NUM@cdisMOSEIVtwofive\endcsname{39}
\expandafter\gdef\csname NUM@charmaAFEWAVTfive\endcsname{0.1}
\expandafter\gdef\csname NUM@charmaAFEWAVTtwofive\endcsname{0.1}
\expandafter\gdef\csname NUM@charmaAFEWAfive\endcsname{2.3}
\expandafter\gdef\csname NUM@charmaAFEWAtwofive\endcsname{2.7}
\expandafter\gdef\csname NUM@charmaAFEWTfive\endcsname{3.9}
\expandafter\gdef\csname NUM@charmaAFEWTtwofive\endcsname{2.9}
\expandafter\gdef\csname NUM@charmaAFEWVfive\endcsname{14.4}
\expandafter\gdef\csname NUM@charmaAFEWVtwofive\endcsname{10.5}
\expandafter\gdef\csname NUM@charmaIEMOCAPAVTfive\endcsname{0.7}
\expandafter\gdef\csname NUM@charmaIEMOCAPAVTtwofive\endcsname{0.8}
\expandafter\gdef\csname NUM@charmaIEMOCAPAfive\endcsname{5.6}
\expandafter\gdef\csname NUM@charmaIEMOCAPAtwofive\endcsname{5.1}
\expandafter\gdef\csname NUM@charmaIEMOCAPTfive\endcsname{5.9}
\expandafter\gdef\csname NUM@charmaIEMOCAPTtwofive\endcsname{5.6}
\expandafter\gdef\csname NUM@charmaIEMOCAPVfive\endcsname{7.5}
\expandafter\gdef\csname NUM@charmaIEMOCAPVtwofive\endcsname{5.8}
\expandafter\gdef\csname NUM@charmaMELDAVTfive\endcsname{1.4}
\expandafter\gdef\csname NUM@charmaMELDAVTtwofive\endcsname{1.1}
\expandafter\gdef\csname NUM@charmaMELDAfive\endcsname{11.2}
\expandafter\gdef\csname NUM@charmaMELDAtwofive\endcsname{9.9}
\expandafter\gdef\csname NUM@charmaMELDTfive\endcsname{20.4}
\expandafter\gdef\csname NUM@charmaMELDTtwofive\endcsname{21.2}
\expandafter\gdef\csname NUM@charmaMELDVfive\endcsname{10.1}
\expandafter\gdef\csname NUM@charmaMELDVtwofive\endcsname{10.4}
\expandafter\gdef\csname NUM@charmaMOSEIAVTfive\endcsname{1.5}
\expandafter\gdef\csname NUM@charmaMOSEIAVTtwofive\endcsname{1.0}
\expandafter\gdef\csname NUM@charmaMOSEIAfive\endcsname{10.2}
\expandafter\gdef\csname NUM@charmaMOSEIAtwofive\endcsname{10.1}
\expandafter\gdef\csname NUM@charmaMOSEITfive\endcsname{14.3}
\expandafter\gdef\csname NUM@charmaMOSEITtwofive\endcsname{12.5}
\expandafter\gdef\csname NUM@charmaMOSEIVfive\endcsname{15.0}
\expandafter\gdef\csname NUM@charmaMOSEIVtwofive\endcsname{14.5}
\expandafter\gdef\csname NUM@charmdAFEWAVTfive\endcsname{6.1}
\expandafter\gdef\csname NUM@charmdAFEWAVTtwofive\endcsname{5.6}
\expandafter\gdef\csname NUM@charmdAFEWAfive\endcsname{4.5}
\expandafter\gdef\csname NUM@charmdAFEWAtwofive\endcsname{5.0}
\expandafter\gdef\csname NUM@charmdAFEWTfive\endcsname{7.0}
\expandafter\gdef\csname NUM@charmdAFEWTtwofive\endcsname{6.3}
\expandafter\gdef\csname NUM@charmdAFEWVfive\endcsname{15.4}
\expandafter\gdef\csname NUM@charmdAFEWVtwofive\endcsname{11.2}
\expandafter\gdef\csname NUM@charmdIEMOCAPAVTfive\endcsname{8.6}
\expandafter\gdef\csname NUM@charmdIEMOCAPAVTtwofive\endcsname{9.2}
\expandafter\gdef\csname NUM@charmdIEMOCAPAfive\endcsname{10.3}
\expandafter\gdef\csname NUM@charmdIEMOCAPAtwofive\endcsname{9.6}
\expandafter\gdef\csname NUM@charmdIEMOCAPTfive\endcsname{10.0}
\expandafter\gdef\csname NUM@charmdIEMOCAPTtwofive\endcsname{9.3}
\expandafter\gdef\csname NUM@charmdIEMOCAPVfive\endcsname{9.6}
\expandafter\gdef\csname NUM@charmdIEMOCAPVtwofive\endcsname{7.4}
\expandafter\gdef\csname NUM@charmdMELDAVTfive\endcsname{9.8}
\expandafter\gdef\csname NUM@charmdMELDAVTtwofive\endcsname{9.0}
\expandafter\gdef\csname NUM@charmdMELDAfive\endcsname{12.4}
\expandafter\gdef\csname NUM@charmdMELDAtwofive\endcsname{10.5}
\expandafter\gdef\csname NUM@charmdMELDTfive\endcsname{17.8}
\expandafter\gdef\csname NUM@charmdMELDTtwofive\endcsname{17.3}
\expandafter\gdef\csname NUM@charmdMELDVfive\endcsname{10.2}
\expandafter\gdef\csname NUM@charmdMELDVtwofive\endcsname{10.6}
\expandafter\gdef\csname NUM@charmdMOSEIAVTfive\endcsname{9.8}
\expandafter\gdef\csname NUM@charmdMOSEIAVTtwofive\endcsname{8.4}
\expandafter\gdef\csname NUM@charmdMOSEIAfive\endcsname{10.4}
\expandafter\gdef\csname NUM@charmdMOSEIAtwofive\endcsname{10.9}
\expandafter\gdef\csname NUM@charmdMOSEITfive\endcsname{14.2}
\expandafter\gdef\csname NUM@charmdMOSEITtwofive\endcsname{13.0}
\expandafter\gdef\csname NUM@charmdMOSEIVfive\endcsname{13.3}
\expandafter\gdef\csname NUM@charmdMOSEIVtwofive\endcsname{13.5}
\expandafter\gdef\csname NUM@cshareAFEWAVTfive\endcsname{38}
\expandafter\gdef\csname NUM@cshareAFEWAVTtwofive\endcsname{44}
\expandafter\gdef\csname NUM@cshareAFEWAfive\endcsname{52}
\expandafter\gdef\csname NUM@cshareAFEWAtwofive\endcsname{55}
\expandafter\gdef\csname NUM@cshareAFEWTfive\endcsname{51}
\expandafter\gdef\csname NUM@cshareAFEWTtwofive\endcsname{60}
\expandafter\gdef\csname NUM@cshareAFEWVfive\endcsname{60}
\expandafter\gdef\csname NUM@cshareAFEWVtwofive\endcsname{65}
\expandafter\gdef\csname NUM@cshareIEMOCAPAVTfive\endcsname{39}
\expandafter\gdef\csname NUM@cshareIEMOCAPAVTtwofive\endcsname{41}
\expandafter\gdef\csname NUM@cshareIEMOCAPAfive\endcsname{47}
\expandafter\gdef\csname NUM@cshareIEMOCAPAtwofive\endcsname{49}
\expandafter\gdef\csname NUM@cshareIEMOCAPTfive\endcsname{58}
\expandafter\gdef\csname NUM@cshareIEMOCAPTtwofive\endcsname{61}
\expandafter\gdef\csname NUM@cshareIEMOCAPVfive\endcsname{61}
\expandafter\gdef\csname NUM@cshareIEMOCAPVtwofive\endcsname{61}
\expandafter\gdef\csname NUM@cshareMELDAVTfive\endcsname{24}
\expandafter\gdef\csname NUM@cshareMELDAVTtwofive\endcsname{24}
\expandafter\gdef\csname NUM@cshareMELDAfive\endcsname{37}
\expandafter\gdef\csname NUM@cshareMELDAtwofive\endcsname{36}
\expandafter\gdef\csname NUM@cshareMELDTfive\endcsname{43}
\expandafter\gdef\csname NUM@cshareMELDTtwofive\endcsname{42}
\expandafter\gdef\csname NUM@cshareMELDVfive\endcsname{51}
\expandafter\gdef\csname NUM@cshareMELDVtwofive\endcsname{51}
\expandafter\gdef\csname NUM@cshareMOSEIAVTfive\endcsname{23}
\expandafter\gdef\csname NUM@cshareMOSEIAVTtwofive\endcsname{24}
\expandafter\gdef\csname NUM@cshareMOSEIAfive\endcsname{37}
\expandafter\gdef\csname NUM@cshareMOSEIAtwofive\endcsname{37}
\expandafter\gdef\csname NUM@cshareMOSEITfive\endcsname{45}
\expandafter\gdef\csname NUM@cshareMOSEITtwofive\endcsname{45}
\expandafter\gdef\csname NUM@cshareMOSEIVfive\endcsname{46}
\expandafter\gdef\csname NUM@cshareMOSEIVtwofive\endcsname{48}
\expandafter\gdef\csname NUM@daitwoAFEWcoralfive\endcsname{+1.8}
\expandafter\gdef\csname NUM@daitwoAFEWcoraltwofive\endcsname{+0.6}
\expandafter\gdef\csname NUM@daitwoAFEWcoralzero\endcsname{+5.3}
\expandafter\gdef\csname NUM@daitwoAFEWdannfive\endcsname{+1.1}
\expandafter\gdef\csname NUM@daitwoAFEWdanntwofive\endcsname{+0.9}
\expandafter\gdef\csname NUM@daitwoAFEWdannzero\endcsname{+1.6}
\expandafter\gdef\csname NUM@daitwoAFEWentminfive\endcsname{\textminus{}0.5}
\expandafter\gdef\csname NUM@daitwoAFEWentmintwofive\endcsname{\textminus{}0.0}
\expandafter\gdef\csname NUM@daitwoAFEWentminzero\endcsname{\textminus{}4.5}
\expandafter\gdef\csname NUM@daitwoAFEWimfive\endcsname{+0.3}
\expandafter\gdef\csname NUM@daitwoAFEWimtwofive\endcsname{+0.6}
\expandafter\gdef\csname NUM@daitwoAFEWimzero\endcsname{\textminus{}2.5}
\expandafter\gdef\csname NUM@daitwoAFEWlpfive\endcsname{+11.0}
\expandafter\gdef\csname NUM@daitwoAFEWlptwofive\endcsname{+9.8}
\expandafter\gdef\csname NUM@daitwoAFEWmmefive\endcsname{+0.4}
\expandafter\gdef\csname NUM@daitwoAFEWmmetwofive\endcsname{+0.2}
\expandafter\gdef\csname NUM@daitwoAFEWmmezero\endcsname{\textminus{}0.3}
\expandafter\gdef\csname NUM@daitwoAFEWownfive\endcsname{+0.8}
\expandafter\gdef\csname NUM@daitwoAFEWowntwofive\endcsname{+0.9}
\expandafter\gdef\csname NUM@daitwoAFEWownzero\endcsname{+0.2}
\expandafter\gdef\csname NUM@daitwoAFEWselffive\endcsname{\textminus{}1.4}
\expandafter\gdef\csname NUM@daitwoAFEWselftwofive\endcsname{+0.2}
\expandafter\gdef\csname NUM@daitwoAFEWselfzero\endcsname{\textminus{}1.1}
\expandafter\gdef\csname NUM@daitwoAFEWstfive\endcsname{+6.3}
\expandafter\gdef\csname NUM@daitwoAFEWsttwofive\endcsname{+7.0}
\expandafter\gdef\csname NUM@daitwoAFEWstzero\endcsname{+4.6}
\expandafter\gdef\csname NUM@daitwoAFEWtthreeafive\endcsname{\textminus{}1.2}
\expandafter\gdef\csname NUM@daitwoAFEWtthreeatwofive\endcsname{+1.9}
\expandafter\gdef\csname NUM@daitwoAFEWtthreeazero\endcsname{\textminus{}5.1}
\expandafter\gdef\csname NUM@daitwoIEMOCAPcoralfive\endcsname{+2.0}
\expandafter\gdef\csname NUM@daitwoIEMOCAPcoraltwofive\endcsname{+0.3}
\expandafter\gdef\csname NUM@daitwoIEMOCAPcoralzero\endcsname{+5.1}
\expandafter\gdef\csname NUM@daitwoIEMOCAPdannfive\endcsname{+1.4}
\expandafter\gdef\csname NUM@daitwoIEMOCAPdanntwofive\endcsname{+1.2}
\expandafter\gdef\csname NUM@daitwoIEMOCAPdannzero\endcsname{+3.5}
\expandafter\gdef\csname NUM@daitwoIEMOCAPentminfive\endcsname{+10.6}
\expandafter\gdef\csname NUM@daitwoIEMOCAPentmintwofive\endcsname{+26.8}
\expandafter\gdef\csname NUM@daitwoIEMOCAPentminzero\endcsname{+16.2}
\expandafter\gdef\csname NUM@daitwoIEMOCAPimfive\endcsname{+2.1}
\expandafter\gdef\csname NUM@daitwoIEMOCAPimtwofive\endcsname{+0.2}
\expandafter\gdef\csname NUM@daitwoIEMOCAPimzero\endcsname{+12.0}
\expandafter\gdef\csname NUM@daitwoIEMOCAPlpfive\endcsname{+11.7}
\expandafter\gdef\csname NUM@daitwoIEMOCAPlptwofive\endcsname{+4.7}
\expandafter\gdef\csname NUM@daitwoIEMOCAPmmefive\endcsname{+10.0}
\expandafter\gdef\csname NUM@daitwoIEMOCAPmmetwofive\endcsname{+0.2}
\expandafter\gdef\csname NUM@daitwoIEMOCAPmmezero\endcsname{+28.2}
\expandafter\gdef\csname NUM@daitwoIEMOCAPownfive\endcsname{+1.1}
\expandafter\gdef\csname NUM@daitwoIEMOCAPowntwofive\endcsname{+0.3}
\expandafter\gdef\csname NUM@daitwoIEMOCAPownzero\endcsname{+2.0}
\expandafter\gdef\csname NUM@daitwoIEMOCAPselffive\endcsname{+1.4}
\expandafter\gdef\csname NUM@daitwoIEMOCAPselftwofive\endcsname{\textminus{}0.4}
\expandafter\gdef\csname NUM@daitwoIEMOCAPselfzero\endcsname{\textminus{}0.3}
\expandafter\gdef\csname NUM@daitwoIEMOCAPstfive\endcsname{+2.5}
\expandafter\gdef\csname NUM@daitwoIEMOCAPsttwofive\endcsname{\textminus{}0.7}
\expandafter\gdef\csname NUM@daitwoIEMOCAPstzero\endcsname{+2.6}
\expandafter\gdef\csname NUM@daitwoIEMOCAPtthreeafive\endcsname{+7.2}
\expandafter\gdef\csname NUM@daitwoIEMOCAPtthreeatwofive\endcsname{+7.6}
\expandafter\gdef\csname NUM@daitwoIEMOCAPtthreeazero\endcsname{+4.6}
\expandafter\gdef\csname NUM@daitwoMELDcoralfive\endcsname{+3.8}
\expandafter\gdef\csname NUM@daitwoMELDcoraltwofive\endcsname{+2.4}
\expandafter\gdef\csname NUM@daitwoMELDcoralzero\endcsname{+5.4}
\expandafter\gdef\csname NUM@daitwoMELDdannfive\endcsname{+3.2}
\expandafter\gdef\csname NUM@daitwoMELDdanntwofive\endcsname{+3.8}
\expandafter\gdef\csname NUM@daitwoMELDdannzero\endcsname{+4.3}
\expandafter\gdef\csname NUM@daitwoMELDentminfive\endcsname{+13.3}
\expandafter\gdef\csname NUM@daitwoMELDentmintwofive\endcsname{+10.7}
\expandafter\gdef\csname NUM@daitwoMELDentminzero\endcsname{+9.0}
\expandafter\gdef\csname NUM@daitwoMELDimfive\endcsname{+10.5}
\expandafter\gdef\csname NUM@daitwoMELDimtwofive\endcsname{+4.5}
\expandafter\gdef\csname NUM@daitwoMELDimzero\endcsname{+12.1}
\expandafter\gdef\csname NUM@daitwoMELDlpfive\endcsname{+16.5}
\expandafter\gdef\csname NUM@daitwoMELDlptwofive\endcsname{+10.9}
\expandafter\gdef\csname NUM@daitwoMELDmmefive\endcsname{+19.0}
\expandafter\gdef\csname NUM@daitwoMELDmmetwofive\endcsname{+5.5}
\expandafter\gdef\csname NUM@daitwoMELDmmezero\endcsname{+24.7}
\expandafter\gdef\csname NUM@daitwoMELDownfive\endcsname{+2.2}
\expandafter\gdef\csname NUM@daitwoMELDowntwofive\endcsname{+1.8}
\expandafter\gdef\csname NUM@daitwoMELDownzero\endcsname{+3.0}
\expandafter\gdef\csname NUM@daitwoMELDselffive\endcsname{+4.1}
\expandafter\gdef\csname NUM@daitwoMELDselftwofive\endcsname{+3.4}
\expandafter\gdef\csname NUM@daitwoMELDselfzero\endcsname{+3.7}
\expandafter\gdef\csname NUM@daitwoMELDstfive\endcsname{\textminus{}1.7}
\expandafter\gdef\csname NUM@daitwoMELDsttwofive\endcsname{\textminus{}0.9}
\expandafter\gdef\csname NUM@daitwoMELDstzero\endcsname{\textminus{}0.8}
\expandafter\gdef\csname NUM@daitwoMELDtthreeafive\endcsname{+8.4}
\expandafter\gdef\csname NUM@daitwoMELDtthreeatwofive\endcsname{+8.5}
\expandafter\gdef\csname NUM@daitwoMELDtthreeazero\endcsname{+6.9}
\expandafter\gdef\csname NUM@daitwoMOSEIcoralfive\endcsname{\textminus{}1.0}
\expandafter\gdef\csname NUM@daitwoMOSEIcoraltwofive\endcsname{+0.4}
\expandafter\gdef\csname NUM@daitwoMOSEIcoralzero\endcsname{\textminus{}0.8}
\expandafter\gdef\csname NUM@daitwoMOSEIdannfive\endcsname{+0.5}
\expandafter\gdef\csname NUM@daitwoMOSEIdanntwofive\endcsname{+1.6}
\expandafter\gdef\csname NUM@daitwoMOSEIdannzero\endcsname{+0.6}
\expandafter\gdef\csname NUM@daitwoMOSEIentminfive\endcsname{+12.8}
\expandafter\gdef\csname NUM@daitwoMOSEIentmintwofive\endcsname{+5.0}
\expandafter\gdef\csname NUM@daitwoMOSEIentminzero\endcsname{+4.1}
\expandafter\gdef\csname NUM@daitwoMOSEIimfive\endcsname{+1.7}
\expandafter\gdef\csname NUM@daitwoMOSEIimtwofive\endcsname{+3.8}
\expandafter\gdef\csname NUM@daitwoMOSEIimzero\endcsname{+3.0}
\expandafter\gdef\csname NUM@daitwoMOSEIlpfive\endcsname{+6.6}
\expandafter\gdef\csname NUM@daitwoMOSEIlptwofive\endcsname{+5.2}
\expandafter\gdef\csname NUM@daitwoMOSEImmefive\endcsname{+6.3}
\expandafter\gdef\csname NUM@daitwoMOSEImmetwofive\endcsname{+6.6}
\expandafter\gdef\csname NUM@daitwoMOSEImmezero\endcsname{+12.5}
\expandafter\gdef\csname NUM@daitwoMOSEIownfive\endcsname{\textminus{}0.5}
\expandafter\gdef\csname NUM@daitwoMOSEIowntwofive\endcsname{+0.4}
\expandafter\gdef\csname NUM@daitwoMOSEIownzero\endcsname{\textminus{}0.7}
\expandafter\gdef\csname NUM@daitwoMOSEIselffive\endcsname{+0.8}
\expandafter\gdef\csname NUM@daitwoMOSEIselftwofive\endcsname{+1.3}
\expandafter\gdef\csname NUM@daitwoMOSEIselfzero\endcsname{\textminus{}0.0}
\expandafter\gdef\csname NUM@daitwoMOSEIstfive\endcsname{+0.0}
\expandafter\gdef\csname NUM@daitwoMOSEIsttwofive\endcsname{+0.7}
\expandafter\gdef\csname NUM@daitwoMOSEIstzero\endcsname{\textminus{}0.5}
\expandafter\gdef\csname NUM@daitwoMOSEItthreeafive\endcsname{+3.7}
\expandafter\gdef\csname NUM@daitwoMOSEItthreeatwofive\endcsname{+3.8}
\expandafter\gdef\csname NUM@daitwoMOSEItthreeazero\endcsname{+3.4}
\expandafter\gdef\csname NUM@daitwoabovecoral\endcsname{5}
\expandafter\gdef\csname NUM@daitwoabovedann\endcsname{4}
\expandafter\gdef\csname NUM@daitwoaboveentmin\endcsname{9}
\expandafter\gdef\csname NUM@daitwoaboveim\endcsname{8}
\expandafter\gdef\csname NUM@daitwoabovelp\endcsname{8}
\expandafter\gdef\csname NUM@daitwoabovemme\endcsname{8}
\expandafter\gdef\csname NUM@daitwoaboveown\endcsname{4}
\expandafter\gdef\csname NUM@daitwoaboveself\endcsname{4}
\expandafter\gdef\csname NUM@daitwoabovest\endcsname{3}
\expandafter\gdef\csname NUM@daitwoabovetthreea\endcsname{8}
\expandafter\gdef\csname NUM@daitwobelowcoral\endcsname{1}
\expandafter\gdef\csname NUM@daitwobelowdann\endcsname{0}
\expandafter\gdef\csname NUM@daitwobelowentmin\endcsname{1}
\expandafter\gdef\csname NUM@daitwobelowim\endcsname{1}
\expandafter\gdef\csname NUM@daitwobelowlp\endcsname{0}
\expandafter\gdef\csname NUM@daitwobelowmme\endcsname{0}
\expandafter\gdef\csname NUM@daitwobelowown\endcsname{0}
\expandafter\gdef\csname NUM@daitwobelowself\endcsname{0}
\expandafter\gdef\csname NUM@daitwobelowst\endcsname{0}
\expandafter\gdef\csname NUM@daitwobelowtthreea\endcsname{1}
\expandafter\gdef\csname NUM@daitwocellAFEWcoralfive\endcsname{+1.8}
\expandafter\gdef\csname NUM@daitwocellAFEWcoraltwofive\endcsname{+0.6}
\expandafter\gdef\csname NUM@daitwocellAFEWcoralzero\endcsname{\textbf{+5.3}}
\expandafter\gdef\csname NUM@daitwocellAFEWdannfive\endcsname{+1.1}
\expandafter\gdef\csname NUM@daitwocellAFEWdanntwofive\endcsname{+0.9}
\expandafter\gdef\csname NUM@daitwocellAFEWdannzero\endcsname{+1.6}
\expandafter\gdef\csname NUM@daitwocellAFEWentminfive\endcsname{\textminus{}0.5}
\expandafter\gdef\csname NUM@daitwocellAFEWentmintwofive\endcsname{\textminus{}0.0}
\expandafter\gdef\csname NUM@daitwocellAFEWentminzero\endcsname{\textcolor{mBad}{\textminus{}4.5}}
\expandafter\gdef\csname NUM@daitwocellAFEWimfive\endcsname{+0.3}
\expandafter\gdef\csname NUM@daitwocellAFEWimtwofive\endcsname{+0.6}
\expandafter\gdef\csname NUM@daitwocellAFEWimzero\endcsname{\textcolor{mBad}{\textminus{}2.5}}
\expandafter\gdef\csname NUM@daitwocellAFEWlpfive\endcsname{\textbf{+11.0}}
\expandafter\gdef\csname NUM@daitwocellAFEWlptwofive\endcsname{\textbf{+9.8}}
\expandafter\gdef\csname NUM@daitwocellAFEWmmefive\endcsname{+0.4}
\expandafter\gdef\csname NUM@daitwocellAFEWmmetwofive\endcsname{+0.2}
\expandafter\gdef\csname NUM@daitwocellAFEWmmezero\endcsname{\textminus{}0.3}
\expandafter\gdef\csname NUM@daitwocellAFEWownfive\endcsname{+0.8}
\expandafter\gdef\csname NUM@daitwocellAFEWowntwofive\endcsname{\textbf{+0.9}}
\expandafter\gdef\csname NUM@daitwocellAFEWownzero\endcsname{+0.2}
\expandafter\gdef\csname NUM@daitwocellAFEWselffive\endcsname{\textminus{}1.4}
\expandafter\gdef\csname NUM@daitwocellAFEWselftwofive\endcsname{+0.2}
\expandafter\gdef\csname NUM@daitwocellAFEWselfzero\endcsname{\textminus{}1.1}
\expandafter\gdef\csname NUM@daitwocellAFEWstfive\endcsname{\textbf{+6.3}}
\expandafter\gdef\csname NUM@daitwocellAFEWsttwofive\endcsname{\textbf{+7.0}}
\expandafter\gdef\csname NUM@daitwocellAFEWstzero\endcsname{\textbf{+4.6}}
\expandafter\gdef\csname NUM@daitwocellAFEWtthreeafive\endcsname{\textminus{}1.2}
\expandafter\gdef\csname NUM@daitwocellAFEWtthreeatwofive\endcsname{+1.9}
\expandafter\gdef\csname NUM@daitwocellAFEWtthreeazero\endcsname{\textcolor{mBad}{\textminus{}5.1}}
\expandafter\gdef\csname NUM@daitwocellIEMOCAPcoralfive\endcsname{+2.0}
\expandafter\gdef\csname NUM@daitwocellIEMOCAPcoraltwofive\endcsname{+0.3}
\expandafter\gdef\csname NUM@daitwocellIEMOCAPcoralzero\endcsname{\textbf{+5.1}}
\expandafter\gdef\csname NUM@daitwocellIEMOCAPdannfive\endcsname{+1.4}
\expandafter\gdef\csname NUM@daitwocellIEMOCAPdanntwofive\endcsname{+1.2}
\expandafter\gdef\csname NUM@daitwocellIEMOCAPdannzero\endcsname{+3.5}
\expandafter\gdef\csname NUM@daitwocellIEMOCAPentminfive\endcsname{\textbf{+10.6}}
\expandafter\gdef\csname NUM@daitwocellIEMOCAPentmintwofive\endcsname{\textbf{+26.8}}
\expandafter\gdef\csname NUM@daitwocellIEMOCAPentminzero\endcsname{\textbf{+16.2}}
\expandafter\gdef\csname NUM@daitwocellIEMOCAPimfive\endcsname{\textbf{+2.1}}
\expandafter\gdef\csname NUM@daitwocellIEMOCAPimtwofive\endcsname{+0.2}
\expandafter\gdef\csname NUM@daitwocellIEMOCAPimzero\endcsname{\textbf{+12.0}}
\expandafter\gdef\csname NUM@daitwocellIEMOCAPlpfive\endcsname{\textbf{+11.7}}
\expandafter\gdef\csname NUM@daitwocellIEMOCAPlptwofive\endcsname{\textbf{+4.7}}
\expandafter\gdef\csname NUM@daitwocellIEMOCAPmmefive\endcsname{\textbf{+10.0}}
\expandafter\gdef\csname NUM@daitwocellIEMOCAPmmetwofive\endcsname{+0.2}
\expandafter\gdef\csname NUM@daitwocellIEMOCAPmmezero\endcsname{\textbf{+28.2}}
\expandafter\gdef\csname NUM@daitwocellIEMOCAPownfive\endcsname{+1.1}
\expandafter\gdef\csname NUM@daitwocellIEMOCAPowntwofive\endcsname{+0.3}
\expandafter\gdef\csname NUM@daitwocellIEMOCAPownzero\endcsname{+2.0}
\expandafter\gdef\csname NUM@daitwocellIEMOCAPselffive\endcsname{+1.4}
\expandafter\gdef\csname NUM@daitwocellIEMOCAPselftwofive\endcsname{\textminus{}0.4}
\expandafter\gdef\csname NUM@daitwocellIEMOCAPselfzero\endcsname{\textminus{}0.3}
\expandafter\gdef\csname NUM@daitwocellIEMOCAPstfive\endcsname{+2.5}
\expandafter\gdef\csname NUM@daitwocellIEMOCAPsttwofive\endcsname{\textminus{}0.7}
\expandafter\gdef\csname NUM@daitwocellIEMOCAPstzero\endcsname{+2.6}
\expandafter\gdef\csname NUM@daitwocellIEMOCAPtthreeafive\endcsname{\textbf{+7.2}}
\expandafter\gdef\csname NUM@daitwocellIEMOCAPtthreeatwofive\endcsname{\textbf{+7.6}}
\expandafter\gdef\csname NUM@daitwocellIEMOCAPtthreeazero\endcsname{+4.6}
\expandafter\gdef\csname NUM@daitwocellMELDcoralfive\endcsname{\textbf{+3.8}}
\expandafter\gdef\csname NUM@daitwocellMELDcoraltwofive\endcsname{\textbf{+2.4}}
\expandafter\gdef\csname NUM@daitwocellMELDcoralzero\endcsname{\textbf{+5.4}}
\expandafter\gdef\csname NUM@daitwocellMELDdannfive\endcsname{\textbf{+3.2}}
\expandafter\gdef\csname NUM@daitwocellMELDdanntwofive\endcsname{\textbf{+3.8}}
\expandafter\gdef\csname NUM@daitwocellMELDdannzero\endcsname{\textbf{+4.3}}
\expandafter\gdef\csname NUM@daitwocellMELDentminfive\endcsname{\textbf{+13.3}}
\expandafter\gdef\csname NUM@daitwocellMELDentmintwofive\endcsname{\textbf{+10.7}}
\expandafter\gdef\csname NUM@daitwocellMELDentminzero\endcsname{\textbf{+9.0}}
\expandafter\gdef\csname NUM@daitwocellMELDimfive\endcsname{\textbf{+10.5}}
\expandafter\gdef\csname NUM@daitwocellMELDimtwofive\endcsname{\textbf{+4.5}}
\expandafter\gdef\csname NUM@daitwocellMELDimzero\endcsname{\textbf{+12.1}}
\expandafter\gdef\csname NUM@daitwocellMELDlpfive\endcsname{\textbf{+16.5}}
\expandafter\gdef\csname NUM@daitwocellMELDlptwofive\endcsname{\textbf{+10.9}}
\expandafter\gdef\csname NUM@daitwocellMELDmmefive\endcsname{\textbf{+19.0}}
\expandafter\gdef\csname NUM@daitwocellMELDmmetwofive\endcsname{\textbf{+5.5}}
\expandafter\gdef\csname NUM@daitwocellMELDmmezero\endcsname{\textbf{+24.7}}
\expandafter\gdef\csname NUM@daitwocellMELDownfive\endcsname{\textbf{+2.2}}
\expandafter\gdef\csname NUM@daitwocellMELDowntwofive\endcsname{\textbf{+1.8}}
\expandafter\gdef\csname NUM@daitwocellMELDownzero\endcsname{\textbf{+3.0}}
\expandafter\gdef\csname NUM@daitwocellMELDselffive\endcsname{\textbf{+4.1}}
\expandafter\gdef\csname NUM@daitwocellMELDselftwofive\endcsname{\textbf{+3.4}}
\expandafter\gdef\csname NUM@daitwocellMELDselfzero\endcsname{+3.7}
\expandafter\gdef\csname NUM@daitwocellMELDstfive\endcsname{\textminus{}1.7}
\expandafter\gdef\csname NUM@daitwocellMELDsttwofive\endcsname{\textminus{}0.9}
\expandafter\gdef\csname NUM@daitwocellMELDstzero\endcsname{\textminus{}0.8}
\expandafter\gdef\csname NUM@daitwocellMELDtthreeafive\endcsname{\textbf{+8.4}}
\expandafter\gdef\csname NUM@daitwocellMELDtthreeatwofive\endcsname{\textbf{+8.5}}
\expandafter\gdef\csname NUM@daitwocellMELDtthreeazero\endcsname{\textbf{+6.9}}
\expandafter\gdef\csname NUM@daitwocellMOSEIcoralfive\endcsname{\textcolor{mBad}{\textminus{}1.0}}
\expandafter\gdef\csname NUM@daitwocellMOSEIcoraltwofive\endcsname{+0.4}
\expandafter\gdef\csname NUM@daitwocellMOSEIcoralzero\endcsname{\textminus{}0.8}
\expandafter\gdef\csname NUM@daitwocellMOSEIdannfive\endcsname{+0.5}
\expandafter\gdef\csname NUM@daitwocellMOSEIdanntwofive\endcsname{\textbf{+1.6}}
\expandafter\gdef\csname NUM@daitwocellMOSEIdannzero\endcsname{+0.6}
\expandafter\gdef\csname NUM@daitwocellMOSEIentminfive\endcsname{\textbf{+12.8}}
\expandafter\gdef\csname NUM@daitwocellMOSEIentmintwofive\endcsname{\textbf{+5.0}}
\expandafter\gdef\csname NUM@daitwocellMOSEIentminzero\endcsname{\textbf{+4.1}}
\expandafter\gdef\csname NUM@daitwocellMOSEIimfive\endcsname{\textbf{+1.7}}
\expandafter\gdef\csname NUM@daitwocellMOSEIimtwofive\endcsname{\textbf{+3.8}}
\expandafter\gdef\csname NUM@daitwocellMOSEIimzero\endcsname{\textbf{+3.0}}
\expandafter\gdef\csname NUM@daitwocellMOSEIlpfive\endcsname{\textbf{+6.6}}
\expandafter\gdef\csname NUM@daitwocellMOSEIlptwofive\endcsname{\textbf{+5.2}}
\expandafter\gdef\csname NUM@daitwocellMOSEImmefive\endcsname{\textbf{+6.3}}
\expandafter\gdef\csname NUM@daitwocellMOSEImmetwofive\endcsname{\textbf{+6.6}}
\expandafter\gdef\csname NUM@daitwocellMOSEImmezero\endcsname{\textbf{+12.5}}
\expandafter\gdef\csname NUM@daitwocellMOSEIownfive\endcsname{\textminus{}0.5}
\expandafter\gdef\csname NUM@daitwocellMOSEIowntwofive\endcsname{+0.4}
\expandafter\gdef\csname NUM@daitwocellMOSEIownzero\endcsname{\textminus{}0.7}
\expandafter\gdef\csname NUM@daitwocellMOSEIselffive\endcsname{\textbf{+0.8}}
\expandafter\gdef\csname NUM@daitwocellMOSEIselftwofive\endcsname{\textbf{+1.3}}
\expandafter\gdef\csname NUM@daitwocellMOSEIselfzero\endcsname{\textminus{}0.0}
\expandafter\gdef\csname NUM@daitwocellMOSEIstfive\endcsname{+0.0}
\expandafter\gdef\csname NUM@daitwocellMOSEIsttwofive\endcsname{+0.7}
\expandafter\gdef\csname NUM@daitwocellMOSEIstzero\endcsname{\textminus{}0.5}
\expandafter\gdef\csname NUM@daitwocellMOSEItthreeafive\endcsname{\textbf{+3.7}}
\expandafter\gdef\csname NUM@daitwocellMOSEItthreeatwofive\endcsname{\textbf{+3.8}}
\expandafter\gdef\csname NUM@daitwocellMOSEItthreeazero\endcsname{\textbf{+3.4}}
\expandafter\gdef\csname NUM@daitwohiAFEWcoralfive\endcsname{+3.8}
\expandafter\gdef\csname NUM@daitwohiAFEWcoraltwofive\endcsname{+1.8}
\expandafter\gdef\csname NUM@daitwohiAFEWcoralzero\endcsname{+9.6}
\expandafter\gdef\csname NUM@daitwohiAFEWdannfive\endcsname{+3.6}
\expandafter\gdef\csname NUM@daitwohiAFEWdanntwofive\endcsname{+2.5}
\expandafter\gdef\csname NUM@daitwohiAFEWdannzero\endcsname{+5.1}
\expandafter\gdef\csname NUM@daitwohiAFEWentminfive\endcsname{+2.2}
\expandafter\gdef\csname NUM@daitwohiAFEWentmintwofive\endcsname{+1.9}
\expandafter\gdef\csname NUM@daitwohiAFEWentminzero\endcsname{\textminus{}1.1}
\expandafter\gdef\csname NUM@daitwohiAFEWimfive\endcsname{+1.7}
\expandafter\gdef\csname NUM@daitwohiAFEWimtwofive\endcsname{+1.5}
\expandafter\gdef\csname NUM@daitwohiAFEWimzero\endcsname{\textminus{}0.7}
\expandafter\gdef\csname NUM@daitwohiAFEWlpfive\endcsname{+14.1}
\expandafter\gdef\csname NUM@daitwohiAFEWlptwofive\endcsname{+12.7}
\expandafter\gdef\csname NUM@daitwohiAFEWmmefive\endcsname{+2.6}
\expandafter\gdef\csname NUM@daitwohiAFEWmmetwofive\endcsname{+1.9}
\expandafter\gdef\csname NUM@daitwohiAFEWmmezero\endcsname{+3.8}
\expandafter\gdef\csname NUM@daitwohiAFEWownfive\endcsname{+2.4}
\expandafter\gdef\csname NUM@daitwohiAFEWowntwofive\endcsname{+1.8}
\expandafter\gdef\csname NUM@daitwohiAFEWownzero\endcsname{+3.2}
\expandafter\gdef\csname NUM@daitwohiAFEWselffive\endcsname{+1.0}
\expandafter\gdef\csname NUM@daitwohiAFEWselftwofive\endcsname{+1.2}
\expandafter\gdef\csname NUM@daitwohiAFEWselfzero\endcsname{+2.7}
\expandafter\gdef\csname NUM@daitwohiAFEWstfive\endcsname{+9.1}
\expandafter\gdef\csname NUM@daitwohiAFEWsttwofive\endcsname{+9.9}
\expandafter\gdef\csname NUM@daitwohiAFEWstzero\endcsname{+7.6}
\expandafter\gdef\csname NUM@daitwohiAFEWtthreeafive\endcsname{+2.4}
\expandafter\gdef\csname NUM@daitwohiAFEWtthreeatwofive\endcsname{+4.2}
\expandafter\gdef\csname NUM@daitwohiAFEWtthreeazero\endcsname{\textminus{}2.1}
\expandafter\gdef\csname NUM@daitwohiIEMOCAPcoralfive\endcsname{+5.9}
\expandafter\gdef\csname NUM@daitwohiIEMOCAPcoraltwofive\endcsname{+4.6}
\expandafter\gdef\csname NUM@daitwohiIEMOCAPcoralzero\endcsname{+9.7}
\expandafter\gdef\csname NUM@daitwohiIEMOCAPdannfive\endcsname{+3.1}
\expandafter\gdef\csname NUM@daitwohiIEMOCAPdanntwofive\endcsname{+8.9}
\expandafter\gdef\csname NUM@daitwohiIEMOCAPdannzero\endcsname{+9.6}
\expandafter\gdef\csname NUM@daitwohiIEMOCAPentminfive\endcsname{+17.8}
\expandafter\gdef\csname NUM@daitwohiIEMOCAPentmintwofive\endcsname{+41.6}
\expandafter\gdef\csname NUM@daitwohiIEMOCAPentminzero\endcsname{+27.3}
\expandafter\gdef\csname NUM@daitwohiIEMOCAPimfive\endcsname{+3.6}
\expandafter\gdef\csname NUM@daitwohiIEMOCAPimtwofive\endcsname{+1.5}
\expandafter\gdef\csname NUM@daitwohiIEMOCAPimzero\endcsname{+16.5}
\expandafter\gdef\csname NUM@daitwohiIEMOCAPlpfive\endcsname{+17.2}
\expandafter\gdef\csname NUM@daitwohiIEMOCAPlptwofive\endcsname{+8.4}
\expandafter\gdef\csname NUM@daitwohiIEMOCAPmmefive\endcsname{+16.9}
\expandafter\gdef\csname NUM@daitwohiIEMOCAPmmetwofive\endcsname{+3.8}
\expandafter\gdef\csname NUM@daitwohiIEMOCAPmmezero\endcsname{+34.4}
\expandafter\gdef\csname NUM@daitwohiIEMOCAPownfive\endcsname{+2.5}
\expandafter\gdef\csname NUM@daitwohiIEMOCAPowntwofive\endcsname{+1.3}
\expandafter\gdef\csname NUM@daitwohiIEMOCAPownzero\endcsname{+4.4}
\expandafter\gdef\csname NUM@daitwohiIEMOCAPselffive\endcsname{+8.9}
\expandafter\gdef\csname NUM@daitwohiIEMOCAPselftwofive\endcsname{+6.6}
\expandafter\gdef\csname NUM@daitwohiIEMOCAPselfzero\endcsname{+7.2}
\expandafter\gdef\csname NUM@daitwohiIEMOCAPstfive\endcsname{+6.0}
\expandafter\gdef\csname NUM@daitwohiIEMOCAPsttwofive\endcsname{+1.5}
\expandafter\gdef\csname NUM@daitwohiIEMOCAPstzero\endcsname{+6.6}
\expandafter\gdef\csname NUM@daitwohiIEMOCAPtthreeafive\endcsname{+11.3}
\expandafter\gdef\csname NUM@daitwohiIEMOCAPtthreeatwofive\endcsname{+10.8}
\expandafter\gdef\csname NUM@daitwohiIEMOCAPtthreeazero\endcsname{+10.7}
\expandafter\gdef\csname NUM@daitwohiMELDcoralfive\endcsname{+4.7}
\expandafter\gdef\csname NUM@daitwohiMELDcoraltwofive\endcsname{+3.7}
\expandafter\gdef\csname NUM@daitwohiMELDcoralzero\endcsname{+7.5}
\expandafter\gdef\csname NUM@daitwohiMELDdannfive\endcsname{+4.8}
\expandafter\gdef\csname NUM@daitwohiMELDdanntwofive\endcsname{+5.8}
\expandafter\gdef\csname NUM@daitwohiMELDdannzero\endcsname{+6.7}
\expandafter\gdef\csname NUM@daitwohiMELDentminfive\endcsname{+17.2}
\expandafter\gdef\csname NUM@daitwohiMELDentmintwofive\endcsname{+13.6}
\expandafter\gdef\csname NUM@daitwohiMELDentminzero\endcsname{+12.7}
\expandafter\gdef\csname NUM@daitwohiMELDimfive\endcsname{+13.0}
\expandafter\gdef\csname NUM@daitwohiMELDimtwofive\endcsname{+6.0}
\expandafter\gdef\csname NUM@daitwohiMELDimzero\endcsname{+16.1}
\expandafter\gdef\csname NUM@daitwohiMELDlpfive\endcsname{+20.1}
\expandafter\gdef\csname NUM@daitwohiMELDlptwofive\endcsname{+13.6}
\expandafter\gdef\csname NUM@daitwohiMELDmmefive\endcsname{+22.6}
\expandafter\gdef\csname NUM@daitwohiMELDmmetwofive\endcsname{+7.2}
\expandafter\gdef\csname NUM@daitwohiMELDmmezero\endcsname{+31.1}
\expandafter\gdef\csname NUM@daitwohiMELDownfive\endcsname{+3.4}
\expandafter\gdef\csname NUM@daitwohiMELDowntwofive\endcsname{+3.3}
\expandafter\gdef\csname NUM@daitwohiMELDownzero\endcsname{+4.8}
\expandafter\gdef\csname NUM@daitwohiMELDselffive\endcsname{+5.4}
\expandafter\gdef\csname NUM@daitwohiMELDselftwofive\endcsname{+4.7}
\expandafter\gdef\csname NUM@daitwohiMELDselfzero\endcsname{+7.5}
\expandafter\gdef\csname NUM@daitwohiMELDstfive\endcsname{+0.5}
\expandafter\gdef\csname NUM@daitwohiMELDsttwofive\endcsname{+0.8}
\expandafter\gdef\csname NUM@daitwohiMELDstzero\endcsname{+2.3}
\expandafter\gdef\csname NUM@daitwohiMELDtthreeafive\endcsname{+10.3}
\expandafter\gdef\csname NUM@daitwohiMELDtthreeatwofive\endcsname{+10.6}
\expandafter\gdef\csname NUM@daitwohiMELDtthreeazero\endcsname{+10.4}
\expandafter\gdef\csname NUM@daitwohiMOSEIcoralfive\endcsname{\textminus{}0.3}
\expandafter\gdef\csname NUM@daitwohiMOSEIcoraltwofive\endcsname{+0.9}
\expandafter\gdef\csname NUM@daitwohiMOSEIcoralzero\endcsname{+0.2}
\expandafter\gdef\csname NUM@daitwohiMOSEIdannfive\endcsname{+1.4}
\expandafter\gdef\csname NUM@daitwohiMOSEIdanntwofive\endcsname{+2.6}
\expandafter\gdef\csname NUM@daitwohiMOSEIdannzero\endcsname{+1.9}
\expandafter\gdef\csname NUM@daitwohiMOSEIentminfive\endcsname{+15.2}
\expandafter\gdef\csname NUM@daitwohiMOSEIentmintwofive\endcsname{+5.9}
\expandafter\gdef\csname NUM@daitwohiMOSEIentminzero\endcsname{+5.4}
\expandafter\gdef\csname NUM@daitwohiMOSEIimfive\endcsname{+2.4}
\expandafter\gdef\csname NUM@daitwohiMOSEIimtwofive\endcsname{+4.7}
\expandafter\gdef\csname NUM@daitwohiMOSEIimzero\endcsname{+4.2}
\expandafter\gdef\csname NUM@daitwohiMOSEIlpfive\endcsname{+8.4}
\expandafter\gdef\csname NUM@daitwohiMOSEIlptwofive\endcsname{+6.4}
\expandafter\gdef\csname NUM@daitwohiMOSEImmefive\endcsname{+8.7}
\expandafter\gdef\csname NUM@daitwohiMOSEImmetwofive\endcsname{+8.0}
\expandafter\gdef\csname NUM@daitwohiMOSEImmezero\endcsname{+14.8}
\expandafter\gdef\csname NUM@daitwohiMOSEIownfive\endcsname{+0.0}
\expandafter\gdef\csname NUM@daitwohiMOSEIowntwofive\endcsname{+0.9}
\expandafter\gdef\csname NUM@daitwohiMOSEIownzero\endcsname{+0.0}
\expandafter\gdef\csname NUM@daitwohiMOSEIselffive\endcsname{+1.4}
\expandafter\gdef\csname NUM@daitwohiMOSEIselftwofive\endcsname{+2.0}
\expandafter\gdef\csname NUM@daitwohiMOSEIselfzero\endcsname{+0.9}
\expandafter\gdef\csname NUM@daitwohiMOSEIstfive\endcsname{+1.0}
\expandafter\gdef\csname NUM@daitwohiMOSEIsttwofive\endcsname{+1.4}
\expandafter\gdef\csname NUM@daitwohiMOSEIstzero\endcsname{+0.6}
\expandafter\gdef\csname NUM@daitwohiMOSEItthreeafive\endcsname{+4.6}
\expandafter\gdef\csname NUM@daitwohiMOSEItthreeatwofive\endcsname{+4.7}
\expandafter\gdef\csname NUM@daitwohiMOSEItthreeazero\endcsname{+4.6}
\expandafter\gdef\csname NUM@daitwoloAFEWcoralfive\endcsname{\textminus{}0.1}
\expandafter\gdef\csname NUM@daitwoloAFEWcoraltwofive\endcsname{\textminus{}0.6}
\expandafter\gdef\csname NUM@daitwoloAFEWcoralzero\endcsname{+0.9}
\expandafter\gdef\csname NUM@daitwoloAFEWdannfive\endcsname{\textminus{}1.4}
\expandafter\gdef\csname NUM@daitwoloAFEWdanntwofive\endcsname{\textminus{}0.6}
\expandafter\gdef\csname NUM@daitwoloAFEWdannzero\endcsname{\textminus{}2.0}
\expandafter\gdef\csname NUM@daitwoloAFEWentminfive\endcsname{\textminus{}3.2}
\expandafter\gdef\csname NUM@daitwoloAFEWentmintwofive\endcsname{\textminus{}1.9}
\expandafter\gdef\csname NUM@daitwoloAFEWentminzero\endcsname{\textminus{}7.9}
\expandafter\gdef\csname NUM@daitwoloAFEWimfive\endcsname{\textminus{}1.1}
\expandafter\gdef\csname NUM@daitwoloAFEWimtwofive\endcsname{\textminus{}0.2}
\expandafter\gdef\csname NUM@daitwoloAFEWimzero\endcsname{\textminus{}4.2}
\expandafter\gdef\csname NUM@daitwoloAFEWlpfive\endcsname{+7.9}
\expandafter\gdef\csname NUM@daitwoloAFEWlptwofive\endcsname{+6.9}
\expandafter\gdef\csname NUM@daitwoloAFEWmmefive\endcsname{\textminus{}1.8}
\expandafter\gdef\csname NUM@daitwoloAFEWmmetwofive\endcsname{\textminus{}1.6}
\expandafter\gdef\csname NUM@daitwoloAFEWmmezero\endcsname{\textminus{}4.4}
\expandafter\gdef\csname NUM@daitwoloAFEWownfive\endcsname{\textminus{}0.8}
\expandafter\gdef\csname NUM@daitwoloAFEWowntwofive\endcsname{+0.1}
\expandafter\gdef\csname NUM@daitwoloAFEWownzero\endcsname{\textminus{}2.8}
\expandafter\gdef\csname NUM@daitwoloAFEWselffive\endcsname{\textminus{}3.8}
\expandafter\gdef\csname NUM@daitwoloAFEWselftwofive\endcsname{\textminus{}0.7}
\expandafter\gdef\csname NUM@daitwoloAFEWselfzero\endcsname{\textminus{}5.0}
\expandafter\gdef\csname NUM@daitwoloAFEWstfive\endcsname{+3.4}
\expandafter\gdef\csname NUM@daitwoloAFEWsttwofive\endcsname{+4.0}
\expandafter\gdef\csname NUM@daitwoloAFEWstzero\endcsname{+1.7}
\expandafter\gdef\csname NUM@daitwoloAFEWtthreeafive\endcsname{\textminus{}4.8}
\expandafter\gdef\csname NUM@daitwoloAFEWtthreeatwofive\endcsname{\textminus{}0.4}
\expandafter\gdef\csname NUM@daitwoloAFEWtthreeazero\endcsname{\textminus{}8.2}
\expandafter\gdef\csname NUM@daitwoloIEMOCAPcoralfive\endcsname{\textminus{}2.0}
\expandafter\gdef\csname NUM@daitwoloIEMOCAPcoraltwofive\endcsname{\textminus{}3.9}
\expandafter\gdef\csname NUM@daitwoloIEMOCAPcoralzero\endcsname{+0.5}
\expandafter\gdef\csname NUM@daitwoloIEMOCAPdannfive\endcsname{\textminus{}0.3}
\expandafter\gdef\csname NUM@daitwoloIEMOCAPdanntwofive\endcsname{\textminus{}6.6}
\expandafter\gdef\csname NUM@daitwoloIEMOCAPdannzero\endcsname{\textminus{}2.5}
\expandafter\gdef\csname NUM@daitwoloIEMOCAPentminfive\endcsname{+3.5}
\expandafter\gdef\csname NUM@daitwoloIEMOCAPentmintwofive\endcsname{+12.1}
\expandafter\gdef\csname NUM@daitwoloIEMOCAPentminzero\endcsname{+5.1}
\expandafter\gdef\csname NUM@daitwoloIEMOCAPimfive\endcsname{+0.6}
\expandafter\gdef\csname NUM@daitwoloIEMOCAPimtwofive\endcsname{\textminus{}1.2}
\expandafter\gdef\csname NUM@daitwoloIEMOCAPimzero\endcsname{+7.5}
\expandafter\gdef\csname NUM@daitwoloIEMOCAPlpfive\endcsname{+6.3}
\expandafter\gdef\csname NUM@daitwoloIEMOCAPlptwofive\endcsname{+1.0}
\expandafter\gdef\csname NUM@daitwoloIEMOCAPmmefive\endcsname{+3.1}
\expandafter\gdef\csname NUM@daitwoloIEMOCAPmmetwofive\endcsname{\textminus{}3.3}
\expandafter\gdef\csname NUM@daitwoloIEMOCAPmmezero\endcsname{+22.0}
\expandafter\gdef\csname NUM@daitwoloIEMOCAPownfive\endcsname{\textminus{}0.3}
\expandafter\gdef\csname NUM@daitwoloIEMOCAPowntwofive\endcsname{\textminus{}0.7}
\expandafter\gdef\csname NUM@daitwoloIEMOCAPownzero\endcsname{\textminus{}0.3}
\expandafter\gdef\csname NUM@daitwoloIEMOCAPselffive\endcsname{\textminus{}6.1}
\expandafter\gdef\csname NUM@daitwoloIEMOCAPselftwofive\endcsname{\textminus{}7.5}
\expandafter\gdef\csname NUM@daitwoloIEMOCAPselfzero\endcsname{\textminus{}7.7}
\expandafter\gdef\csname NUM@daitwoloIEMOCAPstfive\endcsname{\textminus{}1.0}
\expandafter\gdef\csname NUM@daitwoloIEMOCAPsttwofive\endcsname{\textminus{}3.0}
\expandafter\gdef\csname NUM@daitwoloIEMOCAPstzero\endcsname{\textminus{}1.3}
\expandafter\gdef\csname NUM@daitwoloIEMOCAPtthreeafive\endcsname{+3.2}
\expandafter\gdef\csname NUM@daitwoloIEMOCAPtthreeatwofive\endcsname{+4.3}
\expandafter\gdef\csname NUM@daitwoloIEMOCAPtthreeazero\endcsname{\textminus{}1.5}
\expandafter\gdef\csname NUM@daitwoloMELDcoralfive\endcsname{+2.9}
\expandafter\gdef\csname NUM@daitwoloMELDcoraltwofive\endcsname{+1.2}
\expandafter\gdef\csname NUM@daitwoloMELDcoralzero\endcsname{+3.4}
\expandafter\gdef\csname NUM@daitwoloMELDdannfive\endcsname{+1.6}
\expandafter\gdef\csname NUM@daitwoloMELDdanntwofive\endcsname{+1.9}
\expandafter\gdef\csname NUM@daitwoloMELDdannzero\endcsname{+1.8}
\expandafter\gdef\csname NUM@daitwoloMELDentminfive\endcsname{+9.3}
\expandafter\gdef\csname NUM@daitwoloMELDentmintwofive\endcsname{+7.9}
\expandafter\gdef\csname NUM@daitwoloMELDentminzero\endcsname{+5.3}
\expandafter\gdef\csname NUM@daitwoloMELDimfive\endcsname{+8.1}
\expandafter\gdef\csname NUM@daitwoloMELDimtwofive\endcsname{+2.9}
\expandafter\gdef\csname NUM@daitwoloMELDimzero\endcsname{+8.2}
\expandafter\gdef\csname NUM@daitwoloMELDlpfive\endcsname{+13.0}
\expandafter\gdef\csname NUM@daitwoloMELDlptwofive\endcsname{+8.1}
\expandafter\gdef\csname NUM@daitwoloMELDmmefive\endcsname{+15.4}
\expandafter\gdef\csname NUM@daitwoloMELDmmetwofive\endcsname{+3.8}
\expandafter\gdef\csname NUM@daitwoloMELDmmezero\endcsname{+18.3}
\expandafter\gdef\csname NUM@daitwoloMELDownfive\endcsname{+1.1}
\expandafter\gdef\csname NUM@daitwoloMELDowntwofive\endcsname{+0.4}
\expandafter\gdef\csname NUM@daitwoloMELDownzero\endcsname{+1.2}
\expandafter\gdef\csname NUM@daitwoloMELDselffive\endcsname{+2.7}
\expandafter\gdef\csname NUM@daitwoloMELDselftwofive\endcsname{+2.1}
\expandafter\gdef\csname NUM@daitwoloMELDselfzero\endcsname{\textminus{}0.2}
\expandafter\gdef\csname NUM@daitwoloMELDstfive\endcsname{\textminus{}3.8}
\expandafter\gdef\csname NUM@daitwoloMELDsttwofive\endcsname{\textminus{}2.6}
\expandafter\gdef\csname NUM@daitwoloMELDstzero\endcsname{\textminus{}3.8}
\expandafter\gdef\csname NUM@daitwoloMELDtthreeafive\endcsname{+6.4}
\expandafter\gdef\csname NUM@daitwoloMELDtthreeatwofive\endcsname{+6.4}
\expandafter\gdef\csname NUM@daitwoloMELDtthreeazero\endcsname{+3.4}
\expandafter\gdef\csname NUM@daitwoloMOSEIcoralfive\endcsname{\textminus{}1.7}
\expandafter\gdef\csname NUM@daitwoloMOSEIcoraltwofive\endcsname{\textminus{}0.2}
\expandafter\gdef\csname NUM@daitwoloMOSEIcoralzero\endcsname{\textminus{}1.7}
\expandafter\gdef\csname NUM@daitwoloMOSEIdannfive\endcsname{\textminus{}0.4}
\expandafter\gdef\csname NUM@daitwoloMOSEIdanntwofive\endcsname{+0.7}
\expandafter\gdef\csname NUM@daitwoloMOSEIdannzero\endcsname{\textminus{}0.7}
\expandafter\gdef\csname NUM@daitwoloMOSEIentminfive\endcsname{+10.3}
\expandafter\gdef\csname NUM@daitwoloMOSEIentmintwofive\endcsname{+4.0}
\expandafter\gdef\csname NUM@daitwoloMOSEIentminzero\endcsname{+2.8}
\expandafter\gdef\csname NUM@daitwoloMOSEIimfive\endcsname{+0.9}
\expandafter\gdef\csname NUM@daitwoloMOSEIimtwofive\endcsname{+3.0}
\expandafter\gdef\csname NUM@daitwoloMOSEIimzero\endcsname{+1.8}
\expandafter\gdef\csname NUM@daitwoloMOSEIlpfive\endcsname{+4.8}
\expandafter\gdef\csname NUM@daitwoloMOSEIlptwofive\endcsname{+4.1}
\expandafter\gdef\csname NUM@daitwoloMOSEImmefive\endcsname{+3.8}
\expandafter\gdef\csname NUM@daitwoloMOSEImmetwofive\endcsname{+5.1}
\expandafter\gdef\csname NUM@daitwoloMOSEImmezero\endcsname{+10.3}
\expandafter\gdef\csname NUM@daitwoloMOSEIownfive\endcsname{\textminus{}1.0}
\expandafter\gdef\csname NUM@daitwoloMOSEIowntwofive\endcsname{\textminus{}0.1}
\expandafter\gdef\csname NUM@daitwoloMOSEIownzero\endcsname{\textminus{}1.5}
\expandafter\gdef\csname NUM@daitwoloMOSEIselffive\endcsname{+0.1}
\expandafter\gdef\csname NUM@daitwoloMOSEIselftwofive\endcsname{+0.6}
\expandafter\gdef\csname NUM@daitwoloMOSEIselfzero\endcsname{\textminus{}0.9}
\expandafter\gdef\csname NUM@daitwoloMOSEIstfive\endcsname{\textminus{}0.9}
\expandafter\gdef\csname NUM@daitwoloMOSEIsttwofive\endcsname{\textminus{}0.1}
\expandafter\gdef\csname NUM@daitwoloMOSEIstzero\endcsname{\textminus{}1.6}
\expandafter\gdef\csname NUM@daitwoloMOSEItthreeafive\endcsname{+2.8}
\expandafter\gdef\csname NUM@daitwoloMOSEItthreeatwofive\endcsname{+2.9}
\expandafter\gdef\csname NUM@daitwoloMOSEItthreeazero\endcsname{+2.3}
\expandafter\gdef\csname NUM@daitwoncoral\endcsname{12}
\expandafter\gdef\csname NUM@daitwondann\endcsname{12}
\expandafter\gdef\csname NUM@daitwonentmin\endcsname{12}
\expandafter\gdef\csname NUM@daitwonim\endcsname{12}
\expandafter\gdef\csname NUM@daitwonlp\endcsname{8}
\expandafter\gdef\csname NUM@daitwonmme\endcsname{12}
\expandafter\gdef\csname NUM@daitwonown\endcsname{12}
\expandafter\gdef\csname NUM@daitwonself\endcsname{12}
\expandafter\gdef\csname NUM@daitwonst\endcsname{12}
\expandafter\gdef\csname NUM@daitwontthreea\endcsname{12}
\expandafter\gdef\csname NUM@daitwopropAFEWfive\endcsname{70.8}
\expandafter\gdef\csname NUM@daitwopropAFEWtwofive\endcsname{75.8}
\expandafter\gdef\csname NUM@daitwopropAFEWzero\endcsname{66.2}
\expandafter\gdef\csname NUM@daitwopropIEMOCAPfive\endcsname{72.7}
\expandafter\gdef\csname NUM@daitwopropIEMOCAPtwofive\endcsname{74.7}
\expandafter\gdef\csname NUM@daitwopropIEMOCAPzero\endcsname{71.3}
\expandafter\gdef\csname NUM@daitwopropMELDfive\endcsname{50.1}
\expandafter\gdef\csname NUM@daitwopropMELDtwofive\endcsname{51.9}
\expandafter\gdef\csname NUM@daitwopropMELDzero\endcsname{50.6}
\expandafter\gdef\csname NUM@daitwopropMOSEIfive\endcsname{41.1}
\expandafter\gdef\csname NUM@daitwopropMOSEItwofive\endcsname{44.1}
\expandafter\gdef\csname NUM@daitwopropMOSEIzero\endcsname{39.3}
\expandafter\gdef\csname NUM@daitwouarAFEWcoralfive\endcsname{68.9}
\expandafter\gdef\csname NUM@daitwouarAFEWcoraltwofive\endcsname{75.2}
\expandafter\gdef\csname NUM@daitwouarAFEWcoralzero\endcsname{61.0}
\expandafter\gdef\csname NUM@daitwouarAFEWdannfive\endcsname{69.7}
\expandafter\gdef\csname NUM@daitwouarAFEWdanntwofive\endcsname{75.0}
\expandafter\gdef\csname NUM@daitwouarAFEWdannzero\endcsname{64.7}
\expandafter\gdef\csname NUM@daitwouarAFEWentminfive\endcsname{71.3}
\expandafter\gdef\csname NUM@daitwouarAFEWentmintwofive\endcsname{75.8}
\expandafter\gdef\csname NUM@daitwouarAFEWentminzero\endcsname{70.7}
\expandafter\gdef\csname NUM@daitwouarAFEWimfive\endcsname{70.4}
\expandafter\gdef\csname NUM@daitwouarAFEWimtwofive\endcsname{75.2}
\expandafter\gdef\csname NUM@daitwouarAFEWimzero\endcsname{68.7}
\expandafter\gdef\csname NUM@daitwouarAFEWlpfive\endcsname{59.8}
\expandafter\gdef\csname NUM@daitwouarAFEWlptwofive\endcsname{66.0}
\expandafter\gdef\csname NUM@daitwouarAFEWmmefive\endcsname{70.4}
\expandafter\gdef\csname NUM@daitwouarAFEWmmetwofive\endcsname{75.7}
\expandafter\gdef\csname NUM@daitwouarAFEWmmezero\endcsname{66.6}
\expandafter\gdef\csname NUM@daitwouarAFEWownfive\endcsname{69.9}
\expandafter\gdef\csname NUM@daitwouarAFEWowntwofive\endcsname{74.9}
\expandafter\gdef\csname NUM@daitwouarAFEWownzero\endcsname{66.1}
\expandafter\gdef\csname NUM@daitwouarAFEWselffive\endcsname{72.1}
\expandafter\gdef\csname NUM@daitwouarAFEWselftwofive\endcsname{75.6}
\expandafter\gdef\csname NUM@daitwouarAFEWselfzero\endcsname{67.4}
\expandafter\gdef\csname NUM@daitwouarAFEWstfive\endcsname{64.5}
\expandafter\gdef\csname NUM@daitwouarAFEWsttwofive\endcsname{68.9}
\expandafter\gdef\csname NUM@daitwouarAFEWstzero\endcsname{61.6}
\expandafter\gdef\csname NUM@daitwouarAFEWtthreeafive\endcsname{72.0}
\expandafter\gdef\csname NUM@daitwouarAFEWtthreeatwofive\endcsname{73.9}
\expandafter\gdef\csname NUM@daitwouarAFEWtthreeazero\endcsname{71.4}
\expandafter\gdef\csname NUM@daitwouarIEMOCAPcoralfive\endcsname{70.7}
\expandafter\gdef\csname NUM@daitwouarIEMOCAPcoraltwofive\endcsname{74.3}
\expandafter\gdef\csname NUM@daitwouarIEMOCAPcoralzero\endcsname{66.2}
\expandafter\gdef\csname NUM@daitwouarIEMOCAPdannfive\endcsname{71.6}
\expandafter\gdef\csname NUM@daitwouarIEMOCAPdanntwofive\endcsname{74.2}
\expandafter\gdef\csname NUM@daitwouarIEMOCAPdannzero\endcsname{67.7}
\expandafter\gdef\csname NUM@daitwouarIEMOCAPentminfive\endcsname{62.1}
\expandafter\gdef\csname NUM@daitwouarIEMOCAPentmintwofive\endcsname{47.8}
\expandafter\gdef\csname NUM@daitwouarIEMOCAPentminzero\endcsname{55.1}
\expandafter\gdef\csname NUM@daitwouarIEMOCAPimfive\endcsname{70.6}
\expandafter\gdef\csname NUM@daitwouarIEMOCAPimtwofive\endcsname{74.5}
\expandafter\gdef\csname NUM@daitwouarIEMOCAPimzero\endcsname{59.3}
\expandafter\gdef\csname NUM@daitwouarIEMOCAPlpfive\endcsname{61.0}
\expandafter\gdef\csname NUM@daitwouarIEMOCAPlptwofive\endcsname{70.0}
\expandafter\gdef\csname NUM@daitwouarIEMOCAPmmefive\endcsname{63.0}
\expandafter\gdef\csname NUM@daitwouarIEMOCAPmmetwofive\endcsname{75.2}
\expandafter\gdef\csname NUM@daitwouarIEMOCAPmmezero\endcsname{43.0}
\expandafter\gdef\csname NUM@daitwouarIEMOCAPownfive\endcsname{71.6}
\expandafter\gdef\csname NUM@daitwouarIEMOCAPowntwofive\endcsname{74.4}
\expandafter\gdef\csname NUM@daitwouarIEMOCAPownzero\endcsname{69.3}
\expandafter\gdef\csname NUM@daitwouarIEMOCAPselffive\endcsname{71.3}
\expandafter\gdef\csname NUM@daitwouarIEMOCAPselftwofive\endcsname{75.1}
\expandafter\gdef\csname NUM@daitwouarIEMOCAPselfzero\endcsname{71.5}
\expandafter\gdef\csname NUM@daitwouarIEMOCAPstfive\endcsname{70.2}
\expandafter\gdef\csname NUM@daitwouarIEMOCAPsttwofive\endcsname{75.4}
\expandafter\gdef\csname NUM@daitwouarIEMOCAPstzero\endcsname{68.7}
\expandafter\gdef\csname NUM@daitwouarIEMOCAPtthreeafive\endcsname{65.5}
\expandafter\gdef\csname NUM@daitwouarIEMOCAPtthreeatwofive\endcsname{67.1}
\expandafter\gdef\csname NUM@daitwouarIEMOCAPtthreeazero\endcsname{66.7}
\expandafter\gdef\csname NUM@daitwouarMELDcoralfive\endcsname{46.3}
\expandafter\gdef\csname NUM@daitwouarMELDcoraltwofive\endcsname{49.5}
\expandafter\gdef\csname NUM@daitwouarMELDcoralzero\endcsname{45.1}
\expandafter\gdef\csname NUM@daitwouarMELDdannfive\endcsname{47.1}
\expandafter\gdef\csname NUM@daitwouarMELDdanntwofive\endcsname{48.3}
\expandafter\gdef\csname NUM@daitwouarMELDdannzero\endcsname{46.3}
\expandafter\gdef\csname NUM@daitwouarMELDentminfive\endcsname{36.9}
\expandafter\gdef\csname NUM@daitwouarMELDentmintwofive\endcsname{41.2}
\expandafter\gdef\csname NUM@daitwouarMELDentminzero\endcsname{41.6}
\expandafter\gdef\csname NUM@daitwouarMELDimfive\endcsname{39.6}
\expandafter\gdef\csname NUM@daitwouarMELDimtwofive\endcsname{47.4}
\expandafter\gdef\csname NUM@daitwouarMELDimzero\endcsname{38.5}
\expandafter\gdef\csname NUM@daitwouarMELDlpfive\endcsname{33.6}
\expandafter\gdef\csname NUM@daitwouarMELDlptwofive\endcsname{41.0}
\expandafter\gdef\csname NUM@daitwouarMELDmmefive\endcsname{31.2}
\expandafter\gdef\csname NUM@daitwouarMELDmmetwofive\endcsname{46.6}
\expandafter\gdef\csname NUM@daitwouarMELDmmezero\endcsname{25.9}
\expandafter\gdef\csname NUM@daitwouarMELDownfive\endcsname{47.9}
\expandafter\gdef\csname NUM@daitwouarMELDowntwofive\endcsname{50.1}
\expandafter\gdef\csname NUM@daitwouarMELDownzero\endcsname{47.6}
\expandafter\gdef\csname NUM@daitwouarMELDselffive\endcsname{46.1}
\expandafter\gdef\csname NUM@daitwouarMELDselftwofive\endcsname{48.5}
\expandafter\gdef\csname NUM@daitwouarMELDselfzero\endcsname{46.9}
\expandafter\gdef\csname NUM@daitwouarMELDstfive\endcsname{51.8}
\expandafter\gdef\csname NUM@daitwouarMELDsttwofive\endcsname{52.8}
\expandafter\gdef\csname NUM@daitwouarMELDstzero\endcsname{51.4}
\expandafter\gdef\csname NUM@daitwouarMELDtthreeafive\endcsname{41.7}
\expandafter\gdef\csname NUM@daitwouarMELDtthreeatwofive\endcsname{43.4}
\expandafter\gdef\csname NUM@daitwouarMELDtthreeazero\endcsname{43.7}
\expandafter\gdef\csname NUM@daitwouarMOSEIcoralfive\endcsname{42.1}
\expandafter\gdef\csname NUM@daitwouarMOSEIcoraltwofive\endcsname{43.7}
\expandafter\gdef\csname NUM@daitwouarMOSEIcoralzero\endcsname{40.1}
\expandafter\gdef\csname NUM@daitwouarMOSEIdannfive\endcsname{40.5}
\expandafter\gdef\csname NUM@daitwouarMOSEIdanntwofive\endcsname{42.0}
\expandafter\gdef\csname NUM@daitwouarMOSEIdannzero\endcsname{38.7}
\expandafter\gdef\csname NUM@daitwouarMOSEIentminfive\endcsname{28.3}
\expandafter\gdef\csname NUM@daitwouarMOSEIentmintwofive\endcsname{39.1}
\expandafter\gdef\csname NUM@daitwouarMOSEIentminzero\endcsname{35.2}
\expandafter\gdef\csname NUM@daitwouarMOSEIimfive\endcsname{39.4}
\expandafter\gdef\csname NUM@daitwouarMOSEIimtwofive\endcsname{40.3}
\expandafter\gdef\csname NUM@daitwouarMOSEIimzero\endcsname{36.3}
\expandafter\gdef\csname NUM@daitwouarMOSEIlpfive\endcsname{34.5}
\expandafter\gdef\csname NUM@daitwouarMOSEIlptwofive\endcsname{38.8}
\expandafter\gdef\csname NUM@daitwouarMOSEImmefive\endcsname{34.8}
\expandafter\gdef\csname NUM@daitwouarMOSEImmetwofive\endcsname{37.0}
\expandafter\gdef\csname NUM@daitwouarMOSEImmezero\endcsname{26.8}
\expandafter\gdef\csname NUM@daitwouarMOSEIownfive\endcsname{41.6}
\expandafter\gdef\csname NUM@daitwouarMOSEIowntwofive\endcsname{43.7}
\expandafter\gdef\csname NUM@daitwouarMOSEIownzero\endcsname{40.0}
\expandafter\gdef\csname NUM@daitwouarMOSEIselffive\endcsname{40.4}
\expandafter\gdef\csname NUM@daitwouarMOSEIselftwofive\endcsname{42.8}
\expandafter\gdef\csname NUM@daitwouarMOSEIselfzero\endcsname{39.3}
\expandafter\gdef\csname NUM@daitwouarMOSEIstfive\endcsname{41.1}
\expandafter\gdef\csname NUM@daitwouarMOSEIsttwofive\endcsname{43.4}
\expandafter\gdef\csname NUM@daitwouarMOSEIstzero\endcsname{39.8}
\expandafter\gdef\csname NUM@daitwouarMOSEItthreeafive\endcsname{37.4}
\expandafter\gdef\csname NUM@daitwouarMOSEItthreeatwofive\endcsname{40.3}
\expandafter\gdef\csname NUM@daitwouarMOSEItthreeazero\endcsname{35.9}
\expandafter\gdef\csname NUM@dirAFEWAVTnonefive\endcsname{9.5}
\expandafter\gdef\csname NUM@dirAFEWAVTnonetwofive\endcsname{4.7}
\expandafter\gdef\csname NUM@dirAFEWAVTpropfive\endcsname{8.7}
\expandafter\gdef\csname NUM@dirAFEWAVTproptwofive\endcsname{3.8}
\expandafter\gdef\csname NUM@dirAFEWAVTstart\endcsname{13.3}
\expandafter\gdef\csname NUM@dirAFEWAnonefive\endcsname{5.0}
\expandafter\gdef\csname NUM@dirAFEWAnonetwofive\endcsname{4.4}
\expandafter\gdef\csname NUM@dirAFEWApropfive\endcsname{2.0}
\expandafter\gdef\csname NUM@dirAFEWAproptwofive\endcsname{1.6}
\expandafter\gdef\csname NUM@dirAFEWAstart\endcsname{5.9}
\expandafter\gdef\csname NUM@dirAFEWTnonefive\endcsname{3.8}
\expandafter\gdef\csname NUM@dirAFEWTnonetwofive\endcsname{3.1}
\expandafter\gdef\csname NUM@dirAFEWTpropfive\endcsname{2.1}
\expandafter\gdef\csname NUM@dirAFEWTproptwofive\endcsname{2.5}
\expandafter\gdef\csname NUM@dirAFEWTstart\endcsname{4.8}
\expandafter\gdef\csname NUM@dirAFEWVnonefive\endcsname{3.8}
\expandafter\gdef\csname NUM@dirAFEWVnonetwofive\endcsname{1.8}
\expandafter\gdef\csname NUM@dirAFEWVpropfive\endcsname{3.7}
\expandafter\gdef\csname NUM@dirAFEWVproptwofive\endcsname{1.7}
\expandafter\gdef\csname NUM@dirAFEWVstart\endcsname{10.2}
\expandafter\gdef\csname NUM@dirIEMOCAPAVTnonefive\endcsname{8.6}
\expandafter\gdef\csname NUM@dirIEMOCAPAVTnonetwofive\endcsname{5.5}
\expandafter\gdef\csname NUM@dirIEMOCAPAVTpropfive\endcsname{7.0}
\expandafter\gdef\csname NUM@dirIEMOCAPAVTproptwofive\endcsname{4.9}
\expandafter\gdef\csname NUM@dirIEMOCAPAVTstart\endcsname{11.3}
\expandafter\gdef\csname NUM@dirIEMOCAPAnonefive\endcsname{8.0}
\expandafter\gdef\csname NUM@dirIEMOCAPAnonetwofive\endcsname{5.4}
\expandafter\gdef\csname NUM@dirIEMOCAPApropfive\endcsname{5.0}
\expandafter\gdef\csname NUM@dirIEMOCAPAproptwofive\endcsname{3.1}
\expandafter\gdef\csname NUM@dirIEMOCAPAstart\endcsname{10.8}
\expandafter\gdef\csname NUM@dirIEMOCAPTnonefive\endcsname{10.8}
\expandafter\gdef\csname NUM@dirIEMOCAPTnonetwofive\endcsname{7.5}
\expandafter\gdef\csname NUM@dirIEMOCAPTpropfive\endcsname{3.4}
\expandafter\gdef\csname NUM@dirIEMOCAPTproptwofive\endcsname{2.0}
\expandafter\gdef\csname NUM@dirIEMOCAPTstart\endcsname{12.4}
\expandafter\gdef\csname NUM@dirIEMOCAPVnonefive\endcsname{0.4}
\expandafter\gdef\csname NUM@dirIEMOCAPVnonetwofive\endcsname{0.6}
\expandafter\gdef\csname NUM@dirIEMOCAPVpropfive\endcsname{\textminus{}0.1}
\expandafter\gdef\csname NUM@dirIEMOCAPVproptwofive\endcsname{0.1}
\expandafter\gdef\csname NUM@dirIEMOCAPVstart\endcsname{0.8}
\expandafter\gdef\csname NUM@dirMELDAVTnonefive\endcsname{9.8}
\expandafter\gdef\csname NUM@dirMELDAVTnonetwofive\endcsname{7.7}
\expandafter\gdef\csname NUM@dirMELDAVTpropfive\endcsname{7.7}
\expandafter\gdef\csname NUM@dirMELDAVTproptwofive\endcsname{5.8}
\expandafter\gdef\csname NUM@dirMELDAVTstart\endcsname{9.6}
\expandafter\gdef\csname NUM@dirMELDAnonefive\endcsname{5.7}
\expandafter\gdef\csname NUM@dirMELDAnonetwofive\endcsname{5.0}
\expandafter\gdef\csname NUM@dirMELDApropfive\endcsname{2.9}
\expandafter\gdef\csname NUM@dirMELDAproptwofive\endcsname{1.9}
\expandafter\gdef\csname NUM@dirMELDAstart\endcsname{5.9}
\expandafter\gdef\csname NUM@dirMELDTnonefive\endcsname{7.2}
\expandafter\gdef\csname NUM@dirMELDTnonetwofive\endcsname{6.1}
\expandafter\gdef\csname NUM@dirMELDTpropfive\endcsname{8.1}
\expandafter\gdef\csname NUM@dirMELDTproptwofive\endcsname{6.2}
\expandafter\gdef\csname NUM@dirMELDTstart\endcsname{7.0}
\expandafter\gdef\csname NUM@dirMELDVnonefive\endcsname{0.5}
\expandafter\gdef\csname NUM@dirMELDVnonetwofive\endcsname{0.7}
\expandafter\gdef\csname NUM@dirMELDVpropfive\endcsname{0.1}
\expandafter\gdef\csname NUM@dirMELDVproptwofive\endcsname{\textminus{}0.2}
\expandafter\gdef\csname NUM@dirMELDVstart\endcsname{0.3}
\expandafter\gdef\csname NUM@dirMOSEIAVTnonefive\endcsname{10.0}
\expandafter\gdef\csname NUM@dirMOSEIAVTnonetwofive\endcsname{7.9}
\expandafter\gdef\csname NUM@dirMOSEIAVTpropfive\endcsname{10.5}
\expandafter\gdef\csname NUM@dirMOSEIAVTproptwofive\endcsname{7.3}
\expandafter\gdef\csname NUM@dirMOSEIAVTstart\endcsname{11.6}
\expandafter\gdef\csname NUM@dirMOSEIAnonefive\endcsname{11.4}
\expandafter\gdef\csname NUM@dirMOSEIAnonetwofive\endcsname{8.9}
\expandafter\gdef\csname NUM@dirMOSEIApropfive\endcsname{8.4}
\expandafter\gdef\csname NUM@dirMOSEIAproptwofive\endcsname{5.5}
\expandafter\gdef\csname NUM@dirMOSEIAstart\endcsname{12.0}
\expandafter\gdef\csname NUM@dirMOSEITnonefive\endcsname{7.7}
\expandafter\gdef\csname NUM@dirMOSEITnonetwofive\endcsname{5.1}
\expandafter\gdef\csname NUM@dirMOSEITpropfive\endcsname{8.0}
\expandafter\gdef\csname NUM@dirMOSEITproptwofive\endcsname{4.8}
\expandafter\gdef\csname NUM@dirMOSEITstart\endcsname{8.4}
\expandafter\gdef\csname NUM@dirMOSEIVnonefive\endcsname{8.4}
\expandafter\gdef\csname NUM@dirMOSEIVnonetwofive\endcsname{7.7}
\expandafter\gdef\csname NUM@dirMOSEIVpropfive\endcsname{8.1}
\expandafter\gdef\csname NUM@dirMOSEIVproptwofive\endcsname{5.1}
\expandafter\gdef\csname NUM@dirMOSEIVstart\endcsname{8.9}
\expandafter\gdef\csname NUM@famoneAFEWAVTfive\endcsname{+1.0}
\expandafter\gdef\csname NUM@famoneAFEWAVTtwofive\endcsname{+0.6}
\expandafter\gdef\csname NUM@famoneAFEWAVTzero\endcsname{+1.7}
\expandafter\gdef\csname NUM@famoneAFEWAfive\endcsname{+4.1}
\expandafter\gdef\csname NUM@famoneAFEWAtwofive\endcsname{+3.3}
\expandafter\gdef\csname NUM@famoneAFEWAzero\endcsname{+4.9}
\expandafter\gdef\csname NUM@famoneAFEWTfive\endcsname{+1.6}
\expandafter\gdef\csname NUM@famoneAFEWTtwofive\endcsname{+0.7}
\expandafter\gdef\csname NUM@famoneAFEWTzero\endcsname{+1.6}
\expandafter\gdef\csname NUM@famoneAFEWVfive\endcsname{+0.1}
\expandafter\gdef\csname NUM@famoneAFEWVtwofive\endcsname{\textminus{}0.2}
\expandafter\gdef\csname NUM@famoneAFEWVzero\endcsname{+2.5}
\expandafter\gdef\csname NUM@famoneIEMOCAPAVTfive\endcsname{+0.2}
\expandafter\gdef\csname NUM@famoneIEMOCAPAVTtwofive\endcsname{\textminus{}0.8}
\expandafter\gdef\csname NUM@famoneIEMOCAPAVTzero\endcsname{+1.7}
\expandafter\gdef\csname NUM@famoneIEMOCAPAfive\endcsname{+2.4}
\expandafter\gdef\csname NUM@famoneIEMOCAPAtwofive\endcsname{+1.4}
\expandafter\gdef\csname NUM@famoneIEMOCAPAzero\endcsname{+4.8}
\expandafter\gdef\csname NUM@famoneIEMOCAPTfive\endcsname{+7.3}
\expandafter\gdef\csname NUM@famoneIEMOCAPTtwofive\endcsname{+5.0}
\expandafter\gdef\csname NUM@famoneIEMOCAPTzero\endcsname{+7.5}
\expandafter\gdef\csname NUM@famoneIEMOCAPVfive\endcsname{\textminus{}0.8}
\expandafter\gdef\csname NUM@famoneIEMOCAPVtwofive\endcsname{\textminus{}0.7}
\expandafter\gdef\csname NUM@famoneIEMOCAPVzero\endcsname{\textminus{}2.1}
\expandafter\gdef\csname NUM@famoneMELDAVTfive\endcsname{+2.5}
\expandafter\gdef\csname NUM@famoneMELDAVTtwofive\endcsname{+1.4}
\expandafter\gdef\csname NUM@famoneMELDAVTzero\endcsname{+3.4}
\expandafter\gdef\csname NUM@famoneMELDAfive\endcsname{+3.1}
\expandafter\gdef\csname NUM@famoneMELDAtwofive\endcsname{+2.9}
\expandafter\gdef\csname NUM@famoneMELDAzero\endcsname{+3.9}
\expandafter\gdef\csname NUM@famoneMELDTfive\endcsname{\textminus{}1.2}
\expandafter\gdef\csname NUM@famoneMELDTtwofive\endcsname{\textminus{}0.9}
\expandafter\gdef\csname NUM@famoneMELDTzero\endcsname{+0.2}
\expandafter\gdef\csname NUM@famoneMELDVfive\endcsname{+0.8}
\expandafter\gdef\csname NUM@famoneMELDVtwofive\endcsname{+1.1}
\expandafter\gdef\csname NUM@famoneMELDVzero\endcsname{+0.7}
\expandafter\gdef\csname NUM@famoneMOSEIAVTfive\endcsname{\textminus{}0.4}
\expandafter\gdef\csname NUM@famoneMOSEIAVTtwofive\endcsname{+0.7}
\expandafter\gdef\csname NUM@famoneMOSEIAVTzero\endcsname{\textminus{}0.3}
\expandafter\gdef\csname NUM@famoneMOSEIAfive\endcsname{+3.2}
\expandafter\gdef\csname NUM@famoneMOSEIAtwofive\endcsname{+3.4}
\expandafter\gdef\csname NUM@famoneMOSEIAzero\endcsname{+2.1}
\expandafter\gdef\csname NUM@famoneMOSEITfive\endcsname{\textminus{}0.0}
\expandafter\gdef\csname NUM@famoneMOSEITtwofive\endcsname{+0.0}
\expandafter\gdef\csname NUM@famoneMOSEITzero\endcsname{\textminus{}0.7}
\expandafter\gdef\csname NUM@famoneMOSEIVfive\endcsname{+0.2}
\expandafter\gdef\csname NUM@famoneMOSEIVtwofive\endcsname{+2.6}
\expandafter\gdef\csname NUM@famoneMOSEIVzero\endcsname{\textminus{}0.9}
\expandafter\gdef\csname NUM@famonehiAFEWAVTfive\endcsname{+2.7}
\expandafter\gdef\csname NUM@famonehiAFEWAVTtwofive\endcsname{+1.6}
\expandafter\gdef\csname NUM@famonehiAFEWAVTzero\endcsname{+5.3}
\expandafter\gdef\csname NUM@famonehiAFEWAfive\endcsname{+5.7}
\expandafter\gdef\csname NUM@famonehiAFEWAtwofive\endcsname{+5.3}
\expandafter\gdef\csname NUM@famonehiAFEWAzero\endcsname{+7.4}
\expandafter\gdef\csname NUM@famonehiAFEWTfive\endcsname{+4.8}
\expandafter\gdef\csname NUM@famonehiAFEWTtwofive\endcsname{+3.6}
\expandafter\gdef\csname NUM@famonehiAFEWTzero\endcsname{+5.5}
\expandafter\gdef\csname NUM@famonehiAFEWVfive\endcsname{+3.2}
\expandafter\gdef\csname NUM@famonehiAFEWVtwofive\endcsname{+1.0}
\expandafter\gdef\csname NUM@famonehiAFEWVzero\endcsname{+6.3}
\expandafter\gdef\csname NUM@famonehiIEMOCAPAVTfive\endcsname{+4.8}
\expandafter\gdef\csname NUM@famonehiIEMOCAPAVTtwofive\endcsname{+3.5}
\expandafter\gdef\csname NUM@famonehiIEMOCAPAVTzero\endcsname{+6.2}
\expandafter\gdef\csname NUM@famonehiIEMOCAPAfive\endcsname{+5.3}
\expandafter\gdef\csname NUM@famonehiIEMOCAPAtwofive\endcsname{+3.4}
\expandafter\gdef\csname NUM@famonehiIEMOCAPAzero\endcsname{+7.6}
\expandafter\gdef\csname NUM@famonehiIEMOCAPTfive\endcsname{+10.3}
\expandafter\gdef\csname NUM@famonehiIEMOCAPTtwofive\endcsname{+7.6}
\expandafter\gdef\csname NUM@famonehiIEMOCAPTzero\endcsname{+10.3}
\expandafter\gdef\csname NUM@famonehiIEMOCAPVfive\endcsname{+3.1}
\expandafter\gdef\csname NUM@famonehiIEMOCAPVtwofive\endcsname{+4.8}
\expandafter\gdef\csname NUM@famonehiIEMOCAPVzero\endcsname{+1.0}
\expandafter\gdef\csname NUM@famonehiMELDAVTfive\endcsname{+3.4}
\expandafter\gdef\csname NUM@famonehiMELDAVTtwofive\endcsname{+2.2}
\expandafter\gdef\csname NUM@famonehiMELDAVTzero\endcsname{+5.8}
\expandafter\gdef\csname NUM@famonehiMELDAfive\endcsname{+5.1}
\expandafter\gdef\csname NUM@famonehiMELDAtwofive\endcsname{+4.5}
\expandafter\gdef\csname NUM@famonehiMELDAzero\endcsname{+7.5}
\expandafter\gdef\csname NUM@famonehiMELDTfive\endcsname{+2.1}
\expandafter\gdef\csname NUM@famonehiMELDTtwofive\endcsname{+1.0}
\expandafter\gdef\csname NUM@famonehiMELDTzero\endcsname{+2.9}
\expandafter\gdef\csname NUM@famonehiMELDVfive\endcsname{+2.6}
\expandafter\gdef\csname NUM@famonehiMELDVtwofive\endcsname{+2.4}
\expandafter\gdef\csname NUM@famonehiMELDVzero\endcsname{+4.1}
\expandafter\gdef\csname NUM@famonehiMOSEIAVTfive\endcsname{+0.2}
\expandafter\gdef\csname NUM@famonehiMOSEIAVTtwofive\endcsname{+1.3}
\expandafter\gdef\csname NUM@famonehiMOSEIAVTzero\endcsname{+0.5}
\expandafter\gdef\csname NUM@famonehiMOSEIAfive\endcsname{+4.0}
\expandafter\gdef\csname NUM@famonehiMOSEIAtwofive\endcsname{+4.2}
\expandafter\gdef\csname NUM@famonehiMOSEIAzero\endcsname{+3.2}
\expandafter\gdef\csname NUM@famonehiMOSEITfive\endcsname{+0.8}
\expandafter\gdef\csname NUM@famonehiMOSEITtwofive\endcsname{+0.7}
\expandafter\gdef\csname NUM@famonehiMOSEITzero\endcsname{+0.3}
\expandafter\gdef\csname NUM@famonehiMOSEIVfive\endcsname{+1.6}
\expandafter\gdef\csname NUM@famonehiMOSEIVtwofive\endcsname{+4.0}
\expandafter\gdef\csname NUM@famonehiMOSEIVzero\endcsname{+0.5}
\expandafter\gdef\csname NUM@famonehitwodAFEWAVTfive\endcsname{+2.72}
\expandafter\gdef\csname NUM@famonehitwodAFEWAVTtwofive\endcsname{+1.57}
\expandafter\gdef\csname NUM@famonehitwodAFEWAVTzero\endcsname{+5.30}
\expandafter\gdef\csname NUM@famonehitwodAFEWAfive\endcsname{+5.74}
\expandafter\gdef\csname NUM@famonehitwodAFEWAtwofive\endcsname{+5.25}
\expandafter\gdef\csname NUM@famonehitwodAFEWAzero\endcsname{+7.40}
\expandafter\gdef\csname NUM@famonehitwodAFEWTfive\endcsname{+4.79}
\expandafter\gdef\csname NUM@famonehitwodAFEWTtwofive\endcsname{+3.56}
\expandafter\gdef\csname NUM@famonehitwodAFEWTzero\endcsname{+5.46}
\expandafter\gdef\csname NUM@famonehitwodAFEWVfive\endcsname{+3.17}
\expandafter\gdef\csname NUM@famonehitwodAFEWVtwofive\endcsname{+1.04}
\expandafter\gdef\csname NUM@famonehitwodAFEWVzero\endcsname{+6.31}
\expandafter\gdef\csname NUM@famonehitwodIEMOCAPAVTfive\endcsname{+4.85}
\expandafter\gdef\csname NUM@famonehitwodIEMOCAPAVTtwofive\endcsname{+3.46}
\expandafter\gdef\csname NUM@famonehitwodIEMOCAPAVTzero\endcsname{+6.16}
\expandafter\gdef\csname NUM@famonehitwodIEMOCAPAfive\endcsname{+5.27}
\expandafter\gdef\csname NUM@famonehitwodIEMOCAPAtwofive\endcsname{+3.37}
\expandafter\gdef\csname NUM@famonehitwodIEMOCAPAzero\endcsname{+7.63}
\expandafter\gdef\csname NUM@famonehitwodIEMOCAPTfive\endcsname{+10.31}
\expandafter\gdef\csname NUM@famonehitwodIEMOCAPTtwofive\endcsname{+7.57}
\expandafter\gdef\csname NUM@famonehitwodIEMOCAPTzero\endcsname{+10.29}
\expandafter\gdef\csname NUM@famonehitwodIEMOCAPVfive\endcsname{+3.06}
\expandafter\gdef\csname NUM@famonehitwodIEMOCAPVtwofive\endcsname{+4.82}
\expandafter\gdef\csname NUM@famonehitwodIEMOCAPVzero\endcsname{+1.03}
\expandafter\gdef\csname NUM@famonehitwodMELDAVTfive\endcsname{+3.42}
\expandafter\gdef\csname NUM@famonehitwodMELDAVTtwofive\endcsname{+2.19}
\expandafter\gdef\csname NUM@famonehitwodMELDAVTzero\endcsname{+5.81}
\expandafter\gdef\csname NUM@famonehitwodMELDAfive\endcsname{+5.06}
\expandafter\gdef\csname NUM@famonehitwodMELDAtwofive\endcsname{+4.51}
\expandafter\gdef\csname NUM@famonehitwodMELDAzero\endcsname{+7.51}
\expandafter\gdef\csname NUM@famonehitwodMELDTfive\endcsname{+2.08}
\expandafter\gdef\csname NUM@famonehitwodMELDTtwofive\endcsname{+0.95}
\expandafter\gdef\csname NUM@famonehitwodMELDTzero\endcsname{+2.92}
\expandafter\gdef\csname NUM@famonehitwodMELDVfive\endcsname{+2.61}
\expandafter\gdef\csname NUM@famonehitwodMELDVtwofive\endcsname{+2.40}
\expandafter\gdef\csname NUM@famonehitwodMELDVzero\endcsname{+4.10}
\expandafter\gdef\csname NUM@famonehitwodMOSEIAVTfive\endcsname{+0.19}
\expandafter\gdef\csname NUM@famonehitwodMOSEIAVTtwofive\endcsname{+1.30}
\expandafter\gdef\csname NUM@famonehitwodMOSEIAVTzero\endcsname{+0.54}
\expandafter\gdef\csname NUM@famonehitwodMOSEIAfive\endcsname{+3.99}
\expandafter\gdef\csname NUM@famonehitwodMOSEIAtwofive\endcsname{+4.16}
\expandafter\gdef\csname NUM@famonehitwodMOSEIAzero\endcsname{+3.19}
\expandafter\gdef\csname NUM@famonehitwodMOSEITfive\endcsname{+0.84}
\expandafter\gdef\csname NUM@famonehitwodMOSEITtwofive\endcsname{+0.74}
\expandafter\gdef\csname NUM@famonehitwodMOSEITzero\endcsname{+0.35}
\expandafter\gdef\csname NUM@famonehitwodMOSEIVfive\endcsname{+1.56}
\expandafter\gdef\csname NUM@famonehitwodMOSEIVtwofive\endcsname{+3.98}
\expandafter\gdef\csname NUM@famonehitwodMOSEIVzero\endcsname{+0.54}
\expandafter\gdef\csname NUM@famoneloAFEWAVTfive\endcsname{\textminus{}0.7}
\expandafter\gdef\csname NUM@famoneloAFEWAVTtwofive\endcsname{\textminus{}0.3}
\expandafter\gdef\csname NUM@famoneloAFEWAVTzero\endcsname{\textminus{}1.9}
\expandafter\gdef\csname NUM@famoneloAFEWAfive\endcsname{+2.6}
\expandafter\gdef\csname NUM@famoneloAFEWAtwofive\endcsname{+1.4}
\expandafter\gdef\csname NUM@famoneloAFEWAzero\endcsname{+2.3}
\expandafter\gdef\csname NUM@famoneloAFEWTfive\endcsname{\textminus{}1.5}
\expandafter\gdef\csname NUM@famoneloAFEWTtwofive\endcsname{\textminus{}2.1}
\expandafter\gdef\csname NUM@famoneloAFEWTzero\endcsname{\textminus{}2.3}
\expandafter\gdef\csname NUM@famoneloAFEWVfive\endcsname{\textminus{}3.0}
\expandafter\gdef\csname NUM@famoneloAFEWVtwofive\endcsname{\textminus{}1.4}
\expandafter\gdef\csname NUM@famoneloAFEWVzero\endcsname{\textminus{}1.3}
\expandafter\gdef\csname NUM@famoneloIEMOCAPAVTfive\endcsname{\textminus{}4.5}
\expandafter\gdef\csname NUM@famoneloIEMOCAPAVTtwofive\endcsname{\textminus{}5.1}
\expandafter\gdef\csname NUM@famoneloIEMOCAPAVTzero\endcsname{\textminus{}2.9}
\expandafter\gdef\csname NUM@famoneloIEMOCAPAfive\endcsname{\textminus{}0.4}
\expandafter\gdef\csname NUM@famoneloIEMOCAPAtwofive\endcsname{\textminus{}0.5}
\expandafter\gdef\csname NUM@famoneloIEMOCAPAzero\endcsname{+1.9}
\expandafter\gdef\csname NUM@famoneloIEMOCAPTfive\endcsname{+4.4}
\expandafter\gdef\csname NUM@famoneloIEMOCAPTtwofive\endcsname{+2.4}
\expandafter\gdef\csname NUM@famoneloIEMOCAPTzero\endcsname{+4.7}
\expandafter\gdef\csname NUM@famoneloIEMOCAPVfive\endcsname{\textminus{}4.7}
\expandafter\gdef\csname NUM@famoneloIEMOCAPVtwofive\endcsname{\textminus{}6.3}
\expandafter\gdef\csname NUM@famoneloIEMOCAPVzero\endcsname{\textminus{}5.2}
\expandafter\gdef\csname NUM@famoneloMELDAVTfive\endcsname{+1.7}
\expandafter\gdef\csname NUM@famoneloMELDAVTtwofive\endcsname{+0.6}
\expandafter\gdef\csname NUM@famoneloMELDAVTzero\endcsname{+0.9}
\expandafter\gdef\csname NUM@famoneloMELDAfive\endcsname{+1.2}
\expandafter\gdef\csname NUM@famoneloMELDAtwofive\endcsname{+1.2}
\expandafter\gdef\csname NUM@famoneloMELDAzero\endcsname{+0.3}
\expandafter\gdef\csname NUM@famoneloMELDTfive\endcsname{\textminus{}4.4}
\expandafter\gdef\csname NUM@famoneloMELDTtwofive\endcsname{\textminus{}2.8}
\expandafter\gdef\csname NUM@famoneloMELDTzero\endcsname{\textminus{}2.6}
\expandafter\gdef\csname NUM@famoneloMELDVfive\endcsname{\textminus{}0.9}
\expandafter\gdef\csname NUM@famoneloMELDVtwofive\endcsname{\textminus{}0.2}
\expandafter\gdef\csname NUM@famoneloMELDVzero\endcsname{\textminus{}2.8}
\expandafter\gdef\csname NUM@famoneloMOSEIAVTfive\endcsname{\textminus{}0.9}
\expandafter\gdef\csname NUM@famoneloMOSEIAVTtwofive\endcsname{+0.2}
\expandafter\gdef\csname NUM@famoneloMOSEIAVTzero\endcsname{\textminus{}1.2}
\expandafter\gdef\csname NUM@famoneloMOSEIAfive\endcsname{+2.4}
\expandafter\gdef\csname NUM@famoneloMOSEIAtwofive\endcsname{+2.6}
\expandafter\gdef\csname NUM@famoneloMOSEIAzero\endcsname{+1.1}
\expandafter\gdef\csname NUM@famoneloMOSEITfive\endcsname{\textminus{}0.9}
\expandafter\gdef\csname NUM@famoneloMOSEITtwofive\endcsname{\textminus{}0.7}
\expandafter\gdef\csname NUM@famoneloMOSEITzero\endcsname{\textminus{}1.7}
\expandafter\gdef\csname NUM@famoneloMOSEIVfive\endcsname{\textminus{}1.1}
\expandafter\gdef\csname NUM@famoneloMOSEIVtwofive\endcsname{+1.2}
\expandafter\gdef\csname NUM@famoneloMOSEIVzero\endcsname{\textminus{}2.4}
\expandafter\gdef\csname NUM@famonelotwodAFEWAVTfive\endcsname{\textminus{}0.65}
\expandafter\gdef\csname NUM@famonelotwodAFEWAVTtwofive\endcsname{\textminus{}0.33}
\expandafter\gdef\csname NUM@famonelotwodAFEWAVTzero\endcsname{\textminus{}1.91}
\expandafter\gdef\csname NUM@famonelotwodAFEWAfive\endcsname{+2.56}
\expandafter\gdef\csname NUM@famonelotwodAFEWAtwofive\endcsname{+1.40}
\expandafter\gdef\csname NUM@famonelotwodAFEWAzero\endcsname{+2.32}
\expandafter\gdef\csname NUM@famonelotwodAFEWTfive\endcsname{\textminus{}1.50}
\expandafter\gdef\csname NUM@famonelotwodAFEWTtwofive\endcsname{\textminus{}2.07}
\expandafter\gdef\csname NUM@famonelotwodAFEWTzero\endcsname{\textminus{}2.33}
\expandafter\gdef\csname NUM@famonelotwodAFEWVfive\endcsname{\textminus{}2.99}
\expandafter\gdef\csname NUM@famonelotwodAFEWVtwofive\endcsname{\textminus{}1.40}
\expandafter\gdef\csname NUM@famonelotwodAFEWVzero\endcsname{\textminus{}1.34}
\expandafter\gdef\csname NUM@famonelotwodIEMOCAPAVTfive\endcsname{\textminus{}4.53}
\expandafter\gdef\csname NUM@famonelotwodIEMOCAPAVTtwofive\endcsname{\textminus{}5.12}
\expandafter\gdef\csname NUM@famonelotwodIEMOCAPAVTzero\endcsname{\textminus{}2.86}
\expandafter\gdef\csname NUM@famonelotwodIEMOCAPAfive\endcsname{\textminus{}0.45}
\expandafter\gdef\csname NUM@famonelotwodIEMOCAPAtwofive\endcsname{\textminus{}0.48}
\expandafter\gdef\csname NUM@famonelotwodIEMOCAPAzero\endcsname{+1.89}
\expandafter\gdef\csname NUM@famonelotwodIEMOCAPTfive\endcsname{+4.37}
\expandafter\gdef\csname NUM@famonelotwodIEMOCAPTtwofive\endcsname{+2.44}
\expandafter\gdef\csname NUM@famonelotwodIEMOCAPTzero\endcsname{+4.75}
\expandafter\gdef\csname NUM@famonelotwodIEMOCAPVfive\endcsname{\textminus{}4.75}
\expandafter\gdef\csname NUM@famonelotwodIEMOCAPVtwofive\endcsname{\textminus{}6.30}
\expandafter\gdef\csname NUM@famonelotwodIEMOCAPVzero\endcsname{\textminus{}5.19}
\expandafter\gdef\csname NUM@famonelotwodMELDAVTfive\endcsname{+1.65}
\expandafter\gdef\csname NUM@famonelotwodMELDAVTtwofive\endcsname{+0.63}
\expandafter\gdef\csname NUM@famonelotwodMELDAVTzero\endcsname{+0.90}
\expandafter\gdef\csname NUM@famonelotwodMELDAfive\endcsname{+1.18}
\expandafter\gdef\csname NUM@famonelotwodMELDAtwofive\endcsname{+1.21}
\expandafter\gdef\csname NUM@famonelotwodMELDAzero\endcsname{+0.28}
\expandafter\gdef\csname NUM@famonelotwodMELDTfive\endcsname{\textminus{}4.44}
\expandafter\gdef\csname NUM@famonelotwodMELDTtwofive\endcsname{\textminus{}2.84}
\expandafter\gdef\csname NUM@famonelotwodMELDTzero\endcsname{\textminus{}2.56}
\expandafter\gdef\csname NUM@famonelotwodMELDVfive\endcsname{\textminus{}0.91}
\expandafter\gdef\csname NUM@famonelotwodMELDVtwofive\endcsname{\textminus{}0.18}
\expandafter\gdef\csname NUM@famonelotwodMELDVzero\endcsname{\textminus{}2.76}
\expandafter\gdef\csname NUM@famonelotwodMOSEIAVTfive\endcsname{\textminus{}0.90}
\expandafter\gdef\csname NUM@famonelotwodMOSEIAVTtwofive\endcsname{+0.18}
\expandafter\gdef\csname NUM@famonelotwodMOSEIAVTzero\endcsname{\textminus{}1.17}
\expandafter\gdef\csname NUM@famonelotwodMOSEIAfive\endcsname{+2.38}
\expandafter\gdef\csname NUM@famonelotwodMOSEIAtwofive\endcsname{+2.56}
\expandafter\gdef\csname NUM@famonelotwodMOSEIAzero\endcsname{+1.08}
\expandafter\gdef\csname NUM@famonelotwodMOSEITfive\endcsname{\textminus{}0.89}
\expandafter\gdef\csname NUM@famonelotwodMOSEITtwofive\endcsname{\textminus{}0.69}
\expandafter\gdef\csname NUM@famonelotwodMOSEITzero\endcsname{\textminus{}1.72}
\expandafter\gdef\csname NUM@famonelotwodMOSEIVfive\endcsname{\textminus{}1.14}
\expandafter\gdef\csname NUM@famonelotwodMOSEIVtwofive\endcsname{+1.17}
\expandafter\gdef\csname NUM@famonelotwodMOSEIVzero\endcsname{\textminus{}2.35}
\expandafter\gdef\csname NUM@famonepAFEWAVTfive\endcsname{1.000}
\expandafter\gdef\csname NUM@famonepAFEWAVTtwofive\endcsname{1.000}
\expandafter\gdef\csname NUM@famonepAFEWAVTzero\endcsname{1.000}
\expandafter\gdef\csname NUM@famonepAFEWAfive\endcsname{0.001}
\expandafter\gdef\csname NUM@famonepAFEWAtwofive\endcsname{0.026}
\expandafter\gdef\csname NUM@famonepAFEWAzero\endcsname{0.013}
\expandafter\gdef\csname NUM@famonepAFEWTfive\endcsname{1.000}
\expandafter\gdef\csname NUM@famonepAFEWTtwofive\endcsname{1.000}
\expandafter\gdef\csname NUM@famonepAFEWTzero\endcsname{1.000}
\expandafter\gdef\csname NUM@famonepAFEWVfive\endcsname{1.000}
\expandafter\gdef\csname NUM@famonepAFEWVtwofive\endcsname{1.000}
\expandafter\gdef\csname NUM@famonepAFEWVzero\endcsname{1.000}
\expandafter\gdef\csname NUM@famonepIEMOCAPAVTfive\endcsname{1.000}
\expandafter\gdef\csname NUM@famonepIEMOCAPAVTtwofive\endcsname{1.000}
\expandafter\gdef\csname NUM@famonepIEMOCAPAVTzero\endcsname{1.000}
\expandafter\gdef\csname NUM@famonepIEMOCAPAfive\endcsname{0.598}
\expandafter\gdef\csname NUM@famonepIEMOCAPAtwofive\endcsname{0.678}
\expandafter\gdef\csname NUM@famonepIEMOCAPAzero\endcsname{0.119}
\expandafter\gdef\csname NUM@famonepIEMOCAPTfive\endcsname{0.047}
\expandafter\gdef\csname NUM@famonepIEMOCAPTtwofive\endcsname{0.084}
\expandafter\gdef\csname NUM@famonepIEMOCAPTzero\endcsname{0.039}
\expandafter\gdef\csname NUM@famonepIEMOCAPVfive\endcsname{1.000}
\expandafter\gdef\csname NUM@famonepIEMOCAPVtwofive\endcsname{1.000}
\expandafter\gdef\csname NUM@famonepIEMOCAPVzero\endcsname{0.738}
\expandafter\gdef\csname NUM@famonepMELDAVTfive\endcsname{0.003}
\expandafter\gdef\csname NUM@famonepMELDAVTtwofive\endcsname{0.040}
\expandafter\gdef\csname NUM@famonepMELDAVTzero\endcsname{0.116}
\expandafter\gdef\csname NUM@famonepMELDAfive\endcsname{0.060}
\expandafter\gdef\csname NUM@famonepMELDAtwofive\endcsname{0.046}
\expandafter\gdef\csname NUM@famonepMELDAzero\endcsname{0.267}
\expandafter\gdef\csname NUM@famonepMELDTfive\endcsname{1.000}
\expandafter\gdef\csname NUM@famonepMELDTtwofive\endcsname{1.000}
\expandafter\gdef\csname NUM@famonepMELDTzero\endcsname{1.000}
\expandafter\gdef\csname NUM@famonepMELDVfive\endcsname{1.000}
\expandafter\gdef\csname NUM@famonepMELDVtwofive\endcsname{0.489}
\expandafter\gdef\csname NUM@famonepMELDVzero\endcsname{1.000}
\expandafter\gdef\csname NUM@famonepMOSEIAVTfive\endcsname{1.000}
\expandafter\gdef\csname NUM@famonepMOSEIAVTtwofive\endcsname{0.080}
\expandafter\gdef\csname NUM@famonepMOSEIAVTzero\endcsname{1.000}
\expandafter\gdef\csname NUM@famonepMOSEIAfive\endcsname{0.000}
\expandafter\gdef\csname NUM@famonepMOSEIAtwofive\endcsname{0.000}
\expandafter\gdef\csname NUM@famonepMOSEIAzero\endcsname{0.001}
\expandafter\gdef\csname NUM@famonepMOSEITfive\endcsname{1.000}
\expandafter\gdef\csname NUM@famonepMOSEITtwofive\endcsname{1.000}
\expandafter\gdef\csname NUM@famonepMOSEITzero\endcsname{1.000}
\expandafter\gdef\csname NUM@famonepMOSEIVfive\endcsname{1.000}
\expandafter\gdef\csname NUM@famonepMOSEIVtwofive\endcsname{0.003}
\expandafter\gdef\csname NUM@famonepMOSEIVzero\endcsname{1.000}
\expandafter\gdef\csname NUM@famonetwodAFEWAVTfive\endcsname{+1.03}
\expandafter\gdef\csname NUM@famonetwodAFEWAVTtwofive\endcsname{+0.62}
\expandafter\gdef\csname NUM@famonetwodAFEWAVTzero\endcsname{+1.69}
\expandafter\gdef\csname NUM@famonetwodAFEWAfive\endcsname{+4.15}
\expandafter\gdef\csname NUM@famonetwodAFEWAtwofive\endcsname{+3.32}
\expandafter\gdef\csname NUM@famonetwodAFEWAzero\endcsname{+4.86}
\expandafter\gdef\csname NUM@famonetwodAFEWTfive\endcsname{+1.65}
\expandafter\gdef\csname NUM@famonetwodAFEWTtwofive\endcsname{+0.75}
\expandafter\gdef\csname NUM@famonetwodAFEWTzero\endcsname{+1.56}
\expandafter\gdef\csname NUM@famonetwodAFEWVfive\endcsname{+0.09}
\expandafter\gdef\csname NUM@famonetwodAFEWVtwofive\endcsname{\textminus{}0.18}
\expandafter\gdef\csname NUM@famonetwodAFEWVzero\endcsname{+2.48}
\expandafter\gdef\csname NUM@famonetwodIEMOCAPAVTfive\endcsname{+0.16}
\expandafter\gdef\csname NUM@famonetwodIEMOCAPAVTtwofive\endcsname{\textminus{}0.83}
\expandafter\gdef\csname NUM@famonetwodIEMOCAPAVTzero\endcsname{+1.65}
\expandafter\gdef\csname NUM@famonetwodIEMOCAPAfive\endcsname{+2.41}
\expandafter\gdef\csname NUM@famonetwodIEMOCAPAtwofive\endcsname{+1.45}
\expandafter\gdef\csname NUM@famonetwodIEMOCAPAzero\endcsname{+4.76}
\expandafter\gdef\csname NUM@famonetwodIEMOCAPTfive\endcsname{+7.34}
\expandafter\gdef\csname NUM@famonetwodIEMOCAPTtwofive\endcsname{+5.01}
\expandafter\gdef\csname NUM@famonetwodIEMOCAPTzero\endcsname{+7.52}
\expandafter\gdef\csname NUM@famonetwodIEMOCAPVfive\endcsname{\textminus{}0.84}
\expandafter\gdef\csname NUM@famonetwodIEMOCAPVtwofive\endcsname{\textminus{}0.74}
\expandafter\gdef\csname NUM@famonetwodIEMOCAPVzero\endcsname{\textminus{}2.08}
\expandafter\gdef\csname NUM@famonetwodMELDAVTfive\endcsname{+2.53}
\expandafter\gdef\csname NUM@famonetwodMELDAVTtwofive\endcsname{+1.41}
\expandafter\gdef\csname NUM@famonetwodMELDAVTzero\endcsname{+3.36}
\expandafter\gdef\csname NUM@famonetwodMELDAfive\endcsname{+3.12}
\expandafter\gdef\csname NUM@famonetwodMELDAtwofive\endcsname{+2.86}
\expandafter\gdef\csname NUM@famonetwodMELDAzero\endcsname{+3.89}
\expandafter\gdef\csname NUM@famonetwodMELDTfive\endcsname{\textminus{}1.18}
\expandafter\gdef\csname NUM@famonetwodMELDTtwofive\endcsname{\textminus{}0.94}
\expandafter\gdef\csname NUM@famonetwodMELDTzero\endcsname{+0.18}
\expandafter\gdef\csname NUM@famonetwodMELDVfive\endcsname{+0.85}
\expandafter\gdef\csname NUM@famonetwodMELDVtwofive\endcsname{+1.11}
\expandafter\gdef\csname NUM@famonetwodMELDVzero\endcsname{+0.67}
\expandafter\gdef\csname NUM@famonetwodMOSEIAVTfive\endcsname{\textminus{}0.35}
\expandafter\gdef\csname NUM@famonetwodMOSEIAVTtwofive\endcsname{+0.74}
\expandafter\gdef\csname NUM@famonetwodMOSEIAVTzero\endcsname{\textminus{}0.32}
\expandafter\gdef\csname NUM@famonetwodMOSEIAfive\endcsname{+3.19}
\expandafter\gdef\csname NUM@famonetwodMOSEIAtwofive\endcsname{+3.36}
\expandafter\gdef\csname NUM@famonetwodMOSEIAzero\endcsname{+2.13}
\expandafter\gdef\csname NUM@famonetwodMOSEITfive\endcsname{\textminus{}0.02}
\expandafter\gdef\csname NUM@famonetwodMOSEITtwofive\endcsname{+0.03}
\expandafter\gdef\csname NUM@famonetwodMOSEITzero\endcsname{\textminus{}0.69}
\expandafter\gdef\csname NUM@famonetwodMOSEIVfive\endcsname{+0.21}
\expandafter\gdef\csname NUM@famonetwodMOSEIVtwofive\endcsname{+2.57}
\expandafter\gdef\csname NUM@famonetwodMOSEIVzero\endcsname{\textminus{}0.91}
\expandafter\gdef\csname NUM@famthreeAFEWAVTfive\endcsname{\textminus{}0.2}
\expandafter\gdef\csname NUM@famthreeAFEWAVTtwofive\endcsname{\textminus{}0.1}
\expandafter\gdef\csname NUM@famthreeAFEWAVTzero\endcsname{\textminus{}0.7}
\expandafter\gdef\csname NUM@famthreeAFEWAfive\endcsname{+0.6}
\expandafter\gdef\csname NUM@famthreeAFEWAtwofive\endcsname{+0.1}
\expandafter\gdef\csname NUM@famthreeAFEWAzero\endcsname{+0.9}
\expandafter\gdef\csname NUM@famthreeAFEWTfive\endcsname{\textminus{}0.2}
\expandafter\gdef\csname NUM@famthreeAFEWTtwofive\endcsname{\textminus{}0.3}
\expandafter\gdef\csname NUM@famthreeAFEWTzero\endcsname{+0.1}
\expandafter\gdef\csname NUM@famthreeAFEWVfive\endcsname{\textminus{}0.2}
\expandafter\gdef\csname NUM@famthreeAFEWVtwofive\endcsname{+0.1}
\expandafter\gdef\csname NUM@famthreeAFEWVzero\endcsname{\textminus{}0.8}
\expandafter\gdef\csname NUM@famthreeIEMOCAPAVTfive\endcsname{+0.8}
\expandafter\gdef\csname NUM@famthreeIEMOCAPAVTtwofive\endcsname{+0.2}
\expandafter\gdef\csname NUM@famthreeIEMOCAPAVTzero\endcsname{\textminus{}0.4}
\expandafter\gdef\csname NUM@famthreeIEMOCAPAfive\endcsname{\textminus{}0.2}
\expandafter\gdef\csname NUM@famthreeIEMOCAPAtwofive\endcsname{\textminus{}0.5}
\expandafter\gdef\csname NUM@famthreeIEMOCAPAzero\endcsname{\textminus{}0.8}
\expandafter\gdef\csname NUM@famthreeIEMOCAPTfive\endcsname{\textminus{}0.2}
\expandafter\gdef\csname NUM@famthreeIEMOCAPTtwofive\endcsname{\textminus{}0.4}
\expandafter\gdef\csname NUM@famthreeIEMOCAPTzero\endcsname{\textminus{}0.3}
\expandafter\gdef\csname NUM@famthreeIEMOCAPVfive\endcsname{\textminus{}1.2}
\expandafter\gdef\csname NUM@famthreeIEMOCAPVtwofive\endcsname{\textminus{}1.0}
\expandafter\gdef\csname NUM@famthreeIEMOCAPVzero\endcsname{\textminus{}0.9}
\expandafter\gdef\csname NUM@famthreeMELDAVTfive\endcsname{+2.2}
\expandafter\gdef\csname NUM@famthreeMELDAVTtwofive\endcsname{\textminus{}0.2}
\expandafter\gdef\csname NUM@famthreeMELDAVTzero\endcsname{+1.3}
\expandafter\gdef\csname NUM@famthreeMELDAfive\endcsname{\textminus{}0.0}
\expandafter\gdef\csname NUM@famthreeMELDAtwofive\endcsname{\textminus{}0.3}
\expandafter\gdef\csname NUM@famthreeMELDAzero\endcsname{\textminus{}0.0}
\expandafter\gdef\csname NUM@famthreeMELDTfive\endcsname{\textminus{}0.3}
\expandafter\gdef\csname NUM@famthreeMELDTtwofive\endcsname{\textminus{}1.2}
\expandafter\gdef\csname NUM@famthreeMELDTzero\endcsname{+0.8}
\expandafter\gdef\csname NUM@famthreeMELDVfive\endcsname{+0.2}
\expandafter\gdef\csname NUM@famthreeMELDVtwofive\endcsname{+0.2}
\expandafter\gdef\csname NUM@famthreeMELDVzero\endcsname{\textminus{}0.1}
\expandafter\gdef\csname NUM@famthreeMOSEIAVTfive\endcsname{\textminus{}0.1}
\expandafter\gdef\csname NUM@famthreeMOSEIAVTtwofive\endcsname{\textminus{}0.0}
\expandafter\gdef\csname NUM@famthreeMOSEIAVTzero\endcsname{+0.1}
\expandafter\gdef\csname NUM@famthreeMOSEIAfive\endcsname{\textminus{}0.2}
\expandafter\gdef\csname NUM@famthreeMOSEIAtwofive\endcsname{\textminus{}0.2}
\expandafter\gdef\csname NUM@famthreeMOSEIAzero\endcsname{+0.1}
\expandafter\gdef\csname NUM@famthreeMOSEITfive\endcsname{\textminus{}0.3}
\expandafter\gdef\csname NUM@famthreeMOSEITtwofive\endcsname{\textminus{}0.8}
\expandafter\gdef\csname NUM@famthreeMOSEITzero\endcsname{\textminus{}0.7}
\expandafter\gdef\csname NUM@famthreeMOSEIVfive\endcsname{\textminus{}0.0}
\expandafter\gdef\csname NUM@famthreeMOSEIVtwofive\endcsname{\textminus{}0.1}
\expandafter\gdef\csname NUM@famthreeMOSEIVzero\endcsname{\textminus{}0.7}
\expandafter\gdef\csname NUM@famthreehiAFEWAVTfive\endcsname{+0.5}
\expandafter\gdef\csname NUM@famthreehiAFEWAVTtwofive\endcsname{+0.6}
\expandafter\gdef\csname NUM@famthreehiAFEWAVTzero\endcsname{+0.8}
\expandafter\gdef\csname NUM@famthreehiAFEWAfive\endcsname{+1.7}
\expandafter\gdef\csname NUM@famthreehiAFEWAtwofive\endcsname{+0.9}
\expandafter\gdef\csname NUM@famthreehiAFEWAzero\endcsname{+2.9}
\expandafter\gdef\csname NUM@famthreehiAFEWTfive\endcsname{+0.9}
\expandafter\gdef\csname NUM@famthreehiAFEWTtwofive\endcsname{+0.5}
\expandafter\gdef\csname NUM@famthreehiAFEWTzero\endcsname{+1.8}
\expandafter\gdef\csname NUM@famthreehiAFEWVfive\endcsname{+0.6}
\expandafter\gdef\csname NUM@famthreehiAFEWVtwofive\endcsname{+0.7}
\expandafter\gdef\csname NUM@famthreehiAFEWVzero\endcsname{+1.1}
\expandafter\gdef\csname NUM@famthreehiIEMOCAPAVTfive\endcsname{+1.7}
\expandafter\gdef\csname NUM@famthreehiIEMOCAPAVTtwofive\endcsname{+0.7}
\expandafter\gdef\csname NUM@famthreehiIEMOCAPAVTzero\endcsname{+0.8}
\expandafter\gdef\csname NUM@famthreehiIEMOCAPAfive\endcsname{+0.9}
\expandafter\gdef\csname NUM@famthreehiIEMOCAPAtwofive\endcsname{+0.8}
\expandafter\gdef\csname NUM@famthreehiIEMOCAPAzero\endcsname{+1.3}
\expandafter\gdef\csname NUM@famthreehiIEMOCAPTfive\endcsname{\textminus{}0.0}
\expandafter\gdef\csname NUM@famthreehiIEMOCAPTtwofive\endcsname{+0.2}
\expandafter\gdef\csname NUM@famthreehiIEMOCAPTzero\endcsname{+1.5}
\expandafter\gdef\csname NUM@famthreehiIEMOCAPVfive\endcsname{+0.9}
\expandafter\gdef\csname NUM@famthreehiIEMOCAPVtwofive\endcsname{+0.8}
\expandafter\gdef\csname NUM@famthreehiIEMOCAPVzero\endcsname{+1.3}
\expandafter\gdef\csname NUM@famthreehiMELDAVTfive\endcsname{+3.2}
\expandafter\gdef\csname NUM@famthreehiMELDAVTtwofive\endcsname{+0.0}
\expandafter\gdef\csname NUM@famthreehiMELDAVTzero\endcsname{+3.2}
\expandafter\gdef\csname NUM@famthreehiMELDAfive\endcsname{+0.5}
\expandafter\gdef\csname NUM@famthreehiMELDAtwofive\endcsname{+0.0}
\expandafter\gdef\csname NUM@famthreehiMELDAzero\endcsname{+1.3}
\expandafter\gdef\csname NUM@famthreehiMELDTfive\endcsname{+0.2}
\expandafter\gdef\csname NUM@famthreehiMELDTtwofive\endcsname{\textminus{}0.7}
\expandafter\gdef\csname NUM@famthreehiMELDTzero\endcsname{+2.3}
\expandafter\gdef\csname NUM@famthreehiMELDVfive\endcsname{+0.5}
\expandafter\gdef\csname NUM@famthreehiMELDVtwofive\endcsname{+0.7}
\expandafter\gdef\csname NUM@famthreehiMELDVzero\endcsname{+3.2}
\expandafter\gdef\csname NUM@famthreehiMOSEIAVTfive\endcsname{+0.0}
\expandafter\gdef\csname NUM@famthreehiMOSEIAVTtwofive\endcsname{+0.1}
\expandafter\gdef\csname NUM@famthreehiMOSEIAVTzero\endcsname{+0.3}
\expandafter\gdef\csname NUM@famthreehiMOSEIAfive\endcsname{\textminus{}0.0}
\expandafter\gdef\csname NUM@famthreehiMOSEIAtwofive\endcsname{\textminus{}0.0}
\expandafter\gdef\csname NUM@famthreehiMOSEIAzero\endcsname{+0.4}
\expandafter\gdef\csname NUM@famthreehiMOSEITfive\endcsname{\textminus{}0.2}
\expandafter\gdef\csname NUM@famthreehiMOSEITtwofive\endcsname{\textminus{}0.5}
\expandafter\gdef\csname NUM@famthreehiMOSEITzero\endcsname{\textminus{}0.2}
\expandafter\gdef\csname NUM@famthreehiMOSEIVfive\endcsname{+0.2}
\expandafter\gdef\csname NUM@famthreehiMOSEIVtwofive\endcsname{+0.1}
\expandafter\gdef\csname NUM@famthreehiMOSEIVzero\endcsname{+0.5}
\expandafter\gdef\csname NUM@famthreehitwodAFEWAVTfive\endcsname{+0.50}
\expandafter\gdef\csname NUM@famthreehitwodAFEWAVTtwofive\endcsname{+0.56}
\expandafter\gdef\csname NUM@famthreehitwodAFEWAVTzero\endcsname{+0.76}
\expandafter\gdef\csname NUM@famthreehitwodAFEWAfive\endcsname{+1.65}
\expandafter\gdef\csname NUM@famthreehitwodAFEWAtwofive\endcsname{+0.95}
\expandafter\gdef\csname NUM@famthreehitwodAFEWAzero\endcsname{+2.93}
\expandafter\gdef\csname NUM@famthreehitwodAFEWTfive\endcsname{+0.86}
\expandafter\gdef\csname NUM@famthreehitwodAFEWTtwofive\endcsname{+0.54}
\expandafter\gdef\csname NUM@famthreehitwodAFEWTzero\endcsname{+1.79}
\expandafter\gdef\csname NUM@famthreehitwodAFEWVfive\endcsname{+0.61}
\expandafter\gdef\csname NUM@famthreehitwodAFEWVtwofive\endcsname{+0.70}
\expandafter\gdef\csname NUM@famthreehitwodAFEWVzero\endcsname{+1.09}
\expandafter\gdef\csname NUM@famthreehitwodIEMOCAPAVTfive\endcsname{+1.69}
\expandafter\gdef\csname NUM@famthreehitwodIEMOCAPAVTtwofive\endcsname{+0.65}
\expandafter\gdef\csname NUM@famthreehitwodIEMOCAPAVTzero\endcsname{+0.83}
\expandafter\gdef\csname NUM@famthreehitwodIEMOCAPAfive\endcsname{+0.87}
\expandafter\gdef\csname NUM@famthreehitwodIEMOCAPAtwofive\endcsname{+0.76}
\expandafter\gdef\csname NUM@famthreehitwodIEMOCAPAzero\endcsname{+1.32}
\expandafter\gdef\csname NUM@famthreehitwodIEMOCAPTfive\endcsname{\textminus{}0.02}
\expandafter\gdef\csname NUM@famthreehitwodIEMOCAPTtwofive\endcsname{+0.20}
\expandafter\gdef\csname NUM@famthreehitwodIEMOCAPTzero\endcsname{+1.46}
\expandafter\gdef\csname NUM@famthreehitwodIEMOCAPVfive\endcsname{+0.92}
\expandafter\gdef\csname NUM@famthreehitwodIEMOCAPVtwofive\endcsname{+0.76}
\expandafter\gdef\csname NUM@famthreehitwodIEMOCAPVzero\endcsname{+1.34}
\expandafter\gdef\csname NUM@famthreehitwodMELDAVTfive\endcsname{+3.24}
\expandafter\gdef\csname NUM@famthreehitwodMELDAVTtwofive\endcsname{+0.01}
\expandafter\gdef\csname NUM@famthreehitwodMELDAVTzero\endcsname{+3.20}
\expandafter\gdef\csname NUM@famthreehitwodMELDAfive\endcsname{+0.49}
\expandafter\gdef\csname NUM@famthreehitwodMELDAtwofive\endcsname{+0.01}
\expandafter\gdef\csname NUM@famthreehitwodMELDAzero\endcsname{+1.27}
\expandafter\gdef\csname NUM@famthreehitwodMELDTfive\endcsname{+0.20}
\expandafter\gdef\csname NUM@famthreehitwodMELDTtwofive\endcsname{\textminus{}0.70}
\expandafter\gdef\csname NUM@famthreehitwodMELDTzero\endcsname{+2.30}
\expandafter\gdef\csname NUM@famthreehitwodMELDVfive\endcsname{+0.52}
\expandafter\gdef\csname NUM@famthreehitwodMELDVtwofive\endcsname{+0.65}
\expandafter\gdef\csname NUM@famthreehitwodMELDVzero\endcsname{+3.23}
\expandafter\gdef\csname NUM@famthreehitwodMOSEIAVTfive\endcsname{+0.04}
\expandafter\gdef\csname NUM@famthreehitwodMOSEIAVTtwofive\endcsname{+0.07}
\expandafter\gdef\csname NUM@famthreehitwodMOSEIAVTzero\endcsname{+0.30}
\expandafter\gdef\csname NUM@famthreehitwodMOSEIAfive\endcsname{\textminus{}0.00}
\expandafter\gdef\csname NUM@famthreehitwodMOSEIAtwofive\endcsname{\textminus{}0.04}
\expandafter\gdef\csname NUM@famthreehitwodMOSEIAzero\endcsname{+0.42}
\expandafter\gdef\csname NUM@famthreehitwodMOSEITfive\endcsname{\textminus{}0.16}
\expandafter\gdef\csname NUM@famthreehitwodMOSEITtwofive\endcsname{\textminus{}0.54}
\expandafter\gdef\csname NUM@famthreehitwodMOSEITzero\endcsname{\textminus{}0.24}
\expandafter\gdef\csname NUM@famthreehitwodMOSEIVfive\endcsname{+0.17}
\expandafter\gdef\csname NUM@famthreehitwodMOSEIVtwofive\endcsname{+0.09}
\expandafter\gdef\csname NUM@famthreehitwodMOSEIVzero\endcsname{+0.45}
\expandafter\gdef\csname NUM@famthreeloAFEWAVTfive\endcsname{\textminus{}0.9}
\expandafter\gdef\csname NUM@famthreeloAFEWAVTtwofive\endcsname{\textminus{}0.8}
\expandafter\gdef\csname NUM@famthreeloAFEWAVTzero\endcsname{\textminus{}2.2}
\expandafter\gdef\csname NUM@famthreeloAFEWAfive\endcsname{\textminus{}0.4}
\expandafter\gdef\csname NUM@famthreeloAFEWAtwofive\endcsname{\textminus{}0.7}
\expandafter\gdef\csname NUM@famthreeloAFEWAzero\endcsname{\textminus{}1.2}
\expandafter\gdef\csname NUM@famthreeloAFEWTfive\endcsname{\textminus{}1.2}
\expandafter\gdef\csname NUM@famthreeloAFEWTtwofive\endcsname{\textminus{}1.2}
\expandafter\gdef\csname NUM@famthreeloAFEWTzero\endcsname{\textminus{}1.5}
\expandafter\gdef\csname NUM@famthreeloAFEWVfive\endcsname{\textminus{}1.0}
\expandafter\gdef\csname NUM@famthreeloAFEWVtwofive\endcsname{\textminus{}0.6}
\expandafter\gdef\csname NUM@famthreeloAFEWVzero\endcsname{\textminus{}2.7}
\expandafter\gdef\csname NUM@famthreeloIEMOCAPAVTfive\endcsname{\textminus{}0.1}
\expandafter\gdef\csname NUM@famthreeloIEMOCAPAVTtwofive\endcsname{\textminus{}0.3}
\expandafter\gdef\csname NUM@famthreeloIEMOCAPAVTzero\endcsname{\textminus{}1.6}
\expandafter\gdef\csname NUM@famthreeloIEMOCAPAfive\endcsname{\textminus{}1.3}
\expandafter\gdef\csname NUM@famthreeloIEMOCAPAtwofive\endcsname{\textminus{}1.7}
\expandafter\gdef\csname NUM@famthreeloIEMOCAPAzero\endcsname{\textminus{}2.9}
\expandafter\gdef\csname NUM@famthreeloIEMOCAPTfive\endcsname{\textminus{}0.4}
\expandafter\gdef\csname NUM@famthreeloIEMOCAPTtwofive\endcsname{\textminus{}1.0}
\expandafter\gdef\csname NUM@famthreeloIEMOCAPTzero\endcsname{\textminus{}2.0}
\expandafter\gdef\csname NUM@famthreeloIEMOCAPVfive\endcsname{\textminus{}3.3}
\expandafter\gdef\csname NUM@famthreeloIEMOCAPVtwofive\endcsname{\textminus{}2.8}
\expandafter\gdef\csname NUM@famthreeloIEMOCAPVzero\endcsname{\textminus{}3.1}
\expandafter\gdef\csname NUM@famthreeloMELDAVTfive\endcsname{+1.1}
\expandafter\gdef\csname NUM@famthreeloMELDAVTtwofive\endcsname{\textminus{}0.5}
\expandafter\gdef\csname NUM@famthreeloMELDAVTzero\endcsname{\textminus{}0.7}
\expandafter\gdef\csname NUM@famthreeloMELDAfive\endcsname{\textminus{}0.5}
\expandafter\gdef\csname NUM@famthreeloMELDAtwofive\endcsname{\textminus{}0.6}
\expandafter\gdef\csname NUM@famthreeloMELDAzero\endcsname{\textminus{}1.4}
\expandafter\gdef\csname NUM@famthreeloMELDTfive\endcsname{\textminus{}0.9}
\expandafter\gdef\csname NUM@famthreeloMELDTtwofive\endcsname{\textminus{}1.8}
\expandafter\gdef\csname NUM@famthreeloMELDTzero\endcsname{\textminus{}0.6}
\expandafter\gdef\csname NUM@famthreeloMELDVfive\endcsname{\textminus{}0.1}
\expandafter\gdef\csname NUM@famthreeloMELDVtwofive\endcsname{\textminus{}0.2}
\expandafter\gdef\csname NUM@famthreeloMELDVzero\endcsname{\textminus{}3.5}
\expandafter\gdef\csname NUM@famthreeloMOSEIAVTfive\endcsname{\textminus{}0.2}
\expandafter\gdef\csname NUM@famthreeloMOSEIAVTtwofive\endcsname{\textminus{}0.2}
\expandafter\gdef\csname NUM@famthreeloMOSEIAVTzero\endcsname{\textminus{}0.2}
\expandafter\gdef\csname NUM@famthreeloMOSEIAfive\endcsname{\textminus{}0.3}
\expandafter\gdef\csname NUM@famthreeloMOSEIAtwofive\endcsname{\textminus{}0.4}
\expandafter\gdef\csname NUM@famthreeloMOSEIAzero\endcsname{\textminus{}0.2}
\expandafter\gdef\csname NUM@famthreeloMOSEITfive\endcsname{\textminus{}0.5}
\expandafter\gdef\csname NUM@famthreeloMOSEITtwofive\endcsname{\textminus{}1.0}
\expandafter\gdef\csname NUM@famthreeloMOSEITzero\endcsname{\textminus{}1.1}
\expandafter\gdef\csname NUM@famthreeloMOSEIVfive\endcsname{\textminus{}0.2}
\expandafter\gdef\csname NUM@famthreeloMOSEIVtwofive\endcsname{\textminus{}0.4}
\expandafter\gdef\csname NUM@famthreeloMOSEIVzero\endcsname{\textminus{}1.9}
\expandafter\gdef\csname NUM@famthreelotwodAFEWAVTfive\endcsname{\textminus{}0.85}
\expandafter\gdef\csname NUM@famthreelotwodAFEWAVTtwofive\endcsname{\textminus{}0.82}
\expandafter\gdef\csname NUM@famthreelotwodAFEWAVTzero\endcsname{\textminus{}2.17}
\expandafter\gdef\csname NUM@famthreelotwodAFEWAfive\endcsname{\textminus{}0.41}
\expandafter\gdef\csname NUM@famthreelotwodAFEWAtwofive\endcsname{\textminus{}0.68}
\expandafter\gdef\csname NUM@famthreelotwodAFEWAzero\endcsname{\textminus{}1.18}
\expandafter\gdef\csname NUM@famthreelotwodAFEWTfive\endcsname{\textminus{}1.21}
\expandafter\gdef\csname NUM@famthreelotwodAFEWTtwofive\endcsname{\textminus{}1.15}
\expandafter\gdef\csname NUM@famthreelotwodAFEWTzero\endcsname{\textminus{}1.51}
\expandafter\gdef\csname NUM@famthreelotwodAFEWVfive\endcsname{\textminus{}0.99}
\expandafter\gdef\csname NUM@famthreelotwodAFEWVtwofive\endcsname{\textminus{}0.59}
\expandafter\gdef\csname NUM@famthreelotwodAFEWVzero\endcsname{\textminus{}2.66}
\expandafter\gdef\csname NUM@famthreelotwodIEMOCAPAVTfive\endcsname{\textminus{}0.13}
\expandafter\gdef\csname NUM@famthreelotwodIEMOCAPAVTtwofive\endcsname{\textminus{}0.31}
\expandafter\gdef\csname NUM@famthreelotwodIEMOCAPAVTzero\endcsname{\textminus{}1.65}
\expandafter\gdef\csname NUM@famthreelotwodIEMOCAPAfive\endcsname{\textminus{}1.34}
\expandafter\gdef\csname NUM@famthreelotwodIEMOCAPAtwofive\endcsname{\textminus{}1.74}
\expandafter\gdef\csname NUM@famthreelotwodIEMOCAPAzero\endcsname{\textminus{}2.90}
\expandafter\gdef\csname NUM@famthreelotwodIEMOCAPTfive\endcsname{\textminus{}0.35}
\expandafter\gdef\csname NUM@famthreelotwodIEMOCAPTtwofive\endcsname{\textminus{}0.99}
\expandafter\gdef\csname NUM@famthreelotwodIEMOCAPTzero\endcsname{\textminus{}1.98}
\expandafter\gdef\csname NUM@famthreelotwodIEMOCAPVfive\endcsname{\textminus{}3.25}
\expandafter\gdef\csname NUM@famthreelotwodIEMOCAPVtwofive\endcsname{\textminus{}2.81}
\expandafter\gdef\csname NUM@famthreelotwodIEMOCAPVzero\endcsname{\textminus{}3.08}
\expandafter\gdef\csname NUM@famthreelotwodMELDAVTfive\endcsname{+1.08}
\expandafter\gdef\csname NUM@famthreelotwodMELDAVTtwofive\endcsname{\textminus{}0.48}
\expandafter\gdef\csname NUM@famthreelotwodMELDAVTzero\endcsname{\textminus{}0.65}
\expandafter\gdef\csname NUM@famthreelotwodMELDAfive\endcsname{\textminus{}0.50}
\expandafter\gdef\csname NUM@famthreelotwodMELDAtwofive\endcsname{\textminus{}0.56}
\expandafter\gdef\csname NUM@famthreelotwodMELDAzero\endcsname{\textminus{}1.37}
\expandafter\gdef\csname NUM@famthreelotwodMELDTfive\endcsname{\textminus{}0.90}
\expandafter\gdef\csname NUM@famthreelotwodMELDTtwofive\endcsname{\textminus{}1.77}
\expandafter\gdef\csname NUM@famthreelotwodMELDTzero\endcsname{\textminus{}0.61}
\expandafter\gdef\csname NUM@famthreelotwodMELDVfive\endcsname{\textminus{}0.10}
\expandafter\gdef\csname NUM@famthreelotwodMELDVtwofive\endcsname{\textminus{}0.21}
\expandafter\gdef\csname NUM@famthreelotwodMELDVzero\endcsname{\textminus{}3.52}
\expandafter\gdef\csname NUM@famthreelotwodMOSEIAVTfive\endcsname{\textminus{}0.16}
\expandafter\gdef\csname NUM@famthreelotwodMOSEIAVTtwofive\endcsname{\textminus{}0.17}
\expandafter\gdef\csname NUM@famthreelotwodMOSEIAVTzero\endcsname{\textminus{}0.17}
\expandafter\gdef\csname NUM@famthreelotwodMOSEIAfive\endcsname{\textminus{}0.31}
\expandafter\gdef\csname NUM@famthreelotwodMOSEIAtwofive\endcsname{\textminus{}0.36}
\expandafter\gdef\csname NUM@famthreelotwodMOSEIAzero\endcsname{\textminus{}0.20}
\expandafter\gdef\csname NUM@famthreelotwodMOSEITfive\endcsname{\textminus{}0.50}
\expandafter\gdef\csname NUM@famthreelotwodMOSEITtwofive\endcsname{\textminus{}1.01}
\expandafter\gdef\csname NUM@famthreelotwodMOSEITzero\endcsname{\textminus{}1.07}
\expandafter\gdef\csname NUM@famthreelotwodMOSEIVfive\endcsname{\textminus{}0.23}
\expandafter\gdef\csname NUM@famthreelotwodMOSEIVtwofive\endcsname{\textminus{}0.37}
\expandafter\gdef\csname NUM@famthreelotwodMOSEIVzero\endcsname{\textminus{}1.92}
\expandafter\gdef\csname NUM@famthreepAFEWAVTfive\endcsname{1.000}
\expandafter\gdef\csname NUM@famthreepAFEWAVTtwofive\endcsname{1.000}
\expandafter\gdef\csname NUM@famthreepAFEWAVTzero\endcsname{1.000}
\expandafter\gdef\csname NUM@famthreepAFEWAfive\endcsname{1.000}
\expandafter\gdef\csname NUM@famthreepAFEWAtwofive\endcsname{1.000}
\expandafter\gdef\csname NUM@famthreepAFEWAzero\endcsname{1.000}
\expandafter\gdef\csname NUM@famthreepAFEWTfive\endcsname{1.000}
\expandafter\gdef\csname NUM@famthreepAFEWTtwofive\endcsname{1.000}
\expandafter\gdef\csname NUM@famthreepAFEWTzero\endcsname{1.000}
\expandafter\gdef\csname NUM@famthreepAFEWVfive\endcsname{1.000}
\expandafter\gdef\csname NUM@famthreepAFEWVtwofive\endcsname{1.000}
\expandafter\gdef\csname NUM@famthreepAFEWVzero\endcsname{1.000}
\expandafter\gdef\csname NUM@famthreepIEMOCAPAVTfive\endcsname{0.785}
\expandafter\gdef\csname NUM@famthreepIEMOCAPAVTtwofive\endcsname{1.000}
\expandafter\gdef\csname NUM@famthreepIEMOCAPAVTzero\endcsname{1.000}
\expandafter\gdef\csname NUM@famthreepIEMOCAPAfive\endcsname{1.000}
\expandafter\gdef\csname NUM@famthreepIEMOCAPAtwofive\endcsname{1.000}
\expandafter\gdef\csname NUM@famthreepIEMOCAPAzero\endcsname{1.000}
\expandafter\gdef\csname NUM@famthreepIEMOCAPTfive\endcsname{0.445}
\expandafter\gdef\csname NUM@famthreepIEMOCAPTtwofive\endcsname{1.000}
\expandafter\gdef\csname NUM@famthreepIEMOCAPTzero\endcsname{1.000}
\expandafter\gdef\csname NUM@famthreepIEMOCAPVfive\endcsname{1.000}
\expandafter\gdef\csname NUM@famthreepIEMOCAPVtwofive\endcsname{1.000}
\expandafter\gdef\csname NUM@famthreepIEMOCAPVzero\endcsname{1.000}
\expandafter\gdef\csname NUM@famthreepMELDAVTfive\endcsname{0.023}
\expandafter\gdef\csname NUM@famthreepMELDAVTtwofive\endcsname{0.542}
\expandafter\gdef\csname NUM@famthreepMELDAVTzero\endcsname{1.000}
\expandafter\gdef\csname NUM@famthreepMELDAfive\endcsname{1.000}
\expandafter\gdef\csname NUM@famthreepMELDAtwofive\endcsname{0.542}
\expandafter\gdef\csname NUM@famthreepMELDAzero\endcsname{1.000}
\expandafter\gdef\csname NUM@famthreepMELDTfive\endcsname{1.000}
\expandafter\gdef\csname NUM@famthreepMELDTtwofive\endcsname{0.012}
\expandafter\gdef\csname NUM@famthreepMELDTzero\endcsname{1.000}
\expandafter\gdef\csname NUM@famthreepMELDVfive\endcsname{1.000}
\expandafter\gdef\csname NUM@famthreepMELDVtwofive\endcsname{1.000}
\expandafter\gdef\csname NUM@famthreepMELDVzero\endcsname{1.000}
\expandafter\gdef\csname NUM@famthreepMOSEIAVTfive\endcsname{1.000}
\expandafter\gdef\csname NUM@famthreepMOSEIAVTtwofive\endcsname{1.000}
\expandafter\gdef\csname NUM@famthreepMOSEIAVTzero\endcsname{1.000}
\expandafter\gdef\csname NUM@famthreepMOSEIAfive\endcsname{0.358}
\expandafter\gdef\csname NUM@famthreepMOSEIAtwofive\endcsname{0.145}
\expandafter\gdef\csname NUM@famthreepMOSEIAzero\endcsname{1.000}
\expandafter\gdef\csname NUM@famthreepMOSEITfive\endcsname{0.002}
\expandafter\gdef\csname NUM@famthreepMOSEITtwofive\endcsname{0.000}
\expandafter\gdef\csname NUM@famthreepMOSEITzero\endcsname{0.021}
\expandafter\gdef\csname NUM@famthreepMOSEIVfive\endcsname{1.000}
\expandafter\gdef\csname NUM@famthreepMOSEIVtwofive\endcsname{1.000}
\expandafter\gdef\csname NUM@famthreepMOSEIVzero\endcsname{1.000}
\expandafter\gdef\csname NUM@famthreetwodAFEWAVTfive\endcsname{\textminus{}0.18}
\expandafter\gdef\csname NUM@famthreetwodAFEWAVTtwofive\endcsname{\textminus{}0.13}
\expandafter\gdef\csname NUM@famthreetwodAFEWAVTzero\endcsname{\textminus{}0.71}
\expandafter\gdef\csname NUM@famthreetwodAFEWAfive\endcsname{+0.62}
\expandafter\gdef\csname NUM@famthreetwodAFEWAtwofive\endcsname{+0.13}
\expandafter\gdef\csname NUM@famthreetwodAFEWAzero\endcsname{+0.87}
\expandafter\gdef\csname NUM@famthreetwodAFEWTfive\endcsname{\textminus{}0.18}
\expandafter\gdef\csname NUM@famthreetwodAFEWTtwofive\endcsname{\textminus{}0.31}
\expandafter\gdef\csname NUM@famthreetwodAFEWTzero\endcsname{+0.14}
\expandafter\gdef\csname NUM@famthreetwodAFEWVfive\endcsname{\textminus{}0.19}
\expandafter\gdef\csname NUM@famthreetwodAFEWVtwofive\endcsname{+0.06}
\expandafter\gdef\csname NUM@famthreetwodAFEWVzero\endcsname{\textminus{}0.78}
\expandafter\gdef\csname NUM@famthreetwodIEMOCAPAVTfive\endcsname{+0.78}
\expandafter\gdef\csname NUM@famthreetwodIEMOCAPAVTtwofive\endcsname{+0.17}
\expandafter\gdef\csname NUM@famthreetwodIEMOCAPAVTzero\endcsname{\textminus{}0.41}
\expandafter\gdef\csname NUM@famthreetwodIEMOCAPAfive\endcsname{\textminus{}0.24}
\expandafter\gdef\csname NUM@famthreetwodIEMOCAPAtwofive\endcsname{\textminus{}0.49}
\expandafter\gdef\csname NUM@famthreetwodIEMOCAPAzero\endcsname{\textminus{}0.79}
\expandafter\gdef\csname NUM@famthreetwodIEMOCAPTfive\endcsname{\textminus{}0.19}
\expandafter\gdef\csname NUM@famthreetwodIEMOCAPTtwofive\endcsname{\textminus{}0.40}
\expandafter\gdef\csname NUM@famthreetwodIEMOCAPTzero\endcsname{\textminus{}0.26}
\expandafter\gdef\csname NUM@famthreetwodIEMOCAPVfive\endcsname{\textminus{}1.17}
\expandafter\gdef\csname NUM@famthreetwodIEMOCAPVtwofive\endcsname{\textminus{}1.03}
\expandafter\gdef\csname NUM@famthreetwodIEMOCAPVzero\endcsname{\textminus{}0.87}
\expandafter\gdef\csname NUM@famthreetwodMELDAVTfive\endcsname{+2.16}
\expandafter\gdef\csname NUM@famthreetwodMELDAVTtwofive\endcsname{\textminus{}0.24}
\expandafter\gdef\csname NUM@famthreetwodMELDAVTzero\endcsname{+1.27}
\expandafter\gdef\csname NUM@famthreetwodMELDAfive\endcsname{\textminus{}0.00}
\expandafter\gdef\csname NUM@famthreetwodMELDAtwofive\endcsname{\textminus{}0.28}
\expandafter\gdef\csname NUM@famthreetwodMELDAzero\endcsname{\textminus{}0.05}
\expandafter\gdef\csname NUM@famthreetwodMELDTfive\endcsname{\textminus{}0.35}
\expandafter\gdef\csname NUM@famthreetwodMELDTtwofive\endcsname{\textminus{}1.24}
\expandafter\gdef\csname NUM@famthreetwodMELDTzero\endcsname{+0.84}
\expandafter\gdef\csname NUM@famthreetwodMELDVfive\endcsname{+0.21}
\expandafter\gdef\csname NUM@famthreetwodMELDVtwofive\endcsname{+0.22}
\expandafter\gdef\csname NUM@famthreetwodMELDVzero\endcsname{\textminus{}0.15}
\expandafter\gdef\csname NUM@famthreetwodMOSEIAVTfive\endcsname{\textminus{}0.06}
\expandafter\gdef\csname NUM@famthreetwodMOSEIAVTtwofive\endcsname{\textminus{}0.05}
\expandafter\gdef\csname NUM@famthreetwodMOSEIAVTzero\endcsname{+0.07}
\expandafter\gdef\csname NUM@famthreetwodMOSEIAfive\endcsname{\textminus{}0.16}
\expandafter\gdef\csname NUM@famthreetwodMOSEIAtwofive\endcsname{\textminus{}0.20}
\expandafter\gdef\csname NUM@famthreetwodMOSEIAzero\endcsname{+0.11}
\expandafter\gdef\csname NUM@famthreetwodMOSEITfive\endcsname{\textminus{}0.33}
\expandafter\gdef\csname NUM@famthreetwodMOSEITtwofive\endcsname{\textminus{}0.78}
\expandafter\gdef\csname NUM@famthreetwodMOSEITzero\endcsname{\textminus{}0.66}
\expandafter\gdef\csname NUM@famthreetwodMOSEIVfive\endcsname{\textminus{}0.03}
\expandafter\gdef\csname NUM@famthreetwodMOSEIVtwofive\endcsname{\textminus{}0.14}
\expandafter\gdef\csname NUM@famthreetwodMOSEIVzero\endcsname{\textminus{}0.73}
\expandafter\gdef\csname NUM@favtAFEWfive\endcsname{70.8}
\expandafter\gdef\csname NUM@favtAFEWtwofive\endcsname{75.8}
\expandafter\gdef\csname NUM@favtIEMOCAPfive\endcsname{72.7}
\expandafter\gdef\csname NUM@favtIEMOCAPtwofive\endcsname{74.7}
\expandafter\gdef\csname NUM@favtMELDfive\endcsname{50.1}
\expandafter\gdef\csname NUM@favtMELDtwofive\endcsname{51.9}
\expandafter\gdef\csname NUM@favtMOSEIfive\endcsname{41.1}
\expandafter\gdef\csname NUM@favtMOSEItwofive\endcsname{44.1}
\expandafter\gdef\csname NUM@fclokAFEWfive\endcsname{84}
\expandafter\gdef\csname NUM@fclokAFEWtwofive\endcsname{98}
\expandafter\gdef\csname NUM@fclokIEMOCAPfive\endcsname{68}
\expandafter\gdef\csname NUM@fclokIEMOCAPtwofive\endcsname{98}
\expandafter\gdef\csname NUM@fclokMELDfive\endcsname{65}
\expandafter\gdef\csname NUM@fclokMELDtwofive\endcsname{93}
\expandafter\gdef\csname NUM@fclokMOSEIfive\endcsname{32}
\expandafter\gdef\csname NUM@fclokMOSEItwofive\endcsname{45}
\expandafter\gdef\csname NUM@fheavyAFEWfive\endcsname{2}
\expandafter\gdef\csname NUM@fheavyAFEWtwofive\endcsname{4}
\expandafter\gdef\csname NUM@fheavyIEMOCAPfive\endcsname{4}
\expandafter\gdef\csname NUM@fheavyIEMOCAPtwofive\endcsname{50}
\expandafter\gdef\csname NUM@fheavyMELDfive\endcsname{14}
\expandafter\gdef\csname NUM@fheavyMELDtwofive\endcsname{58}
\expandafter\gdef\csname NUM@fheavyMOSEIfive\endcsname{46}
\expandafter\gdef\csname NUM@fheavyMOSEItwofive\endcsname{2}
\expandafter\gdef\csname NUM@fixAFEWAVTfive\endcsname{5}
\expandafter\gdef\csname NUM@fixAFEWAVTtwofive\endcsname{4}
\expandafter\gdef\csname NUM@fixAFEWAfive\endcsname{25}
\expandafter\gdef\csname NUM@fixAFEWAtwofive\endcsname{28}
\expandafter\gdef\csname NUM@fixAFEWTfive\endcsname{33}
\expandafter\gdef\csname NUM@fixAFEWTtwofive\endcsname{35}
\expandafter\gdef\csname NUM@fixAFEWVfive\endcsname{16}
\expandafter\gdef\csname NUM@fixAFEWVtwofive\endcsname{14}
\expandafter\gdef\csname NUM@fixIEMOCAPAVTfive\endcsname{5}
\expandafter\gdef\csname NUM@fixIEMOCAPAVTtwofive\endcsname{4}
\expandafter\gdef\csname NUM@fixIEMOCAPAfive\endcsname{15}
\expandafter\gdef\csname NUM@fixIEMOCAPAtwofive\endcsname{14}
\expandafter\gdef\csname NUM@fixIEMOCAPTfive\endcsname{30}
\expandafter\gdef\csname NUM@fixIEMOCAPTtwofive\endcsname{28}
\expandafter\gdef\csname NUM@fixIEMOCAPVfive\endcsname{22}
\expandafter\gdef\csname NUM@fixIEMOCAPVtwofive\endcsname{22}
\expandafter\gdef\csname NUM@fixMELDAVTfive\endcsname{5}
\expandafter\gdef\csname NUM@fixMELDAVTtwofive\endcsname{5}
\expandafter\gdef\csname NUM@fixMELDAfive\endcsname{16}
\expandafter\gdef\csname NUM@fixMELDAtwofive\endcsname{17}
\expandafter\gdef\csname NUM@fixMELDTfive\endcsname{14}
\expandafter\gdef\csname NUM@fixMELDTtwofive\endcsname{12}
\expandafter\gdef\csname NUM@fixMELDVfive\endcsname{22}
\expandafter\gdef\csname NUM@fixMELDVtwofive\endcsname{23}
\expandafter\gdef\csname NUM@fixMOSEIAVTfive\endcsname{5}
\expandafter\gdef\csname NUM@fixMOSEIAVTtwofive\endcsname{5}
\expandafter\gdef\csname NUM@fixMOSEIAfive\endcsname{17}
\expandafter\gdef\csname NUM@fixMOSEIAtwofive\endcsname{18}
\expandafter\gdef\csname NUM@fixMOSEITfive\endcsname{20}
\expandafter\gdef\csname NUM@fixMOSEITtwofive\endcsname{20}
\expandafter\gdef\csname NUM@fixMOSEIVfive\endcsname{17}
\expandafter\gdef\csname NUM@fixMOSEIVtwofive\endcsname{18}
\expandafter\gdef\csname NUM@fkshotAFEWfive\endcsname{68.7}
\expandafter\gdef\csname NUM@fkshotAFEWtwofive\endcsname{69.9}
\expandafter\gdef\csname NUM@fkshotIEMOCAPfive\endcsname{71.5}
\expandafter\gdef\csname NUM@fkshotIEMOCAPtwofive\endcsname{73.8}
\expandafter\gdef\csname NUM@fkshotMELDfive\endcsname{50.4}
\expandafter\gdef\csname NUM@fkshotMELDtwofive\endcsname{52.2}
\expandafter\gdef\csname NUM@fkshotMOSEIfive\endcsname{38.7}
\expandafter\gdef\csname NUM@fkshotMOSEItwofive\endcsname{42.5}
\expandafter\gdef\csname NUM@fkshotdAFEWfive\endcsname{\textminus{}2.0}
\expandafter\gdef\csname NUM@fkshotdAFEWtwofive\endcsname{\textminus{}5.9}
\expandafter\gdef\csname NUM@fkshotdIEMOCAPfive\endcsname{\textminus{}1.2}
\expandafter\gdef\csname NUM@fkshotdIEMOCAPtwofive\endcsname{\textminus{}0.9}
\expandafter\gdef\csname NUM@fkshotdMELDfive\endcsname{+0.2}
\expandafter\gdef\csname NUM@fkshotdMELDtwofive\endcsname{+0.3}
\expandafter\gdef\csname NUM@fkshotdMOSEIfive\endcsname{\textminus{}2.4}
\expandafter\gdef\csname NUM@fkshotdMOSEItwofive\endcsname{\textminus{}1.6}
\expandafter\gdef\csname NUM@foracleAFEWfive\endcsname{76.4}
\expandafter\gdef\csname NUM@foracleAFEWtwofive\endcsname{78.9}
\expandafter\gdef\csname NUM@foracleIEMOCAPfive\endcsname{73.4}
\expandafter\gdef\csname NUM@foracleIEMOCAPtwofive\endcsname{75.6}
\expandafter\gdef\csname NUM@foracleMELDfive\endcsname{51.4}
\expandafter\gdef\csname NUM@foracleMELDtwofive\endcsname{53.0}
\expandafter\gdef\csname NUM@foracleMOSEIfive\endcsname{40.0}
\expandafter\gdef\csname NUM@foracleMOSEItwofive\endcsname{43.6}
\expandafter\gdef\csname NUM@foracledAFEWfive\endcsname{+5.6}
\expandafter\gdef\csname NUM@foracledAFEWtwofive\endcsname{+3.0}
\expandafter\gdef\csname NUM@foracledIEMOCAPfive\endcsname{+0.7}
\expandafter\gdef\csname NUM@foracledIEMOCAPtwofive\endcsname{+0.9}
\expandafter\gdef\csname NUM@foracledMELDfive\endcsname{+1.3}
\expandafter\gdef\csname NUM@foracledMELDtwofive\endcsname{+1.1}
\expandafter\gdef\csname NUM@foracledMOSEIfive\endcsname{\textminus{}1.1}
\expandafter\gdef\csname NUM@foracledMOSEItwofive\endcsname{\textminus{}0.4}
\expandafter\gdef\csname NUM@fprodAFEWfive\endcsname{69.1}
\expandafter\gdef\csname NUM@fprodAFEWtwofive\endcsname{73.1}
\expandafter\gdef\csname NUM@fprodIEMOCAPfive\endcsname{72.4}
\expandafter\gdef\csname NUM@fprodIEMOCAPtwofive\endcsname{74.4}
\expandafter\gdef\csname NUM@fprodMELDfive\endcsname{50.5}
\expandafter\gdef\csname NUM@fprodMELDtwofive\endcsname{52.2}
\expandafter\gdef\csname NUM@fprodMOSEIfive\endcsname{39.6}
\expandafter\gdef\csname NUM@fprodMOSEItwofive\endcsname{43.5}
\expandafter\gdef\csname NUM@fproddAFEWfive\endcsname{\textminus{}1.6}
\expandafter\gdef\csname NUM@fproddAFEWtwofive\endcsname{\textminus{}2.7}
\expandafter\gdef\csname NUM@fproddIEMOCAPfive\endcsname{\textminus{}0.3}
\expandafter\gdef\csname NUM@fproddIEMOCAPtwofive\endcsname{\textminus{}0.3}
\expandafter\gdef\csname NUM@fproddMELDfive\endcsname{+0.4}
\expandafter\gdef\csname NUM@fproddMELDtwofive\endcsname{+0.3}
\expandafter\gdef\csname NUM@fproddMOSEIfive\endcsname{\textminus{}1.5}
\expandafter\gdef\csname NUM@fproddMOSEItwofive\endcsname{\textminus{}0.6}
\expandafter\gdef\csname NUM@fruleAFEWfive\endcsname{73.7}
\expandafter\gdef\csname NUM@fruleAFEWtwofive\endcsname{77.8}
\expandafter\gdef\csname NUM@fruleIEMOCAPfive\endcsname{72.0}
\expandafter\gdef\csname NUM@fruleIEMOCAPtwofive\endcsname{74.5}
\expandafter\gdef\csname NUM@fruleMELDfive\endcsname{50.0}
\expandafter\gdef\csname NUM@fruleMELDtwofive\endcsname{52.7}
\expandafter\gdef\csname NUM@fruleMOSEIfive\endcsname{38.3}
\expandafter\gdef\csname NUM@fruleMOSEItwofive\endcsname{43.1}
\expandafter\gdef\csname NUM@fruledAFEWfive\endcsname{+2.9}
\expandafter\gdef\csname NUM@fruledAFEWtwofive\endcsname{+2.0}
\expandafter\gdef\csname NUM@fruledIEMOCAPfive\endcsname{\textminus{}0.7}
\expandafter\gdef\csname NUM@fruledIEMOCAPtwofive\endcsname{\textminus{}0.2}
\expandafter\gdef\csname NUM@fruledMELDfive\endcsname{\textminus{}0.1}
\expandafter\gdef\csname NUM@fruledMELDtwofive\endcsname{+0.8}
\expandafter\gdef\csname NUM@fruledMOSEIfive\endcsname{\textminus{}2.8}
\expandafter\gdef\csname NUM@fruledMOSEItwofive\endcsname{\textminus{}1.0}
\expandafter\gdef\csname NUM@gainrAFEWAfive\endcsname{+4.1}
\expandafter\gdef\csname NUM@gainrAFEWAtwofive\endcsname{+3.3}
\expandafter\gdef\csname NUM@gainrAFEWTfive\endcsname{+1.6}
\expandafter\gdef\csname NUM@gainrAFEWTtwofive\endcsname{+0.7}
\expandafter\gdef\csname NUM@gainrAFEWVfive\endcsname{+0.1}
\expandafter\gdef\csname NUM@gainrAFEWVtwofive\endcsname{\textminus{}0.2}
\expandafter\gdef\csname NUM@gainrIEMOCAPAfive\endcsname{+2.4}
\expandafter\gdef\csname NUM@gainrIEMOCAPAtwofive\endcsname{+1.4}
\expandafter\gdef\csname NUM@gainrIEMOCAPTfive\endcsname{+7.3}
\expandafter\gdef\csname NUM@gainrIEMOCAPTtwofive\endcsname{+5.0}
\expandafter\gdef\csname NUM@gainrIEMOCAPVfive\endcsname{\textminus{}0.8}
\expandafter\gdef\csname NUM@gainrIEMOCAPVtwofive\endcsname{\textminus{}0.7}
\expandafter\gdef\csname NUM@gainrMELDAfive\endcsname{+3.1}
\expandafter\gdef\csname NUM@gainrMELDAtwofive\endcsname{+2.9}
\expandafter\gdef\csname NUM@gainrMELDTfive\endcsname{\textminus{}1.2}
\expandafter\gdef\csname NUM@gainrMELDTtwofive\endcsname{\textminus{}0.9}
\expandafter\gdef\csname NUM@gainrMELDVfive\endcsname{+0.8}
\expandafter\gdef\csname NUM@gainrMELDVtwofive\endcsname{+1.1}
\expandafter\gdef\csname NUM@gainrMOSEIAfive\endcsname{+3.2}
\expandafter\gdef\csname NUM@gainrMOSEIAtwofive\endcsname{+3.4}
\expandafter\gdef\csname NUM@gainrMOSEITfive\endcsname{\textminus{}0.0}
\expandafter\gdef\csname NUM@gainrMOSEITtwofive\endcsname{+0.0}
\expandafter\gdef\csname NUM@gainrMOSEIVfive\endcsname{+0.2}
\expandafter\gdef\csname NUM@gainrMOSEIVtwofive\endcsname{+2.6}
\expandafter\gdef\csname NUM@gapkAFEW\endcsname{4.6}
\expandafter\gdef\csname NUM@gapkIEMOCAP\endcsname{8.1}
\expandafter\gdef\csname NUM@gapkMELD\endcsname{8.4}
\expandafter\gdef\csname NUM@gapkMOSEI\endcsname{8.3}
\expandafter\gdef\csname NUM@gapstartAFEW\endcsname{15.2}
\expandafter\gdef\csname NUM@gapstartIEMOCAP\endcsname{14.0}
\expandafter\gdef\csname NUM@gapstartMELD\endcsname{11.7}
\expandafter\gdef\csname NUM@gapstartMOSEI\endcsname{12.0}
\expandafter\gdef\csname NUM@harmAFEWAVTfive\endcsname{4}
\expandafter\gdef\csname NUM@harmAFEWAVTtwofive\endcsname{3}
\expandafter\gdef\csname NUM@harmAFEWAfive\endcsname{3}
\expandafter\gdef\csname NUM@harmAFEWAtwofive\endcsname{4}
\expandafter\gdef\csname NUM@harmAFEWTfive\endcsname{5}
\expandafter\gdef\csname NUM@harmAFEWTtwofive\endcsname{4}
\expandafter\gdef\csname NUM@harmAFEWVfive\endcsname{15}
\expandafter\gdef\csname NUM@harmAFEWVtwofive\endcsname{11}
\expandafter\gdef\csname NUM@harmIEMOCAPAVTfive\endcsname{6}
\expandafter\gdef\csname NUM@harmIEMOCAPAVTtwofive\endcsname{6}
\expandafter\gdef\csname NUM@harmIEMOCAPAfive\endcsname{8}
\expandafter\gdef\csname NUM@harmIEMOCAPAtwofive\endcsname{7}
\expandafter\gdef\csname NUM@harmIEMOCAPTfive\endcsname{8}
\expandafter\gdef\csname NUM@harmIEMOCAPTtwofive\endcsname{7}
\expandafter\gdef\csname NUM@harmIEMOCAPVfive\endcsname{8}
\expandafter\gdef\csname NUM@harmIEMOCAPVtwofive\endcsname{6}
\expandafter\gdef\csname NUM@harmMELDAVTfive\endcsname{8}
\expandafter\gdef\csname NUM@harmMELDAVTtwofive\endcsname{7}
\expandafter\gdef\csname NUM@harmMELDAfive\endcsname{12}
\expandafter\gdef\csname NUM@harmMELDAtwofive\endcsname{10}
\expandafter\gdef\csname NUM@harmMELDTfive\endcsname{19}
\expandafter\gdef\csname NUM@harmMELDTtwofive\endcsname{19}
\expandafter\gdef\csname NUM@harmMELDVfive\endcsname{10}
\expandafter\gdef\csname NUM@harmMELDVtwofive\endcsname{11}
\expandafter\gdef\csname NUM@harmMOSEIAVTfive\endcsname{8}
\expandafter\gdef\csname NUM@harmMOSEIAVTtwofive\endcsname{7}
\expandafter\gdef\csname NUM@harmMOSEIAfive\endcsname{10}
\expandafter\gdef\csname NUM@harmMOSEIAtwofive\endcsname{11}
\expandafter\gdef\csname NUM@harmMOSEITfive\endcsname{14}
\expandafter\gdef\csname NUM@harmMOSEITtwofive\endcsname{13}
\expandafter\gdef\csname NUM@harmMOSEIVfive\endcsname{14}
\expandafter\gdef\csname NUM@harmMOSEIVtwofive\endcsname{14}
\expandafter\gdef\csname NUM@ktfAFEW\endcsname{75.2}
\expandafter\gdef\csname NUM@ktfIEMOCAP\endcsname{75.5}
\expandafter\gdef\csname NUM@ktfMELD\endcsname{50.5}
\expandafter\gdef\csname NUM@ktfMOSEI\endcsname{43.4}
\expandafter\gdef\csname NUM@lvlnoneAFEWAVTfive\endcsname{69.7}
\expandafter\gdef\csname NUM@lvlnoneAFEWAVTtwofive\endcsname{75.2}
\expandafter\gdef\csname NUM@lvlnoneAFEWAfive\endcsname{50.1}
\expandafter\gdef\csname NUM@lvlnoneAFEWAtwofive\endcsname{51.9}
\expandafter\gdef\csname NUM@lvlnoneAFEWTfive\endcsname{44.9}
\expandafter\gdef\csname NUM@lvlnoneAFEWTtwofive\endcsname{45.7}
\expandafter\gdef\csname NUM@lvlnoneAFEWVfive\endcsname{70.7}
\expandafter\gdef\csname NUM@lvlnoneAFEWVtwofive\endcsname{73.3}
\expandafter\gdef\csname NUM@lvlnoneIEMOCAPAVTfive\endcsname{72.6}
\expandafter\gdef\csname NUM@lvlnoneIEMOCAPAVTtwofive\endcsname{75.5}
\expandafter\gdef\csname NUM@lvlnoneIEMOCAPAfive\endcsname{67.1}
\expandafter\gdef\csname NUM@lvlnoneIEMOCAPAtwofive\endcsname{69.9}
\expandafter\gdef\csname NUM@lvlnoneIEMOCAPTfive\endcsname{50.3}
\expandafter\gdef\csname NUM@lvlnoneIEMOCAPTtwofive\endcsname{54.3}
\expandafter\gdef\csname NUM@lvlnoneIEMOCAPVfive\endcsname{53.3}
\expandafter\gdef\csname NUM@lvlnoneIEMOCAPVtwofive\endcsname{53.5}
\expandafter\gdef\csname NUM@lvlnoneMELDAVTfive\endcsname{47.6}
\expandafter\gdef\csname NUM@lvlnoneMELDAVTtwofive\endcsname{50.5}
\expandafter\gdef\csname NUM@lvlnoneMELDAfive\endcsname{41.5}
\expandafter\gdef\csname NUM@lvlnoneMELDAtwofive\endcsname{42.8}
\expandafter\gdef\csname NUM@lvlnoneMELDTfive\endcsname{49.6}
\expandafter\gdef\csname NUM@lvlnoneMELDTtwofive\endcsname{51.2}
\expandafter\gdef\csname NUM@lvlnoneMELDVfive\endcsname{36.7}
\expandafter\gdef\csname NUM@lvlnoneMELDVtwofive\endcsname{36.6}
\expandafter\gdef\csname NUM@lvlnoneMOSEIAVTfive\endcsname{41.5}
\expandafter\gdef\csname NUM@lvlnoneMOSEIAVTtwofive\endcsname{43.4}
\expandafter\gdef\csname NUM@lvlnoneMOSEIAfive\endcsname{34.9}
\expandafter\gdef\csname NUM@lvlnoneMOSEIAtwofive\endcsname{37.2}
\expandafter\gdef\csname NUM@lvlnoneMOSEITfive\endcsname{33.3}
\expandafter\gdef\csname NUM@lvlnoneMOSEITtwofive\endcsname{36.1}
\expandafter\gdef\csname NUM@lvlnoneMOSEIVfive\endcsname{36.4}
\expandafter\gdef\csname NUM@lvlnoneMOSEIVtwofive\endcsname{37.2}
\expandafter\gdef\csname NUM@lvlpropAFEWAVTfive\endcsname{70.8}
\expandafter\gdef\csname NUM@lvlpropAFEWAVTtwofive\endcsname{75.8}
\expandafter\gdef\csname NUM@lvlpropAFEWAfive\endcsname{54.2}
\expandafter\gdef\csname NUM@lvlpropAFEWAtwofive\endcsname{55.3}
\expandafter\gdef\csname NUM@lvlpropAFEWTfive\endcsname{46.6}
\expandafter\gdef\csname NUM@lvlpropAFEWTtwofive\endcsname{46.4}
\expandafter\gdef\csname NUM@lvlpropAFEWVfive\endcsname{70.8}
\expandafter\gdef\csname NUM@lvlpropAFEWVtwofive\endcsname{73.1}
\expandafter\gdef\csname NUM@lvlpropIEMOCAPAVTfive\endcsname{72.7}
\expandafter\gdef\csname NUM@lvlpropIEMOCAPAVTtwofive\endcsname{74.7}
\expandafter\gdef\csname NUM@lvlpropIEMOCAPAfive\endcsname{69.5}
\expandafter\gdef\csname NUM@lvlpropIEMOCAPAtwofive\endcsname{71.4}
\expandafter\gdef\csname NUM@lvlpropIEMOCAPTfive\endcsname{57.7}
\expandafter\gdef\csname NUM@lvlpropIEMOCAPTtwofive\endcsname{59.3}
\expandafter\gdef\csname NUM@lvlpropIEMOCAPVfive\endcsname{52.5}
\expandafter\gdef\csname NUM@lvlpropIEMOCAPVtwofive\endcsname{52.8}
\expandafter\gdef\csname NUM@lvlpropMELDAVTfive\endcsname{50.1}
\expandafter\gdef\csname NUM@lvlpropMELDAVTtwofive\endcsname{51.9}
\expandafter\gdef\csname NUM@lvlpropMELDAfive\endcsname{44.6}
\expandafter\gdef\csname NUM@lvlpropMELDAtwofive\endcsname{45.6}
\expandafter\gdef\csname NUM@lvlpropMELDTfive\endcsname{48.5}
\expandafter\gdef\csname NUM@lvlpropMELDTtwofive\endcsname{50.2}
\expandafter\gdef\csname NUM@lvlpropMELDVfive\endcsname{37.6}
\expandafter\gdef\csname NUM@lvlpropMELDVtwofive\endcsname{37.7}
\expandafter\gdef\csname NUM@lvlpropMOSEIAVTfive\endcsname{41.1}
\expandafter\gdef\csname NUM@lvlpropMOSEIAVTtwofive\endcsname{44.1}
\expandafter\gdef\csname NUM@lvlpropMOSEIAfive\endcsname{38.1}
\expandafter\gdef\csname NUM@lvlpropMOSEIAtwofive\endcsname{40.6}
\expandafter\gdef\csname NUM@lvlpropMOSEITfive\endcsname{33.3}
\expandafter\gdef\csname NUM@lvlpropMOSEITtwofive\endcsname{36.1}
\expandafter\gdef\csname NUM@lvlpropMOSEIVfive\endcsname{36.6}
\expandafter\gdef\csname NUM@lvlpropMOSEIVtwofive\endcsname{39.7}
\expandafter\gdef\csname NUM@posAFEWAVTnonefive\endcsname{0.5}
\expandafter\gdef\csname NUM@posAFEWAVTnonetwofive\endcsname{\textminus{}0.1}
\expandafter\gdef\csname NUM@posAFEWAVTpropfive\endcsname{0.4}
\expandafter\gdef\csname NUM@posAFEWAVTproptwofive\endcsname{0.1}
\expandafter\gdef\csname NUM@posAFEWAVTstart\endcsname{1.9}
\expandafter\gdef\csname NUM@posAFEWAnonefive\endcsname{1.2}
\expandafter\gdef\csname NUM@posAFEWAnonetwofive\endcsname{\textminus{}0.1}
\expandafter\gdef\csname NUM@posAFEWApropfive\endcsname{\textminus{}0.1}
\expandafter\gdef\csname NUM@posAFEWAproptwofive\endcsname{\textminus{}0.7}
\expandafter\gdef\csname NUM@posAFEWAstart\endcsname{1.8}
\expandafter\gdef\csname NUM@posAFEWTnonefive\endcsname{\textminus{}0.1}
\expandafter\gdef\csname NUM@posAFEWTnonetwofive\endcsname{\textminus{}0.2}
\expandafter\gdef\csname NUM@posAFEWTpropfive\endcsname{\textminus{}0.1}
\expandafter\gdef\csname NUM@posAFEWTproptwofive\endcsname{\textminus{}0.4}
\expandafter\gdef\csname NUM@posAFEWTstart\endcsname{\textminus{}0.4}
\expandafter\gdef\csname NUM@posAFEWVnonefive\endcsname{1.2}
\expandafter\gdef\csname NUM@posAFEWVnonetwofive\endcsname{0.6}
\expandafter\gdef\csname NUM@posAFEWVpropfive\endcsname{1.2}
\expandafter\gdef\csname NUM@posAFEWVproptwofive\endcsname{0.8}
\expandafter\gdef\csname NUM@posAFEWVstart\endcsname{\textminus{}0.5}
\expandafter\gdef\csname NUM@posIEMOCAPAVTnonefive\endcsname{2.5}
\expandafter\gdef\csname NUM@posIEMOCAPAVTnonetwofive\endcsname{2.7}
\expandafter\gdef\csname NUM@posIEMOCAPAVTpropfive\endcsname{4.0}
\expandafter\gdef\csname NUM@posIEMOCAPAVTproptwofive\endcsname{4.1}
\expandafter\gdef\csname NUM@posIEMOCAPAVTstart\endcsname{2.7}
\expandafter\gdef\csname NUM@posIEMOCAPAnonefive\endcsname{2.1}
\expandafter\gdef\csname NUM@posIEMOCAPAnonetwofive\endcsname{1.8}
\expandafter\gdef\csname NUM@posIEMOCAPApropfive\endcsname{2.7}
\expandafter\gdef\csname NUM@posIEMOCAPAproptwofive\endcsname{2.6}
\expandafter\gdef\csname NUM@posIEMOCAPAstart\endcsname{1.9}
\expandafter\gdef\csname NUM@posIEMOCAPTnonefive\endcsname{0.7}
\expandafter\gdef\csname NUM@posIEMOCAPTnonetwofive\endcsname{0.1}
\expandafter\gdef\csname NUM@posIEMOCAPTpropfive\endcsname{0.7}
\expandafter\gdef\csname NUM@posIEMOCAPTproptwofive\endcsname{0.6}
\expandafter\gdef\csname NUM@posIEMOCAPTstart\endcsname{\textminus{}0.2}
\expandafter\gdef\csname NUM@posIEMOCAPVnonefive\endcsname{4.3}
\expandafter\gdef\csname NUM@posIEMOCAPVnonetwofive\endcsname{3.9}
\expandafter\gdef\csname NUM@posIEMOCAPVpropfive\endcsname{5.6}
\expandafter\gdef\csname NUM@posIEMOCAPVproptwofive\endcsname{5.1}
\expandafter\gdef\csname NUM@posIEMOCAPVstart\endcsname{3.3}
\expandafter\gdef\csname NUM@posMELDAVTnonefive\endcsname{1.5}
\expandafter\gdef\csname NUM@posMELDAVTnonetwofive\endcsname{0.7}
\expandafter\gdef\csname NUM@posMELDAVTpropfive\endcsname{1.1}
\expandafter\gdef\csname NUM@posMELDAVTproptwofive\endcsname{1.2}
\expandafter\gdef\csname NUM@posMELDAVTstart\endcsname{2.1}
\expandafter\gdef\csname NUM@posMELDAnonefive\endcsname{0.8}
\expandafter\gdef\csname NUM@posMELDAnonetwofive\endcsname{0.3}
\expandafter\gdef\csname NUM@posMELDApropfive\endcsname{0.4}
\expandafter\gdef\csname NUM@posMELDAproptwofive\endcsname{0.4}
\expandafter\gdef\csname NUM@posMELDAstart\endcsname{0.8}
\expandafter\gdef\csname NUM@posMELDTnonefive\endcsname{0.6}
\expandafter\gdef\csname NUM@posMELDTnonetwofive\endcsname{0.2}
\expandafter\gdef\csname NUM@posMELDTpropfive\endcsname{1.0}
\expandafter\gdef\csname NUM@posMELDTproptwofive\endcsname{1.1}
\expandafter\gdef\csname NUM@posMELDTstart\endcsname{0.2}
\expandafter\gdef\csname NUM@posMELDVnonefive\endcsname{1.2}
\expandafter\gdef\csname NUM@posMELDVnonetwofive\endcsname{1.2}
\expandafter\gdef\csname NUM@posMELDVpropfive\endcsname{0.8}
\expandafter\gdef\csname NUM@posMELDVproptwofive\endcsname{1.0}
\expandafter\gdef\csname NUM@posMELDVstart\endcsname{1.4}
\expandafter\gdef\csname NUM@posMOSEIAVTnonefive\endcsname{0.1}
\expandafter\gdef\csname NUM@posMOSEIAVTnonetwofive\endcsname{0.4}
\expandafter\gdef\csname NUM@posMOSEIAVTpropfive\endcsname{0.0}
\expandafter\gdef\csname NUM@posMOSEIAVTproptwofive\endcsname{0.3}
\expandafter\gdef\csname NUM@posMOSEIAVTstart\endcsname{0.4}
\expandafter\gdef\csname NUM@posMOSEIAnonefive\endcsname{0.0}
\expandafter\gdef\csname NUM@posMOSEIAnonetwofive\endcsname{0.2}
\expandafter\gdef\csname NUM@posMOSEIApropfive\endcsname{\textminus{}0.1}
\expandafter\gdef\csname NUM@posMOSEIAproptwofive\endcsname{0.2}
\expandafter\gdef\csname NUM@posMOSEIAstart\endcsname{\textminus{}0.0}
\expandafter\gdef\csname NUM@posMOSEITnonefive\endcsname{0.2}
\expandafter\gdef\csname NUM@posMOSEITnonetwofive\endcsname{0.0}
\expandafter\gdef\csname NUM@posMOSEITpropfive\endcsname{\textminus{}0.1}
\expandafter\gdef\csname NUM@posMOSEITproptwofive\endcsname{0.3}
\expandafter\gdef\csname NUM@posMOSEITstart\endcsname{0.3}
\expandafter\gdef\csname NUM@posMOSEIVnonefive\endcsname{0.3}
\expandafter\gdef\csname NUM@posMOSEIVnonetwofive\endcsname{0.2}
\expandafter\gdef\csname NUM@posMOSEIVpropfive\endcsname{0.4}
\expandafter\gdef\csname NUM@posMOSEIVproptwofive\endcsname{0.2}
\expandafter\gdef\csname NUM@posMOSEIVstart\endcsname{\textminus{}0.2}
\expandafter\gdef\csname NUM@reproAFEWAfive\endcsname{69}
\expandafter\gdef\csname NUM@reproAFEWAtwofive\endcsname{63}
\expandafter\gdef\csname NUM@reproAFEWTfive\endcsname{58}
\expandafter\gdef\csname NUM@reproAFEWTtwofive\endcsname{55}
\expandafter\gdef\csname NUM@reproAFEWVfive\endcsname{65}
\expandafter\gdef\csname NUM@reproAFEWVtwofive\endcsname{72}
\expandafter\gdef\csname NUM@reproIEMOCAPAfive\endcsname{73}
\expandafter\gdef\csname NUM@reproIEMOCAPAtwofive\endcsname{74}
\expandafter\gdef\csname NUM@reproIEMOCAPTfive\endcsname{62}
\expandafter\gdef\csname NUM@reproIEMOCAPTtwofive\endcsname{62}
\expandafter\gdef\csname NUM@reproIEMOCAPVfive\endcsname{60}
\expandafter\gdef\csname NUM@reproIEMOCAPVtwofive\endcsname{61}
\expandafter\gdef\csname NUM@reproMELDAfive\endcsname{65}
\expandafter\gdef\csname NUM@reproMELDAtwofive\endcsname{66}
\expandafter\gdef\csname NUM@reproMELDTfive\endcsname{57}
\expandafter\gdef\csname NUM@reproMELDTtwofive\endcsname{60}
\expandafter\gdef\csname NUM@reproMELDVfive\endcsname{57}
\expandafter\gdef\csname NUM@reproMELDVtwofive\endcsname{54}
\expandafter\gdef\csname NUM@reproMOSEIAfive\endcsname{69}
\expandafter\gdef\csname NUM@reproMOSEIAtwofive\endcsname{69}
\expandafter\gdef\csname NUM@reproMOSEITfive\endcsname{55}
\expandafter\gdef\csname NUM@reproMOSEITtwofive\endcsname{57}
\expandafter\gdef\csname NUM@reproMOSEIVfive\endcsname{60}
\expandafter\gdef\csname NUM@reproMOSEIVtwofive\endcsname{58}
\expandafter\gdef\csname NUM@rhofive\endcsname{0.51}
\expandafter\gdef\csname NUM@rhotwofive\endcsname{0.52}
\expandafter\gdef\csname NUM@roomAFEWAfive\endcsname{6.1}
\expandafter\gdef\csname NUM@roomAFEWAtwofive\endcsname{4.3}
\expandafter\gdef\csname NUM@roomAFEWTfive\endcsname{3.7}
\expandafter\gdef\csname NUM@roomAFEWTtwofive\endcsname{2.9}
\expandafter\gdef\csname NUM@roomAFEWVfive\endcsname{4.9}
\expandafter\gdef\csname NUM@roomAFEWVtwofive\endcsname{2.4}
\expandafter\gdef\csname NUM@roomIEMOCAPAfive\endcsname{10.1}
\expandafter\gdef\csname NUM@roomIEMOCAPAtwofive\endcsname{7.2}
\expandafter\gdef\csname NUM@roomIEMOCAPTfive\endcsname{11.5}
\expandafter\gdef\csname NUM@roomIEMOCAPTtwofive\endcsname{7.6}
\expandafter\gdef\csname NUM@roomIEMOCAPVfive\endcsname{4.7}
\expandafter\gdef\csname NUM@roomIEMOCAPVtwofive\endcsname{4.5}
\expandafter\gdef\csname NUM@roomMELDAfive\endcsname{6.5}
\expandafter\gdef\csname NUM@roomMELDAtwofive\endcsname{5.2}
\expandafter\gdef\csname NUM@roomMELDTfive\endcsname{7.9}
\expandafter\gdef\csname NUM@roomMELDTtwofive\endcsname{6.3}
\expandafter\gdef\csname NUM@roomMELDVfive\endcsname{1.8}
\expandafter\gdef\csname NUM@roomMELDVtwofive\endcsname{1.9}
\expandafter\gdef\csname NUM@roomMOSEIAfive\endcsname{11.4}
\expandafter\gdef\csname NUM@roomMOSEIAtwofive\endcsname{9.0}
\expandafter\gdef\csname NUM@roomMOSEITfive\endcsname{7.9}
\expandafter\gdef\csname NUM@roomMOSEITtwofive\endcsname{5.1}
\expandafter\gdef\csname NUM@roomMOSEIVfive\endcsname{8.7}
\expandafter\gdef\csname NUM@roomMOSEIVtwofive\endcsname{7.9}
\expandafter\gdef\csname NUM@savtAFEWfive\endcsname{70.8}
\expandafter\gdef\csname NUM@savtAFEWtwofive\endcsname{75.8}
\expandafter\gdef\csname NUM@savtIEMOCAPfive\endcsname{72.7}
\expandafter\gdef\csname NUM@savtIEMOCAPtwofive\endcsname{74.7}
\expandafter\gdef\csname NUM@savtMELDfive\endcsname{50.1}
\expandafter\gdef\csname NUM@savtMELDtwofive\endcsname{51.9}
\expandafter\gdef\csname NUM@savtMOSEIfive\endcsname{41.1}
\expandafter\gdef\csname NUM@savtMOSEItwofive\endcsname{44.1}
\expandafter\gdef\csname NUM@selcorpora\endcsname{AFEW，IEMOCAP，MELD，MOSEI}
\expandafter\gdef\csname NUM@sestaAFEWfive\endcsname{67.8}
\expandafter\gdef\csname NUM@sestaAFEWtwofive\endcsname{75.2}
\expandafter\gdef\csname NUM@sestaIEMOCAPfive\endcsname{70.1}
\expandafter\gdef\csname NUM@sestaIEMOCAPtwofive\endcsname{74.9}
\expandafter\gdef\csname NUM@sestaMELDfive\endcsname{46.3}
\expandafter\gdef\csname NUM@sestaMELDtwofive\endcsname{50.7}
\expandafter\gdef\csname NUM@sestaMOSEIfive\endcsname{38.3}
\expandafter\gdef\csname NUM@sestaMOSEItwofive\endcsname{43.0}
\expandafter\gdef\csname NUM@sestpAFEWfive\endcsname{71.6}
\expandafter\gdef\csname NUM@sestpAFEWtwofive\endcsname{76.6}
\expandafter\gdef\csname NUM@sestpIEMOCAPfive\endcsname{72.3}
\expandafter\gdef\csname NUM@sestpIEMOCAPtwofive\endcsname{75.0}
\expandafter\gdef\csname NUM@sestpMELDfive\endcsname{53.0}
\expandafter\gdef\csname NUM@sestpMELDtwofive\endcsname{54.8}
\expandafter\gdef\csname NUM@sestpMOSEIfive\endcsname{44.0}
\expandafter\gdef\csname NUM@sestpMOSEItwofive\endcsname{43.9}
\expandafter\gdef\csname NUM@sgridAFEWfive\endcsname{76.4}
\expandafter\gdef\csname NUM@sgridAFEWtwofive\endcsname{78.9}
\expandafter\gdef\csname NUM@sgridIEMOCAPfive\endcsname{73.4}
\expandafter\gdef\csname NUM@sgridIEMOCAPtwofive\endcsname{75.6}
\expandafter\gdef\csname NUM@sgridMELDfive\endcsname{51.4}
\expandafter\gdef\csname NUM@sgridMELDtwofive\endcsname{53.0}
\expandafter\gdef\csname NUM@sgridMOSEIfive\endcsname{40.0}
\expandafter\gdef\csname NUM@sgridMOSEItwofive\endcsname{43.6}
\expandafter\gdef\csname NUM@shareAFEWAfive\endcsname{+67}
\expandafter\gdef\csname NUM@shareAFEWAtwofive\endcsname{+78}
\expandafter\gdef\csname NUM@shareAFEWTfive\endcsname{+44}
\expandafter\gdef\csname NUM@shareAFEWTtwofive\endcsname{+26}
\expandafter\gdef\csname NUM@shareAFEWVfive\endcsname{+2}
\expandafter\gdef\csname NUM@shareAFEWVtwofive\endcsname{\textminus{}8}
\expandafter\gdef\csname NUM@shareIEMOCAPAfive\endcsname{+24}
\expandafter\gdef\csname NUM@shareIEMOCAPAtwofive\endcsname{+20}
\expandafter\gdef\csname NUM@shareIEMOCAPTfive\endcsname{+64}
\expandafter\gdef\csname NUM@shareIEMOCAPTtwofive\endcsname{+66}
\expandafter\gdef\csname NUM@shareIEMOCAPVfive\endcsname{\textminus{}18}
\expandafter\gdef\csname NUM@shareIEMOCAPVtwofive\endcsname{\textminus{}17}
\expandafter\gdef\csname NUM@shareMELDAfive\endcsname{+48}
\expandafter\gdef\csname NUM@shareMELDAtwofive\endcsname{+55}
\expandafter\gdef\csname NUM@shareMELDTfive\endcsname{\textminus{}15}
\expandafter\gdef\csname NUM@shareMELDTtwofive\endcsname{\textminus{}15}
\expandafter\gdef\csname NUM@shareMELDVfive\endcsname{+49}
\expandafter\gdef\csname NUM@shareMELDVtwofive\endcsname{+59}
\expandafter\gdef\csname NUM@shareMOSEIAfive\endcsname{+28}
\expandafter\gdef\csname NUM@shareMOSEIAtwofive\endcsname{+37}
\expandafter\gdef\csname NUM@shareMOSEITfive\endcsname{\textminus{}0}
\expandafter\gdef\csname NUM@shareMOSEITtwofive\endcsname{+0}
\expandafter\gdef\csname NUM@shareMOSEIVfive\endcsname{+2}
\expandafter\gdef\csname NUM@shareMOSEIVtwofive\endcsname{+33}
\expandafter\gdef\csname NUM@snAFEWfive\endcsname{50}
\expandafter\gdef\csname NUM@snAFEWtwofive\endcsname{50}
\expandafter\gdef\csname NUM@snIEMOCAPfive\endcsname{50}
\expandafter\gdef\csname NUM@snIEMOCAPtwofive\endcsname{50}
\expandafter\gdef\csname NUM@snMELDfive\endcsname{50}
\expandafter\gdef\csname NUM@snMELDtwofive\endcsname{50}
\expandafter\gdef\csname NUM@snMOSEIfive\endcsname{50}
\expandafter\gdef\csname NUM@snMOSEItwofive\endcsname{50}
\expandafter\gdef\csname NUM@soracleAFEWfive\endcsname{74.3}
\expandafter\gdef\csname NUM@soracleAFEWtwofive\endcsname{78.2}
\expandafter\gdef\csname NUM@soracleIEMOCAPfive\endcsname{73.1}
\expandafter\gdef\csname NUM@soracleIEMOCAPtwofive\endcsname{75.1}
\expandafter\gdef\csname NUM@soracleMELDfive\endcsname{50.6}
\expandafter\gdef\csname NUM@soracleMELDtwofive\endcsname{52.8}
\expandafter\gdef\csname NUM@soracleMOSEIfive\endcsname{41.1}
\expandafter\gdef\csname NUM@soracleMOSEItwofive\endcsname{44.1}
\expandafter\gdef\csname NUM@soracledAFEWfive\endcsname{+3.6}
\expandafter\gdef\csname NUM@soracledAFEWtwofive\endcsname{+2.4}
\expandafter\gdef\csname NUM@soracledIEMOCAPfive\endcsname{+0.4}
\expandafter\gdef\csname NUM@soracledIEMOCAPtwofive\endcsname{+0.5}
\expandafter\gdef\csname NUM@soracledMELDfive\endcsname{+0.5}
\expandafter\gdef\csname NUM@soracledMELDtwofive\endcsname{+0.9}
\expandafter\gdef\csname NUM@soracledMOSEIfive\endcsname{+0.0}
\expandafter\gdef\csname NUM@soracledMOSEItwofive\endcsname{+0.0}
\expandafter\gdef\csname NUM@sruleAFEWfive\endcsname{73.7}
\expandafter\gdef\csname NUM@sruleAFEWtwofive\endcsname{77.8}
\expandafter\gdef\csname NUM@sruleIEMOCAPfive\endcsname{72.0}
\expandafter\gdef\csname NUM@sruleIEMOCAPtwofive\endcsname{74.5}
\expandafter\gdef\csname NUM@sruleMELDfive\endcsname{50.0}
\expandafter\gdef\csname NUM@sruleMELDtwofive\endcsname{52.7}
\expandafter\gdef\csname NUM@sruleMOSEIfive\endcsname{38.3}
\expandafter\gdef\csname NUM@sruleMOSEItwofive\endcsname{43.1}
\expandafter\gdef\csname NUM@sruledAFEWfive\endcsname{+2.9}
\expandafter\gdef\csname NUM@sruledAFEWtwofive\endcsname{+2.0}
\expandafter\gdef\csname NUM@sruledIEMOCAPfive\endcsname{\textminus{}0.7}
\expandafter\gdef\csname NUM@sruledIEMOCAPtwofive\endcsname{\textminus{}0.2}
\expandafter\gdef\csname NUM@sruledMELDfive\endcsname{\textminus{}0.1}
\expandafter\gdef\csname NUM@sruledMELDtwofive\endcsname{+0.8}
\expandafter\gdef\csname NUM@sruledMOSEIfive\endcsname{\textminus{}2.8}
\expandafter\gdef\csname NUM@sruledMOSEItwofive\endcsname{\textminus{}1.0}
\expandafter\gdef\csname NUM@sselAFEWfive\endcsname{72.7}
\expandafter\gdef\csname NUM@sselAFEWtwofive\endcsname{77.2}
\expandafter\gdef\csname NUM@sselIEMOCAPfive\endcsname{72.3}
\expandafter\gdef\csname NUM@sselIEMOCAPtwofive\endcsname{74.7}
\expandafter\gdef\csname NUM@sselMELDfive\endcsname{50.1}
\expandafter\gdef\csname NUM@sselMELDtwofive\endcsname{52.5}
\expandafter\gdef\csname NUM@sselMOSEIfive\endcsname{39.3}
\expandafter\gdef\csname NUM@sselMOSEItwofive\endcsname{43.4}
\expandafter\gdef\csname NUM@sseldAFEWfive\endcsname{+1.9}
\expandafter\gdef\csname NUM@sseldAFEWtwofive\endcsname{+1.3}
\expandafter\gdef\csname NUM@sseldIEMOCAPfive\endcsname{\textminus{}0.4}
\expandafter\gdef\csname NUM@sseldIEMOCAPtwofive\endcsname{\textminus{}0.0}
\expandafter\gdef\csname NUM@sseldMELDfive\endcsname{+0.0}
\expandafter\gdef\csname NUM@sseldMELDtwofive\endcsname{+0.6}
\expandafter\gdef\csname NUM@sseldMOSEIfive\endcsname{\textminus{}1.8}
\expandafter\gdef\csname NUM@sseldMOSEItwofive\endcsname{\textminus{}0.7}
\expandafter\gdef\csname NUM@ssellooAFEWfive\endcsname{73.0}
\expandafter\gdef\csname NUM@ssellooAFEWtwofive\endcsname{76.4}
\expandafter\gdef\csname NUM@ssellooIEMOCAPfive\endcsname{72.2}
\expandafter\gdef\csname NUM@ssellooIEMOCAPtwofive\endcsname{74.7}
\expandafter\gdef\csname NUM@ssellooMELDfive\endcsname{50.2}
\expandafter\gdef\csname NUM@ssellooMELDtwofive\endcsname{52.1}
\expandafter\gdef\csname NUM@ssellooMOSEIfive\endcsname{39.2}
\expandafter\gdef\csname NUM@ssellooMOSEItwofive\endcsname{43.4}
\expandafter\gdef\csname NUM@sshareAFEWfive\endcsname{54}
\expandafter\gdef\csname NUM@sshareAFEWtwofive\endcsname{64}
\expandafter\gdef\csname NUM@sshareIEMOCAPfive\endcsname{48}
\expandafter\gdef\csname NUM@sshareIEMOCAPtwofive\endcsname{46}
\expandafter\gdef\csname NUM@sshareMELDfive\endcsname{68}
\expandafter\gdef\csname NUM@sshareMELDtwofive\endcsname{82}
\expandafter\gdef\csname NUM@sshareMOSEIfive\endcsname{70}
\expandafter\gdef\csname NUM@sshareMOSEItwofive\endcsname{56}
\expandafter\gdef\csname NUM@startAFEW\endcsname{64.5}
\expandafter\gdef\csname NUM@startIEMOCAP\endcsname{69.6}
\expandafter\gdef\csname NUM@startMELD\endcsname{47.2}
\expandafter\gdef\csname NUM@startMOSEI\endcsname{39.6}
\expandafter\gdef\csname NUM@suarAFEWAVTfive\endcsname{71.4}
\expandafter\gdef\csname NUM@suarAFEWAVTtwofive\endcsname{76.8}
\expandafter\gdef\csname NUM@suarAFEWAfive\endcsname{51.0}
\expandafter\gdef\csname NUM@suarAFEWAtwofive\endcsname{53.6}
\expandafter\gdef\csname NUM@suarAFEWTfive\endcsname{45.5}
\expandafter\gdef\csname NUM@suarAFEWTtwofive\endcsname{46.7}
\expandafter\gdef\csname NUM@suarAFEWVfive\endcsname{71.1}
\expandafter\gdef\csname NUM@suarAFEWVtwofive\endcsname{74.7}
\expandafter\gdef\csname NUM@suarIEMOCAPAVTfive\endcsname{74.1}
\expandafter\gdef\csname NUM@suarIEMOCAPAVTtwofive\endcsname{78.7}
\expandafter\gdef\csname NUM@suarIEMOCAPAfive\endcsname{68.1}
\expandafter\gdef\csname NUM@suarIEMOCAPAtwofive\endcsname{71.8}
\expandafter\gdef\csname NUM@suarIEMOCAPTfive\endcsname{50.3}
\expandafter\gdef\csname NUM@suarIEMOCAPTtwofive\endcsname{55.2}
\expandafter\gdef\csname NUM@suarIEMOCAPVfive\endcsname{57.8}
\expandafter\gdef\csname NUM@suarIEMOCAPVtwofive\endcsname{59.6}
\expandafter\gdef\csname NUM@suarMELDAVTfive\endcsname{48.2}
\expandafter\gdef\csname NUM@suarMELDAVTtwofive\endcsname{51.2}
\expandafter\gdef\csname NUM@suarMELDAfive\endcsname{42.3}
\expandafter\gdef\csname NUM@suarMELDAtwofive\endcsname{43.2}
\expandafter\gdef\csname NUM@suarMELDTfive\endcsname{50.3}
\expandafter\gdef\csname NUM@suarMELDTtwofive\endcsname{52.0}
\expandafter\gdef\csname NUM@suarMELDVfive\endcsname{38.0}
\expandafter\gdef\csname NUM@suarMELDVtwofive\endcsname{38.3}
\expandafter\gdef\csname NUM@suarMOSEIAVTfive\endcsname{41.3}
\expandafter\gdef\csname NUM@suarMOSEIAVTtwofive\endcsname{44.1}
\expandafter\gdef\csname NUM@suarMOSEIAfive\endcsname{34.6}
\expandafter\gdef\csname NUM@suarMOSEIAtwofive\endcsname{37.5}
\expandafter\gdef\csname NUM@suarMOSEITfive\endcsname{33.4}
\expandafter\gdef\csname NUM@suarMOSEITtwofive\endcsname{36.2}
\expandafter\gdef\csname NUM@suarMOSEIVfive\endcsname{36.4}
\expandafter\gdef\csname NUM@suarMOSEIVtwofive\endcsname{37.9}
\expandafter\gdef\csname NUM@totAFEWAVTnonefive\endcsname{10.0}
\expandafter\gdef\csname NUM@totAFEWAVTnonetwofive\endcsname{4.6}
\expandafter\gdef\csname NUM@totAFEWAVTpropfive\endcsname{9.0}
\expandafter\gdef\csname NUM@totAFEWAVTproptwofive\endcsname{4.0}
\expandafter\gdef\csname NUM@totAFEWAVTstart\endcsname{15.2}
\expandafter\gdef\csname NUM@totAFEWAnonefive\endcsname{6.1}
\expandafter\gdef\csname NUM@totAFEWAnonetwofive\endcsname{4.3}
\expandafter\gdef\csname NUM@totAFEWApropfive\endcsname{2.0}
\expandafter\gdef\csname NUM@totAFEWAproptwofive\endcsname{0.9}
\expandafter\gdef\csname NUM@totAFEWAstart\endcsname{7.7}
\expandafter\gdef\csname NUM@totAFEWTnonefive\endcsname{3.7}
\expandafter\gdef\csname NUM@totAFEWTnonetwofive\endcsname{2.9}
\expandafter\gdef\csname NUM@totAFEWTpropfive\endcsname{2.0}
\expandafter\gdef\csname NUM@totAFEWTproptwofive\endcsname{2.1}
\expandafter\gdef\csname NUM@totAFEWTstart\endcsname{4.4}
\expandafter\gdef\csname NUM@totAFEWVnonefive\endcsname{4.9}
\expandafter\gdef\csname NUM@totAFEWVnonetwofive\endcsname{2.4}
\expandafter\gdef\csname NUM@totAFEWVpropfive\endcsname{4.8}
\expandafter\gdef\csname NUM@totAFEWVproptwofive\endcsname{2.5}
\expandafter\gdef\csname NUM@totAFEWVstart\endcsname{9.7}
\expandafter\gdef\csname NUM@totIEMOCAPAVTnonefive\endcsname{11.1}
\expandafter\gdef\csname NUM@totIEMOCAPAVTnonetwofive\endcsname{8.1}
\expandafter\gdef\csname NUM@totIEMOCAPAVTpropfive\endcsname{10.9}
\expandafter\gdef\csname NUM@totIEMOCAPAVTproptwofive\endcsname{9.0}
\expandafter\gdef\csname NUM@totIEMOCAPAVTstart\endcsname{14.0}
\expandafter\gdef\csname NUM@totIEMOCAPAnonefive\endcsname{10.1}
\expandafter\gdef\csname NUM@totIEMOCAPAnonetwofive\endcsname{7.2}
\expandafter\gdef\csname NUM@totIEMOCAPApropfive\endcsname{7.7}
\expandafter\gdef\csname NUM@totIEMOCAPAproptwofive\endcsname{5.7}
\expandafter\gdef\csname NUM@totIEMOCAPAstart\endcsname{12.7}
\expandafter\gdef\csname NUM@totIEMOCAPTnonefive\endcsname{11.5}
\expandafter\gdef\csname NUM@totIEMOCAPTnonetwofive\endcsname{7.6}
\expandafter\gdef\csname NUM@totIEMOCAPTpropfive\endcsname{4.2}
\expandafter\gdef\csname NUM@totIEMOCAPTproptwofive\endcsname{2.6}
\expandafter\gdef\csname NUM@totIEMOCAPTstart\endcsname{12.2}
\expandafter\gdef\csname NUM@totIEMOCAPVnonefive\endcsname{4.7}
\expandafter\gdef\csname NUM@totIEMOCAPVnonetwofive\endcsname{4.5}
\expandafter\gdef\csname NUM@totIEMOCAPVpropfive\endcsname{5.5}
\expandafter\gdef\csname NUM@totIEMOCAPVproptwofive\endcsname{5.2}
\expandafter\gdef\csname NUM@totIEMOCAPVstart\endcsname{4.1}
\expandafter\gdef\csname NUM@totMELDAVTnonefive\endcsname{11.3}
\expandafter\gdef\csname NUM@totMELDAVTnonetwofive\endcsname{8.4}
\expandafter\gdef\csname NUM@totMELDAVTpropfive\endcsname{8.8}
\expandafter\gdef\csname NUM@totMELDAVTproptwofive\endcsname{7.0}
\expandafter\gdef\csname NUM@totMELDAVTstart\endcsname{11.7}
\expandafter\gdef\csname NUM@totMELDAnonefive\endcsname{6.5}
\expandafter\gdef\csname NUM@totMELDAnonetwofive\endcsname{5.2}
\expandafter\gdef\csname NUM@totMELDApropfive\endcsname{3.4}
\expandafter\gdef\csname NUM@totMELDAproptwofive\endcsname{2.4}
\expandafter\gdef\csname NUM@totMELDAstart\endcsname{6.7}
\expandafter\gdef\csname NUM@totMELDTnonefive\endcsname{7.9}
\expandafter\gdef\csname NUM@totMELDTnonetwofive\endcsname{6.3}
\expandafter\gdef\csname NUM@totMELDTpropfive\endcsname{9.1}
\expandafter\gdef\csname NUM@totMELDTproptwofive\endcsname{7.3}
\expandafter\gdef\csname NUM@totMELDTstart\endcsname{7.2}
\expandafter\gdef\csname NUM@totMELDVnonefive\endcsname{1.8}
\expandafter\gdef\csname NUM@totMELDVnonetwofive\endcsname{1.9}
\expandafter\gdef\csname NUM@totMELDVpropfive\endcsname{0.9}
\expandafter\gdef\csname NUM@totMELDVproptwofive\endcsname{0.8}
\expandafter\gdef\csname NUM@totMELDVstart\endcsname{1.7}
\expandafter\gdef\csname NUM@totMOSEIAVTnonefive\endcsname{10.2}
\expandafter\gdef\csname NUM@totMOSEIAVTnonetwofive\endcsname{8.3}
\expandafter\gdef\csname NUM@totMOSEIAVTpropfive\endcsname{10.5}
\expandafter\gdef\csname NUM@totMOSEIAVTproptwofive\endcsname{7.6}
\expandafter\gdef\csname NUM@totMOSEIAVTstart\endcsname{12.0}
\expandafter\gdef\csname NUM@totMOSEIAnonefive\endcsname{11.4}
\expandafter\gdef\csname NUM@totMOSEIAnonetwofive\endcsname{9.0}
\expandafter\gdef\csname NUM@totMOSEIApropfive\endcsname{8.2}
\expandafter\gdef\csname NUM@totMOSEIAproptwofive\endcsname{5.7}
\expandafter\gdef\csname NUM@totMOSEIAstart\endcsname{12.0}
\expandafter\gdef\csname NUM@totMOSEITnonefive\endcsname{7.9}
\expandafter\gdef\csname NUM@totMOSEITnonetwofive\endcsname{5.1}
\expandafter\gdef\csname NUM@totMOSEITpropfive\endcsname{7.9}
\expandafter\gdef\csname NUM@totMOSEITproptwofive\endcsname{5.1}
\expandafter\gdef\csname NUM@totMOSEITstart\endcsname{8.8}
\expandafter\gdef\csname NUM@totMOSEIVnonefive\endcsname{8.7}
\expandafter\gdef\csname NUM@totMOSEIVnonetwofive\endcsname{7.9}
\expandafter\gdef\csname NUM@totMOSEIVpropfive\endcsname{8.5}
\expandafter\gdef\csname NUM@totMOSEIVproptwofive\endcsname{5.3}
\expandafter\gdef\csname NUM@totMOSEIVstart\endcsname{8.8}
\expandafter\gdef\csname NUM@tuarAFEWAVTfive\endcsname{72.7}
\expandafter\gdef\csname NUM@tuarAFEWAVTtwofive\endcsname{77.8}
\expandafter\gdef\csname NUM@tuarAFEWAfive\endcsname{72.7}
\expandafter\gdef\csname NUM@tuarAFEWAtwofive\endcsname{77.8}
\expandafter\gdef\csname NUM@tuarAFEWTfive\endcsname{72.7}
\expandafter\gdef\csname NUM@tuarAFEWTtwofive\endcsname{77.8}
\expandafter\gdef\csname NUM@tuarAFEWVfive\endcsname{72.7}
\expandafter\gdef\csname NUM@tuarAFEWVtwofive\endcsname{77.8}
\expandafter\gdef\csname NUM@tuarIEMOCAPAVTfive\endcsname{75.4}
\expandafter\gdef\csname NUM@tuarIEMOCAPAVTtwofive\endcsname{78.8}
\expandafter\gdef\csname NUM@tuarIEMOCAPAfive\endcsname{75.4}
\expandafter\gdef\csname NUM@tuarIEMOCAPAtwofive\endcsname{78.8}
\expandafter\gdef\csname NUM@tuarIEMOCAPTfive\endcsname{75.4}
\expandafter\gdef\csname NUM@tuarIEMOCAPTtwofive\endcsname{78.8}
\expandafter\gdef\csname NUM@tuarIEMOCAPVfive\endcsname{75.4}
\expandafter\gdef\csname NUM@tuarIEMOCAPVtwofive\endcsname{78.8}
\expandafter\gdef\csname NUM@tuarMELDAVTfive\endcsname{50.6}
\expandafter\gdef\csname NUM@tuarMELDAVTtwofive\endcsname{52.7}
\expandafter\gdef\csname NUM@tuarMELDAfive\endcsname{50.6}
\expandafter\gdef\csname NUM@tuarMELDAtwofive\endcsname{52.7}
\expandafter\gdef\csname NUM@tuarMELDTfive\endcsname{50.6}
\expandafter\gdef\csname NUM@tuarMELDTtwofive\endcsname{52.7}
\expandafter\gdef\csname NUM@tuarMELDVfive\endcsname{50.6}
\expandafter\gdef\csname NUM@tuarMELDVtwofive\endcsname{52.7}
\expandafter\gdef\csname NUM@tuarMOSEIAVTfive\endcsname{41.4}
\expandafter\gdef\csname NUM@tuarMOSEIAVTtwofive\endcsname{45.0}
\expandafter\gdef\csname NUM@tuarMOSEIAfive\endcsname{41.4}
\expandafter\gdef\csname NUM@tuarMOSEIAtwofive\endcsname{45.0}
\expandafter\gdef\csname NUM@tuarMOSEITfive\endcsname{41.4}
\expandafter\gdef\csname NUM@tuarMOSEITtwofive\endcsname{45.0}
\expandafter\gdef\csname NUM@tuarMOSEIVfive\endcsname{41.4}
\expandafter\gdef\csname NUM@tuarMOSEIVtwofive\endcsname{45.0}
\expandafter\gdef\csname NUM@upperAFEW\endcsname{79.8}
\expandafter\gdef\csname NUM@upperIEMOCAP\endcsname{83.6}
\expandafter\gdef\csname NUM@upperMELD\endcsname{58.9}
\expandafter\gdef\csname NUM@upperMOSEI\endcsname{51.7}


% ---------- 言語 ----------
\newcommand{\T}[2]{\iflangja #1\else #2\fi}
\newcommand{\N}[1]{\ifcsname NUM@#1\endcsname\csname NUM@#1\endcsname\else\PackageError{numbers}{no number #1}{}\fi}
\renewcommand{\hosoku}[2]{\gdef\hosokutext{\textcolor{mDarkTeal}{\T{補足}{Appendix} p.~\pageref{#1}\T{：}{: }#2}}}
\renewcommand{\toi}[1]{\band{cAsk}{\T{問い}{Question}}{#1}}
\renewcommand{\imi}[1]{\band{mDarkTeal}{\T{この結果の意味}{What this means}}{#1}}
\newcommand{\notreal}{\par{\tiny\color{mMuted}\T{模式図．実験の値ではない}{Schematic; not experimental values}}\par}
\newcommand{\pt}{\T{ポイント}{pt}}
\newcommand{\interim}{\textcolor{mLightBrown}{\textbf{\T{途中}{interim}}}}
\newcommand{\emo}{\T{感情}{emotion}}

% ---------- 色（手法と種類） ----------
\definecolor{cNone}{HTML}{8A8985}    % K 本だけ（比較の基準）
\definecolor{cProp}{HTML}{2A78D6}    % 提案（別のモダリティの予測で学ぶ）
\definecolor{cKD}{HTML}{6250D6}      % 温度つき蒸留
\definecolor{cUpper}{HTML}{1B7F5C}   % 上限
\definecolor{kA}{HTML}{2A78D6}\definecolor{kV}{HTML}{EB6834}\definecolor{kT}{HTML}{1BAF7A}\definecolor{kAVT}{HTML}{6250D6}
\definecolor{cK0}{HTML}{B0B0B0}\definecolor{cK5}{HTML}{2A78D6}\definecolor{cK25}{HTML}{0B3D6E}
\definecolor{cAgree}{HTML}{1B7F5C}\definecolor{cDis}{HTML}{B3312B}
\pgfplotsset{
  sK0/.style={cK0, mark=*, mark size=1.6pt, only marks, error bars/.cd, y dir=both, y explicit, error bar style={cK0, line width=0.7pt}},
  sK5/.style={cK5, mark=square*, mark size=1.6pt, only marks, error bars/.cd, y dir=both, y explicit, error bar style={cK5, line width=0.7pt}},
  sK25/.style={cK25, mark=triangle*, mark size=2.0pt, only marks, error bars/.cd, y dir=both, y explicit, error bar style={cK25, line width=0.7pt}},
}
\tikzset{
  lblbox/.style={draw=mDarkTeal!70, fill=cTgt, rounded corners=2pt, align=center, font=\scriptsize, inner sep=3pt},
  unlbox/.style={draw=mDarkTeal!70, fill=white, rounded corners=2pt, align=center, font=\scriptsize, inner sep=3pt},
  propbox/.style={draw=cProp, fill=cProp!10, rounded corners=2pt, align=center, font=\scriptsize, inner sep=3pt},
  kdbox/.style={draw=cKD, fill=cKD!10, rounded corners=2pt, align=center, font=\scriptsize, inner sep=3pt},
  utt/.style={draw=mDarkTeal!60, fill=white, minimum width=5mm, minimum height=5mm, font=\tiny, inner sep=1pt},
}

% ---------- 発表情報 ----------
\title{\T{少数のターゲットラベルで学ぶマルチモーダル感情認識}{Multimodal Emotion Recognition with Few Target Labels}}
\subtitle{\T{別のモダリティの予測で学ぶ手法の測定と説明 ── 途中経過と，いま進めていること}{Measuring and explaining learning from the other modalities' predictions --- interim results and current work}}
\author{\T{杉浦 伸一}{Shinichi Sugiura}}
\institute[\T{北陸先端大学}{JAIST}]{\T{北陸先端大学}{Japan Advanced Institute of Science and Technology}}
\date{\T{2026年11月5日}{November 5, 2026}}

\begin{document}
\maketitle

%% 1. 目的と設定（骨組みの 1・2）
\setband{\T{\S1 目的と設定}{\S1 Goal and setting}}

\begin{frame}{\T{なぜ重要か：新しいコーパスでは，ラベル付きの発話が感情ごとに数本〜数十本しか用意できない}{Why it matters: a new corpus comes with only a handful of labeled utterances per emotion}}
\rl{\T{目的（骨組み 1）}{Goal (skeleton 1)}}
\begin{columns}[T]
\begin{column}{0.60\textwidth}
\centering
\begin{adjustbox}{max width=\linewidth}
\begin{tikzpicture}[font=\scriptsize, x=1cm, y=1cm]
  % ソース
  \foreach \n/\x in {AFEW/0, IEMOCAP/2.05, MELD/4.1} {
    \node[srcbox, minimum width=1.9cm, minimum height=1.15cm] (s\n) at (\x,0) {\n\\ \T{全発話に}{all utterances}\\ \T{ラベル}{labeled}};
  }
  \node[font=\scriptsize\bfseries, color=mDarkTeal] at (2.05,1.0) {\T{ソース（3 つ）}{Sources (3)}};
  % ターゲット
  \node[draw=mDarkTeal!70, fill=cTgt!50, rounded corners=3pt, minimum width=3.0cm, minimum height=2.6cm] (tgt) at (7.2,-0.35) {};
  \node[font=\scriptsize\bfseries, color=mDarkTeal] at (7.2,1.25) {\T{ターゲット（1 つ）}{Target (1)}: MOSEI};
  \node[lblbox, minimum width=2.6cm] (k) at (7.2,0.45) {\T{K 本 × 4 感情（ラベル付き）}{K × 4 emotions (labeled)}};
  \node[unlbox, minimum width=2.6cm] (u) at (7.2,-0.35) {\T{ラベル無しの発話（候補）}{unlabeled utterances (pool)}};
  \node[testbox, minimum width=2.6cm] (t) at (7.2,-1.2) {\T{テスト（別の話者）}{test (other speakers)}};
  % 矢印
  \draw[arr] (4.1,-0.75) -- node[below, font=\tiny, align=center]{\T{識別器を作って持ち込む}{build classifiers,\\carry them over}} (5.7,-0.75);
  % 回す
  \node[font=\tiny, color=mMuted, align=center] at (3.6,-1.75) {\T{4 コーパスを順にターゲットにする（leave-one-corpus-out）}{each corpus takes the target role in turn (leave-one-corpus-out)}};
\end{tikzpicture}
\end{adjustbox}
\end{column}
\begin{column}{0.40\textwidth}
\small
\T{%
\begin{itemize}
\item 感情認識のラベル付けは高く，新しい場面（映画，対話，動画）ごとにコーパスを作れない
\item そこで，\textbf{ほかのコーパスの識別器}と，\textbf{ターゲットのラベル無しの発話}で，K 本だけの識別器を補いたい
\item 音声・顔・言葉の 3 モダリティがそろう 4 コーパスで測る
\end{itemize}
}{%
\begin{itemize}
\item Emotion labels are expensive; one cannot build a corpus for every new setting (films, dialogues, videos)
\item So we want to improve a classifier trained on K labels using \textbf{classifiers from other corpora} and \textbf{unlabeled target utterances}
\item Measured on 4 corpora where audio, face and text are all available
\end{itemize}
}
\end{column}
\end{columns}
\takeaway{\T{問い：ソースの識別器とラベル無しの発話は，K 本だけの識別器をどれだけ，どこで，なぜ良くするか}{Question: how much, where, and why do source classifiers and unlabeled utterances improve a classifier trained on K labels?}}
\end{frame}

\begin{frame}{\T{設定：凍結した特徴の上の線形の識別器を，4 つのコーパスで順に良くする}{Setting: linear classifiers on frozen features, each corpus in turn as the target}}
\rl{\T{設定}{Setting}}
\hosoku{appdata}{\T{コーパスの詳細（発話数，塊，分割）}{corpus details (utterances, clusters, folds)}}
\begin{columns}[T]
\begin{column}{0.55\textwidth}
\centering
\begin{adjustbox}{max width=\linewidth}
\begin{tikzpicture}[font=\scriptsize, x=1cm, y=1cm]
  \node[encbox, minimum width=1.5cm] (ea) at (0,1.2) {\T{音声}{audio}\\WavLM};
  \node[encbox, minimum width=1.5cm] (ev) at (0,0.4) {\T{顔}{face}\\ResNet-50};
  \node[encbox, minimum width=1.5cm] (et) at (0,-0.4) {\T{言葉}{text}\\RoBERTa};
  \node[font=\tiny, color=mMuted, align=center] at (0,-1.1) {\T{凍結（微調整しない）}{frozen (no fine-tuning)}};
  \node[clf, minimum width=1.7cm] (ca) at (2.6,1.2) {\T{音声の識別器}{audio classifier}};
  \node[clf, minimum width=1.7cm] (cv) at (2.6,0.4) {\T{顔の識別器}{face classifier}};
  \node[clf, minimum width=1.7cm] (ct) at (2.6,-0.4) {\T{言葉の識別器}{text classifier}};
  \node[clf, minimum width=2.1cm, fill=kAVT!10] (cf) at (2.6,-1.35) {AVT\T{の識別器}{ classifier}\\ \T{（3 つの特徴をつないだ上）}{(on concatenated features)}};
  \foreach \e/\c in {ea/ca, ev/cv, et/ct} {\draw[arr] (\e) -- (\c);}
  \draw[arr] (0.8,-0.7) -- (1.5,-1.35) -- (cf);
  \node[font=\scriptsize, align=left, anchor=west] at (-0.9,-2.35) {\T{どれも線形（ロジスティック回帰）}{all linear (logistic regression)}\\ \T{4 感情：Neutral，Happy，Sad，Angry}{4 emotions: Neutral, Happy, Sad, Angry}\\ \T{良くしたい識別器＝この 4 つ（A，V，T，AVT）}{classifiers to improve = these 4 (A, V, T, AVT)}};
\end{tikzpicture}
\end{adjustbox}
\end{column}
\begin{column}{0.45\textwidth}
\scriptsize
\begin{tabular}{@{}lrr@{}}
\toprule
\T{コーパス}{Corpus} & \T{発話}{Utterances} & \T{話者の塊}{Clusters} \\
\midrule
AFEW (\T{映画}{films}) & 790 & 51 \\
IEMOCAP (\T{即興の対話}{improvised}) & 2,943 & 5 \T{セッション}{sessions} \\
MELD (\T{ドラマ}{TV series}) & 11,267 & 272 \\
MOSEI (\T{動画}{online videos}) & 18,634 & 2,614 \\
\bottomrule
\end{tabular}

\vspace{0.4em}
\T{%
\begin{itemize}
\item 話者を 5 つに分け，1 つをテストにする．テストの話者の発話は学習側に入れない
\item K ∈ \{0, 5, 25\}．K 本の抜き方を 10 種類変える → コーパスごとに 50 単位
\item 手法の形と設定は，評価の前に「仮のターゲット」（ソースの 1 つをターゲットの代わりにする回）で選んで固定し，評価の回は 1 回だけ流した
\end{itemize}
}{%
\begin{itemize}
\item Speakers are split into 5 folds; one fold is the test. Test speakers never enter training
\item K ∈ \{0, 5, 25\}; 10 random draws of the K labels → 50 units per corpus
\item Method forms and settings were chosen and frozen before evaluation on ``proxy targets'' (one source standing in for the target); the evaluation round was run once
\end{itemize}
}
\end{column}
\end{columns}
\end{frame}

\begin{frame}{\T{どれほど難しいか：ソースの識別器の平均は，最良の線形識別器より 12〜15 ポイント低い}{How hard: the mean of source classifiers is 12--15 points below the best linear classifier}}
\rl{\T{難しさ（骨組み 2）}{Difficulty (skeleton 2)}}
\hosoku{appdecomp}{\T{損の内訳の表（全生徒）}{loss decomposition tables (all students)}}
\begin{columns}[T]
\begin{column}{0.56\textwidth}
\begin{tikzpicture}
\begin{axis}[width=\linewidth, height=5.6cm, ybar, bar width=8pt, ymin=30, ymax=118, enlarge x limits=0.18,
  xtick={1,2,3,4}, xticklabels={AFEW, IEMOCAP, MELD, MOSEI}, ylabel={\T{UAR（AVT の生徒，\%）}{UAR (AVT student, \%)}},
  legend style={at={(0.98,0.98)}, anchor=north east, font=\tiny, cells={anchor=west}}, legend columns=1,
  nodes near coords={\pgfmathprintnumber[fixed, precision=0]{\pgfplotspointmeta}}, every node near coord/.append style={font=\tiny}, ymajorgrids]
\addplot[fill=cNone, draw=cNone] table[x=ci, y=start]{difficulty.csv};
\addplot[fill=cK25, draw=cK25] table[x=ci, y=k25]{difficulty.csv};
\addplot[fill=cUpper, draw=cUpper] table[x=ci, y=upper]{difficulty.csv};
\legend{\T{出発点（ソースの平均，K = 0）}{start (sources' mean, K = 0)}, \T{K = 25 で学び直し}{retrained, K = 25}, \T{上限（候補の全ラベルで学んだ重み）}{upper bound (all pool labels)}}
\end{axis}
\end{tikzpicture}
\end{column}
\begin{column}{0.44\textwidth}
\footnotesize
\T{%
\begin{itemize}
\item 出発点と上限の差：AFEW \N{gapstartAFEW}，IEMOCAP \N{gapstartIEMOCAP}，MELD \N{gapstartMELD}，MOSEI \N{gapstartMOSEI} ポイント
\item K = 25（100 本）で学び直しても，差は \N{gapkAFEW}〜\N{gapkMELD} ポイント残る
\item 差の大半は \textbf{重みの向き}のずれで，切片のずれは 0.4〜2.7 ポイント（\S4）．少数のラベルでは 2,304 次元の重みを直し切れない
\item 標準のラベル無しの手法 14 種（自己学習，DANN，CORAL など）を足しても，AVT では 1〜3 ポイントしか埋まらない（\S3）
\end{itemize}
}{%
\begin{itemize}
\item Gap between start and upper bound: AFEW \N{gapstartAFEW}, IEMOCAP \N{gapstartIEMOCAP}, MELD \N{gapstartMELD}, MOSEI \N{gapstartMOSEI} points
\item After retraining with K = 25 (100 labels), \N{gapkAFEW}--\N{gapkMELD} points remain
\item Most of the gap is a \textbf{weight-direction} error; the intercept part is 0.4--2.7 points (\S4). A few labels cannot fix 2,304 weights
\item Adding 14 standard unlabeled-data methods (self-training, DANN, CORAL, …) closes only 1--3 points for AVT (\S3)
\end{itemize}
}
\end{column}
\end{columns}
{\tiny\color{mMuted}\T{上限＝候補の全ラベルで学んだ重みに，テストの感情の釣り合いで切片を当て直した識別器（点の値，区間なし）．出どころ：figures/difficulty.csv}{Upper bound = weights trained on all pool labels with the intercept refit on the test emotion balance (point values, no interval). Source: figures/difficulty.csv}}
\end{frame}

\begin{frame}{\T{実験の全体像：4 つの実験で「どれだけ・どこで・なぜ」を分けて測る}{Overview: four experiments separate ``how much'', ``where'' and ``why''}}
\rl{\T{計画}{Plan}}
\scriptsize
\begin{tabular}{@{}p{0.07\linewidth}p{0.42\linewidth}p{0.27\linewidth}p{0.17\linewidth}@{}}
\toprule
& \T{何を測るか}{What is measured} & \T{相手}{Compared against} & \T{状態}{Status} \\
\midrule
\T{実験 1}{Exp.~1} & \T{別のモダリティの予測（疑似ラベル）で学ぶ手法の得}{Gain from learning on the other modalities' predictions (pseudo-labels)} & \T{K 本だけ；同じ情報を受け取る 5 手法（温度つき蒸留など）}{K-only; 5 methods using the same information (e.g.\ temperature distillation)} & \done \\
\T{実験 2}{Exp.~2} & \T{標準のドメイン適応・半教師あり手法との差}{Difference from standard DA / semi-supervised methods} & \T{自己学習，エントロピー最小化，T3A，ラベル伝播，S+T，CORAL，DANN，MME など 14}{self-training, entropy min., T3A, label propagation, S+T, CORAL, DANN, MME, … (14)} & \done \\
\T{実験 3}{Exp.~3} & \T{得の説明：損を切片（位置）と重みの向きに分ける，再現できる度合い}{Explanation: split the loss into intercept (position) and weight direction; reproducibility} & \T{事前に固定した判定の規則}{pre-fixed decision rules} & \done \\
\T{実験 4}{Exp.~4} & \T{蒸留の弱点（全発話を同じ重みで信じる）を直した手法の得}{Gain from fixing a weakness of distillation (uniform trust in all pool utterances)} & \T{実験 1 の提案・蒸留}{Exp.~1 proposal / distillation} & \todo\ \T{設計済み}{designed} \\
\T{探索}{Exploration} & \T{どのモダリティを信じるか（信頼度の見積もり）}{Which modality to trust (reliability estimation)} & --- & \done\ \T{測定済み}{measured} \\
\bottomrule
\end{tabular}

\vspace{0.3em}
\scriptsize
\T{実験 1〜3 は設計・台本・集計の台本を評価の前に固定し，関門（統計・漏れ・比較・定義・台本）と ML・感情認識の査読想定の検査を通してから計算した．実験 4 と探索は評価の結果を見た後の案で，論文では節を分け，限界に由来を書く．}{Experiments 1--3: designs, scripts and aggregation were frozen before evaluation and passed the gate checks (statistics, leakage, fairness, definitions, code) and simulated ML / affective-computing review. Experiment 4 and the exploration arose after seeing the results; the paper separates them and states their origin.}
\end{frame}

%% 2. 手法（骨組みの 3・5・6）
\setband{\T{\S2 手法と比較}{\S2 Methods and comparisons}}

\begin{frame}{\T{既存研究の不足：「上がるか」だけを測り，どのモダリティで・損のどの部分が減るかを測らない}{What existing work misses: it measures ``does it help'' but not for which modality or which part of the loss}}
\rl{\T{既存研究（骨組み 3）}{Prior work (skeleton 3)}}
\centering
\begin{tikzpicture}[font=\scriptsize, x=1cm, y=1cm]
  % (a)
  \node[qbox, text width=3.9cm, minimum height=2.5cm, anchor=north] (a) at (0,0) {\textbf{(a) \T{同じモダリティのラベル無しの手法}{Unlabeled-data methods within one modality}}\\[2pt]
    \T{自己学習，エントロピー最小化，T3A，ラベル伝播，DANN，CORAL，MME …}{self-training, entropy min., T3A, label propagation, DANN, CORAL, MME, …}\\[3pt]
    \T{疑似ラベルは\textbf{生徒自身}の予測か，特徴の分布合わせ．生徒より当たる予測は入らない}{Pseudo-labels come from the \textbf{student itself}, or features are aligned. No prediction better than the student enters}};
  % (b)
  \node[qbox, text width=3.9cm, minimum height=2.5cm, anchor=north] (b) at (4.6,0) {\textbf{(b) \T{別のモダリティ 1 つの予測}{One other modality's prediction}}\\[2pt]
    \T{Albanie ほか 2018：顔の識別器の予測で音声の識別器を学ぶ}{Albanie et al.\ 2018: the face classifier teaches the audio classifier}\\[3pt]
    \T{1 つのモダリティの予測は弱いことがあり，学ぶ側を悪くしうる（TelME 2024：言葉が音声から学ぶと 66.6 → 57.4）}{A single modality can be weak and hurt the learner (TelME 2024: text learning from audio, 66.6 → 57.4)}};
  % (c)
  \node[propbox, text width=3.9cm, minimum height=2.5cm, anchor=north] (c) at (9.2,0) {\textbf{(c) \T{この研究}{This work}}\\[2pt]
    \T{3 モダリティの識別器を学び直し，確率の\textbf{積}を疑似ラベルにする．K 本と出発点への引き戻しを足す}{Retrain the 3 unimodal classifiers and use the \textbf{product} of their probabilities as pseudo-labels, plus K labels and a pull toward the start}\\[3pt]
    \T{そして\textbf{測る}：どのモダリティの生徒で，損のどの部分（切片か重みの向きか）が減るか}{And \textbf{measure}: for which student, and which part of the loss (intercept or weight direction) shrinks}};
\end{tikzpicture}

\vspace{0.4em}
\takeaway{\T{骨組み 3 の文：既存の評価は，ほかのモダリティの予測で学ぶと成績が上がるかだけを測り，どのモダリティで，損のどの部分（切片か重みの向きか）が減るかを測らない}{Skeleton item 3: existing evaluations only ask whether learning from other modalities' predictions helps, not for which modality or which part of the loss (intercept or weight direction) shrinks}}
\end{frame}

\begin{frame}{\T{提案の流れ：ソースの識別器を出発点にし，3 モダリティの積の予測と K 本で生徒を学ぶ}{Proposal: start from the source classifiers; train the student on K labels plus the product of three modalities' predictions}}
\rl{\T{手法}{Method}}
\hosoku{appsolver}{\T{学習器と引き戻しの詳細}{solver and pull details}}
\centering
\begin{tikzpicture}[font=\scriptsize, x=1cm, y=1cm]
  % 行 A：出発点
  \node[srcbox, text width=3.6cm, minimum height=0.9cm] (src) at (0,0) {\T{ソースのコーパスごとの識別器}{Per-source classifiers}\\ \T{（3 コーパス × A・V・T・AVT）}{(3 corpora × A, V, T, AVT)}};
  \node[clf, text width=2.6cm, minimum height=0.9cm] (start) at (4.3,0) {\T{出発点}{Start}\\ \T{＝重みと切片の平均}{= mean of weights and intercepts}};
  \draw[arr] (src) -- node[above, font=\tiny]{\T{平均}{average}} (start);
  \node[font=\tiny, color=mMuted, anchor=west] at (6.0,0) {\T{すべての手法が同じ出発点を使う}{every method uses the same start}};
  % 行 B：教師
  \node[unlbox, text width=2.9cm, minimum height=1.0cm] (avtk) at (0.35,-1.7) {\T{出発点の AVT を}{AVT start, retrained}\\ \T{K 本で学び直す → 候補に予測}{with K labels → predict the pool}};
  \node[unlbox, text width=3.3cm, minimum height=1.0cm] (ret) at (4.3,-1.7) {\T{A・V・T の識別器を出発点から}{Retrain A, V, T from the start on}\\ \T{K 本 ＋ その予測で学び直す}{K labels + those predictions}};
  \node[modbox, text width=3.3cm, minimum height=1.0cm] (q) at (8.6,-1.7) {\T{\textbf{教師}：3 つの確率の積}{\textbf{Teacher}: product of the 3 probabilities}\\ \T{→ 候補の各発話に疑似ラベル Q}{→ pseudo-label Q for each pool utterance}};
  \draw[arr] (avtk) -- (ret); \draw[arr] (ret) -- node[above, font=\tiny]{\T{掛ける}{multiply}} (q);
  \draw[arr] (start.south) -- ++(0,-0.35) -| (ret.north);
  % 行 C：生徒
  \node[propbox, text width=8.6cm, minimum height=1.25cm] (stu) at (4.3,-3.6) {\T{\textbf{生徒}（A／V／T／AVT）を出発点から学ぶ．損失 ＝}{\textbf{Student} (A / V / T / AVT), trained from the start. Loss =}\\
    \colorbox{cTgt}{\T{K 本の正解との交差エントロピー}{cross-entropy on the K labels}} ＋
    \colorbox{cMod}{\T{疑似ラベル Q との交差エントロピー（感情ごとに総量をそろえる）}{cross-entropy to Q (total weight balanced per emotion)}} ＋
    \colorbox{cSrc}{\T{出発点への引き戻し}{pull toward the start}}};
  \draw[arr] (q.south) |- ++(-1.5,-0.55) -| (stu.north east);
  \node[lblbox, text width=1.6cm] (k) at (-1.1,-3.6) {K \T{本}{labels}};
  \draw[arr] (k) -- (stu);
  \node[testbox, text width=1.4cm] (test) at (10.5,-3.6) {\T{テストで UAR}{UAR on test}};
  \draw[arr] (stu) -- (test);
\end{tikzpicture}
\notreal
\end{frame}

\begin{frame}{\T{提案と温度つき蒸留の違いは 2 点だけで，評価ではほぼ同じ値になった}{The proposal differs from temperature distillation in only two points, and scored the same in evaluation}}
\rl{\T{手法}{Method}}
\begin{columns}[T]
\begin{column}{0.58\textwidth}
\small
\begin{tabular}{@{}p{0.30\linewidth}p{0.31\linewidth}p{0.31\linewidth}@{}}
\toprule
& \T{提案}{Proposal} & \T{温度つき蒸留（Hinton 2015，Albanie 2018 の形）}{Temperature distillation (Hinton 2015; Albanie 2018 form)} \\
\midrule
\T{教師}{Teacher} & \T{同じ Q}{same Q} & \T{同じ Q}{same Q} \\
\T{疑似ラベルの形}{Pseudo-label form} & \T{Q そのまま}{Q as is} & \T{Q を温度 2 で柔らげる}{Q softened with temperature 2} \\
\T{感情ごとの総量}{Per-emotion total weight} & \T{そろえる}{balanced} & \T{そろえない}{not balanced} \\
\T{K 本・引き戻し}{K labels, pull} & \T{同じ}{same} & \T{同じ}{same} \\
\T{つまみ（格子）}{Knobs (grid)} & \T{全体の重み ρ，K 本の倍率}{overall weight ρ, label multiplier} & \T{温度，全体の重み ρ}{temperature, overall weight ρ} \\
\bottomrule
\end{tabular}
\end{column}
\begin{column}{0.42\textwidth}
\small
\T{%
\begin{itemize}
\item 評価の回で，提案 − 蒸留 は 48 セル中 40 セルで区間が 0 を含む（\S3）
\item 残りの差も −1.2〜+2.2 ポイントで，向きはコーパスで分かれる
\item したがって解析では\textbf{両者を区別せず}，「別のモダリティの予測（教師）で学ぶ手法」として扱う
\end{itemize}
}{%
\begin{itemize}
\item In evaluation, proposal − distillation has an interval containing 0 in 40 of 48 cells (\S3)
\item The remaining differences are −1.2 to +2.2 points, with signs varying by corpus
\item Hence the analysis \textbf{does not distinguish} them: both are ``learning from the other modalities' predictions (the teacher)''
\end{itemize}
}
\end{column}
\end{columns}
\takeaway{\T{「蒸留が良い理由」を問うときの「蒸留」は，この 2 手法の共通部分＝教師 Q で学ぶことを指す}{When we ask ``why distillation helps'', ``distillation'' means what the two share: learning from the teacher Q}}
\end{frame}

\begin{frame}{\T{評価の尺度と比較：UAR で測り，使う情報ごとに相手を分ける}{Measure and comparisons: UAR, with baselines grouped by the information they use}}
\rl{\T{評価（骨組み 5・6）}{Evaluation (skeleton 5, 6)}}
\hosoku{appmethods}{\T{手法と格子の一覧}{list of methods and grids}}
\begin{columns}[T]
\begin{column}{0.46\textwidth}
\footnotesize
\textbf{\T{尺度（骨組み 5）}{Measure (skeleton 5)}}
\T{%
\begin{itemize}
\item UAR（4 感情の再現率の平均）．MELD と MOSEI は Neutral が半分近くを占め，正解率だと「全部 Neutral」が高く出るため
\item テストの話者は学習側に入れない（MOSEI は話者の代わりにチャンネルか動画）
\end{itemize}
\textbf{除去（骨組み 6）}：提案から 1 つずつ外した版を実験 1 で計算済み
\begin{itemize}
\item (i) 教師を学び直さない　(ii) 教師を 1 モダリティにする
\item (iii) 総量をそろえない（＝温度 1 の蒸留）　(iv) 切片だけを K 本で学ぶ
\end{itemize}
}{%
\begin{itemize}
\item UAR (mean recall over 4 emotions). In MELD and MOSEI nearly half the utterances are Neutral, so accuracy rewards ``always Neutral''
\item Test speakers never enter training (MOSEI: channel or video stands in for speaker)
\end{itemize}
\textbf{Ablations (skeleton 6)}: versions with one part removed, already computed in Exp.~1
\begin{itemize}
\item (i) teacher not retrained\quad (ii) single-modality teacher
\item (iii) no balancing (= distillation at T = 1)\quad (iv) intercept only from K labels
\end{itemize}
}
\end{column}
\begin{column}{0.54\textwidth}
\footnotesize
\textbf{\T{比較の相手（骨組み 6）}{Baselines (skeleton 6)}}
\begin{tabular}{@{}p{0.30\linewidth}p{0.66\linewidth}@{}}
\toprule
\T{使う情報}{Information used} & \T{相手}{Baselines} \\
\midrule
\T{K 本だけ}{K labels only} & \T{比較の基準（出発点から K 本で学び直す）}{the baseline (start retrained on K labels)} \\
\T{K 本 ＋ 同じ教師 Q}{K labels + the same teacher Q} & \T{温度つき蒸留，一致した発話だけ，自分の予測と混ぜる，柔らかい共訓練，共正則化（実験 1）}{temperature distillation, agreed-only, mixing with own prediction, soft co-training, co-regularization (Exp.~1)} \\
\T{K 本 ＋ ラベル無し（同じモダリティ）}{K labels + unlabeled (same modality)} & \T{自己学習，エントロピー最小化，IM，T3A，ラベル伝播，自分の予測を提案の形で（実験 2）}{self-training, entropy min., IM, T3A, label propagation, own prediction in the proposal's form (Exp.~2)} \\
\T{＋ ソースのデータ}{+ source data} & \T{S+T，CORAL，DANN，MME，xMUDA 型（実験 2）}{S+T, CORAL, DANN, MME, xMUDA-style (Exp.~2)} \\
\bottomrule
\end{tabular}
\vspace{0.3em}
{\scriptsize\color{mMuted}\T{全手法が同じ出発点・同じ標準化・同じ K 本を使い，つまみは仮のターゲットで選んで固定した}{All methods share the start, the standardization and the K labels; knobs were selected on proxy targets and frozen}}
\end{column}
\end{columns}
\end{frame}

\begin{frame}{\T{統計の決まり：検定に残るのは 128 セル中 15 セルで，残りは区間つきの記述}{Statistical rules: 15 of 128 cells are tested; the rest are descriptions with intervals}}
\rl{\T{統計}{Statistics}}
\hosoku{appstats}{\T{区間と検定の規則の詳細}{details of intervals and tests}}
\small
\T{%
\begin{itemize}
\item セル ＝（コーパス，生徒，K）．族 1（提案 − K 本だけ），族 3（提案 − 同じ情報の相手）など，コーパスごとに Holm で補正する
\item 区間は，テストの話者の塊の引き直し（jackknife）と K 本の抜き方の揺れを含む．AFEW・IEMOCAP・MELD は実効の塊が 15・5・9 と少なく，1〜2 ポイントの差は有意にならない
\item そこで評価の\textbf{前}に，検出力の見込みが足りるセル（15）だけを検定に残し，残り 113 セルは「区間つきの記述」と決めた．記述のセルでは「差が無い」「効かない」と書かない
\item 手法の形は評価と同じ 4 コーパスの上で作ったので，p 値は確認の検定ではない（論文の限界に書く）
\end{itemize}
}{%
\begin{itemize}
\item A cell = (corpus, student, K). Family 1 (proposal − K-only), family 3 (proposal − same-information baseline), etc., with Holm correction per corpus
\item Intervals include resampling of test speaker clusters (jackknife) and the random draw of the K labels. AFEW, IEMOCAP and MELD have only 15, 5 and 9 effective clusters, so 1--2 point differences cannot be significant
\item Therefore, \textbf{before} evaluation, only the 15 cells with sufficient expected power were kept for testing; the other 113 are ``descriptions with intervals''. For those we never write ``no difference'' or ``does not help''
\item Method forms were built on the same 4 corpora used for evaluation, so p-values are not confirmatory (stated in the limitations)
\end{itemize}
}
\end{frame}

%% 3. 結果（実験 1・2）
\setband{\T{\S3 結果}{\S3 Results}}

% セルの値：#1 = famone|famthree，#2 = コーパス，#3 = 生徒，#4 = K（zero|five|twofive）
\newcommand{\fc}[4]{\N{#1#2#3#4}\,{\color{mMuted}[\N{#1lo#2#3#4}, \N{#1hi#2#3#4}]}}

% 族の図：#1 = fam1|fam3，#2 = ymin，#3 = ymax，#4 = ylabel
\newcommand{\fampanel}[3]{%
  \nextgroupplot[title=#2]
  \addplot[gray!60, no marks, forget plot] coordinates {(0.5,0) (4.5,0)};
  \addplot[sK0] table[x expr=\thisrow{xi}-0.25, y=effect, y error plus expr=\thisrow{hi}-\thisrow{effect}, y error minus expr=\thisrow{effect}-\thisrow{lo}]{#1_#2_K0.csv};
  \addplot[sK5] table[x expr=\thisrow{xi}, y=effect, y error plus expr=\thisrow{hi}-\thisrow{effect}, y error minus expr=\thisrow{effect}-\thisrow{lo}]{#1_#2_K5.csv};
  \addplot[sK25] table[x expr=\thisrow{xi}+0.25, y=effect, y error plus expr=\thisrow{hi}-\thisrow{effect}, y error minus expr=\thisrow{effect}-\thisrow{lo}]{#1_#2_K25.csv};
  #3}
\newcommand{\famchart}[4]{%
\begin{tikzpicture}
\begin{groupplot}[group style={group size=4 by 1, horizontal sep=0.3cm, ylabels at=edge left, yticklabels at=edge left},
  width=0.255\textwidth, height=4.3cm, xmin=0.5, xmax=4.5, ymin=#2, ymax=#3, xtick={1,2,3,4}, xticklabels={A,V,T,AVT},
  ylabel={#4}, ylabel style={font=\tiny}, ymajorgrids, title style={font=\scriptsize, yshift=-4pt}, legend style={at={(0.03,0.97)}, anchor=north west, font=\tiny, cells={anchor=west}}]
\fampanel{#1}{AFEW}{\legend{K = 0, K = 5, K = 25}}
\fampanel{#1}{IEMOCAP}{}
\fampanel{#1}{MELD}{}
\fampanel{#1}{MOSEI}{}
\end{groupplot}
\end{tikzpicture}}

% 実験 2 の図：#1 = K，#2 = xmin，#3 = xmax
\newcommand{\daipanel}[2]{%
  \nextgroupplot[title=#2]
  \addplot[gray!60, no marks, forget plot] coordinates {(0,0.3) (0,10.7)};
  \addplot[xbar, bar width=4.2pt, fill=cProp!70, draw=cProp, error bars/.cd, x dir=both, x explicit, error bar style={mDarkTeal, line width=0.6pt}]
    table[x=effect, y=yi, x error plus expr=\thisrow{hi}-\thisrow{effect}, x error minus expr=\thisrow{effect}-\thisrow{lo}]{dai2_#2_K#1.csv};}
\newcommand{\daichart}[4]{%
\begin{tikzpicture}
\begin{groupplot}[group style={group size=4 by 1, horizontal sep=0.3cm, yticklabels at=edge left},
  width=0.255\textwidth, height=#4, xmin=#2, xmax=#3, ymin=0.3, ymax=10.7, ytick={1,...,10},
  yticklabels={\T{自己学習}{self-training}, \T{エントロピー最小化}{entropy min.}, IM, T3A, \T{ラベル伝播}{label prop.}, \T{自分の予測}{own prediction}, S+T, CORAL, DANN, MME},
  y dir=reverse, xmajorgrids, title style={font=\scriptsize, yshift=-4pt}, tick label style={font=\tiny}]
\daipanel{#1}{AFEW}
\daipanel{#1}{IEMOCAP}
\daipanel{#1}{MELD}
\daipanel{#1}{MOSEI}
\end{groupplot}
\end{tikzpicture}}

% 得の棒グラフ：#1 = fam1，#2 = コーパス
\newcommand{\gainpanel}[2]{%
  \nextgroupplot[title=#2]
  \addplot[gray!60, no marks, forget plot] coordinates {(0.5,0) (4.5,0)};
  \addplot[ybar, bar shift=0pt, bar width=8pt, fill=cK5, draw=cK5, error bars/.cd, y dir=both, y explicit, error bar style={mDarkTeal, line width=0.6pt}]
    table[x expr=\thisrow{xi}-0.19, y=effect, y error plus expr=\thisrow{hi}-\thisrow{effect}, y error minus expr=\thisrow{effect}-\thisrow{lo}]{#1_#2_K5.csv};
  \addplot[ybar, bar shift=0pt, bar width=8pt, fill=cK25, draw=cK25, error bars/.cd, y dir=both, y explicit, error bar style={mDarkTeal, line width=0.6pt}]
    table[x expr=\thisrow{xi}+0.19, y=effect, y error plus expr=\thisrow{hi}-\thisrow{effect}, y error minus expr=\thisrow{effect}-\thisrow{lo}]{#1_#2_K25.csv};}
\newcommand{\gainchart}[1]{%
\begin{tikzpicture}
\begin{groupplot}[group style={group size=4 by 1, horizontal sep=0.35cm, ylabels at=edge left, yticklabels at=edge left},
  width=0.235\textwidth, height=4.2cm, xmin=0.5, xmax=4.5, ymin=-6, ymax=12, xtick={1,2,3,4}, xticklabels={A,V,T,AVT},
  ylabel={\T{得（ポイント）}{gain (points)}}, ylabel style={font=\tiny}, ymajorgrids, title style={font=\scriptsize, yshift=-4pt},
  legend style={at={(0.03,0.97)}, anchor=north west, font=\tiny, cells={anchor=west}}]
\gainpanel{#1}{AFEW}\legend{K = 5, K = 25}
\gainpanel{#1}{IEMOCAP}
\gainpanel{#1}{MELD}
\gainpanel{#1}{MOSEI}
\end{groupplot}
\end{tikzpicture}}

% 絶対の UAR の棒グラフ：#1 = 生徒の記号（A|V|T|AVT），#2 = K，#3 = 題
\newcommand{\levelpanel}[3]{%
  \nextgroupplot[title={#3}]
  \addplot[ybar, bar shift=0pt, bar width=11pt, fill=cNone!60, draw=cNone] table[x expr=\thisrow{ci}-0.2, y=none#2]{levels_#1.csv};
  \addplot[ybar, bar shift=0pt, bar width=11pt, fill=cProp, draw=cProp] table[x expr=\thisrow{ci}+0.2, y=prop#2]{levels_#1.csv};}
\newcommand{\levelchart}[3]{%
\begin{tikzpicture}
\begin{groupplot}[group style={group size=2 by 1, horizontal sep=0.6cm, ylabels at=edge left, yticklabels at=edge left}, width=0.5\textwidth, height=4.1cm,
  xmin=0.4, xmax=4.6, ymin=#2, ymax=#3, xtick={1,2,3,4}, xticklabels={AFEW, IEMOCAP, MELD, MOSEI}, tick label style={font=\scriptsize},
  ylabel={\T{テスト UAR（\%）}{test UAR (\%)}}, ylabel style={font=\scriptsize}, ymajorgrids, title style={font=\small, yshift=-3pt},
  nodes near coords={\pgfmathprintnumber[fixed, precision=1]{\pgfplotspointmeta}}, every node near coord/.append style={font=\tiny},
  ]
\levelpanel{#1}{5}{K = 5\T{（感情ごとに 5 本）}{ (5 labels per emotion)}}
\levelpanel{#1}{25}{K = 25\T{（感情ごとに 25 本）}{ (25 labels per emotion)}}
\end{groupplot}
\end{tikzpicture}}
% 凡例（手で描く）
\newcommand{\levellegend}{{\scriptsize\tikz[baseline=-0.3ex]\fill[cNone!60, draw=cNone] (0,-0.1) rectangle (0.4,0.15);\ \T{K 本だけ}{K labels only}\qquad \tikz[baseline=-0.3ex]\fill[cProp, draw=cProp] (0,-0.1) rectangle (0.4,0.15);\ \T{提案（K 本 ＋ 教師の疑似ラベル）}{proposal (K labels + teacher pseudo-labels)}}}
% 得の表：#1 = 生徒の記号
\newcommand{\gaintable}[1]{\scriptsize\setlength{\tabcolsep}{5pt}\begin{tabular}{@{}lrrrr@{}}\toprule
\T{得（提案 − K 本だけ）[95\% 区間]}{Gain (proposal − K-only) [95\% interval]} & AFEW & IEMOCAP & MELD & MOSEI \\ \midrule
K = 5 & \fc{famone}{AFEW}{#1}{five} & \fc{famone}{IEMOCAP}{#1}{five} & \fc{famone}{MELD}{#1}{five} & \fc{famone}{MOSEI}{#1}{five} \\
K = 25 & \fc{famone}{AFEW}{#1}{twofive} & \fc{famone}{IEMOCAP}{#1}{twofive} & \fc{famone}{MELD}{#1}{twofive} & \fc{famone}{MOSEI}{#1}{twofive} \\ \bottomrule\end{tabular}}

\begin{frame}{\T{結果 A-1（音声の識別器）：別のモダリティの予測を足すと，4 コーパスで上がる}{Result A-1 (audio classifier): the other modalities' predictions lift it in all 4 corpora}}
\rl{\T{実験 1．比べるもの：提案（K 本 ＋ 教師の疑似ラベル）対 K 本だけ．同じ K 本・同じ出発点}{Exp.~1. Compared: proposal (K labels + teacher pseudo-labels) vs K labels only; same labels, same start}}
\hosoku{appfam1}{\T{全セルの数値表（K = 0 を含む）}{full table of cells (incl.\ K = 0)}}
\centering
\levellegend\par\vspace{-0.3em}
\levelchart{A}{30}{85}\par
\vspace{-0.1em}
\footnotesize
\gaintable{A}\par
\takeaway{\T{音声では 12 セルすべてで正，10 セルで区間が 0 より上．MOSEI は検定でも有意（Holm 後 p < 0.001）}{Audio: positive in all 12 cells, interval above 0 in 10; MOSEI also significant (Holm p < 0.001)}}
\end{frame}

\begin{frame}{\T{結果 A-2（言葉の識別器）：IEMOCAP だけ大きく上がる．ほかは揺れの中}{Result A-2 (text classifier): a large gain only in IEMOCAP; elsewhere within the noise}}
\rl{\T{実験 1．比べるもの：提案 対 K 本だけ}{Exp.~1. Compared: proposal vs K labels only}}
\centering
\levellegend\par\vspace{-0.3em}
\levelchart{T}{30}{72}\par
\vspace{-0.1em}
\footnotesize
\gaintable{T}\par
\takeaway{\T{IEMOCAP の言葉は +5.0〜+7.5．AFEW・MELD・MOSEI は −1.2〜+1.6 で区間が 0 をまたぐ}{IEMOCAP text: +5.0 to +7.5. AFEW, MELD, MOSEI: −1.2 to +1.6, intervals crossing 0}}
\end{frame}

\begin{frame}{\T{結果 A-3（顔の識別器）：上がらない．IEMOCAP では下がる}{Result A-3 (face classifier): no gain; a loss in IEMOCAP}}
\rl{\T{実験 1．比べるもの：提案 対 K 本だけ}{Exp.~1. Compared: proposal vs K labels only}}
\centering
\levellegend\par\vspace{-0.3em}
\levelchart{V}{30}{90}\par
\vspace{-0.1em}
\footnotesize
\gaintable{V}\par
\takeaway{\T{顔で区間が 0 より上なのは MOSEI の K = 25 だけ．IEMOCAP は負（理由は \S4：切片が悪くなる）}{Face: the interval is above 0 only for MOSEI at K = 25; IEMOCAP is negative at all K (why: \S4, the intercept gets worse)}}
\end{frame}

\begin{frame}{\T{結果 A-4（AVT：3 つをつないだ識別器）：MELD 以外は上がらない}{Result A-4 (AVT, the concatenated classifier): no gain except in MELD}}
\rl{\T{実験 1．比べるもの：提案 対 K 本だけ}{Exp.~1. Compared: proposal vs K labels only}}
\centering
\levellegend\par\vspace{-0.3em}
\levelchart{AVT}{35}{92}\par
\vspace{-0.1em}
\footnotesize
\gaintable{AVT}\par
\takeaway{\T{AVT は MELD で +1.4〜+3.4，ほかは −0.8〜+1.7 で区間が 0 をまたぐ（理由は \S4：教師が AVT より当たっていない）}{AVT: +1.4 to +3.4 in MELD, −0.8 to +1.7 elsewhere with intervals crossing 0 (why: \S4, the teacher is no better than AVT)}}
\end{frame}

% 実験 2 の表の 1 行：#1 = 手法の記号（daitwo の key），#2 = 手法名
\newcommand{\dairow}[2]{#2 & \N{daitwocellAFEW#1five} & \N{daitwocellIEMOCAP#1five} & \N{daitwocellMELD#1five} & \N{daitwocellMOSEI#1five} \\}

\begin{frame}{\T{結果 B（実験 2）：同じラベル無しの発話を使う標準の手法と比べると，提案は多くに勝つが，全部ではない}{Result B (Exp.~2): against standard methods that use the same unlabeled utterances, the proposal beats most but not all}}
\rl{\T{比べるもの：提案 − 標準の手法（AVT の生徒，K = 5）}{Compared: proposal − standard method (AVT student, K = 5)}}
\hosoku{appdai2}{\T{区間つきの図（K = 0・5・25）}{charts with intervals (K = 0, 5, 25)}}
\begin{columns}[T]
\begin{column}{0.52\textwidth}
\scriptsize\setlength{\tabcolsep}{4pt}
\begin{tabular}{@{}lrrrr@{}}
\toprule
\T{手法（使う情報）}{Method (information used)} & AFEW & IEMOCAP & MELD & MOSEI \\
\midrule
\multicolumn{5}{@{}l}{\textit{\T{K 本 ＋ 同じモダリティのラベル無し}{K labels + same-modality unlabeled data}}} \\
\dairow{self}{\T{自己学習}{self-training}}
\dairow{entmin}{\T{エントロピー最小化}{entropy min.}}
\dairow{im}{IM}
\dairow{tthreea}{T3A}
\dairow{lp}{\T{ラベル伝播}{label propagation}}
\dairow{own}{\T{自分の予測を提案の形で}{own prediction}}
\multicolumn{5}{@{}l}{\textit{\T{＋ ソースのデータも使う}{+ source data}}} \\
\dairow{st}{S+T}
\dairow{coral}{CORAL}
\dairow{dann}{DANN}
\dairow{mme}{MME}
\bottomrule
\end{tabular}
\end{column}
\begin{column}{0.48\textwidth}
\scriptsize
\T{%
\begin{itemize}
\item \textbf{何を比べたか}：結果 A と同じ AVT の識別器を，提案（教師の疑似ラベル）ではなく標準の手法で直した場合との差．正なら提案が高い．\textbf{太字}は区間が 0 より上（提案が確かに高い），\textcolor{mBad}{赤}は区間が 0 より下（手法が確かに高い）
\item 提案は多くの手法に勝つ（12 セル中，エントロピー最小化 \N{daitwoaboveentmin}，T3A \N{daitwoabovetthreea}，MME \N{daitwoabovemme}，CORAL \N{daitwoabovecoral}，DANN \N{daitwoabovedann}）
\item 負けるのは MOSEI の CORAL（\N{daitwoMOSEIcoralfive}）と，AFEW の K = 0 の T3A・エントロピー最小化・IM（補足）
\item 使う情報が違うので，優劣でなく「情報の違い」として読む
\end{itemize}
}{%
\begin{itemize}
\item \textbf{What is compared}: the AVT classifier of Result A, adapted by a standard method instead of the proposal. Positive = proposal higher. \textbf{Bold}: interval above 0; \textcolor{mBad}{red}: interval below 0 (method reliably higher)
\item The proposal beats most methods (of 12 cells: entropy min.\ \N{daitwoaboveentmin}, T3A \N{daitwoabovetthreea}, MME \N{daitwoabovemme}, CORAL \N{daitwoabovecoral}, DANN \N{daitwoabovedann})
\item It loses to CORAL in MOSEI (\N{daitwoMOSEIcoralfive}) and, at K = 0 in AFEW, to T3A, entropy min.\ and IM (appendix)
\item Different information is used, so read as information differences, not a ranking
\end{itemize}
}
\end{column}
\end{columns}
{\tiny\color{mMuted}\T{提案 − 手法（UAR ポイント，50 単位の平均）．ラベル伝播は K ≥ 1 だけ．出どころ：figures/dai2\_*.csv}{Proposal − method (UAR points, mean over 50 units). Label propagation exists only for K ≥ 1. Source: figures/dai2\_*.csv}}
\end{frame}

%% 4. なぜ効くか（実験 3 と探索）
\setband{\T{\S4 なぜ効くか}{\S4 Why it works}}

\begin{frame}{\T{得が出るのは，教師が生徒より当たっているときである}{Gains appear where the teacher is more accurate than the student}}
\rl{\T{探索（評価の後）}{Exploration (post-evaluation)}}
\hosoku{appadv}{\T{K = 25 の図と数値}{chart and numbers for K = 25}}
\begin{columns}[T]
\begin{column}{0.54\textwidth}
\begin{tikzpicture}
\begin{axis}[width=\linewidth, height=6.2cm, xmin=-2, xmax=33, ymin=-3.5, ymax=9.5,
  xlabel={\T{教師の先回り（候補での UAR，教師 − 生徒．ポイント）}{teacher lead (pool UAR, teacher − student, points)}},
  ylabel={\T{提案の得（テスト UAR，提案 − K 本だけ）}{gain (test UAR, proposal − K-only)}},
  legend style={at={(0.98,0.03)}, anchor=south east, font=\tiny, cells={anchor=west}}, legend columns=2, grid=major,
  every node near coord/.append style={font=\tiny, anchor=west, xshift=2pt, color=mMuted}]
\addplot[only marks, kA, mark=*, mark size=2pt, point meta=explicit symbolic, nodes near coords] table[x=adv_pool, y=gain_prop, meta=T]{teacher_adv_A_K5.csv};
\addplot[only marks, kV, mark=square*, mark size=2pt, point meta=explicit symbolic, nodes near coords] table[x=adv_pool, y=gain_prop, meta=T]{teacher_adv_V_K5.csv};
\addplot[only marks, kT, mark=triangle*, mark size=2.4pt, point meta=explicit symbolic, nodes near coords] table[x=adv_pool, y=gain_prop, meta=T]{teacher_adv_T_K5.csv};
\addplot[only marks, kAVT, mark=diamond*, mark size=2.4pt] table[x=adv_pool, y=gain_prop]{teacher_adv_AVT_K5.csv};
\legend{\T{音声}{audio}, \T{顔}{face}, \T{言葉}{text}, AVT}
\end{axis}
\end{tikzpicture}
\end{column}
\begin{column}{0.46\textwidth}
\footnotesize
\T{%
\begin{itemize}
\item 横軸：教師（3 モダリティの積）が候補でどれだけ生徒（K 本だけ）より当たるか．縦軸：提案のテストでの得．K = 5，16 点（4 コーパス × 4 生徒）
\item 順位の相関は \N{rhofive}（K = 25 で \N{rhotwofive}）．先回りが大きい弱いモダリティ（音声，IEMOCAP の言葉）に得が出る
\item \textbf{AVT には学ぶものが無い}（図の菱形 4 点，コーパス名なし）：先回りは \N{advMOSEIAVTfive}〜\N{advAFEWAVTfive} ポイントで，教師と AVT の生徒の判定は候補の \N{agreeMOSEIAVTfive}〜\N{agreeAFEWAVTfive}\% で一致する
\end{itemize}
}{%
\begin{itemize}
\item x: how much the teacher (product of 3 modalities) beats the student (K-only) on the pool. y: the proposal's test gain. K = 5, 16 points (4 corpora × 4 students)
\item Rank correlation \N{rhofive} (\N{rhotwofive} at K = 25). Gains appear for weak modalities with a large lead (audio, IEMOCAP text)
\item \textbf{Nothing to learn for AVT} (the 4 unlabeled diamonds): the lead is \N{advMOSEIAVTfive} to \N{advAFEWAVTfive} points, and the teacher agrees with the AVT student on \N{agreeMOSEIAVTfive}--\N{agreeAFEWAVTfive}\% of the pool
\end{itemize}
}
\end{column}
\end{columns}
{\tiny\color{mMuted}\T{評価の回を見た後の探索．16 点は独立でなく，p 値は付けない．MOSEI は途中の 39 単位．出どころ：figures/teacher\_adv\_*.csv}{Post-evaluation exploration; the 16 points are not independent and no p-value is attached. MOSEI: 39 interim units. Source: figures/teacher\_adv\_*.csv}}
\end{frame}

\begin{frame}{\T{損を「位置」と「向き」に分ける：出発点の損の大半は重みの向きのずれである}{Splitting the loss into ``position'' and ``direction'': most of the start's loss is a weight-direction error}}
\rl{\T{実験 3}{Exp.~3}}
\hosoku{appdecomp}{\T{全生徒・全 K の表}{tables for all students and K}}
\begin{columns}[T]
\begin{column}{0.5\textwidth}
\centering
\begin{tikzpicture}[font=\scriptsize, x=1cm, y=1cm]
  \pgfmathsetseed{7}
  % 2 つの雲
  \foreach \i in {1,...,14} {\pgfmathsetmacro{\px}{0.9+0.75*rand}\pgfmathsetmacro{\py}{1.7+0.5*rand}\fill[kA!70] (\px,\py) circle (1.3pt);}
  \foreach \i in {1,...,14} {\pgfmathsetmacro{\px}{2.6+0.75*rand}\pgfmathsetmacro{\py}{0.5+0.5*rand}\fill[kV!70] (\px,\py) circle (1.3pt);}
  \node[font=\tiny, kA] at (0.9,2.45) {\T{感情 1}{emotion 1}}; \node[font=\tiny, kV] at (2.6,-0.25) {\T{感情 2}{emotion 2}};
  % 正しい境界
  \draw[cUpper, line width=1pt] (0.2,-0.2) -- (3.4,2.6);
  % 位置のずれ（平行移動）
  \draw[cCrit, line width=1pt, densely dashed] (0.9,-0.2) -- (4.1,2.6);
  % 向きのずれ（回転）
  \draw[cDisc, line width=1pt, dotted] (-0.4,0.8) -- (4.0,1.6);
  \node[font=\tiny, anchor=north, align=left] at (1.9,-0.55) {\textcolor{cUpper}{\rule[0.5ex]{2em}{1pt}} \T{正しい境界}{right boundary}\quad \textcolor{cCrit}{- - -} \T{位置のずれ（切片：平行に動く）}{position error (intercept: parallel shift)}\\ \textcolor{cDisc}{$\cdots$} \T{向きのずれ（重み：回る）}{direction error (weights: rotation)}};
\end{tikzpicture}
\notreal
\vspace{0.2em}
\scriptsize
\T{測り方：識別器の切片だけをテストの感情の釣り合いで当て直す．当て直しで減る分＝位置の損，残る分＝向きの損（上限との差）．}{How measured: refit only the intercept on the test emotion balance. The reduction = position loss; what remains = direction loss (gap to the upper bound).}
\end{column}
\begin{column}{0.5\textwidth}
\begin{tikzpicture}
\begin{axis}[width=\linewidth, height=4.4cm, ybar stacked, bar width=14pt, ymin=0, ymax=18, enlarge x limits=0.2,
  xtick={1,2,3,4}, xticklabels={AFEW, IEMOCAP, MELD, MOSEI}, ylabel={\T{出発点の損（AVT，ポイント）}{loss of the start (AVT, points)}},
  legend style={at={(0.98,0.97)}, anchor=north east, font=\tiny, cells={anchor=west}}, ymajorgrids,
  nodes near coords={\pgfmathprintnumber[fixed, precision=1]{\pgfplotspointmeta}}, every node near coord/.append style={font=\tiny, color=white}]
\addplot[fill=cCrit, draw=cCrit] table[x=ci, y=position]{decomp_start_avt.csv};
\addplot[fill=cDisc, draw=cDisc] table[x=ci, y=direction]{decomp_start_avt.csv};
\legend{\T{位置（切片）}{position (intercept)}, \T{向き（重み）}{direction (weights)}}
\end{axis}
\end{tikzpicture}
\footnotesize
\T{%
\begin{itemize}
\item AVT の出発点：向きの損 \N{dirMOSEIAVTstart}〜\N{dirAFEWAVTstart}，位置の損 \N{posMOSEIAVTstart}〜\N{posIEMOCAPAVTstart} ポイント
\item 向きの物差しの確かめに不合格の組（AFEW と IEMOCAP の AVT など 5 組）では，論文で「向き」と呼ばず「当て直した後に残る差」と書く
\end{itemize}
}{%
\begin{itemize}
\item AVT start: direction loss \N{dirMOSEIAVTstart} to \N{dirAFEWAVTstart}, position loss \N{posMOSEIAVTstart} to \N{posIEMOCAPAVTstart} points
\item Where the direction measure failed its validity check (5 pairs incl.\ AFEW and IEMOCAP AVT), the paper says ``gap remaining after the refit'' instead of ``direction''
\end{itemize}
}
\end{column}
\end{columns}
\end{frame}

\begin{frame}{\T{K 本も向きを直していた．提案は，弱いモダリティの向きの損をさらに減らす}{K labels also fix the direction; the proposal further shrinks the direction loss of weak modalities}}
\rl{\T{実験 3}{Exp.~3}}
\begin{columns}[T]
\begin{column}{0.56\textwidth}
\footnotesize
\T{向きの損（当て直した後に残る差，ポイント，K = 5）}{Direction loss (gap remaining after the refit, points, K = 5)}
\begin{tabular}{@{}llrrr@{}}
\toprule
\T{コーパス}{Corpus} & \T{生徒}{Student} & \T{出発点}{Start} & \T{K 本だけ}{K-only} & \T{提案}{Proposal} \\
\midrule
AFEW & \T{音声}{audio} & \N{dirAFEWAstart} & \N{dirAFEWAnonefive} & \textbf{\N{dirAFEWApropfive}} \\
IEMOCAP & \T{音声}{audio} & \N{dirIEMOCAPAstart} & \N{dirIEMOCAPAnonefive} & \textbf{\N{dirIEMOCAPApropfive}} \\
MELD & \T{音声}{audio} & \N{dirMELDAstart} & \N{dirMELDAnonefive} & \textbf{\N{dirMELDApropfive}} \\
MOSEI & \T{音声}{audio} & \N{dirMOSEIAstart} & \N{dirMOSEIAnonefive} & \textbf{\N{dirMOSEIApropfive}} \\
IEMOCAP & \T{言葉}{text} & \N{dirIEMOCAPTstart} & \N{dirIEMOCAPTnonefive} & \textbf{\N{dirIEMOCAPTpropfive}} \\
MELD & \T{言葉}{text} & \N{dirMELDTstart} & \N{dirMELDTnonefive} & \N{dirMELDTpropfive} \\
MOSEI & \T{言葉}{text} & \N{dirMOSEITstart} & \N{dirMOSEITnonefive} & \N{dirMOSEITpropfive} \\
MELD & AVT & \N{dirMELDAVTstart} & \N{dirMELDAVTnonefive} & \N{dirMELDAVTpropfive} \\
MOSEI & AVT & \N{dirMOSEIAVTstart} & \N{dirMOSEIAVTnonefive} & \N{dirMOSEIAVTpropfive} \\
\bottomrule
\end{tabular}
\end{column}
\begin{column}{0.44\textwidth}
\footnotesize
\T{%
\begin{itemize}
\item 事前に立てた説明「K 本は主に切片を直す」は 4 コーパスとも支持されない．K 本は向きも直す（例：IEMOCAP の言葉 \N{dirIEMOCAPTstart} → \N{dirIEMOCAPTnonefive}）
\item 提案は，得の出る生徒で向きの損をさらに減らす（IEMOCAP の言葉 \N{dirIEMOCAPTnonefive} → \N{dirIEMOCAPTpropfive}，音声は 4 コーパスで減る）
\item 得の出ない生徒では減らない，または増える（MELD・MOSEI の言葉，MOSEI の AVT）
\item 「重みは予測で，切片は K 本で」と役割を分けた版は，同時に学ぶ版と差が無い（−0.4〜+1.7）
\end{itemize}
}{%
\begin{itemize}
\item The pre-registered ``K labels mainly fix the intercept'' is not supported in any corpus; K labels also fix the direction (e.g.\ IEMOCAP text \N{dirIEMOCAPTstart} → \N{dirIEMOCAPTnonefive})
\item The proposal further reduces the direction loss where gains appear (IEMOCAP text \N{dirIEMOCAPTnonefive} → \N{dirIEMOCAPTpropfive}; audio in all 4 corpora)
\item Where there is no gain, it does not shrink or grows (MELD and MOSEI text, MOSEI AVT)
\item A version that learns weights from predictions and the intercept from K labels does not differ from learning both together (−0.4 to +1.7)
\end{itemize}
}
\end{column}
\end{columns}
{\tiny\color{mMuted}\T{点の値．区間と判定の規則は補足．出どころ：figures/decomp.csv（版 four）}{Point values; intervals and decision rules in the appendix. Source: figures/decomp.csv (version four)}}
\end{frame}

\begin{frame}{\T{効かない所の理由：顔は学ぶものが無いか切片が悪くなる．MELD・MOSEI は教師そのものが弱い}{Why it fails: the face has nothing to learn or gets a worse intercept; in MELD and MOSEI the teacher itself is weak}}
\rl{\T{実験 3・探索}{Exp.~3, exploration}}
\begin{columns}[T]
\begin{column}{0.5\textwidth}
\footnotesize
\textbf{\T{顔の生徒}{Face student}}
\T{%
\begin{itemize}
\item AFEW：教師の先回りは \N{advAFEWVfive} ポイントしかない（顔が最も強いモダリティ）→ 学ぶものが無い
\item IEMOCAP：損の大半が位置（切片）で（\N{posIEMOCAPVnonefive} 対 向き \N{dirIEMOCAPVnonefive}），提案は位置の損を \N{posIEMOCAPVnonefive} → \N{posIEMOCAPVpropfive} と\textbf{悪くする}
\item MELD：先回り \N{advMELDVfive} に対して得は \N{gainrMELDVfive}．理由は未解析
\item 直す余地（上限 − K 本だけ）は顔で \N{roomMELDVfive}〜\N{roomMOSEIVfive}，音声で \N{roomAFEWAfive}〜\N{roomMOSEIAfive} ポイント
\end{itemize}
}{%
\begin{itemize}
\item AFEW: the teacher's lead is only \N{advAFEWVfive} points (the face is the strongest modality) → nothing to learn
\item IEMOCAP: most of the loss is position (intercept; \N{posIEMOCAPVnonefive} vs direction \N{dirIEMOCAPVnonefive}), and the proposal \textbf{worsens} it (\N{posIEMOCAPVnonefive} → \N{posIEMOCAPVpropfive})
\item MELD: lead \N{advMELDVfive} but gain \N{gainrMELDVfive}; not yet analyzed
\item Room (upper bound − K-only): face \N{roomMELDVfive}--\N{roomMOSEIVfive}, audio \N{roomAFEWAfive}--\N{roomMOSEIAfive} points
\end{itemize}
}
\end{column}
\begin{column}{0.5\textwidth}
\footnotesize
\textbf{\T{教師の強さ（候補での UAR，K = 5）}{Teacher strength (pool UAR, K = 5)}}
\begin{tabular}{@{}lrr@{}}
\toprule
\T{コーパス}{Corpus} & \T{教師}{Teacher} & \T{AVT の生徒（K 本だけ）}{AVT student (K-only)} \\
\midrule
AFEW & \N{tuarAFEWAVTfive} & \N{suarAFEWAVTfive} \\
IEMOCAP & \N{tuarIEMOCAPAVTfive} & \N{suarIEMOCAPAVTfive} \\
MELD & \N{tuarMELDAVTfive} & \N{suarMELDAVTfive} \\
MOSEI & \N{tuarMOSEIAVTfive} & \N{suarMOSEIAVTfive} \\
\bottomrule
\end{tabular}
\T{%
\begin{itemize}
\item MELD と MOSEI では教師が候補の半分も当てない．それでも音声には効く（音声の生徒がもっと弱いから）
\item 教師の強さはコーパスで大きく違う．\textbf{どのモダリティをどれだけ信じるか}がコーパスごとに違うことが，次の問題になる（\S5，\S6）
\end{itemize}
}{%
\begin{itemize}
\item In MELD and MOSEI the teacher gets less than half the pool right, yet it still helps audio (the audio student is weaker still)
\item Teacher strength differs widely across corpora. \textbf{Which modality to trust, and how much}, differs per corpus --- the next problem (\S5, \S6)
\end{itemize}
}
\end{column}
\end{columns}
\end{frame}

%% 5. 普通の蒸留の弱点と実験 4
\setband{\T{\S5 普通の蒸留の弱点}{\S5 Weakness of plain distillation}}

\begin{frame}{\T{弱点：教師はラベル無しの発話すべてに同じ重みで疑似ラベルを付けるが，当たりやすさは発話で大きく違う}{Weakness: the teacher pseudo-labels every pool utterance with the same weight, yet its accuracy varies widely across utterances}}
\rl{\T{探索（評価の後）}{Exploration (post-evaluation)}}
\hosoku{appcons}{\T{全生徒・K = 25 の数値}{numbers for all students and K = 25}}
\begin{columns}[T]
\begin{column}{0.46\textwidth}
\centering
\begin{tikzpicture}[font=\scriptsize, x=1cm, y=1cm]
  \node[font=\scriptsize, anchor=west] at (-0.3,1.9) {\T{候補の発話と教師の疑似ラベル}{pool utterances and the teacher's pseudo-labels}};
  \foreach \i/\lab/\agr in {0/H/1, 1/N/0, 2/S/1, 3/A/0, 4/H/1, 5/N/1, 6/A/0, 7/S/0} {
    \node[utt] (u\i) at (0.7*\i,1.2) {\lab};
    \fill[cMod] (0.7*\i-0.2,0) rectangle (0.7*\i+0.2,0.7);
    \ifnum\agr=1 \node[font=\tiny, cAgree] at (0.7*\i,-0.3) {\checkmark}; \else \node[font=\tiny, cDis] at (0.7*\i,-0.3) {×}; \fi
  }
  \node[font=\tiny, anchor=west, color=mMuted] at (-0.3,0.35) {\rotatebox{90}{\T{重み}{weight}}};
  \draw[mDarkTeal!50] (-0.1,0) -- (5.1,0);
  \node[font=\tiny, anchor=west, align=left] at (-0.3,-0.75) {\T{普通の蒸留・提案：全部に同じ重み（上の帯）}{plain distillation / proposal: same weight for all (bars)}\\ \T{\checkmark{}／×：生徒以外の 2 モダリティの判定が一致するか}{\checkmark / ×: whether the other two modalities agree}};
\end{tikzpicture}
\notreal
\end{column}
\begin{column}{0.54\textwidth}
\begin{tikzpicture}
\begin{axis}[width=\linewidth, height=4.1cm, ybar, bar width=7pt, ymin=20, ymax=105, enlarge x limits=0.15,
  xtick={1,2,3,4}, xticklabels={AFEW, IEMOCAP, MELD, MOSEI}, ylabel={\T{教師の正解率（候補，\%）}{teacher accuracy (pool, \%)}},
  legend style={at={(0.98,0.97)}, anchor=north east, font=\tiny, cells={anchor=west}}, ymajorgrids,
  title={\T{音声の生徒，K = 5}{audio student, K = 5}}, title style={font=\scriptsize, yshift=-4pt},
  nodes near coords={\pgfmathprintnumber[fixed, precision=0]{\pgfplotspointmeta}}, every node near coord/.append style={font=\tiny}]
\addplot[fill=cAgree!80, draw=cAgree] table[x=ci, y=acc_agree]{consensus_A_K5.csv};
\addplot[fill=cDis!70, draw=cDis] table[x=ci, y=acc_dis]{consensus_A_K5.csv};
\legend{\T{顔と言葉が一致する発話}{face and text agree}, \T{一致しない発話}{they disagree}}
\end{axis}
\end{tikzpicture}
\end{column}
\end{columns}
\footnotesize
\T{%
\begin{itemize}
\item 音声の生徒の教師は，顔と言葉の判定が一致する発話では AFEW \N{cagreeAFEWAfive}\%・IEMOCAP \N{cagreeIEMOCAPAfive}\% 当たり，一致しない発話では \N{cdisAFEWAfive}\%・\N{cdisIEMOCAPAfive}\% しか当たらない（言葉の生徒でも同じ：\N{cagreeAFEWTfive}／\N{cdisAFEWTfive}\%，\N{cagreeIEMOCAPTfive}／\N{cdisIEMOCAPTfive}\%）
\item つまり「生徒以外のモダリティが一致するか」は，教師の当たりやすさをラベル無しで予言する．普通の蒸留はこれを使わず，当たらない発話にも同じ重みで学ばせる
\end{itemize}
}{%
\begin{itemize}
\item For the audio student, the teacher is right on \N{cagreeAFEWAfive}\% (AFEW) and \N{cagreeIEMOCAPAfive}\% (IEMOCAP) of utterances where face and text agree, but only \N{cdisAFEWAfive}\% and \N{cdisIEMOCAPAfive}\% where they disagree (text student: \N{cagreeAFEWTfive}/\N{cdisAFEWTfive}\%, \N{cagreeIEMOCAPTfive}/\N{cdisIEMOCAPTfive}\%)
\item So ``do the other modalities agree'' predicts the teacher's accuracy without labels. Plain distillation ignores it and trains equally on the utterances it gets wrong
\end{itemize}
}
\end{frame}

\begin{frame}{\T{実験 4 の案：一致しない発話の重みを下げる（設計と台本は検査を通り，計算は保留）}{Experiment 4: down-weight the utterances where the other modalities disagree (design and code reviewed; computation on hold)}}
\rl{\T{実験 4（未計算）}{Exp.~4 (not computed)}}
\begin{columns}[T]
\begin{column}{0.46\textwidth}
\centering
\begin{tikzpicture}[font=\scriptsize, x=1cm, y=1cm]
  \foreach \i/\lab/\agr in {0/H/1, 1/N/0, 2/S/1, 3/A/0, 4/H/1, 5/N/1, 6/A/0, 7/S/0} {
    \node[utt] (u\i) at (0.7*\i,1.2) {\lab};
    \ifnum\agr=1 \fill[cMod] (0.7*\i-0.2,0) rectangle (0.7*\i+0.2,0.9); \node[font=\tiny, cAgree] at (0.7*\i,-0.3) {\checkmark};
    \else \fill[cMod!50] (0.7*\i-0.2,0) rectangle (0.7*\i+0.2,0.25); \node[font=\tiny, cDis] at (0.7*\i,-0.3) {×}; \fi
  }
  \draw[mDarkTeal!50] (-0.1,0) -- (5.1,0);
  \node[font=\tiny, anchor=north west, align=left, text width=5.6cm] at (-0.3,-0.55) {\T{一致する発話：重み 1．一致しない発話：重み γ（0，1/4，1/2 から仮のターゲットで選ぶ）}{agree: weight 1. disagree: weight γ (0, 1/4 or 1/2, chosen on proxy targets)}\\ \T{感情ごとの総量はそろえたまま（感情の間の重みは変えない）}{per-emotion totals stay balanced (no shift between emotions)}};
\end{tikzpicture}
\notreal
\vspace{0.3em}
\footnotesize
\T{%
対照（印の定義だけを替えた版）も同じ格子で解く：生徒自身との一致，教師の確信度の上位，乱数で選んだ印．これが無いと「他モダリティの一致が効いた」と「確信度の高い発話を重くすれば効く」を区別できない．
}{%
Controls with only the mark changed are solved on the same grid: agreement with the student itself, top teacher confidence, random marks. Without them one cannot tell ``cross-modal agreement helped'' from ``up-weighting confident utterances helps''.
}
\end{column}
\begin{column}{0.54\textwidth}
\footnotesize
\textbf{\T{状態}{Status}}
\T{%
\begin{itemize}
\item 設計（第 3 版）と台本は，関門 4 体・台本の検査役・ML と感情認識の査読想定の両方を通った
\item 見込みの得は \textbf{+1 ポイント以下}（一致しない発話で教師が壊す割合が一致する発話の約 2 倍で，その半分を避けると仮定）
\item 案は評価の回の診断から出たので，\textbf{確認の検定にはならない}（論文では由来と限界を書く）
\item 所要：H200 1 台で 4〜6 時間（約 50〜80 ドル）．研究者の判断で保留中（「まだ違う筋を探そう」）
\end{itemize}
}{%
\begin{itemize}
\item Design (v3) and scripts passed the four gate checkers, the code checker and both simulated reviewers (ML, affective computing)
\item Expected gain \textbf{≤ +1 point} (the teacher harms about twice as often where the others disagree; assume half of that is avoided)
\item The idea came from evaluation-round diagnostics, so it \textbf{cannot be a confirmatory test} (stated in the paper)
\item Cost: one H200, 4--6 h (≈\$50--80). On hold by the researcher's decision
\end{itemize}
}
\takeaway{\T{弱点は分かったが，直せる量は小さい．より大きな得は「どのモダリティを信じるか」にある（\S6）}{The weakness is identified, but the fixable amount is small. The larger gain lies in ``which modality to trust'' (\S6)}}
\end{column}
\end{columns}
\end{frame}

%% 6. いま進めていること（モダリティの信頼度）
\setband{\T{\S6 いま進めていること}{\S6 Current work}}

\begin{frame}{\T{核心の問い：新しいコーパスで，どのモダリティをどれだけ信じるかを少数のラベルから決められるか}{The core question: can a few labels tell us which modality to trust, and how much, in a new corpus?}}
\rl{\T{問い}{Question}}
\hosoku{appfusion}{\T{融合の表（全部）}{all fusion tables}}
\begin{columns}[T]
\begin{column}{0.5\textwidth}
\centering
\begin{adjustbox}{max width=\linewidth}
\begin{tikzpicture}[font=\scriptsize, x=1cm, y=1cm]
  % つなぐ形
  \node[encbox, minimum width=0.8cm] (a1) at (0,0) {A}; \node[encbox, minimum width=0.8cm] (v1) at (0.9,0) {V}; \node[encbox, minimum width=0.8cm] (t1) at (1.8,0) {T};
  \node[clf, text width=2.4cm, fill=kAVT!10] (avt) at (0.9,-1.1) {\T{AVT の生徒}{AVT student}\\ \T{（つないだ特徴の上の 1 つの識別器）}{(one classifier on concatenated features)}};
  \draw[arr] (0.9,-0.3) -- (avt);
  \node[font=\scriptsize\bfseries] at (0.9,0.7) {\T{つなぐ形（今の提案）}{Concat (current proposal)}};
  % 積の形
  \node[clf, minimum width=0.9cm] (ca) at (4.3,0) {\T{直した A}{adapted A}}; \node[clf, minimum width=0.9cm] (cv) at (5.5,0) {\T{直した V}{adapted V}}; \node[clf, minimum width=0.9cm] (ct) at (6.7,0) {\T{直した T}{adapted T}};
  \node[modbox, text width=2.9cm] (prod) at (5.5,-1.1) {\T{確率の積：}{product of probabilities:}\\ $p_A^{\,e_A}\, p_V^{\,e_V}\, p_T^{\,e_T}$};
  \draw[arr] (ca) -- (prod); \draw[arr] (cv) -- (prod); \draw[arr] (ct) -- (prod);
  \node[font=\scriptsize\bfseries] at (5.5,0.7) {\T{積の形（指数 = 信頼度）}{Product (exponents = trust)}};
  \node[font=\tiny, color=mMuted, align=center] at (5.5,-1.95) {\T{指数が大きいモダリティほど判定に強く入る}{a larger exponent gives that modality more say}};
\end{tikzpicture}
\end{adjustbox}
\notreal
\end{column}
\begin{column}{0.5\textwidth}
\footnotesize
\T{%
\begin{itemize}
\item 「積の形」で指数を\textbf{テストで選んだ上限}は，AVT の生徒より AFEW で \N{foracledAFEWfive}（K = 5）・\N{foracledAFEWtwofive}（K = 25），MELD で \N{foracledMELDfive}・\N{foracledMELDtwofive} ポイント高い．IEMOCAP は \N{foracledIEMOCAPfive}・\N{foracledIEMOCAPtwofive}，MOSEI は \N{foracledMOSEIfive}・\N{foracledMOSEItwofive}
\item 最良の指数はコーパスで違う：AFEW は顔，IEMOCAP は音声，MELD は言葉を重くする
\item 教師の積の指数も同じ量で決まる．つまり「どのモダリティがどれだけ信頼できるか」が，この課題の 1 つの隠れた量である
\end{itemize}
}{%
\begin{itemize}
\item With the product form and exponents \textbf{chosen on the test set (upper bound)}, the gain over the AVT student is \N{foracledAFEWfive} (K = 5) and \N{foracledAFEWtwofive} (K = 25) in AFEW, \N{foracledMELDfive} and \N{foracledMELDtwofive} in MELD; IEMOCAP \N{foracledIEMOCAPfive}, \N{foracledIEMOCAPtwofive}; MOSEI \N{foracledMOSEIfive}, \N{foracledMOSEItwofive}
\item The best exponents differ per corpus: face for AFEW, audio for IEMOCAP, text for MELD
\item The teacher's exponents depend on the same quantity. ``How much each modality can be trusted'' is a hidden variable of this task
\end{itemize}
}
\end{column}
\end{columns}
{\tiny\color{mMuted}\T{上限は 63 点の格子でテストの最大を選んだ値なので，選択だけで正になる．手法の性能としては出さない}{The upper bound picks the test maximum over a 63-point grid, so selection alone makes it positive; it is not reported as a method's performance}}
\end{frame}

\begin{frame}{\T{手がかりが無い：一致率でも，学習に使った K 本でも，指数を決められない}{No usable signal: neither agreement rates nor the training labels can set the exponents}}
\rl{\T{探索}{Exploration}}
\scriptsize
\begin{center}
\T{テスト UAR（AVT の生徒との差，ポイント）}{Test UAR (difference from the AVT student, points)}

\begin{tabular}{@{}lrrrrrrrr@{}}
\toprule
& \multicolumn{2}{c}{AFEW} & \multicolumn{2}{c}{IEMOCAP} & \multicolumn{2}{c}{MELD} & \multicolumn{2}{c}{MOSEI} \\
K & 5 & 25 & 5 & 25 & 5 & 25 & 5 & 25 \\
\midrule
\T{積，指数すべて 1}{Product, all exponents 1} & \N{fproddAFEWfive} & \N{fproddAFEWtwofive} & \N{fproddIEMOCAPfive} & \N{fproddIEMOCAPtwofive} & \N{fproddMELDfive} & \N{fproddMELDtwofive} & \N{fproddMOSEIfive} & \N{fproddMOSEItwofive} \\
\T{積，K 本で選んだ指数}{Product, exponents chosen on the K labels} & \N{fkshotdAFEWfive} & \N{fkshotdAFEWtwofive} & \N{fkshotdIEMOCAPfive} & \N{fkshotdIEMOCAPtwofive} & \N{fkshotdMELDfive} & \N{fkshotdMELDtwofive} & \N{fkshotdMOSEIfive} & \N{fkshotdMOSEItwofive} \\
\T{積，テストで選んだ上限}{Product, test-chosen upper bound} & \N{foracledAFEWfive} & \N{foracledAFEWtwofive} & \N{foracledIEMOCAPfive} & \N{foracledIEMOCAPtwofive} & \N{foracledMELDfive} & \N{foracledMELDtwofive} & \N{foracledMOSEIfive} & \N{foracledMOSEItwofive} \\
\midrule
\T{K 本で最強のモダリティを当てた割合}{K labels pick the strongest modality (\%)} & \N{fheavyAFEWfive}\% & \N{fheavyAFEWtwofive}\% & \N{fheavyIEMOCAPfive}\% & \N{fheavyIEMOCAPtwofive}\% & \N{fheavyMELDfive}\% & \N{fheavyMELDtwofive}\% & \N{fheavyMOSEIfive}\% & \N{fheavyMOSEItwofive}\% \\
\bottomrule
\end{tabular}
\end{center}
\scriptsize
\T{%
\begin{itemize}
\item \textbf{一致率}：強いモダリティほど弱いモダリティと意見が合わないので，候補での一致率は強さの手がかりにならない（AFEW で −0.4〜−3.4）
\item \textbf{K 本のラベル}：単独の識別器の学習に使った発話なので，どのモダリティも K 本の上では当たりすぎて差が見えない．AFEW で最強の顔を当てた単位は 50 中 1〜2
\item \textbf{仮のターゲット}で選んだ指数：強いモダリティがコーパスで違うので移らない
\end{itemize}
}{%
\begin{itemize}
\item \textbf{Agreement rate}: the stronger a modality, the more it disagrees with the weak ones, so pool agreement is no cue for strength (AFEW: −0.4 to −3.4)
\item \textbf{The K labels}: used to train the unimodal classifiers, so every modality looks too good on them. In AFEW the strongest modality (face) is picked in only 1--2 of 50 units
\item \textbf{Proxy targets}: the strongest modality differs per corpus, so exponents chosen elsewhere do not transfer
\end{itemize}
}
\takeaway{\T{難しさの核：同じ K 本を学習にも信頼度の測定にも使うと，測定が壊れる}{The crux: using the same K labels both for training and for measuring trust breaks the measurement}}
\end{frame}

\begin{frame}{\T{いまの案：K 本を「1 塊ずつ抜いて」学び直し，抜いた分で当たるかを数える（leave-one-out）}{Current idea: hold out the K labels one cluster at a time, retrain, and score the held-out part (leave-one-out)}}
\rl{\T{方法（探索）}{Method (exploration)}}
\centering
\begin{adjustbox}{max width=\linewidth}
\begin{tikzpicture}[font=\scriptsize, x=1cm, y=1cm]
  % K 本（塊ごと）
  \node[font=\scriptsize\bfseries, anchor=west] at (-0.2,1.6) {\T{K 本 × 4 感情のラベル（話者の塊ごとに色分け）}{K × 4 labels (colored by speaker cluster)}};
  \foreach \i/\c in {0/kA, 1/kA, 2/kV, 3/kV, 4/kV, 5/kT, 6/kAVT, 7/kAVT} {\node[utt, fill=\c!25] at (0.62*\i,0.9) {};}
  \draw[mBad, line width=1pt, rounded corners=2pt] (1.0,0.55) rectangle (2.95,1.25);
  \node[font=\tiny, mBad] at (2.0,0.3) {\T{塊 2 を抜く}{hold out cluster 2}};
  % 学び直し
  \node[clf, text width=2.5cm] (fit) at (6.3,0.9) {\T{残りの K 本だけで}{retrain A, V, T and AVT}\\ \T{A・V・T・AVT を学び直す}{on the remaining labels only}};
  \draw[arr] (4.6,0.9) -- (fit);
  \node[testbox, text width=2.3cm] (held) at (10.2,0.9) {\T{抜いた塊の発話に}{predict the held-out}\\ \T{予測 → 当たるか}{cluster → right or not}};
  \draw[arr] (fit) -- (held);
  \node[font=\tiny, color=mMuted, align=center] at (6.3,-0.1) {\T{これを塊の数だけ繰り返す（AFEW は 1 単位あたり約 12〜30 塊）}{repeat for every cluster (AFEW: ≈12--30 clusters per unit)}};
  % 出力
  \node[modbox, text width=3.6cm] (rel) at (2.0,-1.5) {\T{モダリティごとの信頼度}{per-modality reliability}\\ \T{＝抜いた発話での UAR}{= UAR on held-out labels}\\ \T{→ 積の指数 $e_A, e_V, e_T$}{→ exponents $e_A, e_V, e_T$}};
  \node[modbox, text width=3.6cm] (est) at (6.3,-1.5) {\T{積の見込み／AVT の見込み}{expected UAR of product / of AVT}\\ \T{（同じ抜いた発話で測る）}{(both measured on the held-out labels)}};
  \node[propbox, text width=3.2cm] (sel) at (10.2,-1.5) {\T{単位ごとに選ぶ：}{choose per unit:}\\ \T{見込みが高い方（積 か AVT）}{the form with the higher estimate}\\ \T{をテストで使う}{and use it on the test set}};
  \draw[arr] (held.south) |- ++(-0.5,-0.55) -| (rel.north);
  \draw[arr] (held.south) |- ++(-0.5,-0.55) -| (est.north);
  \draw[arr] (est) -- (sel);
  \draw[arr] (rel.east) -- (est.west);
\end{tikzpicture}
\end{adjustbox}
\notreal
\scriptsize
\T{%
\begin{itemize}
\item 発話 1 本ずつ抜く版は，同じ話者の残りの発話から学ぶので話者が漏れ，信頼度が高く出る．そこで\textbf{塊（話者）ごと}に抜く
\item テストのラベルは使わない．新しい学習は「K 本だけの小さな識別器」で，1 単位あたり数秒〜数十秒
\end{itemize}
}{%
\begin{itemize}
\item Holding out one utterance at a time leaks the speaker (its other labels stay in training) and inflates reliability, so we hold out \textbf{whole clusters (speakers)}
\item No test labels are used. The new training involves only small K-label classifiers (seconds per unit)
\end{itemize}
}
\end{frame}

\begin{frame}{\T{結果 1：信頼度の見積もりは当たる．ただし「常に積」は AFEW でしか得をしない}{Result 1: the reliability estimate is right, but ``always product'' gains only in AFEW}}
\rl{\T{結果 1}{Result 1}}
\scriptsize
\begin{center}
\begin{tabular}{@{}lrrrrrrrr@{}}
\toprule
& \multicolumn{2}{c}{AFEW} & \multicolumn{2}{c}{IEMOCAP} & \multicolumn{2}{c}{MELD} & \multicolumn{2}{c}{MOSEI} \\
K & 5 & 25 & 5 & 25 & 5 & 25 & 5 & 25 \\
\midrule
\T{最強のモダリティを当てた割合（塊ごとの LOO）}{Strongest modality identified (cluster LOO, \%)} & \N{fclokAFEWfive}\% & \N{fclokAFEWtwofive}\% & \N{fclokIEMOCAPfive}\% & \N{fclokIEMOCAPtwofive}\% & \N{fclokMELDfive}\% & \N{fclokMELDtwofive}\% & \N{fclokMOSEIfive}\% & \N{fclokMOSEItwofive}\% \\
\midrule
\T{AVT の生徒（テスト UAR）}{AVT student (test UAR)} & \N{favtAFEWfive} & \N{favtAFEWtwofive} & \N{favtIEMOCAPfive} & \N{favtIEMOCAPtwofive} & \N{favtMELDfive} & \N{favtMELDtwofive} & \N{favtMOSEIfive} & \N{favtMOSEItwofive} \\
\T{常に積，LOO の信頼度の指数（差）}{Always product, LOO exponents (diff.)} & \N{fruledAFEWfive} & \N{fruledAFEWtwofive} & \N{fruledIEMOCAPfive} & \N{fruledIEMOCAPtwofive} & \N{fruledMELDfive} & \N{fruledMELDtwofive} & \N{fruledMOSEIfive} & \N{fruledMOSEItwofive} \\
\T{テストで選んだ上限（差）}{Test-chosen upper bound (diff.)} & \N{foracledAFEWfive} & \N{foracledAFEWtwofive} & \N{foracledIEMOCAPfive} & \N{foracledIEMOCAPtwofive} & \N{foracledMELDfive} & \N{foracledMELDtwofive} & \N{foracledMOSEIfive} & \N{foracledMOSEItwofive} \\
\bottomrule
\end{tabular}
\end{center}
\footnotesize
\T{%
\begin{itemize}
\item K = 25 では，塊ごとの LOO が最強のモダリティを AFEW・IEMOCAP・MELD で \N{fclokMELDtwofive}〜\N{fclokAFEWtwofive}\% 当てる．MOSEI は \N{fclokMOSEItwofive}\%（モダリティの差が小さく，塊も粗い）
\item 常に積を使うと AFEW で \N{fruledAFEWfive}・\N{fruledAFEWtwofive}（上限との差 1〜3），MOSEI で \N{fruledMOSEIfive}・\N{fruledMOSEItwofive}．MOSEI は上限でも AVT に届かず，\textbf{積という形そのものが負ける}
\item そこで「積か AVT か」を単位ごとに同じ LOO で選ぶ（前のページの右端）．AFEW で積，MOSEI で AVT を選べれば，4 コーパスで AVT を下回らない形になる → 次のページ
\end{itemize}
}{%
\begin{itemize}
\item At K = 25, cluster LOO identifies the strongest modality in \N{fclokMELDtwofive}--\N{fclokAFEWtwofive}\% of units for AFEW, IEMOCAP and MELD; MOSEI \N{fclokMOSEItwofive}\% (small modality differences, coarse clusters)
\item Always using the product gives \N{fruledAFEWfive} / \N{fruledAFEWtwofive} in AFEW (1--3 below the upper bound) but \N{fruledMOSEIfive} / \N{fruledMOSEItwofive} in MOSEI, where even the upper bound is below AVT: \textbf{the product form itself loses there}
\item So we choose ``product or AVT'' per unit with the same LOO (right end of the previous page). Picking product in AFEW and AVT in MOSEI would give a form that never falls below AVT → next page
\end{itemize}
}
\end{frame}

\begin{frame}{\T{結果 2：LOO で選ぶ形は AFEW で得をするが，MOSEI の損を消せない．まだ使える手法ではない}{Result 2: the LOO-selected form gains in AFEW but does not remove the MOSEI loss; not yet a usable method}}
\rl{\T{結果 2（4 コーパス）}{Result 2 (4 corpora)}}
\scriptsize
\begin{center}
\T{テスト UAR（AVT の生徒との差，ポイント．50 単位の平均）}{Test UAR (difference from the AVT student, points; mean over 50 units)}

\begin{tabular}{@{}lrrrrrrrr@{}}
\toprule
& \multicolumn{2}{c}{AFEW} & \multicolumn{2}{c}{IEMOCAP} & \multicolumn{2}{c}{MELD} & \multicolumn{2}{c}{MOSEI} \\
K & 5 & 25 & 5 & 25 & 5 & 25 & 5 & 25 \\
\midrule
\T{AVT の生徒（UAR そのもの）}{AVT student (absolute UAR)} & \N{savtAFEWfive} & \N{savtAFEWtwofive} & \N{savtIEMOCAPfive} & \N{savtIEMOCAPtwofive} & \N{savtMELDfive} & \N{savtMELDtwofive} & \N{savtMOSEIfive} & \N{savtMOSEItwofive} \\
\T{常に積（差）}{Always product (diff.)} & \N{sruledAFEWfive} & \N{sruledAFEWtwofive} & \N{sruledIEMOCAPfive} & \N{sruledIEMOCAPtwofive} & \N{sruledMELDfive} & \N{sruledMELDtwofive} & \N{sruledMOSEIfive} & \N{sruledMOSEItwofive} \\
\textbf{\T{LOO で選ぶ（差）}{LOO-selected (diff.)}} & \textbf{\N{sseldAFEWfive}} & \textbf{\N{sseldAFEWtwofive}} & \textbf{\N{sseldIEMOCAPfive}} & \textbf{\N{sseldIEMOCAPtwofive}} & \textbf{\N{sseldMELDfive}} & \textbf{\N{sseldMELDtwofive}} & \textbf{\N{sseldMOSEIfive}} & \textbf{\N{sseldMOSEItwofive}} \\
\T{単位ごとにテストで良い方（差，上限）}{Test-best per unit (diff., upper bound)} & \N{soracledAFEWfive} & \N{soracledAFEWtwofive} & \N{soracledIEMOCAPfive} & \N{soracledIEMOCAPtwofive} & \N{soracledMELDfive} & \N{soracledMELDtwofive} & \N{soracledMOSEIfive} & \N{soracledMOSEItwofive} \\
\midrule
\T{積を選んだ単位の割合}{Units choosing the product} & \N{sshareAFEWfive}\% & \N{sshareAFEWtwofive}\% & \N{sshareIEMOCAPfive}\% & \N{sshareIEMOCAPtwofive}\% & \N{sshareMELDfive}\% & \N{sshareMELDtwofive}\% & \N{sshareMOSEIfive}\% & \N{sshareMOSEItwofive}\% \\
\bottomrule
\end{tabular}
\end{center}
\vspace{-0.6em}
\scriptsize
\T{%
\begin{itemize}
\item MOSEI の損は「常に積」の半分ほどに減ったが消えず，4 コーパスの平均は AVT とほぼ同じ．理由：MOSEI で積がテストで良い単位は 3〜5\% しか無いのに，LOO の見込みは積のほうが高く出て（\N{sestpMOSEIfive} 対 \N{sestaMOSEIfive}，K = 5），\N{sshareMOSEIfive}\% の単位で積を選ぶ
\item 偏りの原因は 2 つ：(1) 指数と見込みを同じ抜いた発話から作る（楽観）．(2) 見込みは K 本だけの識別器で測るが，テストで使うのは直した識別器で，この差は 2,304 次元の AVT で大きい
\item 信頼度の見積もりは当たるが，その上の「積か AVT か」の選択はまだ使える形ではない．続けるなら，偏りを直す版を設計して検査に回す
\end{itemize}
}{%
\begin{itemize}
\item The MOSEI loss is roughly halved but not removed; averaged over corpora the form is about level with AVT. Why: in MOSEI the product is better on test in only 3--5\% of units, yet the LOO estimate favors the product (\N{sestpMOSEIfive} vs \N{sestaMOSEIfive} at K = 5) and picks it in \N{sshareMOSEIfive}\% of units
\item Two sources of the bias: (1) exponents and estimate come from the same held-out labels (optimistic); (2) the estimate uses K-label-only classifiers while the test uses adapted ones, a gap largest for the 2,304-dim AVT
\item The reliability estimate is right, but the product-vs-AVT selection on top of it is not yet usable. If this line continues, a bias-corrected version must be designed and checked
\end{itemize}
}
{\tiny\color{mMuted}\T{出どころ：figures/select.csv（results/explore\_fusion/select\_cells.csv．本番の台本 DAI/exp/explore\_fusion\_select.py，2026-11-05 10:17）．AFEW の値は保存ずみの対数確率から先に出した途中の値と小数 2 桁まで一致}{Source: figures/select.csv (results/explore\_fusion/select\_cells.csv; script DAI/exp/explore\_fusion\_select.py, 2026-11-05 10:17). AFEW values agree to 2 decimals with the interim values computed earlier from the saved log-probabilities}}
\end{frame}

\begin{frame}{\T{いまの状態の一覧}{Where things stand}}
\rl{\T{状態}{Status}}
\footnotesize
\begin{tabular}{@{}p{0.44\linewidth}p{0.22\linewidth}p{0.30\linewidth}@{}}
\toprule
\T{項目}{Item} & \T{状態}{Status} & \T{次}{Next} \\
\midrule
\T{実験 1〜3 の評価の回，集計，結果の文の照合}{Exp.~1--3 evaluation, aggregation, claim check} & \done & \T{論文の結果の節に写す}{move into the paper's results section} \\
\T{得の仕組みの探索（先回り，一致，位置と向き）}{Mechanism exploration (lead, agreement, position/direction)} & \done & \T{論文の「説明」の節（事前の解析と分けて）}{``explanation'' section, separated from pre-fixed analyses} \\
\T{実験 4（一致で重みづけた疑似ラベル）}{Exp.~4 (agreement-weighted pseudo-labels)} & \todo\ \T{設計・台本は検査済み}{design and code checked} & \T{研究者の判断で起動（H200，4〜6 時間）}{launch at the researcher's decision (H200, 4--6 h)} \\
\T{LOO による信頼度と，積／AVT の選択}{LOO reliability and product/AVT selection} & \done\ \T{4 コーパス}{4 corpora} & \T{続けるかを研究者が決める（見込みの偏りを直す版の設計が要る）}{researcher decides whether to continue (a version that removes the estimate bias would be needed)} \\
\T{第 2 のエンコーダでの繰り返し}{Replication with a second encoder set} & \todo & \T{保留（研究者「まだいい」）}{on hold (researcher: ``not yet'')} \\
\T{論文の型（手法の論文 → 測定と説明の論文）と骨組みの項目 3・4・6 の書き直し}{Paper type (method paper → measurement-and-explanation paper), rewriting skeleton items 3, 4, 6} & \wip & \T{\S7 の案を研究者が決める}{researcher decides on the \S7 drafts} \\
\bottomrule
\end{tabular}
\end{frame}


\appendix
\bandheadline
\metroset{numbering=counter}\setbeamertemplate{frame numbering}{\insertpagenumber/\insertdocumentendpage}
%% 補足（前半）：本編から移した語・全体像・比較・蒸留の違い
\setband{\T{補足}{Appendix}}

\begin{frame}{\T{この資料で使う語（8 つ）}{Terms used in this deck (8)}}
\rl{\T{語}{Terms}}
\footnotesize
\begin{tabular}{@{}p{0.27\linewidth}p{0.70\linewidth}@{}}
\toprule
\T{語}{Term} & \T{意味}{Meaning} \\
\midrule
\T{ソース／ターゲット}{Source / target} & \T{ラベル付きの既存コーパス（3 つ）／新しいコーパス（1 つ）．4 つを順にターゲットにする}{The 3 labeled existing corpora / the 1 new corpus. Each of the 4 corpora takes the target role in turn} \\
\T{K 本}{K labels} & \T{ターゲットの感情ごとに K 本のラベル付き発話（K = 25 なら 4 感情 × 25 = 100 本）}{K labeled target utterances per emotion (K = 25 means 4 × 25 = 100)} \\
\T{候補}{Pool} & \T{ターゲットのラベル無しの発話（学習に使える．テストの話者は含まない）}{Unlabeled target utterances usable for training (test speakers excluded)} \\
\T{出発点}{Start} & \T{ソースのコーパスごとの識別器の重みと切片の平均．学習の初期値であり，引き戻す先}{Mean of the per-source classifiers (weights and intercepts). Initial value and the anchor of the pull} \\
\T{引き戻し}{Pull} & \T{出発点から離れすぎないようにする L2 正則化（強さは仮のターゲットで選んで固定）}{L2 regularization toward the start (strength fixed on proxy targets)} \\
\T{教師／生徒}{Teacher / student} & \T{教師＝ラベル無しの発話に疑似ラベルを付けるモデル（3 モダリティの識別器の確率の積）．生徒＝ターゲットで使う識別器（音声・顔・言葉・AVT）}{Teacher = the model that pseudo-labels the pool (product of the three unimodal classifiers' probabilities). Student = the classifier to deploy (audio, face, text, AVT)} \\
\T{単位}{Unit} & \T{（ターゲット，分割，種）の 1 組．コーパスごとに 5 分割 × 10 種 = 50 単位}{One (target, fold, seed) triple. 5 folds × 10 seeds = 50 units per corpus} \\
UAR & \T{4 感情の再現率の平均（でたらめなら 25\%）．差は「ポイント」で書く}{Mean recall over the 4 emotions (chance = 25\%). Differences are in ``points''} \\
\bottomrule
\end{tabular}
\end{frame}

\begin{frame}{\T{実験の全体像：4 つの実験で「どれだけ・どこで・なぜ」を分けて測る}{Overview: four experiments separate ``how much'', ``where'' and ``why''}}
\rl{\T{計画}{Plan}}
\scriptsize
\begin{tabular}{@{}p{0.07\linewidth}p{0.42\linewidth}p{0.27\linewidth}p{0.17\linewidth}@{}}
\toprule
& \T{何を測るか}{What is measured} & \T{相手}{Compared against} & \T{状態}{Status} \\
\midrule
\T{実験 1}{Exp.~1} & \T{別のモダリティの予測（疑似ラベル）で学ぶ手法の得}{Gain from learning on the other modalities' predictions (pseudo-labels)} & \T{K 本だけ；同じ情報を受け取る 5 手法（温度つき蒸留など）}{K-only; 5 methods using the same information (e.g.\ temperature distillation)} & \done \\
\T{実験 2}{Exp.~2} & \T{標準のドメイン適応・半教師あり手法との差}{Difference from standard DA / semi-supervised methods} & \T{自己学習，エントロピー最小化，T3A，ラベル伝播，S+T，CORAL，DANN，MME など 14}{self-training, entropy min., T3A, label propagation, S+T, CORAL, DANN, MME, … (14)} & \done \\
\T{実験 3}{Exp.~3} & \T{得の説明：損を切片（位置）と重みの向きに分ける，再現できる度合い}{Explanation: split the loss into intercept (position) and weight direction; reproducibility} & \T{事前に固定した判定の規則}{pre-fixed decision rules} & \done \\
\T{実験 4}{Exp.~4} & \T{蒸留の弱点（全発話を同じ重みで信じる）を直した手法の得}{Gain from fixing a weakness of distillation (uniform trust in all pool utterances)} & \T{実験 1 の提案・蒸留}{Exp.~1 proposal / distillation} & \todo\ \T{設計済み}{designed} \\
\T{探索}{Exploration} & \T{どのモダリティを信じるか（信頼度の見積もり）}{Which modality to trust (reliability estimation)} & --- & \done\ \T{測定済み}{measured} \\
\bottomrule
\end{tabular}

\vspace{0.3em}
\scriptsize
\T{実験 1〜3 は設計・台本・集計の台本を評価の前に固定し，関門（統計・漏れ・比較・定義・台本）と ML・感情認識の査読想定の検査を通してから計算した．実験 4 と探索は評価の結果を見た後の案で，論文では節を分け，限界に由来を書く．}{Experiments 1--3: designs, scripts and aggregation were frozen before evaluation and passed the gate checks (statistics, leakage, fairness, definitions, code) and simulated ML / affective-computing review. Experiment 4 and the exploration arose after seeing the results; the paper separates them and states their origin.}
\end{frame}

\begin{frame}{\T{評価の尺度と比較：UAR で測り，使う情報ごとに相手を分ける}{Measure and comparisons: UAR, with baselines grouped by the information they use}}
\rl{\T{評価（骨組み 5・6）}{Evaluation (skeleton 5, 6)}}
\hosoku{appmethods}{\T{手法と格子の一覧}{list of methods and grids}}
\begin{columns}[T]
\begin{column}{0.46\textwidth}
\footnotesize
\textbf{\T{尺度（骨組み 5）}{Measure (skeleton 5)}}
\T{%
\begin{itemize}
\item UAR（4 感情の再現率の平均）．MELD と MOSEI は Neutral が半分近くを占め，正解率だと「全部 Neutral」が高く出るため
\item テストの話者は学習側に入れない（MOSEI は話者の代わりにチャンネルか動画）
\end{itemize}
\textbf{除去（骨組み 6）}：提案から 1 つずつ外した版を実験 1 で計算済み
\begin{itemize}
\item (i) 教師を学び直さない　(ii) 教師を 1 モダリティにする
\item (iii) 総量をそろえない（＝温度 1 の蒸留）　(iv) 切片だけを K 本で学ぶ
\end{itemize}
}{%
\begin{itemize}
\item UAR (mean recall over 4 emotions). In MELD and MOSEI nearly half the utterances are Neutral, so accuracy rewards ``always Neutral''
\item Test speakers never enter training (MOSEI: channel or video stands in for speaker)
\end{itemize}
\textbf{Ablations (skeleton 6)}: versions with one part removed, already computed in Exp.~1
\begin{itemize}
\item (i) teacher not retrained\quad (ii) single-modality teacher
\item (iii) no balancing (= distillation at T = 1)\quad (iv) intercept only from K labels
\end{itemize}
}
\end{column}
\begin{column}{0.54\textwidth}
\footnotesize
\textbf{\T{比較の相手（骨組み 6）}{Baselines (skeleton 6)}}
\begin{tabular}{@{}p{0.30\linewidth}p{0.66\linewidth}@{}}
\toprule
\T{使う情報}{Information used} & \T{相手}{Baselines} \\
\midrule
\T{K 本だけ}{K labels only} & \T{比較の基準（出発点から K 本で学び直す）}{the baseline (start retrained on K labels)} \\
\T{K 本 ＋ 同じ教師 Q}{K labels + the same teacher Q} & \T{温度つき蒸留，一致した発話だけ，自分の予測と混ぜる，柔らかい共訓練，共正則化（実験 1）}{temperature distillation, agreed-only, mixing with own prediction, soft co-training, co-regularization (Exp.~1)} \\
\T{K 本 ＋ ラベル無し（同じモダリティ）}{K labels + unlabeled (same modality)} & \T{自己学習，エントロピー最小化，IM，T3A，ラベル伝播，自分の予測を提案の形で（実験 2）}{self-training, entropy min., IM, T3A, label propagation, own prediction in the proposal's form (Exp.~2)} \\
\T{＋ ソースのデータ}{+ source data} & \T{S+T，CORAL，DANN，MME，xMUDA 型（実験 2）}{S+T, CORAL, DANN, MME, xMUDA-style (Exp.~2)} \\
\bottomrule
\end{tabular}
\vspace{0.3em}
{\scriptsize\color{mMuted}\T{全手法が同じ出発点・同じ標準化・同じ K 本を使い，つまみは仮のターゲットで選んで固定した}{All methods share the start, the standardization and the K labels; knobs were selected on proxy targets and frozen}}
\end{column}
\end{columns}
\end{frame}

\begin{frame}{\T{提案と温度つき蒸留の違いは 2 点だけで，評価ではほぼ同じ値になった}{The proposal differs from temperature distillation in only two points, and scored the same in evaluation}}
\rl{\T{手法}{Method}}
\begin{columns}[T]
\begin{column}{0.58\textwidth}
\small
\begin{tabular}{@{}p{0.30\linewidth}p{0.31\linewidth}p{0.31\linewidth}@{}}
\toprule
& \T{提案}{Proposal} & \T{温度つき蒸留（Hinton 2015，Albanie 2018 の形）}{Temperature distillation (Hinton 2015; Albanie 2018 form)} \\
\midrule
\T{教師}{Teacher} & \T{同じ Q}{same Q} & \T{同じ Q}{same Q} \\
\T{疑似ラベルの形}{Pseudo-label form} & \T{Q そのまま}{Q as is} & \T{Q を温度 2 で柔らげる}{Q softened with temperature 2} \\
\T{感情ごとの総量}{Per-emotion total weight} & \T{そろえる}{balanced} & \T{そろえない}{not balanced} \\
\T{K 本・引き戻し}{K labels, pull} & \T{同じ}{same} & \T{同じ}{same} \\
\T{つまみ（格子）}{Knobs (grid)} & \T{全体の重み ρ，K 本の倍率}{overall weight ρ, label multiplier} & \T{温度，全体の重み ρ}{temperature, overall weight ρ} \\
\bottomrule
\end{tabular}
\end{column}
\begin{column}{0.42\textwidth}
\small
\T{%
\begin{itemize}
\item 評価の回で，提案 − 蒸留 は 48 セル中 40 セルで区間が 0 を含む（\S3）
\item 残りの差も −1.2〜+2.2 ポイントで，向きはコーパスで分かれる
\item したがって解析では\textbf{両者を区別せず}，「別のモダリティの予測（教師）で学ぶ手法」として扱う
\end{itemize}
}{%
\begin{itemize}
\item In evaluation, proposal − distillation has an interval containing 0 in 40 of 48 cells (\S3)
\item The remaining differences are −1.2 to +2.2 points, with signs varying by corpus
\item Hence the analysis \textbf{does not distinguish} them: both are ``learning from the other modalities' predictions (the teacher)''
\end{itemize}
}
\end{column}
\end{columns}
\takeaway{\T{「蒸留が良い理由」を問うときの「蒸留」は，この 2 手法の共通部分＝教師 Q で学ぶことを指す}{When we ask ``why distillation helps'', ``distillation'' means what the two share: learning from the teacher Q}}
\end{frame}

%% 補足（Appendix）
\setband{\T{補足}{Appendix}}

\begin{frame}[label=appdata]{\T{補足：コーパスの詳細}{Appendix: corpus details}}
\rl{\T{データ}{Data}}
\scriptsize
\begin{tabular}{@{}lrrrrp{0.33\linewidth}@{}}
\toprule
\T{コーパス}{Corpus} & \T{発話}{Utter.} & \T{塊}{Clusters} & \T{実効の塊}{Eff.\ clusters} & \T{感情ごとの最少}{Min.\ per emotion} & \T{塊の定義と注意}{Cluster definition and caveats} \\
\midrule
AFEW & 790 & 51 & 15.0 & 127 & \T{俳優と作品のつながり．言葉は Whisper large-v3 の書き起こし}{actor–film connectivity; text from Whisper large-v3 transcripts} \\
IEMOCAP\T{（即興）}{ (improvised)} & 2,943 & 5 & 4.9 & 199 & \T{セッション．顔のマーカーの手がかりが入りうる}{sessions; face-marker cues may leak into face features} \\
MELD & 11,267 & 272 & 8.6 & 465 & \T{話者．テストの対話の発話は学習側から除く．分割 0〜3 は主役 1 人が 70〜82\%}{speakers; test dialogues removed from training; one lead actor holds 70--82\% of folds 0--3} \\
MOSEI & 18,634 & 2,614 & 783 & 1,080 & \T{チャンネルか動画（話者の代用．同じ人物の別の動画が分割をまたぎうる）}{channel or video (stand-in for speaker; the same person may cross folds in other videos)} \\
\bottomrule
\end{tabular}

\vspace{0.3em}
\T{%
\begin{itemize}
\item 特徴：音声 WavLM（1,024 次元），顔 ResNet-50 の平均（512），言葉 RoBERTa（768），AVT はつないで 2,304．エンコーダは微調整しない
\item 4 感情（Neutral，Happy，Sad，Angry）．標準化は単位ごとに候補（K 本を含み，テストを含まない）の平均と標準偏差で行い，全手法で同じ
\item K 本は学習側から感情ごとに乱数で抜く（K = 1，5，10，25 で入れ子）．評価の回は K ∈ \{0, 5, 25\} を主に報告
\item 塊の数が少ないので，区間の自由度は AFEW 13，IEMOCAP 3，MELD 7（t 分布）．MOSEI は塊の引き直し 2,000 回
\end{itemize}
}{%
\begin{itemize}
\item Features: audio WavLM (1,024 dims), face ResNet-50 mean (512), text RoBERTa (768); AVT concatenates them (2,304). Encoders are not fine-tuned
\item 4 emotions (Neutral, Happy, Sad, Angry). Standardization per unit with the pool's mean and SD (includes the K labels, excludes the test); identical for all methods
\item K labels are drawn at random per emotion from the training side (nested over K = 1, 5, 10, 25); the evaluation round reports K ∈ \{0, 5, 25\}
\item Few clusters: interval degrees of freedom AFEW 13, IEMOCAP 3, MELD 7 (t distribution); MOSEI uses 2,000 cluster resamples
\end{itemize}
}
\end{frame}

\begin{frame}[label=appsolver]{\T{補足：学習器，出発点，引き戻し，教師の作り方}{Appendix: solver, start, pull, and how the teacher is built}}
\rl{\T{手法}{Method}}
\scriptsize
\T{%
\begin{itemize}
\item \textbf{識別器}：4 感情の多項ロジスティック回帰（重み W，切片 b）．float64 の L-BFGS で，問をまとめて解く
\item \textbf{ソースの識別器}：コーパスごと・種類ごと（A，V，T，AV，AT，VT，AVT）に，感情の釣り合いの重みつき交差エントロピー ＋ L2 で学ぶ．強さ C はそのコーパスの塊の 5 分割の交差検証で選ぶ
\item \textbf{出発点}：ソース 3 つの W と b の平均（等しい重み）
\item \textbf{引き戻し}：W については $(C'/2)\,\|W - W_0\|^2$，b については $(C_b/2)\,\|b - b_0\|^2$．$(C', C_b)$ は 7 × 2 の格子から，仮のターゲットの「K 本だけ」の UAR で (コーパス，種類，K) ごとに選んで固定．全手法が同じ値を使う
\item \textbf{教師 Q}：(1) 出発点の AVT を K 本で学び直して候補に予測を付ける．(2) A・V・T の識別器を出発点から「K 本 ＋ その予測（感情ごとに総量をそろえる）」で学び直す．(3) 3 つの確率を掛けて正規化する（指数は (1,1,1)．仮のターゲットで指数を選ぶ版も用意したが，主の版は 4 コーパス全 K で (1,1,1) が選ばれた）
\item \textbf{生徒の損失}：K 本の交差エントロピー × a ＋ 疑似ラベル Q の交差エントロピー（重みの合計 ρn，感情ごとに総量をそろえる）＋ 引き戻し．ρ ∈ \{1/16, 1/4, 1, 4\}，a ∈ \{1, 4, 16\} から仮のターゲットで選ぶ（80 セル中 75 で ρ = 1，a = 1）
\end{itemize}
}{%
\begin{itemize}
\item \textbf{Classifier}: 4-class multinomial logistic regression (weights W, intercept b), solved in float64 by batched L-BFGS
\item \textbf{Source classifiers}: per corpus and per feature type (A, V, T, AV, AT, VT, AVT), class-balanced cross-entropy + L2; the strength C is chosen by 5-fold cluster cross-validation within that corpus
\item \textbf{Start}: mean of the three sources' W and b (equal weights)
\item \textbf{Pull}: $(C'/2)\,\|W - W_0\|^2$ for the weights and $(C_b/2)\,\|b - b_0\|^2$ for the intercept. $(C', C_b)$ chosen from a 7 × 2 grid on proxy targets by the K-only UAR, per (corpus, type, K), then frozen; every method uses the same values
\item \textbf{Teacher Q}: (1) retrain the AVT start on the K labels and predict the pool; (2) retrain A, V, T from the start on ``K labels + those predictions (balanced per emotion)''; (3) multiply the three probabilities and normalize (exponents (1,1,1); a version with proxy-selected exponents exists, but (1,1,1) was selected for all corpora and K)
\item \textbf{Student loss}: a × cross-entropy on K labels + cross-entropy to Q (total weight ρn, balanced per emotion) + pull. ρ ∈ \{1/16, 1/4, 1, 4\}, a ∈ \{1, 4, 16\}, chosen on proxy targets (ρ = 1, a = 1 in 75 of 80 cells)
\end{itemize}
}
\end{frame}

\begin{frame}[label=appmethods]{\T{補足：比較の手法と格子}{Appendix: baseline methods and grids}}
\rl{\T{手法}{Method}}
\tiny
\begin{tabular}{@{}p{0.26\linewidth}p{0.50\linewidth}p{0.20\linewidth}@{}}
\toprule
\T{手法}{Method} & \T{形（すべて K 本と引き戻しを含む）}{Form (all include the K labels and the pull)} & \T{格子（点の数）}{Grid (points)} \\
\midrule
\T{温度つき蒸留}{Temperature distillation} & \T{Q を温度 T で柔らげ，総量はそろえない（Hinton 2015，Albanie 2018）}{Q softened at temperature T, no balancing (Hinton 2015; Albanie 2018)} & T ∈ \{1,2,4\}, ρ 4 (12) \\
\T{一致した発話だけ}{Agreed-only} & \T{Q と生徒自身の判定が一致し確信度 ≥ τ の発話に硬いラベル}{hard labels where Q agrees with the student and confidence ≥ τ} & τ 3 × f 3 (9) \\
\T{自分の予測と混ぜる}{Mixing with own prediction} & \T{α × 生徒の確率 ＋ (1−α) × Q の最大が 0.9 超の発話に硬いラベル（Zhang 2021 の形）}{hard label where α × own + (1−α) × Q exceeds 0.9 (Zhang 2021 form)} & α 3 × f 3 (9) \\
\T{柔らかい共訓練}{Soft co-training} & \T{A・V・T を 3 巡，各々「ほかの 2 つの積」で学び直す}{A, V, T retrained for 3 rounds, each on the product of the other two} & ρ (4) \\
\T{共正則化}{Co-regularization} & \T{3 つの識別器の点数の差の二乗を罰する（凸）}{penalize squared score differences among the three classifiers (convex)} & ρ (5) \\
\midrule
\T{自己学習}{Self-training} & \T{生徒自身の確信度 ≥ τ の候補に硬いラベル，3 巡}{hard labels on own confident pool utterances, 3 rounds} & τ 3 × ρ 4 (12) \\
\T{エントロピー最小化／IM}{Entropy min.\ / IM} & \T{候補の予測のエントロピーを下げる（IM は周辺の項を足す）}{lower the pool prediction entropy (IM adds the marginal term)} & ρ 4 / α 3 × ρ 4 \\
T3A & \T{支えの集合で分類する試験時適応（Iwasawa \& Matsuo 2021）}{test-time adaptation with a support set (Iwasawa \& Matsuo 2021)} & M (6) \\
\T{ラベル伝播}{Label propagation} & \T{k 近傍グラフで K 本のラベルを伝播し疑似ラベルに}{propagate the K labels on a k-NN graph as pseudo-labels} & k 3 × α 2 × 2 (12) \\
\T{自分の予測を提案の形で}{Own prediction, proposal form} & \T{Q を生徒自身の確率に替えた提案}{the proposal with Q replaced by the student's own probabilities} & ρ 4 × a 3 (12) \\
S+T & \T{ソースの全ラベル付き発話と K 本（倍率 β）を合わせて学ぶ}{train on all labeled source utterances plus the K labels (multiplier β)} & β (3) \\
CORAL & \T{ソースの共分散を候補にそろえ，出発点を作り直す}{align source covariance to the pool and rebuild the start} & r (3) \\
DANN / MME / xMUDA\T{型}{-style} & \T{線形の写像 ＋ 領域の判別器／エントロピーの minimax／相互の模倣（勾配法，種 3）}{linear map + domain discriminator / entropy minimax / mutual mimicry (gradient methods, 3 seeds)} & 24 / 24 / 24 \\
\bottomrule
\end{tabular}
\vspace{0.2em}
{\tiny\T{すべて凍結した特徴の上の線形の形に移した版で，原論文と同じ手法とは書かない．つまみは仮のターゲットの 120 単位で選び，固定した．}{All are versions transposed to linear classifiers on frozen features and are not claimed to equal the original methods. Knobs were selected on 120 proxy-target units and frozen.}}
\end{frame}

\begin{frame}[label=appstats]{\T{補足：区間と検定の規則}{Appendix: interval and test rules}}
\rl{\T{統計}{Statistics}}
\scriptsize
\T{%
\begin{itemize}
\item \textbf{族}：族 1 ＝ 提案 − K 本だけ，族 2 ＝ 役割を分ける版 − 同時に解く版，族 3 ＝ 提案 − 同じ情報の相手（仮のターゲットで最良だった 1 手法），族 4 ＝ 選んだ設定 − 選択なし．各族はコーパスごとに 12 セル（生徒 4 × K 3）で Holm
\item \textbf{区間}：AFEW・IEMOCAP・MELD は塊の jackknife の分散 ＋ 種（K 本の抜き方）の分散に t 分布（自由度 13・3・7）．MOSEI は塊の引き直し 2,000 回 ＋ 種の分散に正規分布．学習側の揺れとソースの識別器の揺れは含まない
\item \textbf{検出力による振り分け}：評価の前に，仮のターゲットの標準誤差から「検出できる差」を見積もり，見込みが足りるセルだけ検定に残した（族 1〜3 で 15 セル，族 4 で 40 セル）．残り 113 セルは区間つきの記述
\item \textbf{族 4}：検定に残った 40 セルは全部「選んだ設定 ＝ 選択なし」で差が 0．選ぶ効果は測れていない
\item \textbf{族 3 のコーパスの文}：「このコーパスで提案は相手より良い」は，事前の条件（全生徒で下端が 0 より上など）をどのコーパスも満たさず，書かない
\item \textbf{p 値の意味}：手法の形は評価と同じ 4 コーパスで作ったので，確認の検定ではない
\end{itemize}
}{%
\begin{itemize}
\item \textbf{Families}: 1 = proposal − K-only; 2 = role-separated − joint; 3 = proposal − same-information baseline (the best of 5 on proxy targets); 4 = selected − unselected setting. Each family has 12 cells per corpus (4 students × 3 K), Holm-corrected
\item \textbf{Intervals}: AFEW, IEMOCAP, MELD use cluster jackknife variance + seed variance with a t distribution (df 13, 3, 7); MOSEI uses 2,000 cluster resamples + seed variance with a normal. Training-side and source-classifier variability are not included
\item \textbf{Power screening}: before evaluation, the detectable difference was estimated from proxy-target standard errors and only cells with sufficient expected power were kept for testing (15 cells in families 1--3, 40 in family 4); the other 113 cells are descriptions
\item \textbf{Family 4}: all 40 tested cells have ``selected = unselected'' and a difference of exactly 0; the selection effect is not measured
\item \textbf{Corpus-level sentence for family 3}: ``the proposal beats the baseline on this corpus'' is not written because no corpus meets the pre-fixed condition
\item \textbf{Meaning of p-values}: not confirmatory, since method forms were built on the same 4 corpora
\end{itemize}
}
\end{frame}

% セルの値：#1 = famone|famthree，#2 = コーパス，#3 = 生徒，#4 = K（zero|five|twofive）
\newcommand{\fc}[4]{\N{#1#2#3#4}\,{\color{mMuted}[\N{#1lo#2#3#4}, \N{#1hi#2#3#4}]}}
\newcommand{\famtable}[1]{%
\begin{tabular}{@{}llrrrr@{}}
\toprule
\T{生徒}{Student} & K & AFEW & IEMOCAP & MELD & MOSEI \\
\midrule
\T{音声}{audio} & 0 & \fc{#1}{AFEW}{A}{zero} & \fc{#1}{IEMOCAP}{A}{zero} & \fc{#1}{MELD}{A}{zero} & \fc{#1}{MOSEI}{A}{zero} \\
 & 5 & \fc{#1}{AFEW}{A}{five} & \fc{#1}{IEMOCAP}{A}{five} & \fc{#1}{MELD}{A}{five} & \fc{#1}{MOSEI}{A}{five} \\
 & 25 & \fc{#1}{AFEW}{A}{twofive} & \fc{#1}{IEMOCAP}{A}{twofive} & \fc{#1}{MELD}{A}{twofive} & \fc{#1}{MOSEI}{A}{twofive} \\
\T{顔}{face} & 0 & \fc{#1}{AFEW}{V}{zero} & \fc{#1}{IEMOCAP}{V}{zero} & \fc{#1}{MELD}{V}{zero} & \fc{#1}{MOSEI}{V}{zero} \\
 & 5 & \fc{#1}{AFEW}{V}{five} & \fc{#1}{IEMOCAP}{V}{five} & \fc{#1}{MELD}{V}{five} & \fc{#1}{MOSEI}{V}{five} \\
 & 25 & \fc{#1}{AFEW}{V}{twofive} & \fc{#1}{IEMOCAP}{V}{twofive} & \fc{#1}{MELD}{V}{twofive} & \fc{#1}{MOSEI}{V}{twofive} \\
\T{言葉}{text} & 0 & \fc{#1}{AFEW}{T}{zero} & \fc{#1}{IEMOCAP}{T}{zero} & \fc{#1}{MELD}{T}{zero} & \fc{#1}{MOSEI}{T}{zero} \\
 & 5 & \fc{#1}{AFEW}{T}{five} & \fc{#1}{IEMOCAP}{T}{five} & \fc{#1}{MELD}{T}{five} & \fc{#1}{MOSEI}{T}{five} \\
 & 25 & \fc{#1}{AFEW}{T}{twofive} & \fc{#1}{IEMOCAP}{T}{twofive} & \fc{#1}{MELD}{T}{twofive} & \fc{#1}{MOSEI}{T}{twofive} \\
AVT & 0 & \fc{#1}{AFEW}{AVT}{zero} & \fc{#1}{IEMOCAP}{AVT}{zero} & \fc{#1}{MELD}{AVT}{zero} & \fc{#1}{MOSEI}{AVT}{zero} \\
 & 5 & \fc{#1}{AFEW}{AVT}{five} & \fc{#1}{IEMOCAP}{AVT}{five} & \fc{#1}{MELD}{AVT}{five} & \fc{#1}{MOSEI}{AVT}{five} \\
 & 25 & \fc{#1}{AFEW}{AVT}{twofive} & \fc{#1}{IEMOCAP}{AVT}{twofive} & \fc{#1}{MELD}{AVT}{twofive} & \fc{#1}{MOSEI}{AVT}{twofive} \\
\bottomrule
\end{tabular}}

\begin{frame}[label=appfam1]{\T{補足：族 1（提案 − K 本だけ）の全セル}{Appendix: family 1 (proposal − K-only), all cells}}
\rl{\T{実験 1}{Exp.~1}}
\scriptsize\centering
\famtable{famone}

\vspace{0.3em}
\T{UAR ポイント，[95\% 区間]．K = 0 は 5 単位，K = 5・25 は 50 単位．検定に残ったセルは MELD の音声 K = 0・5，MOSEI の音声 K = 0・5・25，MOSEI の言葉 K = 0・5，MOSEI の AVT K = 0・5・25（Holm 後に有意なのは MOSEI の音声の 3 セル）．}{UAR points, [95\% interval]. K = 0: 5 units; K = 5, 25: 50 units. Tested cells: MELD audio K = 0, 5; MOSEI audio K = 0, 5, 25; MOSEI text K = 0, 5; MOSEI AVT K = 0, 5, 25 (Holm-significant: the three MOSEI audio cells).}
\end{frame}

\begin{frame}[label=appfam3]{\T{補足：族 3（提案 − 同じ情報の相手）の全セル}{Appendix: family 3 (proposal − same-information baseline), all cells}}
\rl{\T{実験 1}{Exp.~1}}
\scriptsize\centering
\famtable{famthree}

\vspace{0.3em}
\T{相手は 80 セル中 75 セルで温度つき蒸留（版は仮のターゲットで選ぶ），5 セルで「一致した発話だけ」．相手の版が提案の版（指数 (1,1,1)）と違う 10 セルでは，版をそろえた値も論文に並べる（例：MOSEI の言葉 K = 0 は −0.43 [−0.76, −0.11]）．検定に残ったのは MOSEI の AVT K = 0・5（有意でない）と MOSEI の言葉 K = 0（蒸留が高い）．}{The baseline is temperature distillation in 75 of 80 cells (version chosen on proxy targets) and ``agreed-only'' in 5. In the 10 cells where the baseline's version differs from the proposal's (exponents (1,1,1)), version-aligned values are also reported (e.g.\ MOSEI text K = 0: −0.43 [−0.76, −0.11]). Tested cells: MOSEI AVT K = 0, 5 (not significant) and MOSEI text K = 0 (distillation higher).}
\end{frame}

\begin{frame}[label=appdai2]{\T{補足：実験 2（AVT の生徒）の K = 0 と K = 25}{Appendix: Exp.~2 (AVT student) at K = 0 and K = 25}}
\rl{\T{実験 2}{Exp.~2}}
\centering
{\scriptsize K = 0}\par
\daichart{0}{-9}{30}{3.3cm}
{\scriptsize K = 25}\par
\daichart{25}{-9}{30}{3.3cm}
\end{frame}

\begin{frame}[label=appverdicts]{\T{補足：実験 3 の事前に固定した判定の結果}{Appendix: outcomes of the pre-fixed decision rules of Exp.~3}}
\rl{\T{実験 3}{Exp.~3}}
\scriptsize
\begin{tabular}{@{}p{0.40\linewidth}p{0.18\linewidth}p{0.38\linewidth}@{}}
\toprule
\T{説明}{Explanation} & \T{判定}{Verdict} & \T{根拠}{Basis} \\
\midrule
\T{1. K 本は主に切片を直し，別のモダリティの予測は向きを直す}{1. K labels mainly fix the intercept; the other modalities' predictions fix the direction} & \T{判定できない（4 コーパス）}{undetermined (4 corpora)} & \T{支持の条件を満たすセル 0．切片だけの版は K 本だけを 18 セルで下回る．「向きの差」が上回るのは IEMOCAP 5／8，MOSEI 5／8 のみ}{0 cells meet the support condition; the intercept-only version is below K-only in 18 cells; the direction difference is positive in only IEMOCAP 5/8 and MOSEI 5/8} \\
\T{2a. 積が強いのは学び直しのため／2b. 確率を掛けること自体が効く}{2a. The product is strong because of retraining / 2b. multiplying probabilities itself helps} & \T{判定できない}{undetermined} & \T{2b の条件を満たすのは MELD（5／8）だけ}{only MELD (5/8) meets the 2b condition} \\
\T{3. 生徒は再現できる部分だけ受け取る}{3. The student takes only what it can reproduce} & \T{判定できない}{undetermined} & \T{Freedman–Lane の並べ替えで 4 コーパスとも}{Freedman–Lane permutation, all 4 corpora} \\
\T{5. IEMOCAP の言葉の得は筋書きの話題では説明できない}{5. IEMOCAP text gain is not explained by scripted topics} & \T{判定できない}{undetermined} & \T{筋書きの対の区間が 0 を含む（幅は最大 72 ポイント）．言葉の文は IEMOCAP を含めた数と除いた数の両方で書く}{paired intervals contain 0 (width up to 72 points); text sentences report counts with and without IEMOCAP} \\
\T{族 5：感情の対ごとの差（切片を当て直した後，K = 5）}{Family 5: per-emotion-pair differences after refit (K = 5)} & \T{5 セルが有意}{5 significant cells} & \T{MOSEI の音声・言葉の Neutral–Happy と Happy–Sad，MOSEI の AVT の Neutral–Happy}{MOSEI audio and text Neutral–Happy and Happy–Sad; MOSEI AVT Neutral–Happy} \\
\bottomrule
\end{tabular}
\vspace{0.3em}
{\scriptsize\T{説明 4（顔の切片）と 6 は記述のみ．向きの物差しの確かめ 2 に不合格の 5 組（AFEW の顔・言葉・AVT，IEMOCAP の顔・AVT）では「向き」と書かない．}{Explanations 4 (face intercept) and 6 are descriptive only. In the 5 pairs failing validity check 2 (AFEW face, text, AVT; IEMOCAP face, AVT), ``direction'' is not used.}}
\end{frame}

% 位置と向きの表：#1 = K（five|twofive），#2 = 見出し
\newcommand{\dr}[3]{\N{pos#1#2#3} / \N{dir#1#2#3}}
\newcommand{\decomptable}[2]{%
\setlength{\tabcolsep}{3pt}\begin{tabular}{@{}llrrr@{}}
\toprule
\T{コーパス}{Corpus} & \T{生徒}{Student} & \T{出発点}{Start} & \T{K 本}{K-only} & \T{提案}{Proposal} \\
\midrule
AFEW & A & \dr{AFEW}{A}{start} & \dr{AFEW}{A}{none#1} & \dr{AFEW}{A}{prop#1} \\
 & V & \dr{AFEW}{V}{start} & \dr{AFEW}{V}{none#1} & \dr{AFEW}{V}{prop#1} \\
 & T & \dr{AFEW}{T}{start} & \dr{AFEW}{T}{none#1} & \dr{AFEW}{T}{prop#1} \\
 & AVT & \dr{AFEW}{AVT}{start} & \dr{AFEW}{AVT}{none#1} & \dr{AFEW}{AVT}{prop#1} \\
IEMOCAP & A & \dr{IEMOCAP}{A}{start} & \dr{IEMOCAP}{A}{none#1} & \dr{IEMOCAP}{A}{prop#1} \\
 & V & \dr{IEMOCAP}{V}{start} & \dr{IEMOCAP}{V}{none#1} & \dr{IEMOCAP}{V}{prop#1} \\
 & T & \dr{IEMOCAP}{T}{start} & \dr{IEMOCAP}{T}{none#1} & \dr{IEMOCAP}{T}{prop#1} \\
 & AVT & \dr{IEMOCAP}{AVT}{start} & \dr{IEMOCAP}{AVT}{none#1} & \dr{IEMOCAP}{AVT}{prop#1} \\
MELD & A & \dr{MELD}{A}{start} & \dr{MELD}{A}{none#1} & \dr{MELD}{A}{prop#1} \\
 & V & \dr{MELD}{V}{start} & \dr{MELD}{V}{none#1} & \dr{MELD}{V}{prop#1} \\
 & T & \dr{MELD}{T}{start} & \dr{MELD}{T}{none#1} & \dr{MELD}{T}{prop#1} \\
 & AVT & \dr{MELD}{AVT}{start} & \dr{MELD}{AVT}{none#1} & \dr{MELD}{AVT}{prop#1} \\
MOSEI & A & \dr{MOSEI}{A}{start} & \dr{MOSEI}{A}{none#1} & \dr{MOSEI}{A}{prop#1} \\
 & V & \dr{MOSEI}{V}{start} & \dr{MOSEI}{V}{none#1} & \dr{MOSEI}{V}{prop#1} \\
 & T & \dr{MOSEI}{T}{start} & \dr{MOSEI}{T}{none#1} & \dr{MOSEI}{T}{prop#1} \\
 & AVT & \dr{MOSEI}{AVT}{start} & \dr{MOSEI}{AVT}{none#1} & \dr{MOSEI}{AVT}{prop#1} \\
\bottomrule
\end{tabular}}

\begin{frame}[label=appdecomp]{\T{補足：位置の損 / 向きの損（ポイント，点の値）}{Appendix: position loss / direction loss (points, point values)}}
\rl{\T{実験 3}{Exp.~3}}
\begin{columns}[T]
\begin{column}{0.5\textwidth}\centering\tiny K = 5\par\decomptable{five}{5}\end{column}
\begin{column}{0.5\textwidth}\centering\tiny K = 25\par\decomptable{twofive}{25}\end{column}
\end{columns}
\vspace{0.1em}
{\tiny\T{位置 ＝ テストの感情の釣り合いで切片を当て直すと減る UAR，向き ＝ 当て直した後に残る上限との差．版 four（4 感情すべてで当て直す）．位置は同じテストの話者で当て直した量なので，大きく見せる向きの偏りがある．区間は results/dai3/eval/decomp.csv．}{Position = UAR recovered by refitting the intercept on the test emotion balance; direction = gap to the upper bound remaining after the refit. Version four (all 4 emotions refit). Position is refit on the same test speakers and is therefore biased upward. Intervals in results/dai3/eval/decomp.csv.}}
\end{frame}

\begin{frame}[label=approom]{\T{補足：直す余地，提案の得，再現できる度合い}{Appendix: room, gain, and reproducibility}}
\rl{\T{実験 3}{Exp.~3}}
\scriptsize\centering
\begin{tabular}{@{}llrrrrrr@{}}
\toprule
& & \multicolumn{3}{c}{K = 5} & \multicolumn{3}{c}{K = 25} \\
\T{コーパス}{Corpus} & \T{生徒}{Student} & \T{余地}{Room} & \T{得}{Gain} & \T{割合}{Share} & \T{余地}{Room} & \T{得}{Gain} & \T{割合}{Share} \\
\midrule
AFEW & A & \N{roomAFEWAfive} & \N{gainrAFEWAfive} & \N{shareAFEWAfive}\% & \N{roomAFEWAtwofive} & \N{gainrAFEWAtwofive} & \N{shareAFEWAtwofive}\% \\
 & V & \N{roomAFEWVfive} & \N{gainrAFEWVfive} & \N{shareAFEWVfive}\% & \N{roomAFEWVtwofive} & \N{gainrAFEWVtwofive} & \N{shareAFEWVtwofive}\% \\
 & T & \N{roomAFEWTfive} & \N{gainrAFEWTfive} & \N{shareAFEWTfive}\% & \N{roomAFEWTtwofive} & \N{gainrAFEWTtwofive} & \N{shareAFEWTtwofive}\% \\
IEMOCAP & A & \N{roomIEMOCAPAfive} & \N{gainrIEMOCAPAfive} & \N{shareIEMOCAPAfive}\% & \N{roomIEMOCAPAtwofive} & \N{gainrIEMOCAPAtwofive} & \N{shareIEMOCAPAtwofive}\% \\
 & V & \N{roomIEMOCAPVfive} & \N{gainrIEMOCAPVfive} & \N{shareIEMOCAPVfive}\% & \N{roomIEMOCAPVtwofive} & \N{gainrIEMOCAPVtwofive} & \N{shareIEMOCAPVtwofive}\% \\
 & T & \N{roomIEMOCAPTfive} & \N{gainrIEMOCAPTfive} & \N{shareIEMOCAPTfive}\% & \N{roomIEMOCAPTtwofive} & \N{gainrIEMOCAPTtwofive} & \N{shareIEMOCAPTtwofive}\% \\
MELD & A & \N{roomMELDAfive} & \N{gainrMELDAfive} & \N{shareMELDAfive}\% & \N{roomMELDAtwofive} & \N{gainrMELDAtwofive} & \N{shareMELDAtwofive}\% \\
 & V & \N{roomMELDVfive} & \N{gainrMELDVfive} & \N{shareMELDVfive}\% & \N{roomMELDVtwofive} & \N{gainrMELDVtwofive} & \N{shareMELDVtwofive}\% \\
 & T & \N{roomMELDTfive} & \N{gainrMELDTfive} & \N{shareMELDTfive}\% & \N{roomMELDTtwofive} & \N{gainrMELDTtwofive} & \N{shareMELDTtwofive}\% \\
MOSEI & A & \N{roomMOSEIAfive} & \N{gainrMOSEIAfive} & \N{shareMOSEIAfive}\% & \N{roomMOSEIAtwofive} & \N{gainrMOSEIAtwofive} & \N{shareMOSEIAtwofive}\% \\
 & V & \N{roomMOSEIVfive} & \N{gainrMOSEIVfive} & \N{shareMOSEIVfive}\% & \N{roomMOSEIVtwofive} & \N{gainrMOSEIVtwofive} & \N{shareMOSEIVtwofive}\% \\
 & T & \N{roomMOSEITfive} & \N{gainrMOSEITfive} & \N{shareMOSEITfive}\% & \N{roomMOSEITtwofive} & \N{gainrMOSEITtwofive} & \N{shareMOSEITtwofive}\% \\
\bottomrule
\end{tabular}

\vspace{0.2em}
{\tiny\T{余地 ＝ 上限（候補の全ラベルで学んだ重み，切片を当て直し）− K 本だけ．得 ＝ 提案 − K 本だけ．割合 ＝ 得 / 余地．点の値，区間なし．AVT は表に無い（台本の列名の衝突で壊れた表を行の並びから復元したため，A・V・T のみ）．}{Room = upper bound (weights trained on all pool labels, intercept refit) − K-only. Gain = proposal − K-only. Share = gain / room. Point values, no intervals. AVT is absent (the table was restored from row order after a column-name collision in the script; only A, V, T are available).}}
\end{frame}

\begin{frame}[label=appadv]{\T{補足：教師の先回りと得（K = 25），一致の割合}{Appendix: teacher lead and gain (K = 25), agreement shares}}
\rl{\T{探索}{Exploration}}
\begin{columns}[T]
\begin{column}{0.52\textwidth}
\begin{tikzpicture}
\begin{axis}[width=\linewidth, height=5.0cm, xmin=-2, xmax=33, ymin=-3.5, ymax=9.5,
  xlabel={\T{教師の先回り（ポイント）}{teacher lead (points)}}, ylabel={\T{提案の得（ポイント）}{gain (points)}},
  legend style={at={(0.98,0.03)}, anchor=south east, font=\tiny, cells={anchor=west}}, legend columns=2, grid=major,
  every node near coord/.append style={font=\tiny, anchor=west, xshift=2pt, color=mMuted}, point meta=explicit symbolic, nodes near coords]
\addplot[only marks, kA, mark=*, mark size=2pt] table[x=adv_pool, y=gain_prop, meta=T]{teacher_adv_A_K25.csv};
\addplot[only marks, kV, mark=square*, mark size=2pt] table[x=adv_pool, y=gain_prop, meta=T]{teacher_adv_V_K25.csv};
\addplot[only marks, kT, mark=triangle*, mark size=2.4pt] table[x=adv_pool, y=gain_prop, meta=T]{teacher_adv_T_K25.csv};
\addplot[only marks, kAVT, mark=diamond*, mark size=2.4pt] table[x=adv_pool, y=gain_prop, meta=T]{teacher_adv_AVT_K25.csv};
\legend{\T{音声}{audio}, \T{顔}{face}, \T{言葉}{text}, AVT}
\end{axis}
\end{tikzpicture}
\end{column}
\begin{column}{0.48\textwidth}
\scriptsize
\T{候補での教師と生徒（K 本だけ）の判定の一致の割合（\%，K = 5）}{Share of pool utterances where teacher and K-only student agree (\%, K = 5)}

\begin{tabular}{@{}lrrrr@{}}
\toprule
\T{コーパス}{Corpus} & \T{音声}{audio} & \T{顔}{face} & \T{言葉}{text} & AVT \\
\midrule
AFEW & \N{agreeAFEWAfive} & \N{agreeAFEWVfive} & \N{agreeAFEWTfive} & \N{agreeAFEWAVTfive} \\
IEMOCAP & \N{agreeIEMOCAPAfive} & \N{agreeIEMOCAPVfive} & \N{agreeIEMOCAPTfive} & \N{agreeIEMOCAPAVTfive} \\
MELD & \N{agreeMELDAfive} & \N{agreeMELDVfive} & \N{agreeMELDTfive} & \N{agreeMELDAVTfive} \\
MOSEI & \N{agreeMOSEIAfive} & \N{agreeMOSEIVfive} & \N{agreeMOSEITfive} & \N{agreeMOSEIAVTfive} \\
\bottomrule
\end{tabular}

\vspace{0.4em}
\T{教師が生徒の誤りを直す／正解を壊す割合（\%，K = 5，音声の生徒）}{Teacher fixes a student error / breaks a correct answer (\%, K = 5, audio student)}

\begin{tabular}{@{}lrr@{}}
\toprule
\T{コーパス}{Corpus} & \T{直す}{fixes} & \T{壊す}{breaks} \\
\midrule
AFEW & \N{fixAFEWAfive} & \N{harmAFEWAfive} \\
IEMOCAP & \N{fixIEMOCAPAfive} & \N{harmIEMOCAPAfive} \\
MELD & \N{fixMELDAfive} & \N{harmMELDAfive} \\
MOSEI & \N{fixMOSEIAfive} & \N{harmMOSEIAfive} \\
\bottomrule
\end{tabular}
\end{column}
\end{columns}
{\tiny\color{mMuted}\T{順位の相関：K = 5 で \N{rhofive}，K = 25 で \N{rhotwofive}（16 点，独立でない，p 値なし）．MOSEI は 39 単位．}{Rank correlation: \N{rhofive} at K = 5, \N{rhotwofive} at K = 25 (16 non-independent points, no p-value). MOSEI: 39 units.}}
\end{frame}

\begin{frame}[label=appcons]{\T{補足：生徒以外のモダリティの一致と教師の正解率（全生徒）}{Appendix: agreement of the other modalities vs teacher accuracy (all students)}}
\rl{\T{探索}{Exploration}}
\scriptsize\centering
\T{教師の正解率（\%）：一致する発話 / 一致しない発話（一致する発話の割合）}{Teacher accuracy (\%): agreeing / disagreeing utterances (share of agreeing utterances)}

\begin{tabular}{@{}llrrrr@{}}
\toprule
K & \T{生徒}{Student} & AFEW & IEMOCAP & MELD & MOSEI \\
\midrule
5 & \T{音声}{audio} & \N{cagreeAFEWAfive} / \N{cdisAFEWAfive} (\N{cshareAFEWAfive}\%) & \N{cagreeIEMOCAPAfive} / \N{cdisIEMOCAPAfive} (\N{cshareIEMOCAPAfive}\%) & \N{cagreeMELDAfive} / \N{cdisMELDAfive} (\N{cshareMELDAfive}\%) & \N{cagreeMOSEIAfive} / \N{cdisMOSEIAfive} (\N{cshareMOSEIAfive}\%) \\
 & \T{顔}{face} & \N{cagreeAFEWVfive} / \N{cdisAFEWVfive} (\N{cshareAFEWVfive}\%) & \N{cagreeIEMOCAPVfive} / \N{cdisIEMOCAPVfive} (\N{cshareIEMOCAPVfive}\%) & \N{cagreeMELDVfive} / \N{cdisMELDVfive} (\N{cshareMELDVfive}\%) & \N{cagreeMOSEIVfive} / \N{cdisMOSEIVfive} (\N{cshareMOSEIVfive}\%) \\
 & \T{言葉}{text} & \N{cagreeAFEWTfive} / \N{cdisAFEWTfive} (\N{cshareAFEWTfive}\%) & \N{cagreeIEMOCAPTfive} / \N{cdisIEMOCAPTfive} (\N{cshareIEMOCAPTfive}\%) & \N{cagreeMELDTfive} / \N{cdisMELDTfive} (\N{cshareMELDTfive}\%) & \N{cagreeMOSEITfive} / \N{cdisMOSEITfive} (\N{cshareMOSEITfive}\%) \\
 & AVT & \N{cagreeAFEWAVTfive} / \N{cdisAFEWAVTfive} (\N{cshareAFEWAVTfive}\%) & \N{cagreeIEMOCAPAVTfive} / \N{cdisIEMOCAPAVTfive} (\N{cshareIEMOCAPAVTfive}\%) & \N{cagreeMELDAVTfive} / \N{cdisMELDAVTfive} (\N{cshareMELDAVTfive}\%) & \N{cagreeMOSEIAVTfive} / \N{cdisMOSEIAVTfive} (\N{cshareMOSEIAVTfive}\%) \\
25 & \T{音声}{audio} & \N{cagreeAFEWAtwofive} / \N{cdisAFEWAtwofive} (\N{cshareAFEWAtwofive}\%) & \N{cagreeIEMOCAPAtwofive} / \N{cdisIEMOCAPAtwofive} (\N{cshareIEMOCAPAtwofive}\%) & \N{cagreeMELDAtwofive} / \N{cdisMELDAtwofive} (\N{cshareMELDAtwofive}\%) & \N{cagreeMOSEIAtwofive} / \N{cdisMOSEIAtwofive} (\N{cshareMOSEIAtwofive}\%) \\
 & \T{顔}{face} & \N{cagreeAFEWVtwofive} / \N{cdisAFEWVtwofive} (\N{cshareAFEWVtwofive}\%) & \N{cagreeIEMOCAPVtwofive} / \N{cdisIEMOCAPVtwofive} (\N{cshareIEMOCAPVtwofive}\%) & \N{cagreeMELDVtwofive} / \N{cdisMELDVtwofive} (\N{cshareMELDVtwofive}\%) & \N{cagreeMOSEIVtwofive} / \N{cdisMOSEIVtwofive} (\N{cshareMOSEIVtwofive}\%) \\
 & \T{言葉}{text} & \N{cagreeAFEWTtwofive} / \N{cdisAFEWTtwofive} (\N{cshareAFEWTtwofive}\%) & \N{cagreeIEMOCAPTtwofive} / \N{cdisIEMOCAPTtwofive} (\N{cshareIEMOCAPTtwofive}\%) & \N{cagreeMELDTtwofive} / \N{cdisMELDTtwofive} (\N{cshareMELDTtwofive}\%) & \N{cagreeMOSEITtwofive} / \N{cdisMOSEITtwofive} (\N{cshareMOSEITtwofive}\%) \\
 & AVT & \N{cagreeAFEWAVTtwofive} / \N{cdisAFEWAVTtwofive} (\N{cshareAFEWAVTtwofive}\%) & \N{cagreeIEMOCAPAVTtwofive} / \N{cdisIEMOCAPAVTtwofive} (\N{cshareIEMOCAPAVTtwofive}\%) & \N{cagreeMELDAVTtwofive} / \N{cdisMELDAVTtwofive} (\N{cshareMELDAVTtwofive}\%) & \N{cagreeMOSEIAVTtwofive} / \N{cdisMOSEIAVTtwofive} (\N{cshareMOSEIAVTtwofive}\%) \\
\bottomrule
\end{tabular}

\vspace{0.3em}
\T{「一致」＝生徒以外の 2 モダリティの学び直した識別器の判定が一致（AVT の生徒は 3 つすべて）．候補の全発話，評価の回．MELD の言葉は一致の有無で差が無い．}{``Agree'' = the retrained classifiers of the two modalities other than the student give the same label (all three for the AVT student). All pool utterances, evaluation round. For MELD text there is no difference.}
\end{frame}

\begin{frame}[label=appfusion]{\T{補足：融合の形の比べ（テスト UAR，50 単位の平均）}{Appendix: fusion forms compared (test UAR, mean over 50 units)}}
\rl{\T{探索}{Exploration}}
\scriptsize\centering
\begin{tabular}{@{}lrrrrrrrr@{}}
\toprule
& \multicolumn{2}{c}{AFEW} & \multicolumn{2}{c}{IEMOCAP} & \multicolumn{2}{c}{MELD} & \multicolumn{2}{c}{MOSEI} \\
K & 5 & 25 & 5 & 25 & 5 & 25 & 5 & 25 \\
\midrule
\T{AVT の生徒（提案）}{AVT student (proposal)} & \N{favtAFEWfive} & \N{favtAFEWtwofive} & \N{favtIEMOCAPfive} & \N{favtIEMOCAPtwofive} & \N{favtMELDfive} & \N{favtMELDtwofive} & \N{favtMOSEIfive} & \N{favtMOSEItwofive} \\
\T{積，指数 (1,1,1)}{Product, exponents (1,1,1)} & \N{fprodAFEWfive} & \N{fprodAFEWtwofive} & \N{fprodIEMOCAPfive} & \N{fprodIEMOCAPtwofive} & \N{fprodMELDfive} & \N{fprodMELDtwofive} & \N{fprodMOSEIfive} & \N{fprodMOSEItwofive} \\
\T{積，K 本で選んだ指数}{Product, exponents from the K labels} & \N{fkshotAFEWfive} & \N{fkshotAFEWtwofive} & \N{fkshotIEMOCAPfive} & \N{fkshotIEMOCAPtwofive} & \N{fkshotMELDfive} & \N{fkshotMELDtwofive} & \N{fkshotMOSEIfive} & \N{fkshotMOSEItwofive} \\
\T{積，塊ごとの LOO の指数}{Product, cluster-LOO exponents} & \N{fruleAFEWfive} & \N{fruleAFEWtwofive} & \N{fruleIEMOCAPfive} & \N{fruleIEMOCAPtwofive} & \N{fruleMELDfive} & \N{fruleMELDtwofive} & \N{fruleMOSEIfive} & \N{fruleMOSEItwofive} \\
\T{積，テストで選んだ上限}{Product, test-chosen upper bound} & \N{foracleAFEWfive} & \N{foracleAFEWtwofive} & \N{foracleIEMOCAPfive} & \N{foracleIEMOCAPtwofive} & \N{foracleMELDfive} & \N{foracleMELDtwofive} & \N{foracleMOSEIfive} & \N{foracleMOSEItwofive} \\
\bottomrule
\end{tabular}

\vspace{0.4em}
\T{%
\begin{itemize}
\item 積の中身：提案で直した A・V・T の生徒の確率．指数 e は信頼度 r（塊ごとの LOO の UAR）から e = max(r − 25, 0) / max の形
\item 上限は 63 点の格子（各指数 ∈ \{0, 0.5, 1, 2\}）でテストの平均が最大の点．手法の性能ではない
\item 発話ごとの LOO は話者が漏れる（塊ごとの LOO と比べて最強のモダリティを当てる割合が同じか少し高いが，MOSEI では 32〜45\% にとどまる）
\end{itemize}
}{%
\begin{itemize}
\item The product uses the probabilities of the adapted A, V, T students. Exponents e come from reliability r (cluster-LOO UAR) as e = max(r − 25, 0) / max
\item The upper bound is the test-mean maximum over a 63-point grid (each exponent ∈ \{0, 0.5, 1, 2\}); not a method's performance
\item Utterance-level LOO leaks speakers (its hit rate for the strongest modality is equal to or slightly above cluster LOO, but only 32--45\% in MOSEI)
\end{itemize}
}
\end{frame}

\begin{frame}[label=applimits]{\T{補足：限界として論文に書くこと}{Appendix: limitations to be stated in the paper}}
\rl{\T{限界}{Limitations}}
\scriptsize
\T{%
\begin{itemize}
\item 手法の形・格子・選び方は，評価に使った同じ 4 コーパスの仮のターゲットの回で作った．各コーパスのテストの分割は，ほかの 3 つのターゲットの設定を選ぶ仮のターゲットとしても使った
\item 評価の回の途中と集計の前に，著者は提案・相手・実験 2 の手法の UAR と差，実験 3 の表を見た（台帳 docs/seen-data.md）．集計の台本と設計は見る前に固定し，変えていない．実験 4 と \S6 の案は評価の回の解析から出た
\item 検定に残るセルは 128 セル中 15．族 4 の 40 セルは構造的に差が 0 で情報を持たない
\item 向きの物差しの確かめに不合格の 5 組では「向き」と呼ばない．IEMOCAP の顔にはマーカーの手がかりが入りうる．MOSEI の話者の代用は不完全．AFEW の言葉は Whisper の書き起こし
\item float32 で保存した識別器から作り直した判定は，8 問で 1〜2 本（最大 0.6\%）保存の判定と違う
\item エンコーダは 1 組（第 2 の組での繰り返しは保留）．凍結した特徴の上の線形の識別器に限る
\end{itemize}
}{%
\begin{itemize}
\item Method forms, grids and selection rules were built on proxy-target rounds over the same 4 corpora used for evaluation. Each corpus's test folds also served as proxy targets for selecting the other three targets' settings
\item During the evaluation round and before aggregation, the author looked at UARs and differences of the proposal, the baseline and Exp.~2 methods, and at Exp.~3 tables (ledger docs/seen-data.md). Aggregation scripts and designs were frozen beforehand and not changed. Exp.~4 and the \S6 idea arose from these analyses
\item 15 of 128 cells are tested; the 40 family-4 cells are structurally 0 and carry no information
\item ``Direction'' is not used in the 5 pairs failing the validity check. IEMOCAP face features may carry marker cues. MOSEI's speaker stand-in is imperfect. AFEW text is a Whisper transcript
\item Predictions rebuilt from float32-saved classifiers differ from the saved predictions by 1--2 utterances (≤0.6\%) in 8 problems
\item One encoder set (replication with a second set is on hold). Conclusions are limited to linear classifiers on frozen features
\end{itemize}
}
\end{frame}

\begin{frame}[label=appterms]{\T{補足：この資料の語と文書・台本の記号の対応}{Appendix: terms in this deck vs symbols in the documents and scripts}}
\rl{\T{語}{Terms}}
\scriptsize
\begin{tabular}{@{}p{0.33\linewidth}p{0.30\linewidth}p{0.30\linewidth}@{}}
\toprule
\T{この資料の語}{Term here} & \T{文書・台本の記号}{Symbol in documents / scripts} & \T{場所}{Where} \\
\midrule
\T{K 本だけ（比較の基準）}{K-only (baseline)} & \texttt{none[K]} & \T{実験 1 の設計 5 節}{Exp.~1 design §5} \\
\T{提案（同時に解く版）}{Proposal (joint version)} & \texttt{joint}\T{，主の版}{, main version} \texttt{one} & \T{同 7 節}{§7} \\
\T{温度つき蒸留}{Temperature distillation} & \texttt{kd} & \T{同 7 節}{§7} \\
\T{役割を分ける版／切片だけ}{Role-separated / intercept-only} & \texttt{rolesep} / \texttt{kbias} & \T{同 5・7 節}{§5, §7} \\
\T{教師（3 つの積）}{Teacher (product of three)} & \texttt{ret\_A}, \texttt{ret\_V}, \texttt{ret\_T} \T{の積}{product} & \T{同 6 節}{§6} \\
\T{引き戻し}{Pull} & \texttt{(C′, C\_b)}, \texttt{tether} & \T{dai1\_frozen.json}{dai1\_frozen.json} \\
\T{族 1〜5}{Families 1--5} & \texttt{families.csv}, \texttt{family5.csv} & results/dai1/stats, results/dai3/eval \\
\T{位置／向きの損}{Position / direction loss} & \texttt{position}, \texttt{direction} (\texttt{decomp.csv}) & results/dai3/eval \\
\T{教師の先回り}{Teacher lead} & \texttt{adv\_pool} & results/explore\_consensus \\
\T{一致で重みづけた疑似ラベル（実験 4）}{Agreement-weighted pseudo-labels (Exp.~4)} & \texttt{cons}, \T{印}{mark} \texttt{c\_i}, γ & DAI/docs/実験4\_設計.md \\
\T{塊ごとの LOO の信頼度／選択}{Cluster-LOO reliability / selection} & \texttt{cl\_uar\_m}, \texttt{select\_cl} & DAI/exp/explore\_fusion\_select.py \\
\bottomrule
\end{tabular}
\end{frame}

%% 補足（後半）：本編から移した区間つきの図，信頼度の筋，まとめ
\setband{\T{補足}{Appendix}}

\begin{frame}{\T{補足：提案 − K 本だけ（K = 0・5・25，区間つき）}{Appendix: proposal − K-only (K = 0, 5, 25, with intervals)}}
\rl{\T{実験 1，族 1}{Exp.~1, family 1}}
\hosoku{appfam1}{\T{全セルの数値表}{full table of cells}}
\centering
\famchart{fam1}{-8}{13}{\T{提案 − K 本だけ（ポイント）}{proposal − K-only (points)}}
\vspace{-0.3em}
\begin{columns}[T]
\begin{column}{0.5\textwidth}
\footnotesize
\T{%
\begin{itemize}
\item \textbf{音声}：12 セルすべて正，10 セルで区間の下端が 0 より上（AFEW \N{famoneAFEWAfive}，MELD \N{famoneMELDAfive}，MOSEI \N{famoneMOSEIAfive} ポイント，K = 5）
\item \textbf{言葉}：IEMOCAP だけ \N{famoneIEMOCAPTfive}（K = 5）．ほかは区間が 0 を含む
\end{itemize}
}{%
\begin{itemize}
\item \textbf{Audio}: all 12 cells positive; 10 have the interval above 0 (AFEW \N{famoneAFEWAfive}, MELD \N{famoneMELDAfive}, MOSEI \N{famoneMOSEIAfive} points at K = 5)
\item \textbf{Text}: only IEMOCAP, \N{famoneIEMOCAPTfive} (K = 5); elsewhere the interval contains 0
\end{itemize}
}
\end{column}
\begin{column}{0.5\textwidth}
\footnotesize
\T{%
\begin{itemize}
\item \textbf{顔}：下端が 0 より上なのは MOSEI の K = 25 だけ．IEMOCAP は負
\item \textbf{AVT}：MELD だけ \N{famoneMELDAVTfive}（K = 5）．ほかは −0.8〜+1.7
\item 検定で言えるのは MOSEI の音声（3 つの K，Holm 後 p < 0.001）だけ
\end{itemize}
}{%
\begin{itemize}
\item \textbf{Face}: the interval is above 0 only for MOSEI at K = 25; IEMOCAP is negative
\item \textbf{AVT}: only MELD, \N{famoneMELDAVTfive} (K = 5); elsewhere −0.8 to +1.7
\item Only MOSEI audio survives the tests (all three K, Holm p < 0.001)
\end{itemize}
}
\end{column}
\end{columns}
\end{frame}

\begin{frame}{\T{補足：提案 − 温度つき蒸留（同じ教師を受け取る相手）は同じ値になる}{Appendix: proposal − temperature distillation (same-teacher baseline) scores the same}}
\rl{\T{実験 1，族 3}{Exp.~1, family 3}}
\hosoku{appfam3}{\T{全セルの数値表}{full table of cells}}
\centering
\famchart{fam3}{-4}{4}{\T{提案 − 相手（ポイント）}{proposal − baseline (points)}}
\vspace{-0.3em}
\footnotesize
\T{%
\begin{itemize}
\item 48 セル中 40 セルで区間が 0 を含む．相手は 80 セル中 75 セルで温度つき蒸留（残りは「一致した発話だけ」）
\item 含まない 8 セル：提案が高いのは MELD の AVT K = 5（\N{famthreetwodMELDAVTfive}）だけ．相手が高いのは MOSEI の言葉（\N{famthreetwodMOSEITzero}〜\N{famthreetwodMOSEITtwofive}），MELD の言葉 K = 25（\N{famthreetwodMELDTtwofive}）など
\item 検定で言えるのは MOSEI の言葉 K = 0 で蒸留が高いこと（Holm 後 p = \N{famthreepMOSEITzero}）だけ．「このコーパスで提案は相手より良い」の文は書かない（事前の条件を満たさない）
\end{itemize}
}{%
\begin{itemize}
\item In 40 of 48 cells the interval contains 0. The baseline is temperature distillation in 75 of 80 cells (``agreed-only'' otherwise)
\item Of the 8 other cells, the proposal is higher only for MELD AVT K = 5 (\N{famthreetwodMELDAVTfive}); the baseline is higher for MOSEI text (\N{famthreetwodMOSEITzero} to \N{famthreetwodMOSEITtwofive}), MELD text K = 25 (\N{famthreetwodMELDTtwofive}), etc.
\item The only tested difference: distillation is higher for MOSEI text at K = 0 (Holm p = \N{famthreepMOSEITzero}). No ``proposal beats the baseline on this corpus'' sentence is written (pre-fixed condition not met)
\end{itemize}
}
\end{frame}

\begin{frame}{\T{補足：提案 − 標準の適応手法（AVT の生徒，K = 5，区間つきの図）}{Appendix: proposal − standard adaptation methods (AVT student, K = 5, chart with intervals)}}
\rl{\T{実験 2（K = 5，AVT の生徒）}{Exp.~2 (K = 5, AVT student)}}
\hosoku{appdai2}{\T{K = 0・25 の図}{charts for K = 0 and 25}}
\centering
\daichart{5}{-9}{22}{4.5cm}
\vspace{-0.2em}
\footnotesize
\T{%
\begin{itemize}
\item 横軸は「提案 − 手法」（ポイント，95\% 区間．手法ごとに使う情報が違うので優劣でなく情報の違いとして読む）．右が提案に有利．12 セル（4 コーパス × K ∈ \{0,5,25\}）で下端が 0 より上の数：エントロピー最小化 \N{daitwoaboveentmin}，IM \N{daitwoaboveim}，MME \N{daitwoabovemme}，T3A \N{daitwoabovetthreea}，CORAL \N{daitwoabovecoral}，DANN \N{daitwoabovedann}，自己学習 \N{daitwoaboveself}，S+T \N{daitwoabovest}
\item 提案が負けるセルは 4 つだけ：AFEW の K = 0 で T3A（\N{daitwoAFEWtthreeazero}），エントロピー最小化（\N{daitwoAFEWentminzero}），IM（\N{daitwoAFEWimzero}）と，MOSEI の K = 5 で CORAL（\N{daitwoMOSEIcoralfive}）
\end{itemize}
}{%
\begin{itemize}
\item x-axis: ``proposal − method'' (points, 95\% interval; methods use different information, so read as information differences, not a ranking); right favors the proposal. Cells (of 12 = 4 corpora × K ∈ \{0,5,25\}) with the lower end above 0: entropy min.\ \N{daitwoaboveentmin}, IM \N{daitwoaboveim}, MME \N{daitwoabovemme}, T3A \N{daitwoabovetthreea}, CORAL \N{daitwoabovecoral}, DANN \N{daitwoabovedann}, self-training \N{daitwoaboveself}, S+T \N{daitwoabovest}
\item The proposal loses in only 4 cells: AFEW at K = 0 to T3A (\N{daitwoAFEWtthreeazero}), entropy min.\ (\N{daitwoAFEWentminzero}), IM (\N{daitwoAFEWimzero}); MOSEI at K = 5 to CORAL (\N{daitwoMOSEIcoralfive})
\end{itemize}
}
\end{frame}

\begin{frame}{\T{補足：結果のまとめ（検定で言えるのは 2 つ，記述の傾向は 4 つ）}{Appendix: results in brief (two tested statements, four descriptive trends)}}
\rl{\T{まとめ}{Summary}}
\begin{columns}[T]
\begin{column}{0.5\textwidth}
\footnotesize
\textbf{\T{検定で言えること（Holm 後 p < 0.05）}{Tested statements (Holm p < 0.05)}}
\T{%
\begin{enumerate}
\item MOSEI の音声の生徒で，提案は K 本だけを上回る（K = 0・5・25 で \N{famoneMOSEIAzero}〜\N{famoneMOSEIAtwofive} ポイント）
\item MOSEI の言葉 K = 0 で，温度つき蒸留は提案より高い（\N{famthreetwodMOSEITzero}）
\end{enumerate}
\textbf{事前に立てた説明で支持されなかったもの}
\begin{itemize}
\item 「K 本は主に切片を直し，別のモダリティの予測が重みの向きを直す」→ K 本も向きを直していた（\S4）
\end{itemize}
}{%
\begin{enumerate}
\item For the MOSEI audio student, the proposal beats K-only (\N{famoneMOSEIAzero} to \N{famoneMOSEIAtwofive} points at K = 0, 5, 25)
\item For MOSEI text at K = 0, temperature distillation is higher than the proposal (\N{famthreetwodMOSEITzero})
\end{enumerate}
\textbf{Pre-registered explanation not supported}
\begin{itemize}
\item ``K labels mainly fix the intercept; the other modalities' predictions fix the weight direction'' → K labels also fix the direction (\S4)
\end{itemize}
}
\end{column}
\begin{column}{0.5\textwidth}
\footnotesize
\textbf{\T{区間つきの記述の傾向}{Descriptive trends (with intervals)}}
\T{%
\begin{enumerate}
\item 音声の生徒には 4 コーパスで効く（+1.4〜+4.9）
\item 言葉は IEMOCAP だけ（+5.0〜+7.5）．顔には効かない．AVT は MELD だけ
\item 提案と温度つき蒸留は同じ値
\item 標準の適応手法には多くで勝つが，AFEW の K = 0 ではラベル無しだけを使う手法が 4〜5 ポイント上
\end{enumerate}
}{%
\begin{enumerate}
\item Helps the audio student in all 4 corpora (+1.4 to +4.9)
\item Text only in IEMOCAP (+5.0 to +7.5). No gain for the face. AVT only in MELD
\item Proposal and temperature distillation score the same
\item Beats most standard adaptation methods, but at AFEW K = 0 unlabeled-only methods are 4--5 points higher
\end{enumerate}
}
\end{column}
\end{columns}
\takeaway{\T{次の問い：なぜ弱いモダリティにだけ効くのか．効く仕組みが分かれば，弱点も分かる（\S4，\S5）}{Next question: why only the weak modalities? Understanding the mechanism also reveals the weakness (\S4, \S5)}}
\end{frame}

\begin{frame}{\T{手がかりが無い：一致率でも，学習に使った K 本でも，指数を決められない}{No usable signal: neither agreement rates nor the training labels can set the exponents}}
\rl{\T{探索}{Exploration}}
\scriptsize
\begin{center}
\T{テスト UAR（AVT の生徒との差，ポイント）}{Test UAR (difference from the AVT student, points)}

\begin{tabular}{@{}lrrrrrrrr@{}}
\toprule
& \multicolumn{2}{c}{AFEW} & \multicolumn{2}{c}{IEMOCAP} & \multicolumn{2}{c}{MELD} & \multicolumn{2}{c}{MOSEI} \\
K & 5 & 25 & 5 & 25 & 5 & 25 & 5 & 25 \\
\midrule
\T{積，指数すべて 1}{Product, all exponents 1} & \N{fproddAFEWfive} & \N{fproddAFEWtwofive} & \N{fproddIEMOCAPfive} & \N{fproddIEMOCAPtwofive} & \N{fproddMELDfive} & \N{fproddMELDtwofive} & \N{fproddMOSEIfive} & \N{fproddMOSEItwofive} \\
\T{積，K 本で選んだ指数}{Product, exponents chosen on the K labels} & \N{fkshotdAFEWfive} & \N{fkshotdAFEWtwofive} & \N{fkshotdIEMOCAPfive} & \N{fkshotdIEMOCAPtwofive} & \N{fkshotdMELDfive} & \N{fkshotdMELDtwofive} & \N{fkshotdMOSEIfive} & \N{fkshotdMOSEItwofive} \\
\T{積，テストで選んだ上限}{Product, test-chosen upper bound} & \N{foracledAFEWfive} & \N{foracledAFEWtwofive} & \N{foracledIEMOCAPfive} & \N{foracledIEMOCAPtwofive} & \N{foracledMELDfive} & \N{foracledMELDtwofive} & \N{foracledMOSEIfive} & \N{foracledMOSEItwofive} \\
\midrule
\T{K 本で最強のモダリティを当てた割合}{K labels pick the strongest modality (\%)} & \N{fheavyAFEWfive}\% & \N{fheavyAFEWtwofive}\% & \N{fheavyIEMOCAPfive}\% & \N{fheavyIEMOCAPtwofive}\% & \N{fheavyMELDfive}\% & \N{fheavyMELDtwofive}\% & \N{fheavyMOSEIfive}\% & \N{fheavyMOSEItwofive}\% \\
\bottomrule
\end{tabular}
\end{center}
\scriptsize
\T{%
\begin{itemize}
\item \textbf{一致率}：強いモダリティほど弱いモダリティと意見が合わないので，候補での一致率は強さの手がかりにならない（AFEW で −0.4〜−3.4）
\item \textbf{K 本のラベル}：単独の識別器の学習に使った発話なので，どのモダリティも K 本の上では当たりすぎて差が見えない．AFEW で最強の顔を当てた単位は 50 中 1〜2
\item \textbf{仮のターゲット}で選んだ指数：強いモダリティがコーパスで違うので移らない
\end{itemize}
}{%
\begin{itemize}
\item \textbf{Agreement rate}: the stronger a modality, the more it disagrees with the weak ones, so pool agreement is no cue for strength (AFEW: −0.4 to −3.4)
\item \textbf{The K labels}: used to train the unimodal classifiers, so every modality looks too good on them. In AFEW the strongest modality (face) is picked in only 1--2 of 50 units
\item \textbf{Proxy targets}: the strongest modality differs per corpus, so exponents chosen elsewhere do not transfer
\end{itemize}
}
\takeaway{\T{難しさの核：同じ K 本を学習にも信頼度の測定にも使うと，測定が壊れる}{The crux: using the same K labels both for training and for measuring trust breaks the measurement}}
\end{frame}

\begin{frame}{\T{いまの案：K 本を「1 塊ずつ抜いて」学び直し，抜いた分で当たるかを数える（leave-one-out）}{Current idea: hold out the K labels one cluster at a time, retrain, and score the held-out part (leave-one-out)}}
\rl{\T{方法（探索）}{Method (exploration)}}
\centering
\begin{adjustbox}{max width=\linewidth}
\begin{tikzpicture}[font=\scriptsize, x=1cm, y=1cm]
  % K 本（塊ごと）
  \node[font=\scriptsize\bfseries, anchor=west] at (-0.2,1.6) {\T{K 本 × 4 感情のラベル（話者の塊ごとに色分け）}{K × 4 labels (colored by speaker cluster)}};
  \foreach \i/\c in {0/kA, 1/kA, 2/kV, 3/kV, 4/kV, 5/kT, 6/kAVT, 7/kAVT} {\node[utt, fill=\c!25] at (0.62*\i,0.9) {};}
  \draw[mBad, line width=1pt, rounded corners=2pt] (1.0,0.55) rectangle (2.95,1.25);
  \node[font=\tiny, mBad] at (2.0,0.3) {\T{塊 2 を抜く}{hold out cluster 2}};
  % 学び直し
  \node[clf, text width=2.5cm] (fit) at (6.3,0.9) {\T{残りの K 本だけで}{retrain A, V, T and AVT}\\ \T{A・V・T・AVT を学び直す}{on the remaining labels only}};
  \draw[arr] (4.6,0.9) -- (fit);
  \node[testbox, text width=2.3cm] (held) at (10.2,0.9) {\T{抜いた塊の発話に}{predict the held-out}\\ \T{予測 → 当たるか}{cluster → right or not}};
  \draw[arr] (fit) -- (held);
  \node[font=\tiny, color=mMuted, align=center] at (6.3,-0.1) {\T{これを塊の数だけ繰り返す（AFEW は 1 単位あたり約 12〜30 塊）}{repeat for every cluster (AFEW: ≈12--30 clusters per unit)}};
  % 出力
  \node[modbox, text width=3.6cm] (rel) at (2.0,-1.5) {\T{モダリティごとの信頼度}{per-modality reliability}\\ \T{＝抜いた発話での UAR}{= UAR on held-out labels}\\ \T{→ 積の指数 $e_A, e_V, e_T$}{→ exponents $e_A, e_V, e_T$}};
  \node[modbox, text width=3.6cm] (est) at (6.3,-1.5) {\T{積の見込み／AVT の見込み}{expected UAR of product / of AVT}\\ \T{（同じ抜いた発話で測る）}{(both measured on the held-out labels)}};
  \node[propbox, text width=3.2cm] (sel) at (10.2,-1.5) {\T{単位ごとに選ぶ：}{choose per unit:}\\ \T{見込みが高い方（積 か AVT）}{the form with the higher estimate}\\ \T{をテストで使う}{and use it on the test set}};
  \draw[arr] (held.south) |- ++(-0.5,-0.55) -| (rel.north);
  \draw[arr] (held.south) |- ++(-0.5,-0.55) -| (est.north);
  \draw[arr] (est) -- (sel);
  \draw[arr] (rel.east) -- (est.west);
\end{tikzpicture}
\end{adjustbox}
\notreal
\scriptsize
\T{%
\begin{itemize}
\item 発話 1 本ずつ抜く版は，同じ話者の残りの発話から学ぶので話者が漏れ，信頼度が高く出る．そこで\textbf{塊（話者）ごと}に抜く
\item テストのラベルは使わない．新しい学習は「K 本だけの小さな識別器」で，1 単位あたり数秒〜数十秒
\end{itemize}
}{%
\begin{itemize}
\item Holding out one utterance at a time leaks the speaker (its other labels stay in training) and inflates reliability, so we hold out \textbf{whole clusters (speakers)}
\item No test labels are used. The new training involves only small K-label classifiers (seconds per unit)
\end{itemize}
}
\end{frame}

\begin{frame}{\T{補足（信頼度の筋）結果 1：信頼度の見積もりは当たる．ただし「常に積」は AFEW でしか得をしない}{Appendix (trust line) result 1: the reliability estimate is right, but ``always product'' gains only in AFEW}}
\rl{\T{結果 1}{Result 1}}
\scriptsize
\begin{center}
\begin{tabular}{@{}lrrrrrrrr@{}}
\toprule
& \multicolumn{2}{c}{AFEW} & \multicolumn{2}{c}{IEMOCAP} & \multicolumn{2}{c}{MELD} & \multicolumn{2}{c}{MOSEI} \\
K & 5 & 25 & 5 & 25 & 5 & 25 & 5 & 25 \\
\midrule
\T{最強のモダリティを当てた割合（塊ごとの LOO）}{Strongest modality identified (cluster LOO, \%)} & \N{fclokAFEWfive}\% & \N{fclokAFEWtwofive}\% & \N{fclokIEMOCAPfive}\% & \N{fclokIEMOCAPtwofive}\% & \N{fclokMELDfive}\% & \N{fclokMELDtwofive}\% & \N{fclokMOSEIfive}\% & \N{fclokMOSEItwofive}\% \\
\midrule
\T{AVT の生徒（テスト UAR）}{AVT student (test UAR)} & \N{favtAFEWfive} & \N{favtAFEWtwofive} & \N{favtIEMOCAPfive} & \N{favtIEMOCAPtwofive} & \N{favtMELDfive} & \N{favtMELDtwofive} & \N{favtMOSEIfive} & \N{favtMOSEItwofive} \\
\T{常に積，LOO の信頼度の指数（差）}{Always product, LOO exponents (diff.)} & \N{fruledAFEWfive} & \N{fruledAFEWtwofive} & \N{fruledIEMOCAPfive} & \N{fruledIEMOCAPtwofive} & \N{fruledMELDfive} & \N{fruledMELDtwofive} & \N{fruledMOSEIfive} & \N{fruledMOSEItwofive} \\
\T{テストで選んだ上限（差）}{Test-chosen upper bound (diff.)} & \N{foracledAFEWfive} & \N{foracledAFEWtwofive} & \N{foracledIEMOCAPfive} & \N{foracledIEMOCAPtwofive} & \N{foracledMELDfive} & \N{foracledMELDtwofive} & \N{foracledMOSEIfive} & \N{foracledMOSEItwofive} \\
\bottomrule
\end{tabular}
\end{center}
\footnotesize
\T{%
\begin{itemize}
\item K = 25 では，塊ごとの LOO が最強のモダリティを AFEW・IEMOCAP・MELD で \N{fclokMELDtwofive}〜\N{fclokAFEWtwofive}\% 当てる．MOSEI は \N{fclokMOSEItwofive}\%（モダリティの差が小さく，塊も粗い）
\item 常に積を使うと AFEW で \N{fruledAFEWfive}・\N{fruledAFEWtwofive}（上限との差 1〜3），MOSEI で \N{fruledMOSEIfive}・\N{fruledMOSEItwofive}．MOSEI は上限でも AVT に届かず，\textbf{積という形そのものが負ける}
\item そこで「積か AVT か」を単位ごとに同じ LOO で選ぶ（前のページの右端）．AFEW で積，MOSEI で AVT を選べれば，4 コーパスで AVT を下回らない形になる → 次のページ
\end{itemize}
}{%
\begin{itemize}
\item At K = 25, cluster LOO identifies the strongest modality in \N{fclokMELDtwofive}--\N{fclokAFEWtwofive}\% of units for AFEW, IEMOCAP and MELD; MOSEI \N{fclokMOSEItwofive}\% (small modality differences, coarse clusters)
\item Always using the product gives \N{fruledAFEWfive} / \N{fruledAFEWtwofive} in AFEW (1--3 below the upper bound) but \N{fruledMOSEIfive} / \N{fruledMOSEItwofive} in MOSEI, where even the upper bound is below AVT: \textbf{the product form itself loses there}
\item So we choose ``product or AVT'' per unit with the same LOO (right end of the previous page). Picking product in AFEW and AVT in MOSEI would give a form that never falls below AVT → next page
\end{itemize}
}
\end{frame}

\begin{frame}{\T{補足（信頼度の筋）結果 2：LOO で選ぶ形は AFEW で得をするが，MOSEI の損を消せない．まだ使える手法ではない}{Appendix (trust line) result 2: the LOO-selected form gains in AFEW but does not remove the MOSEI loss; not yet a usable method}}
\rl{\T{結果 2（4 コーパス）}{Result 2 (4 corpora)}}
\scriptsize
\begin{center}
\T{テスト UAR（AVT の生徒との差，ポイント．50 単位の平均）}{Test UAR (difference from the AVT student, points; mean over 50 units)}

\begin{tabular}{@{}lrrrrrrrr@{}}
\toprule
& \multicolumn{2}{c}{AFEW} & \multicolumn{2}{c}{IEMOCAP} & \multicolumn{2}{c}{MELD} & \multicolumn{2}{c}{MOSEI} \\
K & 5 & 25 & 5 & 25 & 5 & 25 & 5 & 25 \\
\midrule
\T{AVT の生徒（UAR そのもの）}{AVT student (absolute UAR)} & \N{savtAFEWfive} & \N{savtAFEWtwofive} & \N{savtIEMOCAPfive} & \N{savtIEMOCAPtwofive} & \N{savtMELDfive} & \N{savtMELDtwofive} & \N{savtMOSEIfive} & \N{savtMOSEItwofive} \\
\T{常に積（差）}{Always product (diff.)} & \N{sruledAFEWfive} & \N{sruledAFEWtwofive} & \N{sruledIEMOCAPfive} & \N{sruledIEMOCAPtwofive} & \N{sruledMELDfive} & \N{sruledMELDtwofive} & \N{sruledMOSEIfive} & \N{sruledMOSEItwofive} \\
\textbf{\T{LOO で選ぶ（差）}{LOO-selected (diff.)}} & \textbf{\N{sseldAFEWfive}} & \textbf{\N{sseldAFEWtwofive}} & \textbf{\N{sseldIEMOCAPfive}} & \textbf{\N{sseldIEMOCAPtwofive}} & \textbf{\N{sseldMELDfive}} & \textbf{\N{sseldMELDtwofive}} & \textbf{\N{sseldMOSEIfive}} & \textbf{\N{sseldMOSEItwofive}} \\
\T{単位ごとにテストで良い方（差，上限）}{Test-best per unit (diff., upper bound)} & \N{soracledAFEWfive} & \N{soracledAFEWtwofive} & \N{soracledIEMOCAPfive} & \N{soracledIEMOCAPtwofive} & \N{soracledMELDfive} & \N{soracledMELDtwofive} & \N{soracledMOSEIfive} & \N{soracledMOSEItwofive} \\
\midrule
\T{積を選んだ単位の割合}{Units choosing the product} & \N{sshareAFEWfive}\% & \N{sshareAFEWtwofive}\% & \N{sshareIEMOCAPfive}\% & \N{sshareIEMOCAPtwofive}\% & \N{sshareMELDfive}\% & \N{sshareMELDtwofive}\% & \N{sshareMOSEIfive}\% & \N{sshareMOSEItwofive}\% \\
\bottomrule
\end{tabular}
\end{center}
\vspace{-0.6em}
\scriptsize
\T{%
\begin{itemize}
\item MOSEI の損は「常に積」の半分ほどに減ったが消えず，4 コーパスの平均は AVT とほぼ同じ．理由：MOSEI で積がテストで良い単位は 3〜5\% しか無いのに，LOO の見込みは積のほうが高く出て（\N{sestpMOSEIfive} 対 \N{sestaMOSEIfive}，K = 5），\N{sshareMOSEIfive}\% の単位で積を選ぶ
\item 偏りの原因は 2 つ：(1) 指数と見込みを同じ抜いた発話から作る（楽観）．(2) 見込みは K 本だけの識別器で測るが，テストで使うのは直した識別器で，この差は 2,304 次元の AVT で大きい
\item 信頼度の見積もりは当たるが，その上の「積か AVT か」の選択はまだ使える形ではない．続けるなら，偏りを直す版を設計して検査に回す
\end{itemize}
}{%
\begin{itemize}
\item The MOSEI loss is roughly halved but not removed; averaged over corpora the form is about level with AVT. Why: in MOSEI the product is better on test in only 3--5\% of units, yet the LOO estimate favors the product (\N{sestpMOSEIfive} vs \N{sestaMOSEIfive} at K = 5) and picks it in \N{sshareMOSEIfive}\% of units
\item Two sources of the bias: (1) exponents and estimate come from the same held-out labels (optimistic); (2) the estimate uses K-label-only classifiers while the test uses adapted ones, a gap largest for the 2,304-dim AVT
\item The reliability estimate is right, but the product-vs-AVT selection on top of it is not yet usable. If this line continues, a bias-corrected version must be designed and checked
\end{itemize}
}
{\tiny\color{mMuted}\T{出どころ：figures/select.csv（results/explore\_fusion/select\_cells.csv．本番の台本 DAI/exp/explore\_fusion\_select.py，2026-11-05 10:17）．AFEW の値は保存ずみの対数確率から先に出した途中の値と小数 2 桁まで一致}{Source: figures/select.csv (results/explore\_fusion/select\_cells.csv; script DAI/exp/explore\_fusion\_select.py, 2026-11-05 10:17). AFEW values agree to 2 decimals with the interim values computed earlier from the saved log-probabilities}}
\end{frame}

\begin{frame}{\T{いまの状態の一覧}{Where things stand}}
\rl{\T{状態}{Status}}
\footnotesize
\begin{tabular}{@{}p{0.44\linewidth}p{0.22\linewidth}p{0.30\linewidth}@{}}
\toprule
\T{項目}{Item} & \T{状態}{Status} & \T{次}{Next} \\
\midrule
\T{実験 1〜3 の評価の回，集計，結果の文の照合}{Exp.~1--3 evaluation, aggregation, claim check} & \done & \T{論文の結果の節に写す}{move into the paper's results section} \\
\T{得の仕組みの探索（先回り，一致，位置と向き）}{Mechanism exploration (lead, agreement, position/direction)} & \done & \T{論文の「説明」の節（事前の解析と分けて）}{``explanation'' section, separated from pre-fixed analyses} \\
\T{実験 4（一致で重みづけた疑似ラベル）}{Exp.~4 (agreement-weighted pseudo-labels)} & \todo\ \T{設計・台本は検査済み}{design and code checked} & \T{研究者の判断で起動（H200，4〜6 時間）}{launch at the researcher's decision (H200, 4--6 h)} \\
\T{LOO による信頼度と，積／AVT の選択}{LOO reliability and product/AVT selection} & \done\ \T{4 コーパス}{4 corpora} & \T{続けるかを研究者が決める（見込みの偏りを直す版の設計が要る）}{researcher decides whether to continue (a version that removes the estimate bias would be needed)} \\
\T{第 2 のエンコーダでの繰り返し}{Replication with a second encoder set} & \todo & \T{保留（研究者「まだいい」）}{on hold (researcher: ``not yet'')} \\
\T{論文の型（手法の論文 → 測定と説明の論文）と骨組みの項目 3・4・6 の書き直し}{Paper type (method paper → measurement-and-explanation paper), rewriting skeleton items 3, 4, 6} & \wip & \T{\S7 の案を研究者が決める}{researcher decides on the \S7 drafts} \\
\bottomrule
\end{tabular}
\end{frame}

\begin{frame}{\T{まとめ：骨組みの 6 項目に，いまの結果で答える}{Summary: answering the six skeleton items with the current results}}
\rl{\T{まとめ}{Summary}}
\scriptsize
\begin{tabular}{@{}cp{0.92\linewidth}@{}}
\toprule
1 & \T{新しいコーパスでは感情ごとに数本〜数十本のラベルしか用意できず，ほかのコーパスの識別器とラベル無しの発話で補う必要がある}{A new corpus offers only a few labels per emotion; classifiers from other corpora and unlabeled utterances must fill the gap} \\
2 & \T{ソースの識別器の平均は最良の線形識別器より \N{gapstartMELD}〜\N{gapstartAFEW} ポイント低く，14 の標準手法とラベル無しの発話を足しても融合の識別器では 1〜3 ポイントしか埋まらない}{The mean of source classifiers is \N{gapstartMELD}--\N{gapstartAFEW} points below the best linear classifier; 14 standard methods plus unlabeled data close only 1--3 points for the fusion classifier} \\
3 & \T{既存の評価は「上がるか」だけを測り，どのモダリティで，損のどの部分（切片か重みの向きか）が減るかを測らない}{Existing evaluations only ask whether it helps, not for which modality or which part of the loss (intercept or weight direction)} \\
4 & \T{\textbf{（書き直し案）}4 コーパスで損を切片と向きに分け，3 モダリティの積の予測で学ぶ得が「教師が生徒より当たる弱いモダリティでの，向きの修正」として現れ，融合と顔では現れないことを示す}{\textbf{(Draft)} Across 4 corpora, split the loss into intercept and direction, and show that the gain from learning on the product of three modalities appears as a direction correction in weak modalities where the teacher leads the student, and not for the fusion or the face} \\
5 & \T{UAR．4 コーパスを順にターゲットにし，テストの話者を学習側に入れない}{UAR; each corpus in turn as the target; test speakers never in training} \\
6 & \T{比較：K 本だけ，温度つき蒸留，自己学習・DANN など 14 手法．除去：教師の学び直し，単一の教師，総量そろえ，切片だけ}{Baselines: K-only, temperature distillation, 14 methods incl.\ self-training and DANN. Ablations: teacher retraining, single-modality teacher, balancing, intercept-only} \\
\bottomrule
\end{tabular}
\vspace{0.3em}
\footnotesize
\T{%
\begin{itemize}
\item 手法の新しさ（提案 ≈ 蒸留）は結果に支えられない．説明の新しさ（どこで・なぜ効くか）は支えられる．論文の型を「測定と説明」にすることを推す
\item 「どのモダリティを信じるか」の見積もりは当たるが，その上の積／AVT の選択は MOSEI の損を消せなかった（\S6）．続けるかは研究者の判断
\end{itemize}
}{%
\begin{itemize}
\item Method novelty (proposal ≈ distillation) is not supported by the results; explanatory novelty (where and why it helps) is. We recommend a ``measurement and explanation'' paper
\item The trust estimate is right, but the product/AVT selection built on it did not remove the MOSEI loss (\S6). Whether to continue is the researcher's call
\end{itemize}
}
\end{frame}

\end{document}
